% Generated mechanically from the frozen varieties.tex target.
% Do not edit this generated file; edit the bound locale source instead.

% OK, start here.
%

\setcounter{chapter}{32}
\chapter{簇}



\phantomsection
\label{varieties-section-phantom}


\section{引言}
\label{varieties-section-introduction}

\noindent
本章开始研究簇，并更一般地研究域上的概形。一个基本参考文献是
\cite{EGA}。








\section{记号}
\label{varieties-section-notation}

\noindent
本章始终用字母 $k$ 表示基域。








\section{簇}
\label{varieties-section-varieties}

\noindent
在 Stacks 计划中，我们采用如下簇的定义。

\begin{definition}
\label{varieties-definition-variety}
设 $k$ 为域。所谓一个{\it 簇}，是指 $X$ 是域 $k$ 上的概形，且
$X$ 为整概形，且结构态射
$X \to \Spec(k)$ 分离并且是有限型的。
\end{definition}

\noindent
这个定义有如下缺点。设 $k'/k$ 是域扩张，并设 $X$ 是域 $k$ 上的簇。
则基变换
$X_{k'} = X \times_{\Spec(k)} \Spec(k')$
未必是域 $k'$ 上的簇。后续各节将以更一般的形式详细讨论这一现象。
两个簇的积未必是簇（本质上这正是同一现象）。下面给出一个例子。

\begin{example}
\label{varieties-example-product-not-a-variety}
令 $k = \mathbf{Q}$。令 $X = \Spec(\mathbf{Q}(i))$ 且
$Y = \Spec(\mathbf{Q}(i))$。则积
$X \times_{\Spec(k)} Y$，亦即簇 $X$ 与 $Y$ 的积，不是簇，因为它是可约的。（它同构于两个
$X$ 的不交并。）
\end{example}

\noindent
然而，如果基域是代数闭域，那么簇的积仍是簇。这可由《代数》章中的
结果推出，不过那里处理的是远为一般的情形。这里也有一个简单的直接
证明，我们将其写出。

\begin{lemma}
\label{varieties-lemma-product-varieties}
\begin{slogan}
代数闭域上簇的积仍是簇。
\end{slogan}
设 $k$ 为代数闭域。设 $X$, $Y$ 为域 $k$ 上的簇。则
$X \times_{\Spec(k)} Y$ 是域 $k$ 上的簇。
\end{lemma}

\begin{proof}
态射 $X \times_{\Spec(k)} Y \to \Spec(k)$ 是有限型且分离的，
因为它是态射
$X \times_{\Spec(k)} Y \to Y \to \Spec(k)$
的复合，而这两个态射都分离且为有限型；见《态射》引理
\ref{morphisms-lemma-base-change-finite-type} 和
\ref{morphisms-lemma-composition-finite-type}，以及《概形》引理
\ref{schemes-lemma-separated-permanence}。
为完成证明，只需证明 $X \times_{\Spec(k)} Y$ 是整概形。
取有限仿射开覆盖
$X = \bigcup_{i = 1, \ldots, n} U_i$、
$Y = \bigcup_{j = 1, \ldots, m} V_j$。如果能证明每个
$U_i \times_{\Spec(k)} V_j$ 都是整概形，则由《性质》引理
\ref{properties-lemma-characterize-reduced}、
\ref{properties-lemma-characterize-irreducible} 和
\ref{properties-lemma-characterize-integral} 即得结论。
所以问题约化到仿射情形。

\medskip\noindent
仿射情形可转化为如下代数命题：设 $A$, $B$ 是整环且为有限生成
$k$-代数，则 $A \otimes_k B$ 是整环。反设不然，并设
$$
(\sum\nolimits_{i = 1, \ldots, n} a_i \otimes b_i)
(\sum\nolimits_{j = 1, \ldots, m} c_j \otimes d_j) = 0
$$
在 $A \otimes_k B$ 中成立，且两个因子在 $A \otimes_k B$ 中均非零。
不妨设 $b_1, \ldots, b_n$ 关于 $k$ 在 $B$ 中线性无关，并且
$d_1, \ldots, d_m$ 关于 $k$ 在 $B$ 中线性无关。当然，也可设
$a_1$ 与 $c_1$ 在 $A$ 中非零。因此
$D(a_1c_1) \subset \Spec(A)$ 非空。由 Hilbert 零点定理
（《代数》定理 \ref{algebra-theorem-nullstellensatz}），可找到极大理想
$\mathfrak m \subset A$，它包含于 $D(a_1c_1)$，且
$A/\mathfrak m = k$；这是因为 $k$ 代数闭。用
$\overline{a}_i, \overline{c}_j$ 表示 $a_i, c_j$ 在
$A/\mathfrak m = k$ 中的剩余类。上述等式变为
$$
(\sum\nolimits_{i = 1, \ldots, n} \overline{a}_i b_i)
(\sum\nolimits_{j = 1, \ldots, m} \overline{c}_j d_j) = 0
$$
，这与 $\mathfrak m \in D(a_1c_1)$、
$b_1, \ldots, b_n$ 和 $d_1, \ldots, d_m$ 的线性无关性以及
$B$ 是整环这一事实矛盾。
\end{proof}






\section{簇与有理映射}
\label{varieties-section-varieties-rational-maps}

\noindent
设 $k$ 为域。设 $X$ 和 $Y$ 为域 $k$ 上的簇。我们所说的
{\it 从 $X$ 到 $Y$ 的簇有理映射}，是指按照《态射》定义
\ref{morphisms-definition-rational-map} 的 $\Spec(k)$-有理映射，
其源概形为 $X$，目标概形为 $Y$。依照惯例，“簇有理映射”这一说法并不指明这些簇
（共同）的基域；但对于一般概形，我们会区分有理映射和给定基上的
有理映射。

\medskip\noindent
本节标题所指的是下面这个基本定理。

\begin{theorem}
\label{varieties-theorem-varieties-rational-maps}
设 $k$ 为域。由簇及占优有理映射构成的范畴等价于有限生成域扩张
$K/k$ 构成的范畴。
\end{theorem}

\begin{proof}
设 $X$ 和 $Y$ 为簇，其一般点分别为 $x \in X$ 和 $y \in Y$。
回忆，从 $X$ 到 $Y$ 的占优有理映射，恰好是把 $x$ 映到 $y$ 的那些
有理映射（见《态射》定义 \ref{morphisms-definition-dominant-rational}
及其后的讨论）。因此，给定占优有理映射 $X \supset U \to Y$，便得到
函数域的映射
$$
k(Y) = \kappa(y) = \mathcal{O}_{Y, y}
\longrightarrow
\mathcal{O}_{X, x} = \kappa(x) = k(X)
$$
反之，这样的 $k$-代数映射（由于源和目标都是域，它自动是局部的）
由《态射》引理 \ref{morphisms-lemma-rational-map-finite-presentation}
唯一确定一个占优有理映射。由此得到一个全忠实函子。为完成证明，
只需证明每个有限生成域扩张 $K/k$ 都在其本质像中。由于 $K/k$ 有限生成，
存在有限型 $k$-代数 $A \subset K$，使 $K$ 是 $A$ 的分式域。于是
$X = \Spec(A)$ 是函数域为 $K$ 的簇。
\end{proof}

\noindent
设 $k$ 为域。设 $X$ 和 $Y$ 为域 $k$ 上的簇。我们所说的
{\it $X$ 与 $Y$ 是双有理簇}，是指按照《态射》定义
\ref{morphisms-definition-birational}，$X$ 与 $Y$ 是
$\Spec(k)$-双有理的。依照惯例，“双有理簇”这一说法并不指明这些簇
（共同）的基域；但对于一般不可约概形，我们会区分双有理和给定基上
双有理。

\begin{lemma}
\label{varieties-lemma-birational-varieties}
设 $X$ 与 $Y$ 是域 $k$ 上的簇。下列条件等价：
\begin{enumerate}
\item $X$ 与 $Y$ 是双有理簇；
\item 函数域 $k(X)$ 与 $k(Y)$ 同构；
\item $X$ 与 $Y$ 各有一个非空开集，二者作为簇同构；
\item 存在开集 $U \subset X$ 以及簇的双有理态射 $U \to Y$。
\end{enumerate}
\end{lemma}

\begin{proof}
这是《态射》引理
\ref{morphisms-lemma-criterion-birational-finite-presentation} 的特殊情形。
\end{proof}






\section{域变换与局部环}
\label{varieties-section-local-rings}

\noindent
先给出基域扩张时局部环所发生变化的一些预备结果。

\begin{lemma}
\label{varieties-lemma-change-fields-flat}
设 $K/k$ 为域扩张。设 $X$ 为域 $k$ 上的概形，并令 $Y = X_K$。
若 $y \in Y$ 的像为 $x \in X$，则：
\begin{enumerate}
\item $\mathcal{O}_{X, x} \to \mathcal{O}_{Y, y}$ 是忠实平坦的局部环同态；
\item 令 $\mathfrak p_0 = \Ker(\kappa(x) \otimes_k K \to \kappa(y))$，
则 $\kappa(y) = \kappa(\mathfrak p_0)$；
\item $\mathcal{O}_{Y, y} = (\mathcal{O}_{X, x} \otimes_k K)_\mathfrak p$，
其中 $\mathfrak p \subset \mathcal{O}_{X, x} \otimes_k K$ 是
$\mathfrak p_0$ 的逆像；
\item 有
$\mathcal{O}_{Y, y}/\mathfrak m_x\mathcal{O}_{Y, y} =
(\kappa(x) \otimes_k K)_{\mathfrak p_0}$。
\end{enumerate}
\end{lemma}

\begin{proof}
不妨设 $X = \Spec(A)$ 为仿射概形。于是 $Y = \Spec(A \otimes_k K)$。
由于 $K$ 在 $k$ 上平坦，$A \to A \otimes_k K$ 平坦。因此
$Y \to X$ 平坦；再用《代数》引理
\ref{algebra-lemma-local-flat-ff}，便得到第一项。第二项由《概形》引理
\ref{schemes-lemma-points-fibre-product} 得出。现在 $y$ 对应于素理想
$\mathfrak q \subset A \otimes_k K$，而 $x$ 对应于
$\mathfrak r = A \cap \mathfrak q$。于是 $\mathfrak p_0$ 是诱导映射
$\kappa(\mathfrak r) \otimes_k K \to \kappa(\mathfrak q)$ 的核。
局部环上的映射为
$$
A_\mathfrak r \longrightarrow (A \otimes_k K)_\mathfrak q
$$
我们可将此映射分解经过
$A_\mathfrak r \otimes_k K = (A \otimes_k K)_{\mathfrak r}$，得到
$$
A_\mathfrak r \longrightarrow A_\mathfrak r \otimes_k K
\longrightarrow (A \otimes_k K)_\mathfrak q
$$
而第二个箭头是关于某个素理想的局部化。这个素理想正是
$\mathfrak p_0$ 的逆像（略去细节），这就证明了 (3)。为证明 (4)，
使用 (3)，以及局部化和 $- \otimes_k K$ 都是正合函子这一事实。
\end{proof}

\begin{lemma}
\label{varieties-lemma-change-fields-algebraic-dim}
采用引理 \ref{varieties-lemma-change-fields-flat} 的记号。设 $X$ 在 $k$ 上局部有限型。
则
$$
\dim(\mathcal{O}_{Y, y}/\mathfrak m_x\mathcal{O}_{Y, y}) =
\text{trdeg}_k(\kappa(x)) - \text{trdeg}_K(\kappa(y)) =
\dim(\mathcal{O}_{Y, y}) - \dim(\mathcal{O}_{X, x})
$$
\end{lemma}

\begin{proof}
这只是《代数》引理
\ref{algebra-lemma-inequalities-under-field-extension} 的改述。
\end{proof}

\begin{lemma}
\label{varieties-lemma-change-fields-algebraic-unramified}
采用引理 \ref{varieties-lemma-change-fields-flat} 的记号。设 $X$ 在 $k$ 上局部有限型，
并且 $\dim(\mathcal{O}_{X, x}) = \dim(\mathcal{O}_{Y, y})$，且
$\kappa(x) \otimes_k K$ 约化（例如，当 $\kappa(x)/k$ 可分或
$K/k$ 可分时）。则 $\mathfrak m_x \mathcal{O}_{Y, y} = \mathfrak m_y$。
\end{lemma}

\begin{proof}
（括号内的断言由《代数》引理
\ref{algebra-lemma-separable-extension-preserves-reducedness} 得出。）
结合引理 \ref{varieties-lemma-change-fields-flat} 与
\ref{varieties-lemma-change-fields-algebraic-dim}，可知
$\mathcal{O}_{Y, y}/\mathfrak m_x \mathcal{O}_{Y, y}$ 的维数为 $0$，
并且它约化。因此它是域。
\end{proof}





\section{几何约化概形}
\label{varieties-section-geometrically-reduced}

\noindent
如果 $X$ 是域上的约化概形，那么扩张基域之后，$X$ 可能变成非约化的。
对于几何约化概形，这种情况不会发生。

\begin{definition}
\label{varieties-definition-geometrically-reduced}
设 $k$ 为域。设 $X$ 为 $k$ 上的概形。
\begin{enumerate}
\item 设 $x \in X$ 为一点。称 $X$ 在 $x$ 处{\it 几何约化}，如果对于
任意域扩张 $k'/k$ 以及任意点 $x' \in X_{k'}$，只要它位于 $x$ 上方，
局部环 $\mathcal{O}_{X_{k'}, x'}$ 就约化。
\item 称 $X$ 在 $k$ 上{\it 几何约化}，如果 $X$ 在 $X$ 的每一点处
都几何约化。
\end{enumerate}
\end{definition}

\noindent
初看起来，这个概念可能略显神秘，但它实际上正是《代数》章中讨论的
概念。下面给出一些说明其联系的基本结果。另见引理
\ref{varieties-lemma-geometrically-reduced-dense-smooth-open}；该引理说明，对域上的簇，
“几何约化”等价于“一般光滑”。

\begin{lemma}
\label{varieties-lemma-geometrically-reduced-at-point}
\begin{slogan}
几何约化性可在局部环上检验。
\end{slogan}
设 $k$ 为域。设 $X$ 为 $k$ 上的概形。设 $x \in X$。下列条件等价：
\begin{enumerate}
\item $X$ 在 $x$ 处几何约化；
\item 环 $\mathcal{O}_{X, x}$ 在 $k$ 上几何约化（见《代数》定义
\ref{algebra-definition-geometrically-reduced}）。
\end{enumerate}
\end{lemma}

\begin{proof}
设 (1) 成立。这尤其说明 $\mathcal{O}_{X, x}$ 约化。设 $k'/k$ 为有限
纯不可分域扩张。考虑环 $\mathcal{O}_{X, x} \otimes_k k'$。由《代数》引理
\ref{algebra-lemma-p-ring-map}，它的谱与 $\mathcal{O}_{X, x}$ 的谱相同，
故它也是局部环（《代数》引理
\ref{algebra-lemma-characterize-local-ring}）。因此存在唯一的点
$x' \in X_{k'}$ 位于 $x$ 上方，并且
$\mathcal{O}_{X_{k'}, x'} \cong \mathcal{O}_{X, x} \otimes_k k'$；见引理
\ref{varieties-lemma-change-fields-flat}。由假设，这是约化环。因此，由《代数》引理
\ref{algebra-lemma-geometrically-reduced-finite-purely-inseparable-extension}
推出 (2)。

\medskip\noindent
设 (2) 成立。设 $k'/k$ 为域扩张。由于
$\Spec(k') \to \Spec(k)$ 满射，故
$X_{k'} \to X$ 也满射（《态射》引理
\ref{morphisms-lemma-base-change-surjective}）。设 $x' \in X_{k'}$ 是位于
$x$ 上方的任意点。由引理 \ref{varieties-lemma-change-fields-flat}，局部环
$\mathcal{O}_{X_{k'}, x'}$ 是环
$\mathcal{O}_{X, x} \otimes_k k'$ 的局部化。因此由假设它约化，(1) 得证。
\end{proof}

\noindent
这个概念在特征零情形没有新的内容。

\begin{lemma}
\label{varieties-lemma-perfect-reduced}
设 $X$ 是完全域 $k$（例如，$k$ 的特征为零）上的概形。设 $x \in X$。
若 $\mathcal{O}_{X, x}$ 约化，则 $X$ 在 $x$ 处几何约化。
若 $X$ 约化，则 $X$ 在 $k$ 上几何约化。
\end{lemma}

\begin{proof}
第一条断言由引理 \ref{varieties-lemma-geometrically-reduced-at-point}、《代数》引理
\ref{algebra-lemma-separable-extension-preserves-reducedness} 以及完全域的定义
（《代数》定义 \ref{algebra-definition-perfect}）得出。第二条断言由第一条得出。
\end{proof}

\begin{lemma}
\label{varieties-lemma-geometrically-reduced}
设 $k$ 为特征 $p > 0$ 的域。设 $X$ 为 $k$ 上的概形。下列条件等价：
\begin{enumerate}
\item $X$ 几何约化；
\item $X_{k'}$ 对每个域扩张 $k'/k$ 都约化；
\item $X_{k'}$ 对每个有限纯不可分域扩张 $k'/k$ 都约化；
\item $X_{k^{1/p}}$ 约化；
\item $X_{k^{perf}}$ 约化；
\item $X_{\bar k}$ 约化；
\item 对每个仿射开集 $U \subset X$，环 $\mathcal{O}_X(U)$ 都几何约化
（见《代数》定义 \ref{algebra-definition-geometrically-reduced}）。
\end{enumerate}
\end{lemma}

\begin{proof}
设 (1) 成立。于是，对每个域扩张 $k'/k$ 和每个点
$x' \in X_{k'}$，$X_{k'}$ 在 $x'$ 处的局部环都约化。换言之，
$X_{k'}$ 约化。因此 (2) 成立。

\medskip\noindent
设 (2) 成立。令 $U \subset X$ 为仿射开集。于是，对每个域扩张
$k'/k$，概形 $X_{k'}$ 都约化，因而
$U_{k'} = \Spec(\mathcal{O}(U)\otimes_k k')$ 约化，继而
$\mathcal{O}(U)\otimes_k k'$ 约化（见《性质》第
\ref{properties-section-integral} 节）。换言之，$\mathcal{O}(U)$ 几何约化，
故 (7) 成立。

\medskip\noindent
设 (7) 成立。对任意域扩张 $k'/k$，基变换 $X_{k'}$ 可由环
$\mathcal{O}_X(U) \otimes_k k'$ 的谱粘合而得，其中 $U$ 遍历 $X$ 中的
仿射开集（见《概形》第 \ref{schemes-section-fibre-products} 节）。因此
$X_{k'}$ 约化，故 (1) 成立。

\medskip\noindent
这证明了 (1)、(2) 与 (7) 等价。它们又等价于 (3)、(4)、(5) 与 (6)，
因为可以对 $\mathcal{O}_X(U)$ 应用《代数》引理，其中 $U \subset X$
为仿射开集：
\ref{algebra-lemma-geometrically-reduced-finite-purely-inseparable-extension}。
\end{proof}

\begin{lemma}
\label{varieties-lemma-check-only-finite-inseparable-extensions}
设 $k$ 为特征 $p > 0$ 的域。设 $X$ 为 $k$ 上的概形。设 $x \in X$。
下列条件等价：
\begin{enumerate}
\item $X$ 在 $x$ 处几何约化；
\item 局部环 $\mathcal{O}_{X_{k'}, x'}$ 约化，其中 $k'$ 是 $k$ 的任一
有限纯不可分域扩张，且 $x' \in X_{k'}$ 是位于 $x$ 上方的唯一点；
\item 局部环 $\mathcal{O}_{X_{k^{1/p}}, x'}$ 约化，其中
$x' \in X_{k^{1/p}}$ 是位于 $x$ 上方的唯一点；
\item 局部环 $\mathcal{O}_{X_{k^{perf}}, x'}$ 约化，其中
$x' \in X_{k^{perf}}$ 是位于 $x$ 上方的唯一点。
\end{enumerate}
\end{lemma}

\begin{proof}
注意，若 $k'/k$ 纯不可分，则 $X_{k'} \to X$ 在底拓扑空间上诱导同胚；
见《代数》引理 \ref{algebra-lemma-p-ring-map}。这说明了命题中 $x'$ 位于
$x$ 上方时的唯一性。此外，在这种情况下，
$\mathcal{O}_{X_{k'}, x'} = \mathcal{O}_{X, x} \otimes_k k'$。
所以本引理由上面的引理 \ref{varieties-lemma-geometrically-reduced-at-point} 以及
《代数》引理
\ref{algebra-lemma-geometrically-reduced-finite-purely-inseparable-extension}
得出。
\end{proof}

\begin{lemma}
\label{varieties-lemma-geometrically-reduced-upstairs}
设 $k$ 为域。设 $X$ 为 $k$ 上的概形。设 $k'/k$ 为域扩张。
设 $x \in X$ 为一点，并设 $x' \in X_{k'}$ 为位于 $x$ 上方的一点。
下列条件等价：
\begin{enumerate}
\item $X$ 在 $x$ 处几何约化；
\item $X_{k'}$ 在 $x'$ 处几何约化。
\end{enumerate}
特别地，$X$ 在 $k$ 上几何约化，当且仅当 $X_{k'}$ 在 $k'$ 上几何约化。
\end{lemma}

\begin{proof}
显然 (1) 蕴含 (2)。设 (2) 成立。令 $k''/k$ 为有限纯不可分域扩张，
并令 $x'' \in X_{k''}$ 为位于 $x$ 上方的一点（事实上，它是唯一的）。
选取 $k'''/k$，使它是 $k''/k$ 与 $k'/k$ 的公共域扩张；见《域》引理
\ref{fields-lemma-common-extension-field}。选取一点 $x''' \in X_{k'''}$，
它同时位于 $x'$ 与 $x''$ 上方。考虑局部环映射
$$
\mathcal{O}_{X_{k''}, x''} \longrightarrow \mathcal{O}_{X_{k'''}, x'''}.
$$
这是平坦局部环同态，因而忠实平坦。由 (2)，右侧局部环约化。
于是由《代数》引理 \ref{algebra-lemma-descent-reduced} 可知
$\mathcal{O}_{X_{k''}, x''}$ 约化。因此，由引理
\ref{varieties-lemma-check-only-finite-inseparable-extensions} 可知 $X$ 在 $x$ 处
几何约化。
\end{proof}

\begin{lemma}
\label{varieties-lemma-geometrically-reduced-any-base-change}
设 $k$ 为域。设 $X$、$Y$ 为 $k$ 上的概形。
\begin{enumerate}
\item 若 $X$ 在 $x$ 处几何约化，且 $Y$ 约化，则 $X \times_k Y$ 在每个
位于 $x$ 上方的点处都约化。
\item 若 $X$ 在 $k$ 上几何约化，且 $Y$ 约化，则 $X \times_k Y$ 约化。
\item 若 $X$ 与 $Y$ 在 $k$ 上几何约化，则 $X \times_k Y$ 几何约化。
\item 若 $k$ 为完全域，且 $X$ 与 $Y$ 约化，则 $X \times_k Y$ 约化。
\item 此处补充更多内容。
\end{enumerate}
\end{lemma}

\begin{proof}
为证明 (1)，结合引理 \ref{varieties-lemma-geometrically-reduced-at-point} 与《代数》引理
\ref{algebra-lemma-geometrically-reduced-any-reduced-base-change}。
为证明 (2)，结合引理 \ref{varieties-lemma-geometrically-reduced} 与《代数》引理
\ref{algebra-lemma-geometrically-reduced-any-reduced-base-change}。
为证明 (3)，注意
$(X \times_k Y)_{\overline{k}} =
X_{\overline{k}} \times_{\overline{k}} Y_{\overline{k}}$，
并使用 (2) 及引理 \ref{varieties-lemma-geometrically-reduced}。
为证明 (4)，结合使用 (3) 与引理 \ref{varieties-lemma-perfect-reduced}。
\end{proof}

\begin{lemma}
\label{varieties-lemma-generic-points-geometrically-reduced}
设 $k$ 为域。设 $X$ 为 $k$ 上的概形。
\begin{enumerate}
\item 若 $x' \leadsto x$ 是一次特化，且 $X$ 在 $x$ 处几何约化，则 $X$
在 $x'$ 处几何约化。
\item 设 $x \in X$ 满足：(a) $\mathcal{O}_{X, x}$ 约化；并且 (b) 对每个
特化 $x' \leadsto x$，若 $x'$ 是 $X$ 的一个不可约分支的泛点，则 $X$
在 $x'$ 处几何约化。那么 $X$ 在 $x$ 处几何约化。
\item 若 $X$ 约化，且在 $X$ 的不可约分支的所有泛点处都几何约化，
则 $X$ 几何约化。
\item 若 $X$ 是 $k$ 上的簇，则 $X$ 几何约化，当且仅当其函数域是 $k$
的可分扩张。
\end{enumerate}
\end{lemma}

\begin{proof}
第 (1) 部分由引理 \ref{varieties-lemma-geometrically-reduced-at-point} 及以下事实得出：
若 $A$ 为几何约化 $k$-代数，则 $S^{-1}A$ 是几何约化 $k$-代数，
其中 $S$ 是 $A$ 的任意乘法子集；见《代数》引理
\ref{algebra-lemma-geometrically-reduced-permanence}。

\medskip\noindent
令 $A = \mathcal{O}_{X, x}$。第 (2) 部分的假设 (a) 与 (b) 蕴含 $A$
约化，并且 $A_{\mathfrak q}$ 在 $k$ 上几何约化，其中 $\mathfrak q$
遍历 $A$ 的极小素理想。因此 $A$ 在 $k$ 上几何约化；见《代数》引理
\ref{algebra-lemma-generic-points-geometrically-reduced}。由此，$X$ 在 $x$
处几何约化；见引理 \ref{varieties-lemma-geometrically-reduced-at-point}。

\medskip\noindent
第 (3) 部分显然由第 (2) 部分得出。最后，第 (4) 部分由第 (3) 部分及
《代数》引理 \ref{algebra-lemma-characterize-separable-field-extensions}
得出。
\end{proof}

\begin{lemma}
\label{varieties-lemma-Noetherian-geometrically-reduced-at-point}
设 $k$ 为域。设 $X$ 为 $k$ 上的概形。设 $x \in X$。假设 $X$ 局部
Noether，且在 $x$ 处几何约化。那么存在开邻域 $U \subset X$，它包含
$x$，并且在 $k$ 上几何约化。
\end{lemma}

\begin{proof}
假设 $X$ 局部 Noether，且在 $x$ 处几何约化。由《性质》引理
\ref{properties-lemma-ring-affine-open-injective-local-ring}，可以找到
仿射开邻域 $U \subset X$，它包含 $x$，并且使得
$R = \mathcal{O}_X(U) \to \mathcal{O}_{X, x}$ 为单射。由引理
\ref{varieties-lemma-geometrically-reduced-at-point}，假设意味着
$\mathcal{O}_{X, x}$ 在 $k$ 上几何约化。由《代数》引理
\ref{algebra-lemma-subalgebra-separable}，这蕴含 $R$ 在 $k$ 上几何约化，
进而蕴含 $U$ 几何约化。
\end{proof}

\begin{example}
\label{varieties-example-not-geometrically-reduced}
设 $k = \mathbf{F}_p(s, t)$，即素域的纯超越扩张。考虑簇
$X = \Spec(k[x, y]/(1 + sx^p + ty^p))$。设 $k'/k$ 为任一扩张，
使 $s$ 与 $t$ 都有 $p$ 次根，且这些根属于 $k'$。那么基变换 $X_{k'}$
非约化。
事实上，环 $k'[x, y]/(1 + s x^p + ty^p)$ 包含元素
$1 + s^{1/p}x + t^{1/p}y$；它的 $p$ 次幂为零，但它本身不为零
（因为理想 $(1 + sx^p + ty^p)$ 显然不含任何次数 $< p$ 的非零元素）。
\end{example}

\begin{lemma}
\label{varieties-lemma-finite-extension-geometrically-reduced}
设 $k$ 为域。设 $X \to \Spec(k)$ 局部有限型。假设 $X$ 只有有限多个
不可约分支。那么存在有限纯不可分扩张 $k'/k$，使得 $(X_{k'})_{red}$
在 $k'$ 上几何约化。
\end{lemma}

\begin{proof}
为证明本引理，可以将 $X$ 替换为其约化 $X_{red}$。因此可以假设 $X$
约化且在 $k$ 上局部有限型。令 $x_1, \ldots, x_n \in X$ 为 $X$ 的
各不可约分支的泛点。注意，对每个纯不可分代数扩张 $k'/k$，态射
$(X_{k'})_{red} \to X$ 都是同胚；见《代数》引理
\ref{algebra-lemma-p-ring-map}。因此，点 $x'_1, \ldots, x'_n$（分别位于
$x_1, \ldots, x_n$ 上方）是 $(X_{k'})_{red}$ 的各不可约分支的泛点。
由于 $X$ 约化，局部环 $K_i = \mathcal{O}_{X, x_i}$ 都是域；见《代数》引理
\ref{algebra-lemma-minimal-prime-reduced-ring}。由于 $X$ 在 $k$ 上局部有限型，
域扩张 $K_i/k$ 都有限生成。最后，局部环
$\mathcal{O}_{(X_{k'})_{red}, x'_i}$ 是域 $(K_i \otimes_k k')_{red}$。
由《代数》引理 \ref{algebra-lemma-make-separable}，可以找到有限纯不可分扩张
$k'/k$，使得 $(K_i \otimes_k k')_{red}$ 都是 $k'$ 的可分域扩张。
特别地，由《代数》引理
\ref{algebra-lemma-characterize-separable-field-extensions}，每个
$(K_i \otimes_k k')_{red}$ 都在 $k'$ 上几何约化。此时，引理
\ref{varieties-lemma-generic-points-geometrically-reduced} 的第 (3) 部分蕴含
$(X_{k'})_{red}$ 几何约化。
\end{proof}

\section{几何连通概形}
\label{varieties-section-geometrically-connected}

\noindent
如果 $X$ 是域上的连通概形，那么扩张基域之后，$X$ 可能变成不连通的。
对于几何连通概形，这种情况不会发生。

\begin{definition}
\label{varieties-definition-geometrically-connected}
设 $X$ 为域 $k$ 上的概形。称 $X$ 在 $k$ 上{\it 几何连通}，如果概形
$X_{k'}$ 对每个域 $k'$（作为 $k$ 的扩张）都连通。
\end{definition}

\noindent
按照约定，连通拓扑空间非空；因此更不必说，几何连通概形非空。
下面给出一个并非几何连通的簇的例子。

\begin{example}
\label{varieties-example-not-geometrically-irreducible}
设 $k = \mathbf{Q}$。概形
$X = \Spec(\mathbf{Q}(i))$ 是 $\Spec(\mathbf{Q})$ 上的簇。
但是基变换 $X_{\mathbf{C}}$ 是
$\mathbf{C} \otimes_{\mathbf{Q}} \mathbf{Q}(i) \cong
\mathbf{C} \times \mathbf{C}$ 的谱，也就是两个 $\Spec(\mathbf{C})$
的不交并。因此，这确实是一个非几何连通簇的例子。
\end{example}

\begin{lemma}
\label{varieties-lemma-geometrically-connected-check-after-extension}
设 $X$ 为域 $k$ 上的概形。设 $k'/k$ 为域扩张。那么 $X$ 在 $k$ 上
几何连通，当且仅当 $X_{k'}$ 在 $k'$ 上几何连通。
\end{lemma}

\begin{proof}
如果 $X$ 在 $k$ 上几何连通，那么显然 $X_{k'}$ 在 $k'$ 上几何连通。
反之，对任意域扩张 $k''/k$，存在公共域扩张 $k'''/k'$ 与 $k'''/k''$；
见《域》引理 \ref{fields-lemma-common-extension-field}。由于态射
$X_{k'''} \to X_{k''}$ 满射（它是域的谱之间满射的基变换），
$X_{k'''}$ 的连通性蕴含 $X_{k''}$ 的连通性。因此，如果 $X_{k'}$ 在
$k'$ 上几何连通，那么 $X$ 在 $k$ 上几何连通。
\end{proof}

\begin{lemma}
\label{varieties-lemma-bijection-connected-components}
设 $k$ 为域。设 $X$、$Y$ 为 $k$ 上的概形。假设 $X$ 在 $k$ 上
几何连通。那么投影态射
$$
p : X \times_k Y \longrightarrow Y
$$
诱导连通分支之间的双射。
\end{lemma}

\begin{proof}
$p$ 的概形论纤维连通，因为它们是几何连通概形 $X$ 沿域扩张的
基变换。此外，概形论纤维与集合论纤维同胚；见《概形》引理
\ref{schemes-lemma-fibre-topological}。由《态射》引理
\ref{morphisms-lemma-scheme-over-field-universally-open}，映射 $p$ 是开的。
因此可以应用《拓扑学》引理
\ref{topology-lemma-connected-fibres-connected-components} 得出结论。
\end{proof}

\begin{lemma}
\label{varieties-lemma-affine-geometrically-connected}
设 $k$ 为域。设 $A$ 为 $k$-代数。那么 $X = \Spec(A)$ 在 $k$ 上
几何连通，当且仅当 $A$ 在 $k$ 上几何连通（见《代数》定义
\ref{algebra-definition-geometrically-connected}）。
\end{lemma}

\begin{proof}
这直接来自定义。
\end{proof}

\begin{lemma}
\label{varieties-lemma-separably-closed-field-connected-components}
设 $k'/k$ 为域扩张。设 $X$ 为 $k$ 上的概形。假设 $k$ 可分代数闭。
那么态射 $X_{k'} \to X$ 诱导连通分支之间的双射。特别地，$X$ 在 $k$
上几何连通，当且仅当 $X$ 连通。
\end{lemma}

\begin{proof}
由于 $k$ 可分代数闭，可知 $k'$ 在 $k$ 上几何连通；见《代数》引理
\ref{algebra-lemma-separably-closed-connected-implies-geometric}。
因此，由上面的引理 \ref{varieties-lemma-affine-geometrically-connected}，
$Z = \Spec(k')$ 在 $k$ 上几何连通。由于
$X_{k'} = Z \times_k X$，结论是引理
\ref{varieties-lemma-bijection-connected-components} 的特殊情形。
\end{proof}

\begin{lemma}
\label{varieties-lemma-characterize-geometrically-connected}
设 $k$ 为域。设 $X$ 为 $k$ 上的概形。设 $\overline{k}$ 为 $k$ 的
可分代数闭包。那么 $X$ 几何连通，当且仅当基变换
$X_{\overline{k}}$ 连通。
\end{lemma}

\begin{proof}
假设 $X_{\overline{k}}$ 连通。设 $k'/k$ 为域扩张。存在域扩张
$\overline{k}'/\overline{k}$，使得 $k'$ 嵌入 $\overline{k}'$，并且该嵌入
是 $k$-扩张之间的嵌入。由引理
\ref{varieties-lemma-separably-closed-field-connected-components} 可知
$X_{\overline{k}'}$ 连通。由于 $X_{\overline{k}'} \to X_{k'}$ 满射，
可知 $X_{k'}$ 连通，如所需。
\end{proof}

\begin{lemma}
\label{varieties-lemma-descend-open}
设 $k$ 为域。设 $X$ 为 $k$ 上的概形。设 $A$ 为 $k$-代数。
设 $V \subset X_A$ 为拟紧开集。那么存在有限生成 $k$-子代数
$A' \subset A$ 以及拟紧开集 $V' \subset X_{A'}$，使得 $V = V'_A$。
\end{lemma}

\begin{proof}
顺带指出，若 $X$ 还拟分离，则该结论由《极限》引理
\ref{limits-lemma-descend-opens} 得出。令 $U_1, \ldots, U_n$ 为 $X$ 的
有限多个仿射开集，使得 $V \subset \bigcup U_{i, A}$。记
$U_i = \Spec(R_i)$。由于 $V$ 拟紧，可以找到有限多个
$f_{ij} \in R_i \otimes_k A$，$j = 1, \ldots, n_i$，使得
$V = \bigcup_i \bigcup_{j = 1, \ldots, n_i} D(f_{ij})$，其中
$D(f_{ij}) \subset U_{i, A}$ 是相应的标准开集。（我们并未断言
$V \cap U_{i, A}$ 是 $D(f_{ij})$（$j = 1, \ldots, n_i$）的并。）
显然可以找到有限生成 $k$-子代数 $A' \subset A$，使得 $f_{ij}$ 是某个
$f_{ij}' \in R_i \otimes_k A'$ 的像。置
$V' = \bigcup D(f_{ij}')$；这是 $X_{A'}$ 的拟紧开集。用
$\pi : X_A \to X_{A'}$ 表示典范态射。有 $\pi(V) \subset V'$，因为
$\pi(D(f_{ij})) \subset D(f_{ij}')$。若 $x \in X_A$ 且
$\pi(x) \in V'$，则 $\pi(x) \in D(f_{ij}')$ 对某个 $i, j$ 成立；于是
$x \in D(f_{ij})$，因为 $f_{ij}'$ 映到 $f_{ij}$。因此
$V = \pi^{-1}(V')$，如所需。
\end{proof}

\noindent
设 $k$ 为域。设 $\overline{k}/k$ 为（可能无限的）Galois 扩张。例如，
$\overline{k}$ 可以是 $k$ 的可分代数闭包。对任意
$\sigma \in \text{Gal}(\overline{k}/k)$，得到相应的自同构
$
\Spec(\sigma) :
\Spec(\overline{k})
\longrightarrow
\Spec(\overline{k})
$。
注意
$\Spec(\sigma) \circ \Spec(\tau) = \Spec(\tau \circ \sigma)$。
因此得到作用
\begingroup\small
$$
\text{Gal}(\overline{k}/k)^{opp} \times \Spec(\overline{k})
\longrightarrow
\Spec(\overline{k})
$$
\endgroup
即反群在概形 $\Spec(\overline{k})$ 上的作用。\par
设 $X$ 为 $k$ 上的概形。
按照定义，
$X_{\overline{k}} =
\Spec(\overline{k}) \times_{\Spec(k)} X$，
所以以上作用诱导典范作用
\begin{equation}
\label{varieties-equation-galois-action-base-change-kbar}
\text{Gal}(\overline{k}/k)^{opp} \times X_{\overline{k}}
\longrightarrow
X_{\overline{k}}.
\end{equation}

\begin{lemma}
\label{varieties-lemma-Galois-action-quasi-compact-open}
设 $k$ 为域。设 $X$ 为 $k$ 上的概形。设 $\overline{k}$ 为 $k$ 的
（可能无限的）Galois 扩张。设 $V \subset X_{\overline{k}}$ 为拟紧开集。
那么：
\begin{enumerate}
\item 存在有限子扩张 $\overline{k}/k'/k$ 以及拟紧开集
$V' \subset X_{k'}$，使得 $V = (V')_{\overline{k}}$；
\item 存在开子群 $H \subset \text{Gal}(\overline{k}/k)$，使得
$\sigma(V) = V$ 对所有 $\sigma \in H$ 都成立。
\end{enumerate}
\end{lemma}

\begin{proof}
由引理 \ref{varieties-lemma-descend-open}，存在有限子扩张
$k'/k \subset \overline{k}$ 以及开集 $V' \subset X_{k'}$，其拉回为 $V$。
这证明了 (1)。由于 $\text{Gal}(\overline{k}/k')$ 是
$\text{Gal}(\overline{k}/k)$ 的开子群，第 (2) 部分也很清楚。
\end{proof}

\begin{lemma}
\label{varieties-lemma-closed-fixed-by-Galois}
设 $k$ 为域。设 $\overline{k}/k$ 为（可能无限的）Galois 扩张。
设 $X$ 为 $k$ 上的概形。设 $\overline{T} \subset X_{\overline{k}}$ 满足：
\begin{enumerate}
\item $\overline{T}$ 是 $X_{\overline{k}}$ 的闭子集；
\item 对每个 $\sigma \in \text{Gal}(\overline{k}/k)$，都有
$\sigma(\overline{T}) = \overline{T}$。
\end{enumerate}
那么存在闭子集 $T \subset X$，它在 $X_{\overline{k}}$ 中的逆像是
$\overline{T}$。
\end{lemma}

\begin{proof}
本引理立即归约到 $X = \Spec(A)$ 为仿射的情形。在这种情况下，令
$\overline{I} \subset A \otimes_k \overline{k}$ 为对应于 $\overline{T}$ 的
根理想。假设 (2) 蕴含 $\sigma(\overline{I}) = \overline{I}$ 对所有
$\sigma \in \text{Gal}(\overline{k}/k)$ 都成立。取 $x \in \overline{I}$。
存在有限 Galois 扩张 $k'/k$，它包含于 $\overline{k}$ 中，并且使得
$x \in A \otimes_k k'$。置 $G = \text{Gal}(k'/k)$。置
$$
P(T) = \prod\nolimits_{\sigma \in G} (T - \sigma(x)) \in (A \otimes_k k')[T]
$$
显然，$P(T)$ 首一，并且它实际上是
$(A \otimes_k k')^G[T] = A[T]$ 的元素（由基本 Galois 理论）。此外，
若写成 $P(T) = T^d + a_1T^{d - 1} + \ldots + a_d$，则可见
$a_i \in I := A \cap \overline{I}$。结合 $P(x) = 0$ 与 $a_i \in I$，得到
$x^d = - a_1 x^{d - 1} - \ldots  - a_d \in I(A \otimes_k \overline{k})$。
因此 $x$ 属于 $I(A \otimes_k \overline{k})$ 的根。故
$\overline{I}$ 是 $I(A \otimes_k \overline{k})$ 的根，而置 $T = V(I)$
便得到一个解。
\end{proof}

\begin{lemma}
\label{varieties-lemma-characterize-geometrically-disconnected}
设 $k$ 为域。设 $X$ 为 $k$ 上的概形。下列条件等价：
\begin{enumerate}
\item $X$ 几何连通；
\item 对每个有限可分域扩张 $k'/k$，概形 $X_{k'}$ 都连通。
\end{enumerate}
\end{lemma}

\begin{proof}
由定义立即可知 (1) 蕴含 (2)。假设 $X$ 不几何连通。设
$k \subset \overline{k}$ 为 $k$ 的可分代数闭包。由引理
\ref{varieties-lemma-characterize-geometrically-connected}，可知
$X_{\overline{k}}$ 不连通。设
$X_{\overline{k}} = \overline{U} \amalg \overline{V}$，其中
$\overline{U}$ 与 $\overline{V}$ 均开、闭且非空。

\medskip\noindent
设 $W \subset X$ 是任意拟紧开集。那么
$W_{\overline{k}} \cap \overline{U}$ 与
$W_{\overline{k}} \cap \overline{V}$ 在 $W_{\overline{k}}$ 中都既开又闭。
特别地，$W_{\overline{k}} \cap \overline{U}$ 与
$W_{\overline{k}} \cap \overline{V}$ 都拟紧，并且由引理
\ref{varieties-lemma-Galois-action-quasi-compact-open}，
$W_{\overline{k}} \cap \overline{U}$ 与
$W_{\overline{k}} \cap \overline{V}$ 都可在一个有限子扩张上定义，
并且在 $\text{Gal}(\overline{k}/k)$ 的一个开子群下不变。下文将不再
另行说明而直接使用这一点。

\medskip\noindent
选取拟紧开集 $W_0 \subset X$，使
$W_{0, \overline{k}} \cap \overline{U}$ 与
$W_{0, \overline{k}} \cap \overline{V}$ 都非空。选取有限子扩张
$\overline{k}/k'/k$ 以及分解 $W_{0, k'} = U_0' \amalg V_0'$，其两项
都是既开又闭的子集，并且使得
$W_{0, \overline{k}} \cap \overline{U} = (U'_0)_{\overline{k}}$ 且
$W_{0, \overline{k}} \cap \overline{V} = (V'_0)_{\overline{k}}$。
令 $H = \text{Gal}(\overline{k}/k') \subset \text{Gal}(\overline{k}/k)$。
特别地，
$\sigma(W_{0, \overline{k}} \cap \overline{U}) =
W_{0, \overline{k}} \cap \overline{U}$，对 $\overline{V}$ 也类似。

\medskip\noindent
如上选定 $W_0$、$k'$ 后，对每个拟紧开集 $W \subset X$，置
$$
U_W =
\bigcap\nolimits_{\sigma \in H} \sigma(W_{\overline{k}} \cap \overline{U}),
\quad
V_W =
\bigcup\nolimits_{\sigma \in H} \sigma(W_{\overline{k}} \cap \overline{V}).
$$
由于 $W_{\overline{k}} \cap \overline{U}$ 与
$W_{\overline{k}} \cap \overline{V}$ 在
$\text{Gal}(\overline{k}/k)$ 的一个开子群下固定，上面的并与交都是
有限的。因此 $U_W$ 与 $V_W$ 都既开又闭。此外，按构造有
$W_{\bar k} = U_W \amalg V_W$。

\medskip\noindent
我们断言，若 $W \subset W' \subset X$ 为拟紧开集，则
$W_{\overline{k}} \cap U_{W'} = U_W$ 且
$W_{\overline{k}} \cap V_{W'} = V_W$。验证从略。因此，在定义
$U = \bigcup_{W \subset X} U_W$ 与
$V = \bigcup_{W \subset X} V_W$ 后，得到
$X_{\overline{k}} = U \amalg V$，这是既开又闭子集的不交并。
显然 $V$ 非空，因为它是（局部地）取并构造出的。另一方面，$U$ 非空，
因为按构造它包含 $W_0 \cap \overline{U}$。最后，
$U, V \subset X_{\bar k}$ 按构造闭且 $H$-不变。因此，由引理
\ref{varieties-lemma-closed-fixed-by-Galois}，有 $U = (U')_{\bar k}$ 且
$V = (V')_{\bar k}$，其中 $U', V' \subset X_{k'}$ 为闭子集。显然
$X_{k'} = U' \amalg V'$，由此可见
$X_{k'}$ 不连通，如所需。
\end{proof}

\begin{lemma}
\label{varieties-lemma-tricky}
设 $k$ 为域。设 $\overline{k}/k$ 为（可能无限的）Galois 扩张。
设 $f : T \to X$ 为 $k$ 上概形的态射。假设 $T_{\overline{k}}$ 连通且
$X_{\overline{k}}$ 不连通。那么 $X$ 不连通。
\end{lemma}

\begin{proof}
写成 $X_{\overline{k}} = \overline{U} \amalg \overline{V}$，其中
$\overline{U}$ 与 $\overline{V}$ 既开又闭。记
$\overline{f} : T_{\overline{k}} \to X_{\overline{k}}$ 为 $f$ 的基变换。
由于 $T_{\overline{k}}$ 连通，可知 $T_{\overline{k}}$ 包含于
$\overline{f}^{-1}(\overline{U})$ 或 $\overline{f}^{-1}(\overline{V})$
二者之一。不妨设
$T_{\overline{k}} \subset \overline{f}^{-1}(\overline{U})$。

\medskip\noindent
固定拟紧开集 $W \subset X$。存在有限 Galois 子扩张
$\overline{k}/k'/k$，使得 $\overline{U} \cap W_{\overline{k}}$ 与
$\overline{V} \cap W_{\overline{k}}$ 来自拟紧开集
$U', V' \subset W_{k'}$。于是也有 $W_{k'} = U' \amalg V'$。考虑
$$
U'' = \bigcap\nolimits_{\sigma \in \text{Gal}(k'/k)} \sigma(U'),
\quad
V'' = \bigcup\nolimits_{\sigma \in \text{Gal}(k'/k)} \sigma(V').
$$
它们 Galois 不变、既开又闭，并且 $W_{k'} = U'' \amalg V''$。
由引理 \ref{varieties-lemma-closed-fixed-by-Galois}，得到既开又闭的子集
$U_W, V_W \subset W$，使得
$U'' = (U_W)_{k'}$、$V'' = (V_W)_{k'}$ 且
$W = U_W \amalg V_W$。

\medskip\noindent
我们断言，若 $W \subset W' \subset X$ 为拟紧开集，则
$W \cap U_{W'} = U_W$ 且 $W \cap V_{W'} = V_W$。验证从略。因此，在定义
$U = \bigcup_{W \subset X} U_W$ 与
$V = \bigcup_{W \subset X} V_W$ 后，得到 $X = U \amalg V$。
显然 $V$ 非空，因为它是（局部地）取并构造出的。另一方面，$U$ 非空，
因为按构造它包含 $f(T)$。
\end{proof}

\begin{lemma}
\label{varieties-lemma-geometrically-connected-criterion}
\begin{reference}
\cite[IV Corollary 4.5.13.1(i)]{EGA}
\end{reference}
设 $k$ 为域。设 $T \to X$ 为 $k$ 上概形的态射。假设 $T$ 几何连通且
$X$ 连通。那么 $X$ 几何连通。
\end{lemma}

\begin{proof}
这是引理 \ref{varieties-lemma-tricky} 的改述。
\end{proof}

\begin{lemma}
\label{varieties-lemma-geometrically-connected-if-connected-and-point}
设 $k$ 为域。设 $X$ 为 $k$ 上的概形。假设 $X$ 连通，并且有一点 $x$，
使得 $k$ 在 $\kappa(x)$ 中代数闭。那么 $X$ 几何连通。特别地，如果
$X$ 有 $k$-有理点且 $X$ 连通，那么 $X$ 几何连通。
\end{lemma}

\begin{proof}
置 $T = \Spec(\kappa(x))$。设 $\overline{k}$ 为 $k$ 的可分代数闭包。
关于 $\kappa(x)/k$ 的假设蕴含 $T_{\overline{k}}$ 不可约；见《代数》引理
\ref{algebra-lemma-field-extension-geometrically-irreducible}。因此，由引理
\ref{varieties-lemma-geometrically-connected-criterion} 可知 $X_{\overline{k}}$ 连通。
由引理 \ref{varieties-lemma-characterize-geometrically-connected}，可知 $X$ 几何连通。
\end{proof}

\begin{lemma}
\label{varieties-lemma-inverse-image-connected-component}
设 $K/k$ 为域扩张。设 $X$ 为 $k$ 上的概形。对每个连通分支 $T$
（属于 $X$），逆像 $T_K \subset X_K$ 是 $X_K$ 的若干连通分支之并。
\end{lemma}

\begin{proof}
这是纯拓扑的陈述。记 $p : X_K \to X$ 为投影态射。设
$T \subset X$ 为 $X$ 的一个连通分支。设
$t \in T_K = p^{-1}(T)$。设 $C \subset X_K$ 为包含 $t$ 的连通分支。
那么 $p(C)$ 是 $X$ 的一个与 $T$ 相交的连通子集，故
$p(C) \subset T$。因此 $C \subset T_K$。
\end{proof}

\noindent
下面的引理将被后面更强的引理 \ref{varieties-lemma-image-connected-component}
所取代。

\begin{lemma}
\label{varieties-lemma-image-connected-component-finite-extension}
设 $K/k$ 为有限域扩张，并设 $X$ 为 $k$ 上的概形。用
$p : X_K \to X$ 表示投影态射。对每个连通分支 $T$（属于 $X_K$），像
$p(T)$ 是 $X$ 的一个连通分支。
\end{lemma}

\begin{proof}
像 $p(T)$ 包含于某个连通分支 $X'$ 中，而该分支属于 $X$。以任意方式
把 $X'$ 视为 $X$ 的闭子概形。那么 $T$ 也是
$X'_K = p^{-1}(X')$ 的连通分支，因此可以假设 $X$ 连通。态射 $p$ 是
开映射（《态射》引理
\ref{morphisms-lemma-scheme-over-field-universally-open}）、闭映射
（《态射》引理 \ref{morphisms-lemma-integral-universally-closed}），并且
$p$ 的纤维是有限集（《态射》引理
\ref{morphisms-lemma-finite-quasi-finite}）。所以可以应用《拓扑学》引理
\ref{topology-lemma-finite-fibre-connected-components} 得出结论。
\end{proof}

\begin{lemma}[Gabber]
\label{varieties-lemma-image-connected-component}
\begin{reference}
Ofer Gabber 于 2016 年 6 月 4 日发送的电子邮件
\end{reference}
设 $K/k$ 为域扩张。设 $X$ 为 $k$ 上的概形。记
$p : X_K \to X$ 为投影态射。设 $\overline{T} \subset X_K$ 为连通分支。
那么 $p(\overline{T})$ 是 $X$ 的连通分支。
\end{lemma}

\begin{proof}
当 $K/k$ 有限时，这就是引理
\ref{varieties-lemma-image-connected-component-finite-extension}。一般情况下，
证明较为困难。

\medskip\noindent
令 $T \subset X$ 为 $X$ 的包含 $\overline{T}$ 之像的连通分支。可以将
$X$ 替换为 $T$（赋予诱导的约化子概形结构）。因此可以假设 $X$ 连通。
令 $A = H^0(X, \mathcal{O}_X)$。设 $L \subset A$ 为极大的弱
\'etale $k$-子代数；见《更多代数》引理
\ref{more-algebra-lemma-max-weakly-etale-subalgebra}。由于 $A$ 没有任何
非平凡幂等元，由《更多代数》引理
\ref{more-algebra-lemma-class-weakly-etale-over-field} 可知，$L$ 是域，
并且是 $k$ 的可分代数扩张。注意，$L$ 也是极大弱 \'etale
$L$-子代数；这里所指的母代数为 $A$。这是因为《更多代数》引理
\ref{more-algebra-lemma-composition-weakly-etale} 表明，任意弱 \'etale
$L$-代数在 $k$ 上都是弱 \'etale 的。由《概形》引理
\ref{schemes-lemma-morphism-into-affine}，得到结构态射的分解
$X \to \Spec(L) \to \Spec(k)$。

\medskip\noindent
设 $L'/L$ 为有限可分扩张。由《概形的上同调》引理
\ref{coherent-lemma-finite-locally-free-base-change-cohomology}，有
$$
A \otimes_L L' =
H^0(X \times_{\Spec(L)} \Spec(L'), \mathcal{O}_{X \times_{\Spec(L)} \Spec(L')})
$$
极大弱 \'etale $L'$-子代数（在 $A \otimes_L L'$ 中）是
$L \otimes_L L' = L'$；见《更多代数》引理
\ref{more-algebra-lemma-change-fields-max-weakly-etale-subalgebra}。特别地，
$A \otimes_L L'$ 没有非平凡幂等元（这样的幂等元会生成一个弱
\'etale 子代数），所以 $X \times_{\Spec(L)} \Spec(L')$ 连通。由引理
\ref{varieties-lemma-characterize-geometrically-disconnected}，可知 $X$ 在 $L$ 上
几何连通。

\medskip\noindent
赋予 $\overline{T}$ 诱导的约化概形结构，并考虑复合
$$
\overline{T} \xrightarrow{i} X_K = X \times_{\Spec(k)} \Spec(K)
\xrightarrow{\pi}
\Spec(L \otimes_k K)
$$
其像包含于 $\Spec(L \otimes_k K)$ 的一个连通分支中。由于
$K \to L \otimes_k K$ 整，可知 $\Spec(L \otimes_k K)$ 的各连通分支
都是点，并且所有点都闭；见《代数》引理
\ref{algebra-lemma-integral-over-field}。因此得到商域
$L \otimes_k K \to E$，使 $\overline{T}$ 映入
$\Spec(E) \subset \Spec(L \otimes_k K)$。故
$i(\overline{T}) \subset \pi^{-1}(\Spec(E))$。但是
$$
\pi^{-1}(\Spec(E)) =
(X \times_{\Spec(k)} \Spec(K)) \times_{\Spec(L \otimes_k K)} \Spec(E) =
X \times_{\Spec(L)} \Spec(E)
$$
它连通，因为 $X$ 在 $L$ 上几何连通。于是得到等式
$\overline{T} = X \times_{\Spec(L)} \Spec(E)$（在集合论意义下），
并由此得出 $\overline{T} \to X$ 满射，如所需。
\end{proof}

\noindent
设 $X$ 为概形。用 $\pi_0(X)$ 表示 $X$ 的连通分支集合。

\begin{lemma}
\label{varieties-lemma-galois-action-connected-components}
设 $k$ 为域，并设 $\overline{k}$ 为其可分代数闭包。设 $X$ 为 $k$ 上的
概形。存在作用
$$
\text{Gal}(\overline{k}/k)^{opp} \times \pi_0(X_{\overline{k}})
\longrightarrow
\pi_0(X_{\overline{k}})
$$
并具有下列性质：
\begin{enumerate}
\item 元素 $\overline{T} \in \pi_0(X_{\overline{k}})$ 被该作用固定，
当且仅当存在连通分支 $T \subset X$，它在 $k$ 上几何连通，并且满足
$T_{\overline{k}} = \overline{T}$。
\item 对任意域扩张 $k'/k$ 及其可分代数闭包 $\overline{k}'$，图表
$$
\xymatrix{
\text{Gal}(\overline{k}'/k') \times \pi_0(X_{\overline{k}'})
\ar[r] \ar[d] &
\pi_0(X_{\overline{k}'}) \ar[d] \\
\text{Gal}(\overline{k}/k) \times \pi_0(X_{\overline{k}})
\ar[r] &
\pi_0(X_{\overline{k}})
}
$$
交换（由引理 \ref{varieties-lemma-separably-closed-field-connected-components}，
右侧竖直箭头为双射）。
\end{enumerate}
\end{lemma}

\begin{proof}
作用（\ref{varieties-equation-galois-action-base-change-kbar}）即
$\text{Gal}(\overline{k}/k)$ 在 $X_{\overline{k}}$ 上的作用，它诱导
在连通分支上的作用。连通分支总是闭的（《拓扑学》引理
\ref{topology-lemma-connected-components}）。因此，若 $\overline{T}$ 如
(1) 所述，则由引理 \ref{varieties-lemma-closed-fixed-by-Galois}，存在闭子集
$T \subset X$，使得 $\overline{T} = T_{\overline{k}}$。注意，$T$ 在
$k$ 上几何连通；见引理 \ref{varieties-lemma-characterize-geometrically-connected}。
为说明 $T$ 是 $X$ 的一个连通分支，假设
$T \subset T'$、$T \not = T'$，其中 $T'$ 是 $X$ 的一个连通分支。
在这种情况下，$T'_{k'}$ 严格包含 $\overline{T}$，因而不连通。由引理
\ref{varieties-lemma-tricky}，这意味着 $T'$ 不连通！矛盾。

\medskip\noindent
我们略去 (2) 中函子性的证明。
\end{proof}

\begin{lemma}
\label{varieties-lemma-galois-action-connected-components-continuous}
设 $k$ 为域，并设 $\overline{k}$ 为其可分代数闭包。设 $X$ 为 $k$ 上的
概形。假设：
\begin{enumerate}
\item $X$ 拟紧；
\item $X_{\overline{k}}$ 的连通分支都是开集。
\end{enumerate}
那么：
\begin{enumerate}
\item[(a)] $\pi_0(X_{\overline{k}})$ 有限；
\item[(b)] $\text{Gal}(\overline{k}/k)$ 在
$\pi_0(X_{\overline{k}})$ 上的作用连续。
\end{enumerate}
此外，当 $X$ 在 $k$ 上有限型时，假设 (1) 与 (2) 均成立。
\end{lemma}

\begin{proof}
由于各连通分支都是开集、覆盖 $X_{\overline{k}}$（《拓扑学》引理
\ref{topology-lemma-connected-components}），而且 $X_{\overline{k}}$ 拟紧，
可知它们只有有限多个。因此 (a) 成立。由引理 \ref{varieties-lemma-descend-open}，
这些连通分支中的每一个都可以在 $\overline{k}/k$ 的某个有限子扩张上
定义，故得到 (b)。若 $X$ 在 $k$ 上有限型，则 $X_{\overline{k}}$ 在
$\overline{k}$ 上有限型（《态射》引理
\ref{morphisms-lemma-base-change-finite-type}）。因此
$X_{\overline{k}}$ 是 Noether 概形（《态射》引理
\ref{morphisms-lemma-finite-type-noetherian}）。所以
$X_{\overline{k}}$ 只有有限多个不可约分支（《性质》引理
\ref{properties-lemma-Noetherian-irreducible-components}），更不必说，它只有
有限多个连通分支（因而这些分支都是开的）。
\end{proof}

\section{几何不可约概形}
\label{varieties-section-geometrically-irreducible}

\noindent
若 $X$ 是域上的不可约概形，则在扩张基域之后，$X$ 可能变为可约概形。
几何不可约概形不会发生这种情况。

\begin{definition}
\label{varieties-definition-geometrically-irreducible}
设 $X$ 为域 $k$ 上的概形。称 $X$ 在 $k$ 上{\it 几何不可约}，如果概形
$X_{k'}$ 不可约\footnote{不可约空间非空。}，其中 $k'$ 是 $k$ 的任意域扩张。
\end{definition}

\begin{lemma}
\label{varieties-lemma-geometrically-irreducible-check-after-extension}
设 $X$ 为域 $k$ 上的概形，并设 $k'/k$ 为域扩张。那么，$X$ 在
$k$ 上几何不可约，当且仅当 $X_{k'}$ 在 $k'$ 上几何不可约。
\end{lemma}

\begin{proof}
若 $X$ 在 $k$ 上几何不可约，则 $X_{k'}$ 在 $k'$ 上几何不可约是显然的。
反过来，对任意域扩张 $k''/k$，存在公共域扩张 $k'''/k'$ 和
$k'''/k''$；见《域》引理 \ref{fields-lemma-common-extension-field}。
态射 $X_{k'''} \to X_{k''}$ 是满的（因为它是域谱之间的满态射的基变换），
故 $X_{k'''}$ 不可约蕴含 $X_{k''}$ 不可约。因此，若 $X_{k'}$ 在
$k'$ 上几何不可约，则 $X$ 在 $k$ 上几何不可约。
\end{proof}

\begin{lemma}
\label{varieties-lemma-separably-closed-irreducible}
设 $X$ 为一个概形，它定义在可分代数闭的域 $k$ 上。若 $X$ 不可约，则 $X_K$ 对任意
域扩张 $K/k$ 都不可约。换言之，$X$ 在 $k$ 上几何不可约。
\end{lemma}

\begin{proof}
使用《性质》引理 \ref{properties-lemma-characterize-irreducible}
和《代数》引理 \ref{algebra-lemma-separably-closed-irreducible}。
\end{proof}

\begin{lemma}
\label{varieties-lemma-bijection-irreducible-components}
设 $k$ 为域，设 $X$、$Y$ 为 $k$ 上的概形。假设 $X$ 在 $k$ 上几何不可约。
那么投影态射
$$
p : X \times_k Y \longrightarrow Y
$$
诱导不可约分支之间的双射。
\end{lemma}

\begin{proof}
首先，$p$ 的概形论纤维不可约，因为它们是几何不可约概形 $X$ 沿域扩张
的基变换。此外，概形论纤维与集合论纤维同胚；见《概形》引理
\ref{schemes-lemma-fibre-topological}。由《态射》引理
\ref{morphisms-lemma-scheme-over-field-universally-open}，映射 $p$ 是开的。
因此可以应用《拓扑学》引理
\ref{topology-lemma-irreducible-fibres-irreducible-components} 得出结论。
\end{proof}

\begin{lemma}
\label{varieties-lemma-geometrically-irreducible-local}
\begin{slogan}
除连通性外，几何不可约性是 Zariski 局部的。
\end{slogan}
设 $k$ 为域，设 $X$ 为 $k$ 上的概形。下列条件等价：
\begin{enumerate}
\item $X$ 在 $k$ 上几何不可约；
\item 对每个非空仿射开集 $U$，$k$-代数 $\mathcal{O}_X(U)$ 在
$k$ 上几何不可约（见《代数》定义
\ref{algebra-definition-geometrically-irreducible}）；
\item $X$ 不可约，并且存在仿射开覆盖 $X = \bigcup U_i$，使得每个
$k$-代数 $\mathcal{O}_X(U_i)$ 都几何不可约；
\item 存在开覆盖 $X = \bigcup_{i \in I} X_i$，其中
$I \not = \emptyset$，使得 $X_i$ 对每个 $i$ 都几何不可约，并且
$X_i \cap X_j \not = \emptyset$ 对所有 $i, j \in I$ 都成立。
\end{enumerate}
此外，若 $X$ 几何不可约，则 $X$ 的每个非空开子概形也几何不可约。
\end{lemma}

\begin{proof}
仿射概形 $\Spec(A)$ 在 $k$ 上几何不可约，当且仅当 $A$ 在 $k$ 上
几何不可约；这直接来自定义。回忆，若 $X$ 不可约，则它的每个非空开子概形
也不可约，而且任意两个非空开子集的交非空。又若每个仿射开集都不可约，
则该概形不可约；见《性质》引理
\ref{properties-lemma-characterize-irreducible}。因此，引理的最后一个断言，
以及蕴含 (1) $\Rightarrow$ (2)、(2) $\Rightarrow$ (3) 和
(3) $\Rightarrow$ (4) 都是显然的。若 (4) 成立，则对任意域扩张
$k'/k$，概形 $X_{k'}$ 有一个由两两相交的不可约开集组成的覆盖。因此
$X_{k'}$ 不可约，故 (4) 蕴含 (1)。
\end{proof}

\begin{lemma}
\label{varieties-lemma-geometrically-irreducible-function-field}
设 $X$ 为域 $k$ 上的不可约概形，设 $\xi \in X$ 为其泛点。下列条件等价：
\begin{enumerate}
\item $X$ 在 $k$ 上几何不可约；
\item $\kappa(\xi)$ 在 $k$ 上几何不可约；
\item $k$ 在 $\kappa(\xi)$ 中的可分代数闭包等于 $k$。
\end{enumerate}
\end{lemma}

\begin{proof}
由《代数》引理
\ref{algebra-lemma-geometrically-irreducible-separable-elements}，
(2) 与 (3) 等价。

\medskip\noindent
假设 (1)。回忆，$\mathcal{O}_{X, \xi}$ 是 $\mathcal{O}_X(U)$ 的滤过
余极限，其中 $U$ 遍历 $X$ 的非空仿射开子概形。结合引理
\ref{varieties-lemma-geometrically-irreducible-local} 与《代数》引理
\ref{algebra-lemma-subalgebra-geometrically-irreducible}，可知
$\mathcal{O}_{X, \xi}$ 在 $k$ 上几何不可约。由于
$\mathcal{O}_{X, \xi} \to \kappa(\xi)$ 是具有局部幂零核的满射
（见《代数》引理 \ref{algebra-lemma-minimal-prime-reduced-ring}），由《代数》
引理 \ref{algebra-lemma-p-ring-map} 可知 $\kappa(\xi)$ 几何不可约。

\medskip\noindent
假设 (2)。可以假设 $X$ 约化。设 $U \subset X$ 为非空仿射开集。则
$U = \Spec(A)$，其中 $A$ 是以 $\kappa(\xi)$ 为分式域的整环。因此，
$A$ 是一个几何不可约 $k$-代数的 $k$-子代数。于是由《代数》引理
\ref{algebra-lemma-subalgebra-geometrically-irreducible}，$A$ 在 $k$ 上
几何不可约。由引理 \ref{varieties-lemma-geometrically-irreducible-local} 得出
$X$ 在 $k$ 上几何不可约。
\end{proof}

\begin{lemma}
\label{varieties-lemma-geometrically-irreducible-self-product}
概形 $X$ 定义在域 $k$ 上时，它在 $k$ 上几何不可约，当且仅当
$X\times_k X$ 不可约。
\end{lemma}

\begin{proof}
若 $X$ 在 $k$ 上几何不可约，则由引理
\ref{varieties-lemma-bijection-irreducible-components}，$X\times_k X$ 不可约。

\medskip\noindent
反过来，假设 $X \times_k X$ 不可约。把 $X$ 替换为其底层约化概形，可以
在不改变问题的前提下假设 $X$ 约化。由于第一投影
$X \times_k X \to X$ 满射，$X$ 也不可约，因而是整概形。假设存在元素
$\alpha$，它属于函数域 $k(X)$，是 $k$ 上的可分代数元且不属于 $k$。
令 $E = k(\alpha)$。

\medskip\noindent
那么，存在非空仿射开子概形 $U$，它是 $X$ 的开子概形，并且使得
$k \subset E \subset O(U)$，因而
$E \otimes_k E \subset
O(U)\otimes_k O(U) = O(U\times_k U)$。
于是有支配态射 $U \times_k U \to \Spec(E\otimes_k E)$。由于
$X\times_k X$ 不可约，$U \times_k U$ 也不可约，因而
$\Spec(E \otimes_k E)$ 连通。若 $E$ 严格包含 $k$，则
$E \otimes_k E$ 是若干域的直积，其中一个因子为 $E$。它的维数为
$(\dim_k(E))^2>\dim_k(E)$；这里维数是作为 $k$-向量空间计算的。故 $E\otimes_k E$
是至少两个域的直积，这与 $\Spec(E \otimes_k E)$ 连通矛盾。因此事实上
$E$ 等于 $k$。由引理 \ref{varieties-lemma-geometrically-irreducible-function-field}，
$X$ 在 $k$ 上几何不可约。
\end{proof}

\begin{lemma}
\label{varieties-lemma-separably-closed-field-irreducible-components}
设 $k'/k$ 为域扩张，设 $X$ 为 $k$ 上的概形。置 $X' = X_{k'}$。
假设 $k$ 可分代数闭。那么态射 $X' \to X$ 诱导不可约分支之间的双射。
\end{lemma}

\begin{proof}
由于 $k$ 可分代数闭，可知 $k'$ 在 $k$ 上几何不可约；见《代数》引理
\ref{algebra-lemma-separably-closed-irreducible-implies-geometric}。
因此 $Z = \Spec(k')$ 在 $k$ 上几何不可约；见上面的引理
\ref{varieties-lemma-geometrically-irreducible-local}。由于
$X' = Z \times_k X$，结论是引理
\ref{varieties-lemma-bijection-irreducible-components} 的特例。
\end{proof}

\begin{lemma}
\label{varieties-lemma-characterize-geometrically-irreducible}
\begin{slogan}
几何不可约性可以在基域的一个可分代数闭包上检验。
\end{slogan}
设 $k$ 为域，设 $X$ 为 $k$ 上的概形。下列条件等价：
\begin{enumerate}
\item $X$ 在 $k$ 上几何不可约；
\item 对每个有限可分域扩张 $k'/k$，概形 $X_{k'}$ 都不可约；
\item $X_{\overline{k}}$ 不可约，其中 $k \subset \overline{k}$ 是
$k$ 的一个可分代数闭包。
\end{enumerate}
\end{lemma}

\begin{proof}
假设 $X_{\overline{k}}$ 不可约，即假设 (3)。设 $k'/k$ 为域扩张。
存在域扩张 $\overline{k}'/\overline{k}$，使得 $k'$ 嵌入
$\overline{k}'$，并且该嵌入是 $k$-扩张的嵌入。由引理
\ref{varieties-lemma-separably-closed-field-irreducible-components} 可知
$X_{\overline{k}'}$ 不可约。由于 $X_{\overline{k}'} \to X_{k'}$ 满射，
可知 $X_{k'}$ 不可约。因此 (1) 成立。

\medskip\noindent
设 $k \subset \overline{k}$ 为 $k$ 的一个可分代数闭包。假设 (3) 不成立，
即假设 $X_{\overline{k}}$ 可约。我们的目标是证明，$X_{k'}$ 对某个有限
子扩张 $\overline{k}/k'/k$ 也可约。设
$X = \bigcup_{i \in I} U_i$ 为仿射开覆盖，其中 $U_i$ 非空。若对某个
$i$，概形 $U_i$ 可约，或者对某对 $i \not = j$，交
$U_i \cap U_j$ 为空，则 $X$ 可约（《性质》引理
\ref{properties-lemma-characterize-irreducible}），于是结论成立。特别地，
可以假设 $U_{i, \overline{k}} \cap U_{j, \overline{k}}$ 对所有
$i, j \in I$ 都非空；由此可知，$U_{i, \overline{k}}$ 对某个 $i$
必须可约。根据《代数》引理
\ref{algebra-lemma-geometrically-irreducible}，这意味着 $U_{i, k'}$ 对某个
有限可分域扩张 $k'/k$ 可约。
因此 $X_{k'}$ 也可约。于是 (2) 蕴含 (3)。

\medskip\noindent
蕴含 (1) $\Rightarrow$ (2) 是直接的。这就证明了引理。
\end{proof}

\begin{lemma}
\label{varieties-lemma-inverse-image-irreducible}
设 $K/k$ 为域扩张，设 $X$ 为 $k$ 上的概形。对每个不可约分支 $T$
（它是 $X$ 的不可约分支），逆像 $T_K \subset X_K$ 是 $X_K$
的若干不可约分支之并。
\end{lemma}

\begin{proof}
设 $T \subset X$ 为 $X$ 的一个不可约分支。态射 $T_K \to T$ 平坦，
故泛化可沿 $T_K \to T$ 提升。因此，每个 $\xi \in T_K$，只要它是
$T_K$ 的一个不可约分支的泛点，就映到泛点 $\eta$；后者是 $T$ 的泛点。若
$\xi' \leadsto \xi$ 是 $X_K$ 中的一次特化，则 $\xi'$ 映到 $\eta$，
因为没有特化到 $\eta$ 的点存在于 $X$ 中。故 $\xi' \in T_K$，由此得出
$\xi = \xi'$。换言之，$\xi$ 是 $X_K$ 的一个不可约分支的泛点。
这意味着 $T_K$ 的所有不可约分支也都是 $X_K$ 的不可约分支。
\end{proof}

\noindent
对概形 $X$，用 $\text{IrredComp}(X)$ 表示 $X$ 的不可约分支集合。

\begin{lemma}
\label{varieties-lemma-image-irreducible}
设 $K/k$ 为域扩张，设 $X$ 为 $k$ 上的概形。对每个不可约分支
$\overline{T} \subset X_K$，$\overline{T}$ 在 $X$ 中的像是 $X$ 的一个
不可约分支。这定义了典范映射
$$
\text{IrredComp}(X_K)
\longrightarrow
\text{IrredComp}(X)
$$
并且该映射满射。
\end{lemma}

\begin{proof}
考虑图表
$$
\xymatrix{
X_K \ar[d] & X_{\overline{K}} \ar[d] \ar[l] \\
X & X_{\overline{k}} \ar[l]
}
$$
其中 $\overline{K}$ 是 $K$ 的可分代数闭包，而 $\overline{k}$ 是
$k$ 的可分代数闭包。由引理
\ref{varieties-lemma-separably-closed-field-irreducible-components}，态射
$X_{\overline{K}} \to X_{\overline{k}}$ 诱导不可约分支之间的双射。因此，
只需对态射 $X_{\overline{k}} \to X$ 和
$X_{\overline{K}} \to X_K$ 证明本引理。换言之，可以假设
$K = \overline{k}$。

\medskip\noindent
态射 $p : X_{\overline{k}} \to X$ 是整、平坦且满的。平坦性蕴含泛化可沿
$p$ 提升；见《态射》引理
\ref{morphisms-lemma-generalizations-lift-flat}。因此，
$X_{\overline{k}}$ 的不可约分支的泛点映到 $X$ 的不可约分支的泛点。
整性蕴含 $p$ 万有闭；见《态射》引理
\ref{morphisms-lemma-integral-universally-closed}。所以，不可约分支的像
$p(\overline{T})$ 是包含 $X$ 的某个不可约分支的泛点的闭不可约子集；
因而 $p(\overline{T})$ 是 $X$ 的一个不可约分支。这证明了第一个断言。
若 $T \subset X$ 是不可约分支，则 $p^{-1}(T) =T_K$ 是不可约分支的
非空并；见引理 \ref{varieties-lemma-inverse-image-irreducible}。根据第一部分，
这些分支中的每一个必然满射到 $T$。故该映射是满的。
\end{proof}

\begin{lemma}
\label{varieties-lemma-irreducible-dense-rational-points-geometrically-irreducible}
设 $k$ 为域，设 $X$ 为 $k$ 上的概形。若 $X$ 不可约且具有稠密的
$k$-有理点集，则 $X$ 几何不可约。
\end{lemma}

\begin{proof}
设 $k'/k$ 为有限域扩张，并设 $Z, Z' \subset X_{k'}$ 为不可约分支。
只需证明 $Z = Z'$；见引理
\ref{varieties-lemma-characterize-geometrically-irreducible}。由引理
\ref{varieties-lemma-image-irreducible}，有 $p(Z) = p(Z') = X$，其中
$p : X_{k'} \to X$ 是投影。若 $Z \not = Z'$，则
$Z \cap Z'$ 在 $X_{k'}$ 中无处稠密，因而由《态射》引理
\ref{morphisms-lemma-image-nowhere-dense-finite}，$p(Z \cap Z')$
不稠密；这里还使用了 $p$ 是有限态射，因为它是有限态射
$\Spec(k') \to \Spec(k)$ 的基变换；见《态射》引理
\ref{morphisms-lemma-base-change-finite}。因此，可以选取一个
$k$-有理点 $x \in X$，使得 $x \not \in p(Z \cap Z')$。由于 $x$ 的
剩余域是 $k$，可知 $p^{-1}(\{x\}) = \{x'\}$，其中
$x' \in X_{k'}$ 是剩余域为 $k'$ 的点。由于
$x \in p(Z) = p(Z')$，可知 $x' \in Z \cap Z'$；这就是所需的矛盾。
\end{proof}

\begin{lemma}
\label{varieties-lemma-galois-action-irreducible-components}
设 $k$ 为域，并设 $\overline{k}$ 为其可分代数闭包。设 $X$ 为 $k$ 上的
概形。存在作用
$$
\text{Gal}(\overline{k}/k)^{opp} \times \text{IrredComp}(X_{\overline{k}})
\longrightarrow
\text{IrredComp}(X_{\overline{k}})
$$
并具有下列性质：
\begin{enumerate}
\item 元素 $\overline{T} \in \text{IrredComp}(X_{\overline{k}})$ 被该作用
固定，当且仅当存在不可约分支 $T \subset X$，它在 $k$ 上几何不可约，
并且满足 $T_{\overline{k}} = \overline{T}$。
\item 对任意域扩张 $k'/k$ 及其可分代数闭包 $\overline{k}'$，图表
$$
\xymatrix{
\text{Gal}(\overline{k}'/k') \times \text{IrredComp}(X_{\overline{k}'})
\ar[r] \ar[d] &
\text{IrredComp}(X_{\overline{k}'}) \ar[d] \\
\text{Gal}(\overline{k}/k) \times \text{IrredComp}(X_{\overline{k}})
\ar[r] &
\text{IrredComp}(X_{\overline{k}})
}
$$
交换（由引理 \ref{varieties-lemma-separably-closed-field-irreducible-components}，
右侧竖直箭头为双射）。
\end{enumerate}
\end{lemma}

\begin{proof}
$\text{Gal}(\overline{k}/k)$ 在 $X_{\overline{k}}$ 上的作用
（\ref{varieties-equation-galois-action-base-change-kbar}）诱导其在不可约分支上的
作用。不可约分支总是闭的（《拓扑学》引理
\ref{topology-lemma-connected-components}）。因此，若 $\overline{T}$
如 (1) 所述，则由引理 \ref{varieties-lemma-closed-fixed-by-Galois}，存在闭子集
$T \subset X$，使得 $\overline{T} = T_{\overline{k}}$。注意，$T$ 在
$k$ 上几何不可约；见引理
\ref{varieties-lemma-characterize-geometrically-irreducible}。为说明 $T$ 是
$X$ 的一个不可约分支，假设 $T \subset T'$、$T \not = T'$，其中
$T'$ 是 $X$ 的一个不可约分支。设 $\overline{\eta}$ 为
$\overline{T}$ 的泛点。它映到泛点 $\eta$，后者是 $T$ 的泛点。于是泛点
$\xi \in T'$ 特化到 $\eta$。由于 $X_{\overline{k}} \to X$ 平坦，
存在点 $\overline{\xi} \in X_{\overline{k}}$，它映到 $\xi$ 并特化到
$\overline{\eta}$。于是，单点集 $\{\overline{\xi}\}$ 在
$X_{\overline{\xi}}$ 中的闭包是一个严格包含 $\overline{T}$ 的
不可约闭子集。这就是所需的矛盾。

\medskip\noindent
我们略去 (2) 中函子性的证明。
\end{proof}

\begin{lemma}
\label{varieties-lemma-orbit-irreducible-components}
设 $k$ 为域，并设 $\overline{k}$ 为其可分代数闭包。设 $X$ 为 $k$ 上的
概形。引理 \ref{varieties-lemma-image-irreducible} 的映射
$$
\text{IrredComp}(X_{\overline{k}})
\longrightarrow
\text{IrredComp}(X)
$$
的各纤维，恰好是 $\text{Gal}(\overline{k}/k)$ 在引理
\ref{varieties-lemma-galois-action-irreducible-components} 的作用下的轨道。
\end{lemma}

\begin{proof}
设 $T \subset X$ 为 $X$ 的一个不可约分支，设 $\eta \in T$ 为其泛点。
由引理 \ref{varieties-lemma-inverse-image-irreducible} 和
\ref{varieties-lemma-image-irreducible}，$\overline{T}$ 的各不可约分支中映入
$T$ 的那些分支的泛点都映到 $\eta$。由《代数》引理
\ref{algebra-lemma-Galois-orbit}，Galois 群在 $X_{\overline{k}}$ 中所有
映到 $\eta$ 的点上可迁地作用。因此引理成立。
\end{proof}

\begin{lemma}
\label{varieties-lemma-galois-action-irreducible-components-locally-finite-type}
设 $k$ 为域。假设 $X \to \Spec(k)$ 局部有限型。在这种情况下：
\begin{enumerate}
\item 若赋予 $\text{IrredComp}(X_{\overline{k}})$ 离散拓扑，则作用
$$
\text{Gal}(\overline{k}/k)^{opp} \times \text{IrredComp}(X_{\overline{k}})
\longrightarrow
\text{IrredComp}(X_{\overline{k}})
$$
是连续的；
\item $X_{\overline{k}}$ 的每个不可约分支都可以在 $k$ 的某个有限扩张上
定义；
\item 给定任意不可约分支 $T \subset X$，概形
$T_{\overline{k}}$ 是 $X_{\overline{k}}$ 的有限多个不可约分支之并，
而且这些分支都属于同一个 $\text{Gal}(\overline{k}/k)$-轨道。
\end{enumerate}
\end{lemma}

\begin{proof}
设 $\overline{T}$ 为 $X_{\overline{k}}$ 的一个不可约分支。可以选取仿射
开集 $U \subset X$，使得 $\overline{T} \cap U_{\overline{k}}$ 非空。
写成 $U = \Spec(A)$；于是 $A$ 是有限型 $k$-代数，见《态射》引理
\ref{morphisms-lemma-locally-finite-type-characterize}。因此
$A_{\overline{k}}$ 是有限型 $\overline{k}$-代数，特别地，它是 Noether
环。设 $\mathfrak p = (f_1, \ldots, f_n)$ 为对应于
$\overline{T} \cap U_{\overline{k}}$ 的素理想。由于
$A_{\overline{k}} = A \otimes_k \overline{k}$，可知存在有限子扩张
$\overline{k}/k'/k$，使得每个 $f_i \in A_{k'}$。显然，
$\text{Gal}(\overline{k}/k')$ 固定 $\overline{T}$，这证明了 (1)。

\medskip\noindent
把引理 \ref{varieties-lemma-galois-action-irreducible-components} (1) 应用于
$k'$ 上的情形，可知不可约分支 $\overline{T}$ 形如
$T'_{\overline{k}}$，其中 $T' \subset X_{k'}$ 不可约，故 (2) 成立。

\medskip\noindent
为证明 (3)，设 $T \subset X$ 为一个不可约分支。选取不可约分支
$\overline{T} \subset X_{\overline{k}}$，它映到 $T$；见引理
\ref{varieties-lemma-image-irreducible}。由上可知，$\overline{T}$ 的轨道有限，
记为 $\overline{T}_1, \ldots, \overline{T}_n$。那么
$\overline{T}_1 \cup \ldots \cup \overline{T}_n$ 是一个
$\text{Gal}(\overline{k}/k)$-不变闭子集，它位于 $X_{\overline{k}}$ 中，故由引理
\ref{varieties-lemma-closed-fixed-by-Galois}，它形如 $W_{\overline{k}}$，其中
$W \subset X$ 是闭子集。显然 $W = T$，故结论成立。
\end{proof}

\begin{lemma}
\label{varieties-lemma-finite-extension-geometrically-irreducible-components}
设 $k$ 为域，设 $X \to \Spec(k)$ 局部有限型。假设 $X$ 只有有限多个
不可约分支。那么存在有限可分扩张 $k'/k$，使得 $X_{k'}$ 的每个
不可约分支都在 $k'$ 上几何不可约。
\end{lemma}

\begin{proof}
设 $\overline{k}$ 为 $k$ 的可分代数闭包。$X$ 只有有限多个不可约分支
这一假设，与引理
\ref{varieties-lemma-galois-action-irreducible-components-locally-finite-type} (3)
合起来表明，$X_{\overline{k}}$ 只有有限多个不可约分支
$\overline{T}_1, \ldots, \overline{T}_n$。由引理
\ref{varieties-lemma-galois-action-irreducible-components-locally-finite-type} (2)，
存在有限扩张 $\overline{k}/k'/k$ 和不可约分支 $T_i \subset X_{k'}$，
使得 $\overline{T}_i = T_{i, \overline{k}}$，故结论成立。
\end{proof}

\begin{lemma}
\label{varieties-lemma-irreducible-components-geometrically-irreducible}
设 $X$ 为域 $k$ 上的概形。假设 $X$ 只有有限多个不可约分支，并且它们
全都几何不可约。那么 $X$ 只有有限多个连通分支，并且每个连通分支都
几何连通。
\end{lemma}

\begin{proof}
这是显然的，因为一个连通分支是若干不可约分支之并。细节略去。
\end{proof}


\section{几何整概形}
\label{varieties-section-geometrically-integral}

\noindent
若 $X$ 是域上的整概形，则在扩张基域后，$X$ 可能变得非约化或可约。
几何整概形不会发生这种情况。

\begin{definition}
\label{varieties-definition-geometrically-integral}
设 $X$ 为域 $k$ 上的概形。
\begin{enumerate}
\item 设 $x \in X$。称 $X$ 在 $x$ 处{\it 几何逐点整}，如果对每个域扩张
$k'/k$，以及每个 $x' \in X_{k'}$，只要它位于 $x$ 上方，局部环
$\mathcal{O}_{X_{k'}, x'}$ 都是整环。
\item 若 $X$ 在每一点处都几何逐点整，则称 $X$ {\it 几何逐点整}。
\item 称 $X$ 在 $k$ 上{\it 几何整}，如果概形 $X_{k'}$ 对每个域扩张
$k'$（它是 $k$ 的扩张）都整。
\end{enumerate}
\end{definition}

\noindent
概念 (2) 与 (3) 之间的区别是必要的。例如，若
$k = \mathbf{R}$ 且 $X = \Spec(\mathbf{C}[x])$，则 $X$ 在
$\mathbf{R}$ 上几何逐点整，但当然不是几何整的。

\begin{lemma}
\label{varieties-lemma-geometrically-integral}
设 $k$ 为域，设 $X$ 为 $k$ 上的概形。那么，$X$ 在 $k$ 上几何整，
当且仅当 $X$ 在 $k$ 上既几何约化又几何不可约。
\end{lemma}

\begin{proof}
见《性质》引理 \ref{properties-lemma-characterize-integral}。
\end{proof}

\begin{lemma}
\label{varieties-lemma-proper-geometrically-reduced-global-sections}
设 $k$ 为域，设 $X$ 为 $k$ 上的固有概形。
\begin{enumerate}
\item $A = H^0(X, \mathcal{O}_X)$ 是有限维 $k$-代数；
\item $A = \prod_{i = 1, \ldots, n} A_i$ 是 Artin 局部
$k$-代数的直积，$X$ 的每个连通分支对应一个因子；
\item 若 $X$ 约化，则 $A = \prod_{i = 1, \ldots, n} k_i$
是若干域的直积，其中每个域都是 $k$ 的有限扩张；
\item 若 $X$ 几何约化，则 $k_i$ 是 $k$ 的有限可分扩张；
\item 若 $X$ 几何连通，则 $A$ 在 $k$ 上几何不可约；
\item 若 $X$ 几何不可约，则 $A$ 在 $k$ 上几何不可约；
\item 若 $X$ 几何约化且几何连通，则 $A = k$；
\item 若 $X$ 几何整，则 $A = k$。
\end{enumerate}
\end{lemma}

\begin{proof}
由《概形的上同调》引理
\ref{coherent-lemma-proper-over-affine-cohomology-finite}，可知
$A = H^0(X, \mathcal{O}_X)$ 是有限维 $k$-代数。这证明了 (1)。

\medskip\noindent
于是，由《代数》引理 \ref{algebra-lemma-finite-dimensional-algebra}
和命题 \ref{algebra-proposition-dimension-zero-ring}，$A$ 是若干 Artin
局部 $k$-代数的直积。若 $X = Y \amalg Z$，其中 $Y$ 与 $Z$ 在 $X$
中都是开集，则可得到幂等元 $e \in A$：取 $\mathcal{O}_X$ 的一个截面，
使其取值为 $1$（在 $Y$ 上）、取值为 $0$（在 $Z$ 上）。反过来，若 $e \in A$
是幂等元，则得到 $X$ 的相应分解。最后，由于 $X$ 的底层拓扑空间是
Noether 空间，它的连通分支都是开的。因此，$X$ 的连通分支以
$1$-to-$1$ 的方式对应于 $A$ 的本原幂等元。这证明了 (2)。

\medskip\noindent
若 $X$ 约化，则 $A$ 约化。因此，局部环 $A_i = k_i$ 约化，从而是域
（例如可用《代数》引理
\ref{algebra-lemma-minimal-prime-reduced-ring}）。这证明了 (3)。

\medskip\noindent
若 $X$ 几何约化，则
$A \otimes_k \overline{k} =
H^0(X_{\overline{k}}, \mathcal{O}_{X_{\overline{k}}})$
（该等式来自《概形的上同调》引理
\ref{coherent-lemma-flat-base-change-cohomology}）是约化的。这蕴含
$k_i \otimes_k \overline{k}$ 是若干域的直积，因而 $k_i/k$ 可分；
例如可用《代数》引理
\ref{algebra-lemma-characterize-separable-field-extensions} 和
\ref{algebra-lemma-geometrically-reduced-finite-purely-inseparable-extension}。
这证明了 (4)。

\medskip\noindent
若 $X$ 几何连通，则
$A \otimes_k \overline{k} =
H^0(X_{\overline{k}}, \mathcal{O}_{X_{\overline{k}}})$
由 (2) 是零维局部环，故它的谱只有一个点，特别地，它不可约。因此
$A$ 几何不可约。这证明了 (5)。当然，(5) 蕴含 (6)。

\medskip\noindent
若 $X$ 几何约化且几何连通，则 $A = k_1$ 是域，而扩张 $k_1/k$
有限可分且几何不可约。然而，此时 $k_1 \otimes_k \overline{k}$
是 $[k_1 : k]$ 份 $\overline{k}$ 的直积，故可知 $k_1 = k$。
这证明了 (7)。当然，(7) 蕴含 (8)。
\end{proof}

\noindent
下面给出 Stein 分解的一个初步版本；真正的 Stein 分解将在《态射续篇》第
\ref{more-morphisms-section-stein-factorization} 节讨论。

\begin{lemma}
\label{varieties-lemma-baby-stein}
设 $X$ 为域 $k$ 上的固有概形。置
$A = H^0(X, \mathcal{O}_X)$。典范态射 $X \to \Spec(A)$ 的纤维几何连通。
\end{lemma}

\begin{proof}
置 $S = \Spec(A)$。典范态射 $X \to S$ 是通过《概形》引理
\ref{schemes-lemma-morphism-into-affine} 与等式
$\Gamma(S, \mathcal{O}_S) = A = \Gamma(X, \mathcal{O}_X)$
相对应的态射。$k$-代数 $A$ 是有限直积
$A = \prod A_i$，其因子是 $k$ 上有限的 Artin 局部 $k$-代数；见引理
\ref{varieties-lemma-proper-geometrically-reduced-global-sections}。用
$s_i \in S$ 表示对应于 $A_i$ 的极大理想的点。选取代数闭包
$\overline{k}$，它是 $k$ 的一个代数闭包，并置
$\overline{A} = A \otimes_k \overline{k}$。
选取嵌入 $\kappa(s_i) \to \overline{k}$，它在 $k$ 上；该嵌入确定
$\overline{k}$-代数映射
$$
\sigma_i : \overline{A} = A \otimes_k \overline{k} \to
\kappa(s_i) \otimes_k \overline{k} \to \overline{k}
$$
考虑基变换
$$
\xymatrix{
\overline{X} \ar[r] \ar[d] & X \ar[d] \\
\overline{S} \ar[r] & S
}
$$
它把 $X$ 基变换到 $\overline{S} = \Spec(\overline{A})$。由
《概形的上同调》引理
\ref{coherent-lemma-flat-base-change-cohomology}，有
$\Gamma(\overline{X}, \mathcal{O}_{\overline{X}}) = \overline{A}$。
若 $\overline{s}_i \in \Spec(\overline{A})$ 表示一个
$\overline{k}$-有理点，它对应于 $\sigma_i$，则 $\overline{s}_i$ 映到
$s_i \in S$，并且
$\overline{X}_{\overline{s}_i}$ 是 $X_{s_i}$ 沿 $\Spec(\sigma_i)$
的基变换。因此，只需在 $k$ 代数闭的情形下证明本引理。

\medskip\noindent
假设 $k$ 代数闭。在这种情况下，$\kappa(s_i)$ 代数闭，而我们需要证明
$X_{s_i}$ 连通。直积分解 $A = \prod A_i$ 对应于不交并分解
$\Spec(A) = \coprod \Spec(A_i)$；见《代数》引理
\ref{algebra-lemma-spec-product}。用 $X_i$ 表示 $\Spec(A_i)$ 的逆像。
由引理 \ref{varieties-lemma-proper-geometrically-reduced-global-sections} (2)，有
$A_i = \Gamma(X_i, \mathcal{O}_{X_i})$。注意，
$X_{s_i} \to X_i$ 是闭浸入，并在底层拓扑空间上诱导同构（因为
$\Spec(A_i)$ 是单点集）。因此，若 $X_{s_i}$ 不连通，则 $X_i$
也不连通。所以，要么 $X_i$ 为空且 $A_i = 0$，要么 $X_i$
可写成 $U \amalg V$，其中 $U$ 与 $V$ 是非空开集；这将蕴含 $A_i$
有非平凡幂等元。由于 $A_i$ 是局部环，这构成矛盾，证明完成。
\end{proof}

\begin{lemma}
\label{varieties-lemma-geometrically-reduced-stein}
设 $k$ 为域，设 $X$ 为 $k$ 上固有且几何约化的概形。下列条件等价：
\begin{enumerate}
\item $H^0(X, \mathcal{O}_X) = k$；
\item $X$ 几何连通。
\end{enumerate}
\end{lemma}

\begin{proof}
由引理 \ref{varieties-lemma-baby-stein}，有 (1) $\Rightarrow$ (2)。
由引理 \ref{varieties-lemma-proper-geometrically-reduced-global-sections}，
有 (2) $\Rightarrow$ (1)。
\end{proof}


\section{几何正规概形}
\label{varieties-section-geometrically-normal}

\noindent
我们已在《性质》定义 \ref{properties-definition-normal} 中定义了正规概形的
概念。这个概念甚至对非 Noether 概形也有定义。因此，与我们对
“几何正则”概形的讨论不同，这里要考虑基域的所有域扩张。

\begin{definition}
\label{varieties-definition-geometrically-normal}
设 $X$ 为域 $k$ 上的概形。
\begin{enumerate}
\item 设 $x \in X$。称 $X$ 在 $x$ 处{\it 几何正规}，如果对每个域扩张
$k'/k$，以及每个 $x' \in X_{k'}$，只要它位于 $x$ 上方，局部环
$\mathcal{O}_{X_{k'}, x'}$ 都正规。
\item 称 $X$ 在 $k$ 上{\it 几何正规}，如果 $X$ 在每个 $x \in X$
处都几何正规。
\end{enumerate}
\end{definition}

\begin{lemma}
\label{varieties-lemma-geometrically-normal-at-point}
设 $k$ 为域，设 $X$ 为 $k$ 上的概形，并设 $x \in X$。下列条件等价：
\begin{enumerate}
\item $X$ 在 $x$ 处几何正规；
\item 对每个有限纯不可分域扩张 $k'$（它是 $k$ 的扩张），以及每个
$x' \in X_{k'}$，只要它位于 $x$ 上方，局部环
$\mathcal{O}_{X_{k'}, x'}$ 就正规；
\item 环 $\mathcal{O}_{X, x}$ 在 $k$ 上几何正规（见《代数》定义
\ref{algebra-definition-geometrically-normal}）。
\end{enumerate}
\end{lemma}

\begin{proof}
(1) 蕴含 (2) 是显然的。假设 (2)。设 $k'/k$ 为有限纯不可分域扩张
（例如 $k = k'$）。考虑环 $\mathcal{O}_{X, x} \otimes_k k'$。
由《代数》引理 \ref{algebra-lemma-p-ring-map}，它的谱与
$\mathcal{O}_{X, x}$ 的谱相同。因此，它也是局部环（《代数》引理
\ref{algebra-lemma-characterize-local-ring}）。所以，存在唯一的点
$x' \in X_{k'}$ 位于 $x$ 上方，并且
$\mathcal{O}_{X_{k'}, x'} \cong \mathcal{O}_{X, x} \otimes_k k'$。
由假设，这是正规环。因此，由《代数》引理
\ref{algebra-lemma-geometrically-normal} 得到 (3)。

\medskip\noindent
假设 (3)。设 $k'/k$ 为域扩张。由于
$\Spec(k') \to \Spec(k)$ 满射，$X_{k'} \to X$ 也满射
（《态射》引理 \ref{morphisms-lemma-base-change-surjective}）。
设 $x' \in X_{k'}$ 为位于 $x$ 上方的任意一点。局部环
$\mathcal{O}_{X_{k'}, x'}$ 是环
$\mathcal{O}_{X, x} \otimes_k k'$ 的一个局部化。因此，它由假设是正规
的，从而证明了 (1)。
\end{proof}

\begin{lemma}
\label{varieties-lemma-geometrically-normal}
设 $k$ 为域，设 $X$ 为 $k$ 上的概形。下列条件等价：
\begin{enumerate}
\item $X$ 几何正规；
\item $X_{k'}$ 对每个域扩张 $k'/k$ 都是正规概形；
\item $X_{k'}$ 对每个有限生成域扩张 $k'/k$ 都是正规概形；
\item $X_{k'}$ 对每个有限纯不可分域扩张 $k'/k$ 都是正规概形；
\item 对每个仿射开集 $U \subset X$，环 $\mathcal{O}_X(U)$ 都几何正规
（见《代数》定义 \ref{algebra-definition-geometrically-normal}）；
\item $X_{k^{perf}}$ 是正规概形。
\end{enumerate}
\end{lemma}

\begin{proof}
假设 (1)。那么，对每个域扩张 $k'/k$ 和每个点
$x' \in X_{k'}$，$X_{k'}$ 在 $x'$ 处的局部环都正规。按定义，这意味着
$X_{k'}$ 正规。因此 (2) 成立。

\medskip\noindent
(2) 蕴含 (3)、(3) 蕴含 (4)，这都是显然的。

\medskip\noindent
假设 (4)，并设 $U \subset X$ 为仿射开子概形。那么，$U_{k'}$ 对任意
有限纯不可分扩张 $k'/k$ 都是正规概形（包括 $k = k'$）。这意味着，环
$k' \otimes_k \mathcal{O}(U)$ 对所有有限纯不可分扩张 $k'/k$ 都正规。因此，
$\mathcal{O}(U)$ 按定义是几何正规 $k$-代数。故 (4) 蕴含 (5)。

\medskip\noindent
假设 (5)。对任意域扩张 $k'/k$，基变换 $X_{k'}$ 由各个环
$\mathcal{O}_X(U) \otimes_k k'$ 的谱粘合而得，其中 $U$ 是 $X$ 中的
仿射开集（见《概形》第 \ref{schemes-section-fibre-products} 节）。
因此 $X_{k'}$ 正规，所以 (1) 成立。

\medskip\noindent
(5) 与 (6) 的等价性来自几何正规代数的定义，以及刚刚证明的 (3) 与
(4) 的等价性。
\end{proof}

\begin{lemma}
\label{varieties-lemma-geometrically-normal-upstairs}
设 $k$ 为域，设 $X$ 为 $k$ 上的概形，设 $k'/k$ 为域扩张。设
$x \in X$ 为一点，并设 $x' \in X_{k'}$ 为位于 $x$ 上方的一点。
下列条件等价：
\begin{enumerate}
\item $X$ 在 $x$ 处几何正规；
\item $X_{k'}$ 在 $x'$ 处几何正规。
\end{enumerate}
特别地，$X$ 在 $k$ 上几何正规，当且仅当 $X_{k'}$ 在 $k'$ 上几何正规。
\end{lemma}

\begin{proof}
(1) 蕴含 (2) 是显然的。假设 (2)。设 $k''/k$ 为有限纯不可分域扩张，
并设 $x'' \in X_{k''}$ 为位于 $x$ 上方的一点（实际上它是唯一的）。
选取公共域扩张 $k'''/k$，它同时扩张 $k'/k$ 与 $k''/k$；见《域》引理
\ref{fields-lemma-common-extension-field}。选取点
$x''' \in X_{k'''}$，它同时位于 $x'$ 与 $x''$ 上方。考虑局部环映射
$$
\mathcal{O}_{X_{k''}, x''} \longrightarrow \mathcal{O}_{X_{k'''}, x''''}.
$$
这是平坦局部环同态，因而忠实平坦。由 (2)，右侧的局部环正规。因此由
《代数》引理 \ref{algebra-lemma-descent-normal}，可知
$\mathcal{O}_{X_{k''}, x''}$ 正规。由引理
\ref{varieties-lemma-geometrically-normal-at-point}，可知 $X$ 在 $x$ 处几何正规。
\end{proof}

\begin{lemma}
\label{varieties-lemma-fibre-product-normal}
设 $k$ 为域，设 $X$ 为 $k$ 上几何正规的概形，并设 $Y$ 为 $k$ 上正规的
概形。那么 $X \times_k Y$ 是正规概形。
\end{lemma}

\begin{proof}
这归结为《代数》引理
\ref{algebra-lemma-geometrically-normal-tensor-normal}；所用的归约来自引理
\ref{varieties-lemma-geometrically-normal}。
\end{proof}

\begin{lemma}
\label{varieties-lemma-base-change-normal-by-separable}
设 $k$ 为域，设 $X$ 为 $k$ 上的正规概形。设 $K/k$ 为可分域扩张。
那么 $X_K$ 是正规概形。
\end{lemma}

\begin{proof}
这来自引理 \ref{varieties-lemma-fibre-product-normal} 和《代数》引理
\ref{algebra-lemma-separable-field-extension-geometrically-normal}。
\end{proof}

\begin{lemma}
\label{varieties-lemma-geometrically-normal-stein}
设 $k$ 为域，设 $X$ 为 $k$ 上固有且几何正规的概形。下列条件等价：
\begin{enumerate}
\item $H^0(X, \mathcal{O}_X) = k$；
\item $X$ 几何连通；
\item $X$ 几何不可约；
\item $X$ 几何整。
\end{enumerate}
\end{lemma}

\begin{proof}
由引理 \ref{varieties-lemma-geometrically-reduced-stein}，(1) 与 (2) 等价。局部
Noether 正规概形（例如 $X_{\overline{k}}$）是其不可约分支的不交并
（《性质》引理 \ref{properties-lemma-normal-Noetherian}）。因此 (2) 与
(3) 等价。由于假设 $X_{\overline{k}}$ 约化，(3) 与 (4) 也等价。
\end{proof}


\section{域变换与局部 Noether 概形}
\label{varieties-section-locally-Noetherian}

\noindent
设 $X$ 为域 $k$ 上的局部 Noether 概形。$X_{k'}$ 不一定仍是局部
Noether 概形。例如，若 $X = \Spec(\overline{\mathbf{Q}})$ 且
$k = \mathbf{Q}$，则 $X_{\overline{\mathbf{Q}}}$ 是
$\overline{\mathbf{Q}} \otimes_{\mathbf{Q}} \overline{\mathbf{Q}}$
的谱，而后者不是 Noether 环。（提示：它有太多幂等元。）但是，如果只沿
有限生成域扩张作基变换，则 Noether 性得以保持。（或者，若 $X$ 在
$k$ 上局部有限型，结论也成立，因为这个性质在基变换下保持。）

\begin{lemma}
\label{varieties-lemma-locally-Noetherian-base-change}
设 $k$ 为域，设 $X$ 为 $k$ 上的概形，并设 $k'/k$ 为有限生成域扩张。
那么，$X$ 局部 Noether，当且仅当 $X_{k'}$ 局部 Noether。
\end{lemma}

\begin{proof}
利用《性质》引理 \ref{properties-lemma-locally-Noetherian}，可归约到
$X$ 仿射的情形，比如说 $X = \Spec(A)$。在这种情况下，需要证明
$A$ 是 Noether 环，当且仅当 $A_{k'}$ 是 Noether 环。由于
$A \to A_{k'} = k' \otimes_k A$ 忠实平坦，由《代数》引理
\ref{algebra-lemma-descent-Noetherian}，若 $A_{k'}$ 是 Noether 环，
则 $A$ 也是 Noether 环。反过来，若 $A$ 是 Noether 环，则由《代数》引理
\ref{algebra-lemma-Noetherian-field-extension}，$A_{k'}$ 是 Noether 环。
\end{proof}


\section{几何正则概形}
\label{varieties-section-geometrically-regular}

\noindent
域 $k$ 上的几何正则概形，是域 $k$ 上的局部 Noether 概形，并且在适当的
基域变换后仍为正则概形。域 $k$ 上的有限型概形几何正则，当且仅当它在
$k$ 上光滑（见引理 \ref{varieties-lemma-geometrically-regular-smooth}）。几何正则性
这个概念最有意义的情形，是无法使用光滑性的情形，例如形式纤维（此处以后
插入引用）。

\medskip\noindent
在下面的定义中，我们限于局部 Noether 概形，因为正则局部环这个性质只对
Noether 局部环定义。由上面的引理
\ref{varieties-lemma-locally-Noetherian-base-change}，如果只考虑有限生成域扩张，
这个性质在基域变换下保持。本节将直接使用这一说明，不再另行引用。特别地，
下面的定义是有意义的。

\begin{definition}
\label{varieties-definition-geometrically-regular}
设 $k$ 为域，设 $X$ 为 $k$ 上的局部 Noether 概形。
\begin{enumerate}
\item 设 $x \in X$。称 $X$ 在 $x$ 处相对于 $k$ {it 几何正则}，如果对
每个有限生成域扩张 $k'/k$ 以及任意 $x' \in X_{k'}$，只要后者位于 $x$
上方，局部环 $\mathcal{O}_{X_{k'}, x'}$ 就正则。
\item 称 $X$ {it 在 $k$ 上几何正则}，如果 $X$ 在其所有点处都几何正则。
\end{enumerate}
\end{definition}

\noindent
类似的定义也可用于定义域上的几何 Cohen--Macaulay、$(R_k)$ 与 $(S_k)$
概形。需要时，我们会另设一节讨论这些性质。

\begin{lemma}
\label{varieties-lemma-geometrically-regular-at-point}
设 $k$ 为域，设 $X$ 为 $k$ 上的局部 Noether 概形，并设 $x \in X$。
下列条件等价：
\begin{enumerate}
\item $X$ 在 $x$ 处几何正则；
\item 对每个有限纯不可分扩域 $k'$（其底域为 $k$）以及
$x' \in X_{k'}$，只要后者位于 $x$ 上方，局部环
$\mathcal{O}_{X_{k'}, x'}$ 就正则；
\item 环 $\mathcal{O}_{X, x}$ 在 $k$ 上几何正则（见《代数》定义
\ref{algebra-definition-geometrically-regular}）。
\end{enumerate}
\end{lemma}

\begin{proof}
(1) 蕴含 (2) 是清楚的。假设 (2)。特别地，这蕴含 $\mathcal{O}_{X, x}$
是正则局部环。设 $k'/k$ 为有限纯不可分域扩张。考虑环
$\mathcal{O}_{X, x} \otimes_k k'$。由《代数》引理
\ref{algebra-lemma-p-ring-map}，它的谱与 $\mathcal{O}_{X, x}$ 的谱相同，
因而它也是局部环（《代数》引理
\ref{algebra-lemma-characterize-local-ring}）。所以存在唯一一点
$x' \in X_{k'}$ 位于 $x$ 上方，且
$\mathcal{O}_{X_{k'}, x'} \cong \mathcal{O}_{X, x} \otimes_k k'$。
根据假设，这是正则环。因此，由几何正则环的定义得到 (3)。

\medskip\noindent
假设 (3)。设 $k'/k$ 为域扩张。由于
$\Spec(k') \to \Spec(k)$ 满射，
$X_{k'} \to X$ 也满射（《态射》引理
\ref{morphisms-lemma-base-change-surjective}）。任取一点
$x' \in X_{k'}$，使它位于 $x$ 上方。局部环
$\mathcal{O}_{X_{k'}, x'}$ 是环
$\mathcal{O}_{X, x} \otimes_k k'$ 的局部化。因此，根据假设它正则，
从而 (1) 得证。
\end{proof}

\begin{lemma}
\label{varieties-lemma-geometrically-regular}
设 $k$ 为域，设 $X$ 为 $k$ 上的局部 Noether 概形。下列条件等价：
\begin{enumerate}
\item $X$ 几何正则；
\item $X_{k'}$ 对每个有限生成域扩张 $k'/k$ 都是正则概形；
\item $X_{k'}$ 对每个有限纯不可分域扩张 $k'/k$ 都是正则概形；
\item 对每个仿射开集 $U \subset X$，环 $\mathcal{O}_X(U)$ 几何正则
（见《代数》定义 \ref{algebra-definition-geometrically-regular}）；
\item 存在仿射开覆盖 $X = \bigcup U_i$，使每个 $\mathcal{O}_X(U_i)$
都在 $k$ 上几何正则。
\end{enumerate}
\end{lemma}

\begin{proof}
假设 (1)。那么，对每个有限生成域扩张 $k'/k$ 以及每个点
$x' \in X_{k'}$，$X_{k'}$ 在 $x'$ 处的局部环都正则。由《性质》引理
\ref{properties-lemma-characterize-regular}，这意味着 $X_{k'}$ 正则，
所以 (2) 成立。

\medskip\noindent
(2) 蕴含 (3) 是清楚的。

\medskip\noindent
假设 (3)，并设 $U \subset X$ 为仿射开子概形。那么 $U_{k'}$ 对任意有限
纯不可分扩张 $k'/k$（包括 $k = k'$）都是正则概形。这意味着
$k' \otimes_k \mathcal{O}(U)$ 对所有有限纯不可分扩张 $k'/k$ 都是正则环。因此，
$\mathcal{O}(U)$ 是几何正则 $k$-代数，从而 (4) 成立。

\medskip\noindent
(4) 蕴含 (5) 是清楚的。设 $X = \bigcup U_i$ 为 (5) 中的仿射开覆盖。
对任意域扩张 $k'/k$，基变换 $X_{k'}$ 由各环
$\mathcal{O}_X(U_i) \otimes_k k'$ 的谱粘合而成（见《概形》第
\ref{schemes-section-fibre-products} 节）。因此 $X_{k'}$ 正则，所以 (1)
成立。
\end{proof}

\begin{lemma}
\label{varieties-lemma-geometrically-regular-upstairs}
设 $k$ 为域，设 $X$ 为 $k$ 上的概形。设 $k'/k$ 为有限生成域扩张。
设 $x \in X$，并设 $x' \in X_{k'}$ 为位于 $x$ 上方的点。下列条件等价：
\begin{enumerate}
\item $X$ 在 $x$ 处几何正则；
\item $X_{k'}$ 在 $x'$ 处几何正则。
\end{enumerate}
特别地，$X$ 在 $k$ 上几何正则，当且仅当 $X_{k'}$ 在 $k'$ 上几何正则。
\end{lemma}

\begin{proof}
(1) 蕴含 (2) 是清楚的。假设 (2)。设 $k''/k$ 为有限纯不可分域扩张，
并设 $x'' \in X_{k''}$ 为位于 $x$ 上方的点（事实上它是唯一的）。我们可以
找到一个公共有限生成域扩张 $k'''/k$，它同时扩张 $k'/k$ 和 $k''/k$，见
《域》引理 \ref{fields-lemma-common-extension-field} 及其证明。并且可以找到
一点 $x''' \in X_{k'''}$，它同时位于 $x'$ 与 $x''$ 上方。考虑局部环映射
$$
\mathcal{O}_{X_{k''}, x''} \longrightarrow \mathcal{O}_{X_{k'''}, x''''}.
$$
这是 Noether 局部环之间的平坦局部同态，因而忠实平坦。由 (2)，右侧局部环
正则。于是由《代数》引理 \ref{algebra-lemma-flat-under-regular}，可得
$\mathcal{O}_{X_{k''}, x''}$ 正则。由引理
\ref{varieties-lemma-geometrically-regular-at-point}，$X$ 在 $x$ 处几何正则。
\end{proof}

\noindent
下一个引理是《代数》引理
\ref{algebra-lemma-geometrically-regular-descent} 的几何版本。

\begin{lemma}
\label{varieties-lemma-flat-under-geometrically-regular}
设 $k$ 为域，设 $f : X \to Y$ 为 $k$ 上局部 Noether 概形之间的态射。
设 $x \in X$，并置 $y = f(x)$。如果 $X$ 在 $x$ 处几何正则，且 $f$
在 $x$ 处平坦，那么 $Y$ 在 $y$ 处几何正则。特别地，如果 $X$ 在 $k$
上几何正则，且 $f$ 平坦满射，那么 $Y$ 在 $k$ 上几何正则。
\end{lemma}

\begin{proof}
设 $k'$ 为 $k$ 的有限纯不可分扩张。设 $f' : X_{k'} \to Y_{k'}$ 为 $f$
的基变换，并设 $x' \in X_{k'}$ 为位于 $x$ 上方的唯一一点。如果证明
$Y_{k'}$ 在 $y' = f'(x')$ 处正则，那么由引理
\ref{varieties-lemma-geometrically-regular}，$Y$ 在 $k$ 上于 $y'$ 处几何正则。
由《态射》引理 \ref{morphisms-lemma-base-change-module-flat}，态射
$X_{k'} \to Y_{k'}$ 在 $x'$ 处平坦。因此环映射
$$
\mathcal{O}_{Y_{k'}, y'}
\longrightarrow
\mathcal{O}_{X_{k'}, x'}
$$
是局部 Noether 环之间的平坦局部同态，而其右侧根据假设正则。因此由
《代数》引理 \ref{algebra-lemma-flat-under-regular}，其左侧是正则局部环。
\end{proof}

\begin{lemma}
\label{varieties-lemma-geometrically-regular-smooth}
设 $k$ 为域，设 $X$ 为 $k$ 上局部有限型的概形，并设 $x \in X$。
那么，$X$ 在 $x$ 处几何正则，当且仅当 $X \to \Spec(k)$ 在 $x$ 处光滑
（《态射》定义 \ref{morphisms-definition-smooth}）。
\end{lemma}

\begin{proof}
问题在 $x$ 附近是局部的，因而可以假设 $X = \Spec(A)$，其中所用代数是
某个有限型 $k$-代数。设 $x$ 对应于素理想 $\mathfrak p$。

\medskip\noindent
如果 $A$ 在 $k$ 上于 $\mathfrak p$ 处光滑，那么可以局部化 $A$，并假设
$A$ 在 $k$ 上光滑。在这种情况下，$k' \otimes_k A$ 对所有扩域都在
$k'$ 上光滑，其中 $k'/k$ 任意；由《代数》引理
\ref{algebra-lemma-characterize-smooth-over-field}，这些 Noether 环都正则。

\medskip\noindent
假设 $X$ 在 $x$ 处几何正则。考虑剩余域
$K := \kappa(x) = \kappa(\mathfrak p)$，它属于点 $x$。这是 $k$ 的有限生成扩张。由
《代数》引理 \ref{algebra-lemma-make-separable}，存在有限纯不可分扩张
$k'/k$，使合成域 $k'K$ 是 $k'$ 的可分域扩张。设
$\mathfrak p' \subset A' = k' \otimes_k A$ 为位于 $\mathfrak p$ 上方的
素理想。由《代数》引理 \ref{algebra-lemma-p-ring-map}，它是位于
$\mathfrak p$ 上方的唯一素理想。因此剩余域
$K' := \kappa(\mathfrak p')$ 就是合成域 $k'K$。根据假设，局部环
$(A')_{\mathfrak p'}$ 正则。于是由《代数》引理
\ref{algebra-lemma-separable-smooth}，$k' \to A'$ 在 $\mathfrak p'$ 处光滑。
进而由《代数》引理 \ref{algebra-lemma-smooth-field-change-local}，
$k \to A$ 在 $\mathfrak p$ 处光滑。引理得证。
\end{proof}


\begin{example}
\label{varieties-example-geometrically-reduced-not-normal}
设 $k =\mathbf{F}_p(t)$。很容易给出一个正则但非几何约化的簇 $V$，它在
$k$ 上定义。例如可取 $\Spec(k[x]/(x^p - t))$。事实上，还存在正则且几何约化、
但甚至不几何正规的簇 $V$。具体地，当 $p > 2$ 时取
$V = \Spec(k[x, y]/(y^2 - x^p + t))$。由于多项式
$y^2 - x^p + t \in k[x, y]$ 不可约，这是一个簇。态射
$V \to \Spec(k)$ 在所有点处都光滑，唯独可能在点 $v_0 \in V$ 处不光滑；
该点对应于极大理想 $(y, x^p - t)$（因为 $2y$ 可逆）。特别地，
我们看到 $V$ 在所有点处都（几何）正则，唯独 $v_0$ 处尚有可能例外。
局部环
$$
\mathcal{O}_{V, v_0} = \left(k[x, y]/(y^2 - x^p + t)\right)_{(y, x^p - t)}
$$
是维数为 $1$ 的整环。它的极大理想由 $1$ 个元素生成，即由 $y$ 生成。
所以它是离散赋值环，从而正则。设 $k' = k[t^{1/p}]$。记
$t' = t^{1/p} \in k'$、$V' = V_{k'}$，并记 $v'_0 \in V'$ 为位于 $v_0$
上方的唯一一点。在 $k'$ 上，可写 $x^p - t = (x - t')^p$，但多项式
$y^2 - (x - t')^p$ 仍不可约，且 $V'$ 仍是簇。然而元素
$$
\frac{y}{x - t'} \in (\text{fraction field of }\mathcal{O}_{V', v'_0})
$$
在 $\mathcal{O}_{V', v'_0}$ 上整（只须计算它的平方），却不属于该环；
所以 $V'$ 在 $v'_0$ 处不正规。这就完成了这个例子。
\end{example}


\section{域变换与 Cohen--Macaulay 性}
\label{varieties-section-CM}

\noindent
下一个引理说明，定义几何 Cohen--Macaulay 概形没有意义，因为这样的概形
恰好就是 Cohen--Macaulay 概形。

\begin{lemma}
\label{varieties-lemma-CM-base-change}
设 $X$ 为域 $k$ 上的局部 Noether 概形。设 $k'/k$ 为有限生成域扩张。
设 $x \in X$，并设 $x' \in X_{k'}$ 为位于 $x$ 上方的点。那么
$$
\mathcal{O}_{X, x}\text{ 为 Cohen--Macaulay}
\Leftrightarrow
\mathcal{O}_{X_{k'}, x'}\text{ 为 Cohen--Macaulay}
$$
如果 $X$ 在 $k$ 上局部有限型，那么对任意域扩张 $k'/k$，同一结论成立。
\end{lemma}

\begin{proof}
引理的第一种情形来自《代数》引理
\ref{algebra-lemma-CM-geometrically-CM}。引理的第二种情形等价于《代数》引理
\ref{algebra-lemma-extend-field-CM-locus}。
\end{proof}


\section{域变换与 Jacobson 性}
\label{varieties-section-overfield}

\noindent
域上局部有限型的概形有充足的闭点，也就是说，它是 Jacobson 概形。此外，
这些剩余域都是基域的有限扩张。

\begin{lemma}
\label{varieties-lemma-locally-finite-type-Jacobson}
设 $X$ 为在 $k$ 上局部有限型的概形。那么：
\begin{enumerate}
\item 对任意闭点 $x \in X$，扩张 $\kappa(x)/k$ 是代数扩张；
\item $X$ 是 Jacobson 概形（《性质》定义
\ref{properties-definition-jacobson}）。
\end{enumerate}
\end{lemma}

\begin{proof}
一个概形是 Jacobson 概形，当且仅当它有一个由 Jacobson 概形组成的仿射
开覆盖，见《性质》引理 \ref{properties-lemma-locally-jacobson}。闭点处剩余域
的性质对 $X$ 也是局部的。因此可以假设 $X$ 仿射。在这种情况下，结论是
Hilbert 零点定理的推论，见《代数》定理
\ref{algebra-theorem-nullstellensatz}。它也可由《态射》引理
\ref{morphisms-lemma-jacobson-finite-type-points}、
\ref{morphisms-lemma-Jacobson-universally-Jacobson} 和
\ref{morphisms-lemma-ubiquity-Jacobson-schemes} 合并推出。
\end{proof}

\noindent
事实证明，如果 $X$ 不局部有限型，那么在作足够大的基域扩张后，可以得到
同样的结论。

\begin{lemma}
\label{varieties-lemma-make-Jacobson}
设 $X$ 为域 $k$ 上的概形。对任意基数足够大的域扩张 $K/k$，有：
\begin{enumerate}
\item 对任意闭点 $x \in X_K$，扩张 $\kappa(x)/K$ 是代数扩张；
\item $X_K$ 是 Jacobson 概形（《性质》定义
\ref{properties-definition-jacobson}）。
\end{enumerate}
\end{lemma}

\begin{proof}
选取仿射开覆盖 $X = \bigcup U_i$。由《代数》引理
\ref{algebra-lemma-base-change-Jacobson} 和《性质》引理
\ref{properties-lemma-affine-jacobson}，存在基数 $\kappa_i$，使得只要
$U_{i, K}$ 是在 $K$ 上考虑的，且 $\#(K) \geq \kappa_i$，它就具有所需性质。
置 $\kappa = \max\{\kappa_i\}$。那么，如果 $K$ 的基数大于 $\kappa$，
每个 $U_{i, K}$ 都满足引理的结论。因此，由《性质》引理
\ref{properties-lemma-locally-jacobson}，$X_K$ 是 Jacobson 概形。关于
$X_K$ 的闭点处剩余域的断言，来自 $U_{i, K}$ 的闭点处剩余域的相应断言。
\end{proof}


\section{域变换与丰沛可逆层}
\label{varieties-section-change-fields-ample}

\noindent
下述结果是本节各结果的典型例子。

\begin{lemma}
\label{varieties-lemma-ample-after-field-extension}
设 $k$ 为域，设 $X$ 为 $k$ 上的概形。如果 $X_K$ 对某个域扩张 $K/k$
有丰沛可逆层，那么 $X$ 有丰沛可逆层。
\end{lemma}

\begin{proof}
设 $K/k$ 为域扩张，使得 $X_K$ 有丰沛可逆层 $\mathcal{L}$。态射
$X_K \to X$ 满射。因此，$X$ 作为拟紧概形的像而拟紧（《性质》定义
\ref{properties-definition-ample}）。由于 $X_K$ 拟分离（由《性质》引理
\ref{properties-lemma-affine-s-opens-cover-quasi-separated}），可知 $X$
拟分离：如果 $U, V \subset X$ 是仿射开集，那么
$(U \cap V)_K = U_K \cap V_K$ 拟紧，而
$(U \cap V)_K \to U \cap V$ 满射。因此《概形》引理
\ref{schemes-lemma-characterize-quasi-separated} 适用。

\medskip\noindent
将 $K = \colim A_i$ 写成 $K$ 的各有限型 $k$-子代数的余极限。记
$X_i = X \times_{\Spec(k)} \Spec(A_i)$。由于 $X_K = \lim X_i$，可找到
一个 $i$ 和一个可逆层 $\mathcal{L}_i$，后者定义在 $X_i$ 上，使它在 $X_K$ 上的拉回
是 $\mathcal{L}$（《极限》引理 \ref{limits-lemma-descend-invertible-modules}；
这里以及下文都使用刚刚证明的 $X$ 拟紧且拟分离）。由《极限》引理
\ref{limits-lemma-limit-ample}，可以假设 $\mathcal{L}_i$ 丰沛，只需在必要时
增大 $i$。固定这样的 $i$，并设
$\mathfrak m \subset A_i$ 为极大理想。由 Hilbert 零点定理（《代数》定理
\ref{algebra-theorem-nullstellensatz}），剩余域
$k' = A_i/\mathfrak m$ 是 $k$ 的有限扩张。因此
$X_{k'} \subset X_i$ 是闭子概形，从而有丰沛可逆层（《性质》引理
\ref{properties-lemma-ample-on-closed}）。由于 $X_{k'} \to X$ 有限局部自由，
由《除子》命题 \ref{divisors-proposition-push-down-ample}，可得 $X$ 有丰沛
可逆层。
\end{proof}

\begin{lemma}
\label{varieties-lemma-quasi-affine-after-field-extension}
设 $k$ 为域，设 $X$ 为 $k$ 上的概形。如果 $X_K$ 对某个域扩张 $K/k$
拟仿射，那么 $X$ 拟仿射。
\end{lemma}

\begin{proof}
设 $K/k$ 为域扩张，使得 $X_K$ 拟仿射。态射 $X_K \to X$ 满射。因此，
$X$ 作为拟紧概形的像而拟紧（《性质》定义
\ref{properties-definition-quasi-affine}）。由于 $X_K$ 拟分离（因为它是
仿射概形的开子概形），可知 $X$ 拟分离：如果 $U, V \subset X$ 是仿射
开集，那么 $(U \cap V)_K = U_K \cap V_K$ 拟紧，而
$(U \cap V)_K \to U \cap V$ 满射。因此《概形》引理
\ref{schemes-lemma-characterize-quasi-separated} 适用。

\medskip\noindent
将 $K = \colim A_i$ 写成 $K$ 的各有限型 $k$-子代数的余极限。记
$X_i = X \times_{\Spec(k)} \Spec(A_i)$。由于 $X_K = \lim X_i$，可找到
一个 $i$，使 $X_i$ 拟仿射（《极限》引理
\ref{limits-lemma-limit-quasi-affine}；这里使用刚刚证明的 $X$ 拟紧且拟分离）。
由 Hilbert 零点定理（《代数》定理
\ref{algebra-theorem-nullstellensatz}），剩余域 $k' = A_i/\mathfrak m$ 是
$k$ 的有限扩张。因此 $X_{k'} \subset X_i$ 是闭子概形，从而拟仿射
（《性质》引理 \ref{properties-lemma-quasi-affine-locally-closed}）。由于
$X_{k'} \to X$ 有限局部自由，结论来自《除子》引理
\ref{divisors-lemma-push-down-quasi-affine}。
\end{proof}

\begin{lemma}
\label{varieties-lemma-quasi-projective-after-field-extension}
设 $k$ 为域，设 $X$ 为 $k$ 上的概形。如果 $X_K$ 在 $K$ 上拟射影，其中
$K/k$ 是某个域扩张，那么 $X$ 在 $k$ 上拟射影。
\end{lemma}

\begin{proof}
根据定义，概形态射 $g : Y \to T$ 拟射影，是指它局部有限型、拟紧，并且
存在一个相对于 $g$ 丰沛的可逆层，定义在 $Y$ 上。设 $K/k$ 为域扩张，
使得 $X_K$ 在 $K$ 上拟射影。设 $\Spec(A) \subset X$ 为仿射开集。那么
$U_K$ 是 $X_K$ 的仿射开子概形，因而 $A_K$ 是有限型 $K$-代数。由《代数》
引理 \ref{algebra-lemma-finite-type-descends}，$A$ 是有限型 $k$-代数。因此
$X \to \Spec(k)$ 局部有限型。由于 $X_K \to \Spec(K)$ 拟紧，可知
$X_K$ 拟紧，进而 $X$ 拟紧，进而 $X \to \Spec(k)$ 有限型。由《态射》引理
\ref{morphisms-lemma-finite-type-over-affine-ample-very-ample}，可知 $X_K$
有丰沛可逆层。再由引理 \ref{varieties-lemma-ample-after-field-extension}，$X$ 有丰沛
可逆层。因此，由《态射》引理
\ref{morphisms-lemma-finite-type-over-affine-ample-very-ample}，
$X \to \Spec(k)$ 拟射影。
\end{proof}

\noindent
下一个引理是《下降》引理
\ref{descent-lemma-descending-property-proper} 的特殊情形。

\begin{lemma}
\label{varieties-lemma-proper-after-field-extension}
设 $k$ 为域，设 $X$ 为 $k$ 上的概形。如果 $X_K$ 在 $K$ 上固有，其中
$K/k$ 是某个域扩张，那么 $X$ 在 $k$ 上固有。
\end{lemma}

\begin{proof}
设 $K/k$ 为域扩张，使得 $X_K$ 在 $K$ 上固有。回忆，这蕴含 $X_K$
分离且拟紧（《态射》定义 \ref{morphisms-definition-proper}）。态射
$X_K \to X$ 满射。因此，$X$ 作为拟紧概形的像而拟紧（《性质》定义
\ref{properties-definition-ample}）。由于 $X_K$ 分离，可知 $X$ 拟分离：
如果 $U, V \subset X$ 是仿射开集，那么
$(U \cap V)_K = U_K \cap V_K$ 拟紧，而
$(U \cap V)_K \to U \cap V$ 满射。因此《概形》引理
\ref{schemes-lemma-characterize-quasi-separated} 适用。

\medskip\noindent
将 $K = \colim A_i$ 写成 $K$ 的各有限型 $k$-子代数的余极限。记
$X_i = X \times_{\Spec(k)} \Spec(A_i)$。由《极限》引理
\ref{limits-lemma-eventually-proper}，存在一个 $i$，使得
$X_i \to \Spec(A_i)$ 固有。这里使用刚刚证明的 $X$ 拟紧且拟分离。选取极大
理想 $\mathfrak m \subset A_i$。由 Hilbert 零点定理（《代数》定理
\ref{algebra-theorem-nullstellensatz}），剩余域 $k' = A_i/\mathfrak m$ 是
$k$ 的有限扩张。基变换 $X_{k'} \to \Spec(k')$ 固有（《态射》引理
\ref{morphisms-lemma-base-change-proper}）。由于 $k'/k$ 有限，
$X_{k'} \to X$ 以及复合态射 $X_{k'} \to \Spec(k)$ 也都固有（《态射》引理
\ref{morphisms-lemma-finite-proper}、
\ref{morphisms-lemma-base-change-proper} 和
\ref{morphisms-lemma-composition-proper}）。第一个断言蕴含 $X$ 在 $k$ 上分离，
因为 $X_{k'}$ 分离（《态射》引理
\ref{morphisms-lemma-image-universally-closed-separated}）。第二个断言由《态射》
引理 \ref{morphisms-lemma-image-proper-is-proper} 蕴含
$X \to \Spec(k)$ 固有。
\end{proof}

\begin{lemma}
\label{varieties-lemma-projective-after-field-extension}
设 $k$ 为域，设 $X$ 为 $k$ 上的概形。如果 $X_K$ 在 $K$ 上射影，其中
$K/k$ 是某个域扩张，那么 $X$ 在 $k$ 上射影。
\end{lemma}

\begin{proof}
域 $k$ 上的概形在 $k$ 上射影，当且仅当它在 $k$ 上拟射影且固有。见
《态射》引理 \ref{morphisms-lemma-projective-is-quasi-projective-proper}。
因此，引理由引理 \ref{varieties-lemma-quasi-projective-after-field-extension} 和
\ref{varieties-lemma-proper-after-field-extension} 推出。
\end{proof}


\section{切空间}
\label{varieties-section-tangent-spaces}

\noindent
本节利用取值于对偶数的点，定义概形态射在源概形一点处的切空间。

\begin{definition}
\label{varieties-definition-dual-numbers}
对任意环 $R$，$R$ 上的{it 对偶数}是一个 $R$-代数，记作
$R[\epsilon]$。作为 $R$-模，它是自由的，基为 $1$、$\epsilon$；其
$R$-代数结构来自规定 $\epsilon^2 = 0$。
\end{definition}

\noindent
设 $f : X \to S$ 为概形态射。设 $x \in X$ 为一点，它的像为
$s = f(x)$，这是 $S$ 中的一点。考虑实线组成的交换图
\begin{equation}
\label{varieties-equation-tangent-space}
\vcenter{
\xymatrix{
\Spec(\kappa(x)) \ar[r] \ar[dr] \ar@/^1pc/[rr] &
\Spec(\kappa(x)[\epsilon]) \ar@{.>}[r] \ar[d]&
X \ar[d] \\
&
\Spec(\kappa(s)) \ar[r] &
S
}
}
\end{equation}
其中弯曲箭头是从 $\Spec(\kappa(x))$ 到 $X$ 的典范态射。

\begin{lemma}
\label{varieties-lemma-tangent-space}
使图 (\ref{varieties-equation-tangent-space}) 交换的虚线箭头之集合，具有典范的
$\kappa(x)$-向量空间结构。
\end{lemma}

\begin{proof}
置 $\kappa = \kappa(x)$。注意，在概形范畴中有推出
$$
\Spec(\kappa[\epsilon]) \amalg_{\Spec(\kappa)} \Spec(\kappa[\epsilon])
= \Spec(\kappa[\epsilon_1, \epsilon_2])
$$
其中 $\kappa[\epsilon_1, \epsilon_2]$ 是一个 $\kappa$-代数，以
$1, \epsilon_1, \epsilon_2$ 为基，并且
$\epsilon_1^2 = \epsilon_1\epsilon_2 = \epsilon_2^2 = 0$。这立即来自环的相应
结论，以及《概形》引理 \ref{schemes-lemma-morphism-from-spec-local-ring}
中从局部环的谱到概形之态射的描述。给定两个箭头
$\theta_1, \theta_2 : \Spec(\kappa[\epsilon]) \to X$，可以考虑态射
$$
\theta_1 + \theta_2 :
\Spec(\kappa[\epsilon]) \to
\Spec(\kappa[\epsilon_1, \epsilon_2])
\xrightarrow{\theta_1, \theta_2} X
$$
其中第一个箭头由 $\epsilon_i \mapsto \epsilon$ 给出。另一方面，给定
$\lambda \in \kappa$，存在 $\Spec(\kappa[\epsilon])$ 的一个自映射，它对应于
$\kappa$-代数 $\kappa[\epsilon]$ 的自同态，而该自同态将 $\epsilon$ 映为
$\lambda \epsilon$。用这个自映射预合成
$\theta : \Spec(\kappa[\epsilon]) \to X$，便得到 $\lambda \theta$。读者可以
通过验证适当概形交换图的存在性来验证向量空间公理。细节从略。（另一种证明
是把一切都用局部环表述，再在环映射层面验证向量空间公理。）
\end{proof}

\begin{definition}
\label{varieties-definition-tangent-space}
设 $f : X \to S$ 为概形态射，并设 $x \in X$。使图
(\ref{varieties-equation-tangent-space}) 交换的虚线箭头之集合，连同其典范的
$\kappa(x)$-向量空间结构，称为{it $X$ 在 $S$ 上于 $x$ 处的切空间}，
记作 $T_{X/S, x}$。这个空间的元素称为 $X/S$ 在 $x$ 处的{it 切向量}。
\end{definition}

\noindent
由于 $x \in X$ 处的切向量存在于概形论纤维 $X_s$ 中，而这个纤维来自
$f : X \to S$ 且位于 $s = f(x)$ 上，我们得到典范等同
\begin{equation}
\label{varieties-equation-tangent-space-fibre}
T_{X/S, x} = T_{X_s/s, x}
\end{equation}
这个涉及点函子的优美定义有如下代数描述；它提示我们把{it $X$ 在 $S$ 上于
$x$ 处的余切空间}定义为 $\kappa(x)$-向量空间
$$
T^*_{X/S, x} = \Omega_{X/S, x} \otimes_{\mathcal{O}_{X, x}} \kappa(x)
$$
其理由就在于，它典范地是 $\kappa(x)$ 上的对偶，而所对偶的是 $X$ 在 $S$
上于 $x$ 处的切空间。

\begin{lemma}
\label{varieties-lemma-tangent-space-cotangent-space}
设 $f : X \to S$ 为概形态射，并设 $x \in X$。存在典范同构
$$
T_{X/S, x} = \Hom_{\mathcal{O}_{X, x}}(\Omega_{X/S, x}, \kappa(x))
$$
它是 $\kappa(x)$ 上的向量空间同构。
\end{lemma}

\begin{proof}
置 $\kappa = \kappa(x)$。给定 $\theta \in T_{X/S, x}$，得到映射
$$
\theta^*\Omega_{X/S} \to
\Omega_{\Spec(\kappa[\epsilon])/\Spec(\kappa(s))} \to
\Omega_{\Spec(\kappa[\epsilon])/\Spec(\kappa)}
$$
取截面得到一个 $\mathcal{O}_{X, x}$-线性映射
$\xi_\theta : \Omega_{X/S, x} \to \kappa \text{d}\epsilon$，也就是引理公式
右侧的一个元素。为了证明 $\theta \mapsto \xi_\theta$ 是同构，可以把 $S$
替换为 $s$，并把 $X$ 替换为概形论纤维 $X_s$。事实上，公式两侧都只依赖于
概形论纤维；对 $T_{X/S, x}$ 这是清楚的，右侧则见《态射》引理
\ref{morphisms-lemma-base-change-differentials}。还可以把 $X$ 替换为
$\mathcal{O}_{X, x}$ 的谱，因为这既不改变 $T_{X/S, x}$（《概形》
引理 \ref{schemes-lemma-morphism-from-spec-local-ring}），也不改变
$\Omega_{X/S, x}$（《模》引理
\ref{modules-lemma-stalk-module-differentials}）。

\medskip\noindent
设 $(A, \mathfrak m, \kappa)$ 为域 $k$ 上的局部环。为完成证明，需要说明
任意 $A$-线性映射 $\xi : \Omega_{A/k} \to \kappa$ 都来自唯一的 $k$-代数映射
$\varphi : A \to \kappa[\epsilon]$，且后者与典范映射
$c : A \to \kappa$ 在模 $\epsilon$ 后一致。写成
$\varphi(a) = c(a) + D(a) \epsilon$，读者可见
$a \mapsto D(a)$ 是 $k$-导子。利用 $\Omega_{A/k}$ 的泛性质，可知每个
$D$ 对应唯一的 $\xi$，反之亦然。证明完成。
\end{proof}

\begin{lemma}
\label{varieties-lemma-tangent-space-rational-point}
设 $f : X \to S$ 为概形态射。设 $x \in X$，并设 $s = f(x) \in S$。
假设 $\kappa(x) = \kappa(s)$。那么存在典范同构
$$
\mathfrak m_x/(\mathfrak m_x^2 + \mathfrak m_s\mathcal{O}_{X, x})
=
\Omega_{X/S, x} \otimes_{\mathcal{O}_{X, x}} \kappa(x)
$$
和
$$
T_{X/S, x} =
\Hom_{\kappa(x)}(
\mathfrak m_x/(\mathfrak m_x^2 + \mathfrak m_s\mathcal{O}_{X, x}),
\kappa(x))
$$
更一般地，当 $\kappa(x)/\kappa(s)$ 是可分代数扩张时，结论仍成立。
\end{lemma}

\begin{proof}
第二个同构由第一个同构和引理
\ref{varieties-lemma-tangent-space-cotangent-space} 推出。对于第一个同构，可以把 $S$
替换为 $s$，并把 $X$ 替换为 $X_s$，见《态射》引理
\ref{morphisms-lemma-base-change-differentials}。还可以把 $X$ 替换为
$\mathcal{O}_{X, x}$ 的谱，见《模》引理
\ref{modules-lemma-stalk-module-differentials}。因此，只须证明如下代数事实：设
$(A, \mathfrak m, \kappa)$ 为域 $k$ 上的局部环，并且 $\kappa/k$ 是可分
代数扩张。那么典范映射
$$
\mathfrak m/\mathfrak m^2
\longrightarrow
\Omega_{A/k} \otimes \kappa
$$
是同构。注意，$\mathfrak m/\mathfrak m^2 = H_1(\NL_{\kappa/A})$。由《代数》
引理 \ref{algebra-lemma-exact-sequence-NL}，只须证明
$\Omega_{\kappa/k} = 0$ 和 $H_1(\NL_{\kappa/k}) = 0$。由于 $\kappa$ 是
$k$ 中各有限可分扩张的并，只须证明 $\kappa$ 是 $k$ 的有限可分扩张时的情形
（《代数》引理 \ref{algebra-lemma-colimits-NL}）。在这种情况下，环映射
$k \to \kappa$ 是 \'etale 的，因而 $\NL_{\kappa/k} = 0$（这基本上就是定义，
见《代数》第 \ref{algebra-section-etale} 节）。
\end{proof}

\begin{lemma}
\label{varieties-lemma-map-tangent-spaces}
设 $f : X \to Y$ 为基概形 $S$ 上的概形态射。设 $x \in X$，并置
$y = f(x)$。如果 $\kappa(y) = \kappa(x)$，那么 $f$ 诱导自然线性映射
$$
\text{d}f : T_{X/S, x} \longrightarrow T_{Y/S, y}
$$
经由引理 \ref{varieties-lemma-tangent-space-cotangent-space} 的等同，它对偶于线性映射
$\Omega_{Y/S, y} \otimes \kappa(y) \to \Omega_{X/S, x} \otimes \kappa(x)$。
\end{lemma}

\begin{proof}
从略。
\end{proof}

\begin{lemma}
\label{varieties-lemma-tangent-space-product}
设 $X$、$Y$ 为基 $S$ 上的概形。设 $x \in X$ 与 $y \in Y$ 有同一
像点 $s \in S$，且 $\kappa(s) = \kappa(x)$ 和
$\kappa(s) = \kappa(y)$。存在典范同构
$$
T_{X \times_S Y/S, (x, y)} = T_{X/S, x} \oplus T_{Y/S, y}
$$
从左到右的映射由投影 $X \times_S Y \to X$ 与
$X \times_S Y \to Y$ 在切空间上诱导的映射给出。从右到左的映射由
$1 \times y : X_s \to X_s \times_s Y_s$ 与
$x \times 1 : Y_s \to X_s \times_s Y_s$ 给出，其中使用等同
(\ref{varieties-equation-tangent-space-fibre}) 将切空间等同于纤维的切空间。
\end{lemma}

\begin{proof}
直和分解来自《态射》引理 \ref{morphisms-lemma-differential-product}，其中
使用引理 \ref{varieties-lemma-tangent-space-rational-point}。它与这些映射的相容性来自引理
\ref{varieties-lemma-map-tangent-spaces}。
\end{proof}

\begin{lemma}
\label{varieties-lemma-injective-tangent-spaces-unramified}
设 $f : X \to Y$ 为基概形 $S$ 上局部有限型的概形态射。设 $x \in X$，
置 $y = f(x)$，并假设 $\kappa(y) = \kappa(x)$。那么下列条件等价：
\begin{enumerate}
\item $\text{d}f : T_{X/S, x} \longrightarrow T_{Y/S, y}$ 是单射；
\item $f$ 在 $x$ 处非分歧。
\end{enumerate}
\end{lemma}

\begin{proof}
由《态射》引理 \ref{morphisms-lemma-permanence-finite-type}，态射 $f$ 局部
有限型。映射 $\text{d}f$ 是单射，当且仅当
$\Omega_{Y/S, y} \otimes \kappa(y) \to \Omega_{X/S, x} \otimes \kappa(x)$
是满射（引理 \ref{varieties-lemma-map-tangent-spaces}）。正合列
$f^*\Omega_{Y/S} \to \Omega_{X/S} \to \Omega_{X/Y} \to 0$
（《态射》引理 \ref{morphisms-lemma-triangle-differentials}）于是说明，发生这种
情况当且仅当 $\Omega_{X/Y, x} \otimes \kappa(x) = 0$。因此结论来自《态射》
引理 \ref{morphisms-lemma-unramified-at-point}。
\end{proof}


\section{一般有限态射}
\label{varieties-section-generically-finite}

\noindent
本节重新讨论《态射》第 \ref{morphisms-section-generically-finite} 节研究过的
概形一般有限态射。

\begin{lemma}
\label{varieties-lemma-quasi-finite-in-codim-1}
设 $f : X \to Y$ 局部有限型。设 $y \in Y$ 为一点，使
$\mathcal{O}_{Y, y}$ 是维数 $\leq 1$ 的 Noether 环。再假设下列条件之一成立：
\begin{enumerate}
\item 对每个泛点 $\eta$，若它属于 $X$ 的某个不可约分支，则域扩张
$\kappa(\eta)/\kappa(f(\eta))$ 有限（或为代数扩张）；
\item 对每个泛点 $\eta$，若它属于 $X$ 的某个不可约分支，且
$f(\eta) \leadsto y$，域扩张 $\kappa(\eta)/\kappa(f(\eta))$ 就有限
（或为代数扩张）；
\item $f$ 在 $X$ 的不可约分支的每个泛点处拟有限；
\item $Y$ 局部 Noether，且 $f$ 在 $X$ 的一个稠密点集上拟有限；
\item 此处添加更多条件。
\end{enumerate}
那么，$f$ 在 $X$ 中位于 $y$ 上方的每一点处都拟有限。
\end{lemma}

\begin{proof}
条件 (4) 蕴含 $X$ 局部 Noether（《态射》引理
\ref{morphisms-lemma-finite-type-noetherian}）。态射拟有限的点集是开集
（《态射》引理 \ref{morphisms-lemma-quasi-finite-points-open}）。局部
Noether 概形的稠密开集包含不可约分支的所有泛点，因而 (4) 蕴含 (3)。
由《态射》引理 \ref{morphisms-lemma-residue-field-quasi-finite}，条件 (3)
蕴含条件 (1)。条件 (1) 蕴含条件 (2)。所以只须在 (2) 成立时证明引理。

\medskip\noindent
假设 (2) 成立。回忆，$\Spec(\mathcal{O}_{Y, y})$ 是 $Y$ 中特化到 $y$
的点集，见《概形》引理 \ref{schemes-lemma-specialize-points}。结合《态射》引理
\ref{morphisms-lemma-base-change-quasi-finite}，这说明可以把 $Y$ 替换为
$\Spec(\mathcal{O}_{Y, y})$。因此可以假设 $Y = \Spec(B)$，其中 $B$ 是
维数 $\leq 1$ 的 Noether 局部环，而 $y$ 是闭点。

\medskip\noindent
设 $X = \bigcup X_i$ 是 $X$ 的不可约分支，并把它们看作约化闭子概形。
如果能证明每个纤维 $X_{i, y}$ 都是离散空间，那么
$X_y = \bigcup X_{i, y}$ 也离散；由《态射》引理
\ref{morphisms-lemma-quasi-finite-at-point-characterize}，可得 $X \to Y$
在 $X_y$ 的所有点处拟有限。因此可以假设 $X$ 是整概形。

\medskip\noindent
如果 $X \to Y$ 把泛点 $\eta$（它是 $X$ 的泛点）映到 $y$，那么 $X$ 是
$\kappa(y)$ 的一个有限扩张的谱，结论成立。假设 $X$ 把 $\eta$ 映到
一个极小素理想 $\mathfrak q$（它属于 $B$）所对应的点，并且该点不同于
$\mathfrak m_B$。我们得到分解
$X \to \Spec(B/\mathfrak q) \to \Spec(B)$。设 $x \in X$ 为位于 $y$
上方的点。由维数公式（《态射》引理
\ref{morphisms-lemma-dimension-formula}），有
$$
\dim(\mathcal{O}_{X, x}) \leq \dim(B/\mathfrak q) +
\text{trdeg}_{\kappa(\mathfrak q)}(R(X)) - \text{trdeg}_{\kappa(y)} \kappa(x)
$$
我们知道 $\dim(B/\mathfrak q) = 1$；$X$ 的泛点不等于 $x$ 且特化到 $x$；
并且 $R(X)$ 是 $\kappa(\mathfrak q)$ 的代数扩张。因此得到
$$
1 \leq 1 - \text{trdeg}_{\kappa(y)} \kappa(x)
$$
于是，每一点 $x$（属于 $X_y$）都在 $X_y$ 中闭，这是由《态射》引理
\ref{morphisms-lemma-algebraic-residue-field-extension-closed-point-fibre} 得到的；
再由《态射》引理 \ref{morphisms-lemma-quasi-finite-at-point-characterize}，
$X \to Y$ 在每一点 $x$（属于 $X_y$）处拟有限（该引理也蕴含 $X_y$
是离散拓扑空间）。
\end{proof}

\begin{lemma}
\label{varieties-lemma-finite-in-codim-1}
设 $f : X \to Y$ 为固有态射。设 $y \in Y$ 为一点，使
$\mathcal{O}_{Y, y}$ 是维数 $\leq 1$ 的 Noether 环。再假设下列条件之一成立：
\begin{enumerate}
\item 对每个泛点 $\eta$，若它属于 $X$ 的某个不可约分支，则域扩张
$\kappa(\eta)/\kappa(f(\eta))$ 有限（或为代数扩张）；
\item 对每个泛点 $\eta$，若它属于 $X$ 的某个不可约分支，且
$f(\eta) \leadsto y$，域扩张 $\kappa(\eta)/\kappa(f(\eta))$ 就有限
（或为代数扩张）；
\item $f$ 在 $X$ 的每个泛点处拟有限；
\item $Y$ 局部 Noether，且 $f$ 在 $X$ 的一个稠密点集上拟有限；
\item 此处添加更多条件。
\end{enumerate}
那么，存在开邻域 $V \subset Y$，它包含 $y$，并使得
$f^{-1}(V) \to V$ 有限。
\end{lemma}

\begin{proof}
由引理 \ref{varieties-lemma-quasi-finite-in-codim-1}，态射 $f$ 在纤维 $X_y$ 的每一点
处拟有限。因此 $X_y$ 是离散拓扑空间（《态射》引理
\ref{morphisms-lemma-quasi-finite-at-point-characterize}）。由于 $f$ 固有，纤维
$X_y$ 拟紧，也就是说，它是有限集。因此可以应用《概形上同调》引理
\ref{coherent-lemma-proper-finite-fibre-finite-in-neighbourhood} 得出结论。
\end{proof}

\begin{lemma}
\label{varieties-lemma-modification-normal-iso-over-codimension-1}
设 $X$ 为 Noether 概形。设 $f : Y \to X$ 为双有理固有态射，且 $Y$ 约化。
设 $U \subset X$ 为 $f$ 在其上为同构的最大开集。那么 $U$ 包含：
\begin{enumerate}
\item 每个余维为 $0$ 的点，只要它属于 $X$；
\item 每个 $x \in X$，只要它余维为 $1$（这是在 $X$ 上计算的），且 $\mathcal{O}_{X, x}$
是离散赋值环；
\item 每个 $x \in X$，只要 $Y \to X$ 在 $x$ 上的纤维有限，且
$\mathcal{O}_{X, x}$ 正规；
\item 每个 $x \in X$，只要 $f$ 在某个 $y \in Y$ 处拟有限，该点
位于 $x$ 上方，并且 $\mathcal{O}_{X, x}$ 正规。
\end{enumerate}
\end{lemma}

\begin{proof}
(1) 来自《态射》引理
\ref{morphisms-lemma-birational-isomorphism-over-dense-open}。(2) 来自 (3)
以及引理 \ref{varieties-lemma-finite-in-codim-1}（并使用有限态射有有限纤维这一事实）。

\medskip\noindent
(3) 来自 (4) 和《态射》引理 \ref{morphisms-lemma-finite-fibre}，但我们也给出
一个直接证明。设 $x \in X$ 如 (3) 所述。由《概形上同调》引理
\ref{coherent-lemma-proper-finite-fibre-finite-in-neighbourhood}，可以假设 $f$
有限。可以假设 $X$ 仿射。这把问题归约为 Noether 仿射概形之间的有限双有理
态射 $Y \to X$ 的情形，其中 $x \in X$ 且 $\mathcal{O}_{X, x}$ 是正规整环。
由于 $\mathcal{O}_{X, x}$ 是整环且 $X$ Noether，可以把 $X$ 替换为 $x$
的某个仿射开集，使之整。于是，由于 $Y \to X$ 双有理且 $Y$ 约化，可知
$Y$ 整。写成 $X = \Spec(A)$ 和 $Y = \Spec(B)$，可知 $A \subset B$ 是
分式域相同的整环之间的有限包含。如果 $\mathfrak p \subset A$ 是对应于
$x$ 的素理想，那么 $A_\mathfrak p$ 正规蕴含
$A_\mathfrak p \subset B_\mathfrak p$ 为等式。由于 $B$ 是有限 $A$-模，
可知存在 $a \in A$、$a \not \in \mathfrak p$，使 $A_a \to B_a$ 是同构。

\medskip\noindent
设 $x \in X$ 与 $y \in Y$ 如 (4) 所述。用仿射开邻域替换 $X$ 后，
可以假设 $X = \Spec(A)$ 且 $A \subset \mathcal{O}_{X, x}$，见《性质》引理
\ref{properties-lemma-ring-affine-open-injective-local-ring}。于是 $A$ 是整环，
从而 $X$ 整。由于 $f$ 双有理且 $Y$ 约化，$Y$ 也整。考虑环映射
$\mathcal{O}_{X, x} \to \mathcal{O}_{Y, y}$。这是本质有限型的环映射，剩余
域扩张有限，并且
$\dim(\mathcal{O}_{Y, y}/\mathfrak m_x\mathcal{O}_{Y, y}) = 0$（要看出这一点，
可逐项追溯《态射》定义 \ref{morphisms-definition-quasi-finite} 与《代数》定义
\ref{algebra-definition-quasi-finite} 中拟有限映射的定义）。由《代数》引理
\ref{algebra-lemma-essentially-finite-type-fibre-dim-zero}，
$\mathcal{O}_{Y, y}$ 是某个有限 $\mathcal{O}_{X, x}$-代数 $B$ 的局部化。
当然，可以用 $B$ 的像替换 $B$，该像取在 $\mathcal{O}_{Y, y}$ 中；并假设 $B$
是与 $\mathcal{O}_{Y, y}$ 有相同分式域的整环。那么
$\mathcal{O}_{X, x} \subset B$ 与之有相同分式域，因为 $f$ 双有理。由于
$\mathcal{O}_{X, x}$ 正规，可得 $\mathcal{O}_{X, x} = B$（因为有限蕴含整），
特别地，$\mathcal{O}_{X, x} = \mathcal{O}_{Y, y}$。由《态射》引理
\ref{morphisms-lemma-morphism-defined-local-ring}，缩小 $X$ 后，可以假设存在
截面 $X \to Y$，它是 $f$ 的截面，把 $x$ 映到 $y$，并在局部环上诱导给定同构。
由于 $X \to Y$ 是闭的（由《概形》引理
\ref{schemes-lemma-section-immersion}），并且必然把 $X$ 的泛点映到 $Y$
的泛点，可知 $X \to Y$ 的像就是 $Y$。于是 $Y = X$，所需结论得证。
\end{proof}


\section{Noether 正规化的变式}
\label{varieties-section-noether-normalization}

\noindent
Noether 正规化定理断言：若 $k$ 是域，$A$ 是有限型 $k$-代数，且维数为
$d$，那么存在有限单射 $k$-代数同态
$k[x_1, \ldots, x_d] \to A$。见《代数》引理
\ref{algebra-lemma-Noether-normalization}。在几何上，这意味着存在有限满态射
$\Spec(A) \to \mathbf{A}^d_k$，它定义在 $\Spec(k)$ 上。

\begin{lemma}
\label{varieties-lemma-noether-normalization}
设 $f : X \to S$ 为概形态射。设 $x \in X$，其像为 $s \in S$。设
$V \subset S$ 是 $s$ 的仿射开邻域。若 $f$ 局部有限型且
$\dim_x(X_s) = d$，那么存在仿射开集 $U \subset X$，使得
$x \in U$ 且 $f(U) \subset V$，并且 $f|_U : U \to V$ 有分解
$$
U \xrightarrow{\pi} \mathbf{A}^d_V \to V
$$
其中 $\pi$ 拟有限。
\end{lemma}

\begin{proof}
这来自《代数》引理
\ref{algebra-lemma-quasi-finite-over-polynomial-algebra}。
\end{proof}

\begin{lemma}
\label{varieties-lemma-noether-normalization-affine}
设 $f : X \to S$ 为仿射概形之间的有限型态射。设 $s \in S$。若
$\dim(X_s) = d$，那么 $f$ 存在分解
$$
X \xrightarrow{\pi} \mathbf{A}^d_S \to S
$$
使得态射 $\pi_s : X_s \to \mathbf{A}^d_{\kappa(s)}$（即在 $s$ 上的纤维态射）有限。
\end{lemma}

\begin{proof}
写成 $S = \Spec(A)$ 与 $X = \Spec(B)$，并设 $A \to B$ 是对应于
$f$ 的环映射。设 $\mathfrak p \subset A$ 是对应于 $s$ 的素理想。可以选取满射
$A[x_1, \ldots, x_r] \to B$。由《代数》引理
\ref{algebra-lemma-Noether-normalization}，存在元
$y_1, \ldots, y_d \in A$，它们位于一个 $\mathbf{Z}$-子代数中；该子代数是在
$A$ 中由 $x_1, \ldots, x_r$ 生成的，并且使 $A$-代数同态
$A[t_1, \ldots, t_d] \to B$（把 $t_i$ 送到 $y_i$）诱导有限
$\kappa(\mathfrak p)$-代数同态
$\kappa(\mathfrak p)[t_1, \ldots, t_d] \to B \otimes_A \kappa(\mathfrak p)$。
这证明了引理。
\end{proof}

\begin{lemma}
\label{varieties-lemma-geometric-structure-unramified}
设 $f : X \to S$ 为概形态射。设 $x \in X$。设
$V = \Spec(A)$ 是 $f(x)$ 的仿射开邻域，并且位于 $S$ 中。若 $f$ 在 $x$
处非分歧，那么存在仿射开集 $U \subset X$，使 $x \in U$ 且
$f(U) \subset V$，并且存在交换图
$$
\xymatrix{
X \ar[d] & U \ar[l] \ar[rd] \ar[r]^-j &
\Spec(A[t]_{g'}/(g)) \ar[d] \ar[r] &
\Spec(A[t]) = \mathbf{A}^1_V \ar[ld] \\
Y & & V \ar[ll]
}
$$
其中 $j$ 是浸入，$g \in A[t]$ 是首一多项式，而 $g'$ 是 $g$ 对 $t$
的导数。若 $f$ 在 $x$ 处是 \'etale 的，那么可选取该图，使 $j$
是开浸入。
\end{lemma}

\begin{proof}
非分歧情形是《代数》命题
\ref{algebra-proposition-unramified-locally-standard} 的几何翻译。在 \'etale
情形，这是《代数》命题
\ref{algebra-proposition-etale-locally-standard} 的几何翻译；或者等价地，它来自
《态射》引理 \ref{morphisms-lemma-etale-locally-standard-etale}，尽管两处的表述略有不同。
\end{proof}

\begin{lemma}
\label{varieties-lemma-unramfied-over-affine}
设 $f : X \to S$ 为仿射概形之间的有限型态射。设 $x \in X$，其像为
$s \in S$。设
$$
r =
\dim_{\kappa(x)} \Omega_{X/S, x} \otimes_{\mathcal{O}_{X, x}} \kappa(x) =
\dim_{\kappa(x)} \Omega_{X_s/s, x} \otimes_{\mathcal{O}_{X_s, x}} \kappa(x) =
\dim_{\kappa(x)} T_{X/S, x}
$$
那么 $f$ 存在分解
$$
X \xrightarrow{\pi} \mathbf{A}^r_S \to S
$$
使得 $\pi$ 在 $x$ 处非分歧。
\end{lemma}

\begin{proof}
由《态射》引理 \ref{morphisms-lemma-finite-type-differentials}，第一个维数是有限的。
第一个等式来自《态射》引理
\ref{morphisms-lemma-base-change-differentials}：$\Omega_{X/S}$ 到纤维上的限制就是微分模。
最后一个等式来自引理 \ref{varieties-lemma-tangent-space-cotangent-space}。因此该表述有意义。

\medskip\noindent
为证明引理，写成 $S = \Spec(A)$ 与 $X = \Spec(B)$，并设 $A \to B$
是对应于 $f$ 的环映射。设 $\mathfrak q \subset B$ 是对应于 $x$ 的素理想。
选取 $A$-代数的满射 $A[x_1, \ldots, x_t] \to B$。由于
$\Omega_{B/A}$ 由 $\text{d}x_1, \ldots, \text{d}x_t$ 生成，它们在
$\Omega_{X/S, x} \otimes_{\mathcal{O}_{X, x}} \kappa(x)$ 中的像生成该
$\kappa(x)$-向量空间。重新编号后，可以假设
$\text{d}x_1, \ldots, \text{d}x_r$ 的像构成
$\Omega_{X/S, x} \otimes_{\mathcal{O}_{X, x}} \kappa(x)$ 的一组基。我们断言
$P = A[x_1, \ldots, x_r] \to B$ 在 $\mathfrak q$ 处非分歧。为此，只须证明
$\Omega_{B/P, \mathfrak q} = 0$（《代数》引理 \ref{algebra-lemma-unramified}）。
注意，$\Omega_{B/P}$ 是 $\Omega_{B/A}$ 对由
$\text{d}x_1, \ldots, \text{d}x_r$ 生成的子模之商。因此有
$\Omega_{B/P, \mathfrak q} \otimes_{B_\mathfrak q} \kappa(\mathfrak q) = 0$；这是由我们对
$x_1, \ldots, x_r$ 的选取得到的。
由 Nakayama 引理，更精确地说，由《代数》引理 \ref{algebra-lemma-NAK}
的 (2)（它可用，因为 $\Omega_{B/P}$ 是有限模，见上面的引用），可得
$\Omega_{B/P, \mathfrak q} = 0$。
\end{proof}

\begin{lemma}
\label{varieties-lemma-immersion-into-affine}
设 $f : X \to S$ 为概形态射。设 $x \in X$，其像为 $s \in S$。设
$V \subset S$ 是 $s$ 的仿射开邻域。若 $f$ 局部有限型且
$$
r =
\dim_{\kappa(x)} \Omega_{X/S, x} \otimes_{\mathcal{O}_{X, x}} \kappa(x) =
\dim_{\kappa(x)} \Omega_{X_s/s, x} \otimes_{\mathcal{O}_{X_s, x}} \kappa(x) =
\dim_{\kappa(x)} T_{X/S, x}
$$
那么存在下列两种情形：
\begin{enumerate}
\item 存在仿射开集 $U \subset X$，使 $x \in U$ 且 $f(U) \subset V$，并且
$f|_U$ 存在分解
$$
U \xrightarrow{j} \mathbf{A}^{r + 1}_V \to V
$$
其中 $j$ 是浸入；或者
\item 存在仿射开集 $U \subset X$，使 $x \in U$ 且 $f(U) \subset V$，并且
$f|_U$ 存在分解
$$
U \xrightarrow{j} D \to V
$$
其中 $j$ 是闭浸入，而 $D \to V$ 光滑且相对维数为 $r$。
\end{enumerate}
\end{lemma}

\begin{proof}
任取仿射开集 $U \subset X$，使 $x \in U$ 且 $f(U) \subset V$。对
$U \to V$ 应用引理 \ref{varieties-lemma-unramfied-over-affine}，得到该引理中的
$U \to \mathbf{A}^r_V \to V$。由引理 \ref{varieties-lemma-geometric-structure-unramified}，
得到分解
$$
U \xrightarrow{j} D \xrightarrow{j'} \mathbf{A}^{r + 1}_V
\xrightarrow{p} \mathbf{A}^r_V \to V
$$
其中 $j$ 与 $j'$ 是浸入，$p$ 是投影，而
$p \circ j'$ 是标准 \'etale 态射。特别地，这说明 (1) 与 (2) 成立。
\end{proof}


\section{纤维的维数}
\label{varieties-section-dimension-fibres}

\noindent
我们已经看到，有限型态射的纤维维数通常会跳升。本节讨论余维
$1$ 中不会发生这种现象，并更一般地讨论纤维维数能跳升多少。
下面列出一些相关结果：
\begin{enumerate}
\item 对有限型态射 $X \to S$，若点 $x \in X$ 满足
$\dim_x(X_{f(x)}) \leq d$，则所有这类点组成开集，见《代数》引理
\ref{algebra-lemma-dimension-fibres-bounded-open-upstairs} 和《态射》引理
\ref{morphisms-lemma-openness-bounded-dimension-fibres}。
\item 我们有维数公式，见《代数》引理
\ref{algebra-lemma-dimension-formula} 和《态射》引理
\ref{morphisms-lemma-dimension-formula}。
\item 对支配赋值环的整有限型概形，纤维维数为常数，见《代数》引理
\ref{algebra-lemma-finite-type-domain-over-valuation-ring-dim-fibres}。
\item 若 $X \to S$ 有限型，并且在 $X$ 的每个泛点处拟有限，那么
$X \to S$ 在余维 $1$ 中拟有限，见《代数》引理
\ref{algebra-lemma-finite-in-codim-1} 和引理
\ref{varieties-lemma-quasi-finite-in-codim-1}。
\end{enumerate}
上述最后一个结果可如下推广。

\begin{lemma}
\label{varieties-lemma-dimension-fibre-in-codim-1}
设 $f : X \to Y$ 局部有限型。设 $x \in X$ 是一点，其像为
$y \in Y$，并且 $\mathcal{O}_{Y, y}$ 是维数 $\leq 1$ 的 Noether 环。设
$d \geq 0$ 是整数，使得对每个泛点 $\eta$，若它属于 $X$ 的一个包含
$x$ 的不可约分支，则 $\dim_\eta(X_{f(\eta)}) = d$。那么
$\dim_x(X_y) = d$。
\end{lemma}

\begin{proof}
回忆，$\Spec(\mathcal{O}_{Y, y})$ 是 $Y$ 中特化到 $y$ 的点集，见《概形》引理
\ref{schemes-lemma-specialize-points}。因此可以替换 $Y$，把它取为
$\Spec(\mathcal{O}_{Y, y})$，并假设 $Y = \Spec(B)$，其中 $B$ 是维数 $\leq 1$ 的
Noether 局部环，而 $y$ 是闭点。还可以替换 $X$，把它取为 $x$ 的一个仿射邻域。

\medskip\noindent
设 $X = \bigcup X_i$ 是 $X$ 的不可约分支，并把它们视为约化闭子概形。
若能证明每个纤维 $X_{i, y}$ 的维数都为 $d$，那么
$X_y = \bigcup X_{i, y}$ 的维数也为 $d$。因此可以假设 $X$ 是整概形。

\medskip\noindent
若 $X \to Y$ 把泛点 $\eta$（它是 $X$ 的泛点）映到 $y$，那么 $X$ 是
$\kappa(y)$ 上的概形，结论由假设立即成立。假设 $X$ 把 $\eta$
映到点 $\xi \in Y$，它对应于极小素理想 $\mathfrak q$（它属于 $B$），
且它不同于 $\mathfrak m_B$。我们得到分解
$X \to \Spec(B/\mathfrak q) \to \Spec(B)$。由维数公式（《态射》引理
\ref{morphisms-lemma-dimension-formula}），有
$$
\dim(\mathcal{O}_{X, x}) + \text{trdeg}_{\kappa(y)} \kappa(x) \leq
\dim(B/\mathfrak q) + \text{trdeg}_{\kappa(\mathfrak q)}(R(X))
$$
我们有 $\dim(B/\mathfrak q) = 1$。又有
$\text{trdeg}_{\kappa(\mathfrak q)}(R(X)) = d$，这来自假设
$\dim_\eta(X_\xi) = d$，见《态射》引理
\ref{morphisms-lemma-dimension-fibre-at-a-point}。由于
$\mathcal{O}_{X, x} \to \mathcal{O}_{X_s, x}$ 有核（因为
$\eta \mapsto \xi \not = y$），并且 $\mathcal{O}_{X, x}$ 是 Noether 整环，可知
$\dim(\mathcal{O}_{X, x}) > \dim(\mathcal{O}_{X_y, x})$。从而
$$
\dim_x(X_s) =
\dim(\mathcal{O}_{X_s, x}) + \text{trdeg}_{\kappa(y)} \kappa(x) \leq d
$$
（《态射》引理 \ref{morphisms-lemma-dimension-fibre-at-a-point}）。另一方面，由《态射》引理
\ref{morphisms-lemma-openness-bounded-dimension-fibres}，有
$\dim_x(X_s) \geq \dim_\eta(X_{f(\eta)}) = d$。
\end{proof}

\begin{lemma}
\label{varieties-lemma-dominate-valuation-ring-dimension-fibres}
设 $f : X \to \Spec(R)$ 是从不可约概形到赋值环之谱的态射。若 $f$
局部有限型且满射，那么特殊纤维等维，且其维数等于一般纤维的维数。
\end{lemma}

\begin{proof}
可以用其约化替换 $X$，因为这不改变 $X$ 或特殊纤维的维数。
于是 $X$ 整，引理来自《代数》引理
\ref{algebra-lemma-finite-type-domain-over-valuation-ring-dim-fibres}。
\end{proof}

\noindent
下一个引理推广引理 \ref{varieties-lemma-dimension-fibre-in-codim-1}。

\begin{lemma}
\label{varieties-lemma-dimension-fibre-in-higher-codimension}
设 $f : X \to Y$ 局部有限型。设 $x \in X$ 是一点，其像为
$y \in Y$，并且 $\mathcal{O}_{Y, y}$ 是 Noether 环。设 $d \geq 0$
是整数，使得对每个泛点 $\eta$，若它属于 $X$ 的一个包含 $x$ 的不可约分支，则有
$f(\eta) \not = y$ 且 $\dim_\eta(X_{f(\eta)}) = d$。那么
$\dim_x(X_y) \leq d + \dim(\mathcal{O}_{Y, y}) - 1$。
\end{lemma}

\begin{proof}
与引理 \ref{varieties-lemma-dimension-fibre-in-codim-1} 的证明完全一样，我们把问题归约到
$X = \Spec(A)$ 的情形，其中 $A$ 是整环，而 $Y = \Spec(B)$，其中
$B$ 是局部 Noether 环，它的极大理想对应于 $y$。把 $B$
替换为 $B/\Ker(B \to A)$ 后，可以假设 $B$ 是整环且
$B \subset A$。然后使用维数公式（《态射》引理
\ref{morphisms-lemma-dimension-formula}）得到
$$
\dim(\mathcal{O}_{X, x}) + \text{trdeg}_{\kappa(y)} \kappa(x) \leq
\dim(B) + \text{trdeg}_B(A)
$$
我们有 $\text{trdeg}_B(A) = d$，这是因为假设 $\dim_\eta(X_\xi) = d$，见《态射》引理
\ref{morphisms-lemma-dimension-fibre-at-a-point}。由于
$\mathcal{O}_{X, x} \to \mathcal{O}_{X_y, x}$ 有核（因为 $f(\eta) \not = y$），
并且 $\mathcal{O}_{X, x}$ 是 Noether 整环，可知
$\dim(\mathcal{O}_{X, x}) > \dim(\mathcal{O}_{X_y, x})$。从而
$$
\dim_x(X_y) =
\dim(\mathcal{O}_{X_y, x}) + \text{trdeg}_{\kappa(y)} \kappa(x)
< \dim(B) + d
$$
（等式来自《态射》引理 \ref{morphisms-lemma-dimension-fibre-at-a-point}），
这就证明了所需结论。
\end{proof}


\section{代数概形}
\label{varieties-section-algebraic-schemes}

\noindent
下列定义取自 \cite[I Definition 6.4.1]{EGA}。

\begin{definition}
\label{varieties-definition-algebraic-scheme}
设 $k$ 为域。称 $k$ 上的概形 $X$ 为{it 代数 $k$-概形}，如果结构态射
$X \to \Spec(k)$ 有限型。称 $k$ 上的概形 $X$ 为{it 局部代数 $k$-概形}，
如果结构态射 $X \to \Spec(k)$ 局部有限型。
\end{definition}

\noindent
注意，每个（局部）代数 $k$-概形都是（局部）Noether 概形，见《态射》引理
\ref{morphisms-lemma-finite-type-noetherian}。代数 $k$-概形的范畴具有所有积和
纤维积（这不同于 $k$ 上簇的范畴）。局部代数 $k$-概形的范畴也一样。

\begin{lemma}
\label{varieties-lemma-algebraic-scheme-dim-0}
设 $k$ 为域。设 $X$ 是局部代数 $k$-概形，且维数为 $0$。那么 $X$ 是若干
局部 Artin $k$-代数 $A$ 的谱的不交并，并且 $\dim_k(A) < \infty$。若 $X$
是代数 $k$-概形，且维数为 $0$，那么此外 $X$ 是仿射的，并且态射
$X \to \Spec(k)$ 是有限态射。
\end{lemma}

\begin{proof}
设 $X$ 是局部代数 $k$-概形，且维数为 $0$。设
$U = \Spec(A) \subset X$ 是仿射开子概形。由于 $\dim(X) = 0$，可知
$\dim(A) = 0$。由 Noether 正规化（见《代数》引理
\ref{algebra-lemma-Noether-normalization}），存在有限单射 $k \to A$，即
$\dim_k(A) < \infty$。因此 $A$ 是 Artin 环，见《代数》引理
\ref{algebra-lemma-finite-dimensional-algebra}。这蕴涵
$A = A_1 \times \ldots \times A_r$ 是有限多个 Artin 局部环的积，见《代数》引理
\ref{algebra-lemma-artinian-finite-length}。当然，$\dim_k(A_i) < \infty$
对每个 $i$ 都成立，因为这些维数之和等于 $\dim_k(A)$。

\medskip\noindent
以上论证说明，$X$ 有一个开覆盖，其成员都是有限离散拓扑空间。因此 $X$
是离散拓扑空间。于是 $X$ 同构于其连通分支的不交并，而每个连通分支都是单点。
单点概形是仿射的，故由上一段的结果，每个这样的单点都是某个局部 Artin
$k$-代数 $A$ 的谱，并且 $\dim_k(A) < \infty$。

\medskip\noindent
最后，若 $X$ 是代数 $k$-概形且维数为 $0$，则 $X$ 拟紧，因而是有限不交并
$X = \Spec(A_1) \amalg \ldots \amalg \Spec(A_r)$，所以是仿射的（见《概形》引理
\ref{schemes-lemma-disjoint-union-affines}）；而证明的第一段已经说明
$X \to \Spec(k)$ 是有限态射。
\end{proof}

\noindent
下列引理汇集了局部代数概形维数论的一些陈述。

\begin{lemma}
\label{varieties-lemma-dimension-locally-algebraic}
设 $k$ 为域。设 $X$ 是局部代数 $k$-概形。
\begin{enumerate}
\item
\label{varieties-item-catenary}
$X$ 的拓扑空间是链状的（《拓扑》定义
\ref{topology-definition-catenary}）。
\item
\label{varieties-item-dimension-at-closed-point}
若 $x \in X$ 是闭点，则 $\dim_x(X) = \dim(\mathcal{O}_{X, x})$。
\item
\label{varieties-item-dimension-irreducible}
若 $X$ 不可约，则有 $\dim(X) = \dim(U)$，其中 $U \subset X$
为任意非空开集；并且有 $\dim(X) = \dim_x(X)$，其中 $x \in X$ 任意。
\item
\label{varieties-item-irreducible-maximal-chains}
若 $X$ 不可约，则不可约闭子集的任意链都可以延拓为极大链，并且不可约闭子集的
每条极大链的长度都等于 $\dim(X)$。
\item
\label{varieties-item-dimension-irreducibles-passing-through}
对 $x \in X$，有
$\dim_x(X) = \max \dim(Z) = \min \dim(\mathcal{O}_{X, x'})$，
其中极大值取遍不可约分支 $Z \subset X$，这些分支包含 $x$；极小值取遍满足
$x \leadsto x'$ 且 $x'$ 在 $X$ 中闭的特化。
\item
\label{varieties-item-dimension-irreducible-trdeg}
若 $X$ 不可约且泛点为 $x$，则
$\dim(X) = \text{trdeg}_k(\kappa(x))$。
\item
\label{varieties-item-immediate-specialization}
若 $x \leadsto x'$ 是 $X$ 中点的直接特化，则
$\text{trdeg}_k(\kappa(x)) = \text{trdeg}_k(\kappa(x')) + 1$。
\item
\label{varieties-item-dimension-sup-trdeg}
$X$ 的维数等于数 $\text{trdeg}_k(\kappa(x))$ 的上确界，其中 $x$
取遍 $X$ 的不可约分支的泛点。
\item
\label{varieties-item-specialization}
若 $x \leadsto x'$ 是 $X$ 中点的非平凡特化，则
\begin{enumerate}
\item $\dim_x(X) \leq \dim_{x'}(X)$，
\item $\dim(\mathcal{O}_{X, x}) < \dim(\mathcal{O}_{X, x'})$，
\item $\text{trdeg}_k(\kappa(x)) > \text{trdeg}_k(\kappa(x'))$，并且
\item 任意非平凡特化的极大链
$x = x_0 \leadsto x_1 \leadsto \ldots \leadsto x_n = x'$ 的长度为
$n = \text{trdeg}_k(\kappa(x)) - \text{trdeg}_k(\kappa(x'))$。
\end{enumerate}
\item
\label{varieties-item-dimension-formula}
对 $x \in X$，有
$\dim_x(X) = \text{trdeg}_k(\kappa(x)) + \dim(\mathcal{O}_{X, x})$。
\item
\label{varieties-item-immediate-specialization-local-ring}
若 $x \leadsto x'$ 是 $X$ 中点的直接特化，并且 $X$ 不可约或等维，则
$\dim(\mathcal{O}_{X, x'}) = \dim(\mathcal{O}_{X, x}) + 1$。
\end{enumerate}
\end{lemma}

\begin{proof}
我们不依赖前面证明的更一般结果，而把这些陈述归约为有限型 $k$-代数的相应陈述，
并引用交换代数章中的结果。

\medskip\noindent
(\ref{varieties-item-catenary}) 的证明。由《拓扑》引理
\ref{topology-lemma-catenary}，这是关于 $X$ 的局部问题。因此可以假设
$X = \Spec(A)$，其中 $A$ 是有限型 $k$-代数。我们必须证明 $A$ 是链环
（《代数》引理 \ref{algebra-lemma-catenary}）。利用《代数》引理
\ref{algebra-lemma-quotient-catenary}，可以归约到 $k[x_1, \ldots, x_n]$，
然后应用《代数》引理 \ref{algebra-lemma-dimension-height-polynomial-ring}。
另一种证明是：$k$ 显然是 Cohen--Macaulay 环，而 Cohen--Macaulay 环是普遍链环
（《代数》引理 \ref{algebra-lemma-CM-ring-catenary}）。

\medskip\noindent
(\ref{varieties-item-dimension-at-closed-point}) 的证明。选取仿射邻域
$U = \Spec(A)$，它包含 $x$。则 $\dim_x(X) = \dim_x(U)$。因此归约到仿射情形，这正是
《代数》引理 \ref{algebra-lemma-dimension-closed-point-finite-type-field}。

\medskip\noindent
(\ref{varieties-item-dimension-irreducible}) 的证明。只需证明任意两个非空仿射开集
$U, U' \subset X$ 具有相同维数（不可约子集的任意有限链都与某个仿射开集相交）。
选取闭点 $x$（它是 $X$ 的闭点），使得 $x \in U \cap U'$。这是可能的，因为 $X$ 不可约，
故 $U \cap U'$ 非空；又由引理 \ref{varieties-lemma-locally-finite-type-Jacobson}，
$X$ 是 Jacobson 概形，所以其中存在这样的闭点。由《代数》引理
\ref{algebra-lemma-dimension-spell-it-out}，有
$\dim(U) = \dim(\mathcal{O}_{X, x}) = \dim(U')$（严格地说，应用该引理前必须
用其约化替换 $X$）。

\medskip\noindent
(\ref{varieties-item-irreducible-maximal-chains}) 的证明。给定不可约闭子集的一条链，
可以找到与其中最小者相交的仿射开集 $U \subset X$。因此，结论来自《代数》引理
\ref{algebra-lemma-dimension-spell-it-out} 以及我们在
(\ref{varieties-item-dimension-irreducible}) 中看到的 $\dim(U) = \dim(X)$。

\medskip\noindent
(\ref{varieties-item-dimension-irreducibles-passing-through}) 的证明。选取仿射邻域
$U = \Spec(A)$，它包含 $x$。则 $\dim_x(X) = \dim_x(U)$。规则 $Z \mapsto Z \cap U$
给出 $X$ 的不可约分支（它们经过 $x$）与 $U$ 的不可约分支（它们经过 $x$）之间的双射。
并且由 (\ref{varieties-item-dimension-irreducible})，有
$\dim(Z \cap U) = \dim(Z)$，这对这样的 $Z$ 成立。因此结论来自《代数》引理
\ref{algebra-lemma-dimension-at-a-point-finite-type-over-field}。

\medskip\noindent
(\ref{varieties-item-dimension-irreducible-trdeg}) 的证明。由
(\ref{varieties-item-dimension-irreducible})，可以归约到 $X = \Spec(A)$ 为仿射的情形。
此时，把《代数》引理 \ref{algebra-lemma-dimension-prime-polynomial-ring}
应用于 $A_{red}$ 即得结论。

\medskip\noindent
(\ref{varieties-item-immediate-specialization}) 的证明。设
$Z = \overline{\{x\}} \supset Z' = \overline{\{x'\}}$。
由 (\ref{varieties-item-irreducible-maximal-chains})，$Z \supset Z'$ 是 $Z$
中不可约闭子概形的一条极大链的开端，因而 $\dim(Z) = \dim(Z') + 1$。
再由 (\ref{varieties-item-dimension-irreducible-trdeg}) 即得结论。

\medskip\noindent
(\ref{varieties-item-dimension-sup-trdeg}) 的证明。一个简单的拓扑论证说明
$\dim(X) = \sup \dim(Z)$，其中上确界取遍 $X$ 的不可约分支（提示：使用《拓扑》引理
\ref{topology-lemma-irreducible}）。因此结论来自
(\ref{varieties-item-dimension-irreducible-trdeg})。

\medskip\noindent
(\ref{varieties-item-specialization}) 的证明。(a) 来自如下事实：每个开集
$U \subset X$，若它包含 $x'$，则也包含 $x$。(b) 来自如下事实：$\mathcal{O}_{X, x}$ 是
$\mathcal{O}_{X, x'}$ 的一个局部化，因而 $\mathcal{O}_{X, x}$ 中的每条素理想链
都对应于 $\mathcal{O}_{X, x'}$ 中的一条素理想链，而且可在末端添加
$\mathfrak m_{x'}$ 来延长它。(c) 和 (d) 都形式地来自
(\ref{varieties-item-immediate-specialization})。

\medskip\noindent
(\ref{varieties-item-dimension-formula}) 的证明。选取仿射邻域
$U = \Spec(A)$，它包含 $x$。则 $\dim_x(X) = \dim_x(U)$。因此归约到仿射情形，这正是
《代数》引理 \ref{algebra-lemma-dimension-at-a-point-finite-type-field}。

\medskip\noindent
(\ref{varieties-item-immediate-specialization-local-ring}) 的证明。若 $X$ 等维
（《拓扑》定义 \ref{topology-definition-equidimensional}），则 $\dim(X)$
等于 $X$ 的每个不可约分支的维数，从而由
(\ref{varieties-item-dimension-irreducibles-passing-through})，有
$\dim_x(X) = \dim(X) = \dim_{x'}(X)$。因此结论来自
(\ref{varieties-item-immediate-specialization})。
\end{proof}

\begin{lemma}
\label{varieties-lemma-dimension-fibres-locally-algebraic}
设 $k$ 为域。设 $f : X \to Y$ 是局部代数 $k$-概形之间的态射。
\begin{enumerate}
\item 对 $y \in Y$，纤维 $X_y$ 是 $\kappa(y)$ 上的局部代数概形，因而引理
\ref{varieties-lemma-dimension-locally-algebraic} 的所有结果都适用。
\item 假设 $X$ 不可约。令 $Z = \overline{f(X)}$，并令
$d = \dim(X) - \dim(Z)$。则
\begin{enumerate}
\item 有 $\dim_x(X_{f(x)}) \geq d$，对所有 $x \in X$ 都成立，
\item 点 $x \in X$ 中满足 $\dim_x(X_{f(x)}) = d$ 的那些点组成稠密开集，
\item 若 $\dim(\mathcal{O}_{Z, f(x)}) \geq 1$，则
$\dim_x(X_{f(x)}) \leq d + \dim(\mathcal{O}_{Z, f(x)}) - 1$，
\item 若 $\dim(\mathcal{O}_{Z, f(x)}) = 1$，则
$\dim_x(X_{f(x)}) = d$。
\end{enumerate}
\item 对 $x \in X$，若 $y = f(x)$，则有
$\dim_x(X_y) \geq \dim_x(X) - \dim_y(Y)$。
\end{enumerate}
\end{lemma}

\begin{proof}
由《态射》引理 \ref{morphisms-lemma-permanence-finite-type}，态射 $f$
局部有限型。因此基变换 $X_y \to \Spec(\kappa(y))$ 局部有限型。这证明了 (1)。
在证明的其余部分，我们将自由地对 $X$、$Y$ 和 $f$ 的纤维使用引理
\ref{varieties-lemma-dimension-locally-algebraic} 的结果。

\medskip\noindent
(2) 的证明。设 $\eta \in X$ 为泛点，并令 $\xi = f(\eta)$。则
$Z = \overline{\{\xi\}}$。因此
$$
d = \dim(X) - \dim(Z) =
\text{trdeg}_k \kappa(\eta) - \text{trdeg}_k \kappa(\xi) =
\text{trdeg}_{\kappa(\xi)} \kappa(\eta) = \dim_\eta(X_\xi)
$$
于是 (2)(a) 和 (2)(b) 来自《态射》引理
\ref{morphisms-lemma-openness-bounded-dimension-fibres}。(2)(c) 和 (2)(d)
来自引理 \ref{varieties-lemma-dimension-fibre-in-higher-codimension} 和
\ref{varieties-lemma-dimension-fibre-in-codim-1}。

\medskip\noindent
(3) 的证明。设 $x \in X$。设 $X' \subset X$ 是 $X$ 的一个不可约分支，
它经过 $x$ 且维数为 $\dim_x(X)$。于是 (2) 蕴涵
$\dim_x(X_y) \geq \dim(X') - \dim(Z')$，其中 $Z' \subset Y$
是 $X'$ 的像的闭包。这证明了 (3)。
\end{proof}

\begin{lemma}
\label{varieties-lemma-dimension-product-locally-algebraic}
\begin{slogan}
积的维数等于各因子的维数之和。
\end{slogan}
设 $k$ 为域。设 $X$、$Y$ 是局部代数 $k$-概形。
\begin{enumerate}
\item 若 $z \in X \times Y$ 位于 $(x, y)$ 上方，则
$\dim_z(X \times Y) = \dim_x(X) + \dim_y(Y)$。
\item 有 $\dim(X \times Y) = \dim(X) + \dim(Y)$。
\end{enumerate}
\end{lemma}

\begin{proof}
(1) 的证明。考虑结构态射的分解
$$
X \times Y \longrightarrow Y \longrightarrow \Spec(k)
$$
第一个态射 $p : X \times Y \to Y$ 是平坦态射，因为它是平坦态射
$X \to \Spec(k)$ 的基变换，见《态射》引理
\ref{morphisms-lemma-base-change-flat}。此外，由《态射》引理
\ref{morphisms-lemma-dimension-fibre-after-base-change}，有
$\dim_z(p^{-1}(y)) = \dim_x(X)$。因此由《态射》引理
\ref{morphisms-lemma-dimension-fibre-at-a-point-additive}，有
$\dim_z(X \times Y) = \dim_x(X) + \dim_y(Y)$。(2) 是 (1) 的直接推论。
\end{proof}


\section{完备局部环}
\label{varieties-section-complete-local-rings}

\noindent
这里给出域上概形的完备局部环的一些结果。

\begin{lemma}
\label{varieties-lemma-complete-local-ring-structure-as-algebra}
设 $k$ 为域。设 $X$ 是 $k$ 上的局部 Noether 概形。设 $x \in X$
是一点，其剩余域为 $\kappa$。存在同构
\begin{equation}
\label{varieties-equation-complete-local-ring}
\kappa[[x_1, \ldots, x_n]]/I \longrightarrow \mathcal{O}_{X, x}^\wedge
\end{equation}
它在剩余域上诱导恒等映射。一般而言，不能选取
(\ref{varieties-equation-complete-local-ring}) 使其成为 $k$-代数同构。然而，若扩张
$\kappa/k$ 可分，则可以选取 (\ref{varieties-equation-complete-local-ring})
使其成为 $k$-代数同构。
\end{lemma}

\begin{proof}
该同构的存在性是 Cohen 结构定理的直接推论\footnote{注意，若 $\kappa$
的特征为 $p$，则该定理只说明存在满射
$\Lambda[[x_1, \ldots, x_n]] \to \mathcal{O}_{X, x}^\wedge$，其中
$\Lambda$ 是 $\kappa$ 的 Cohen 环。但在这种情形下，该映射当然经过
$\Lambda/p\Lambda[[x_1, \ldots, x_n]]$ 分解，并且
$\Lambda/p\Lambda = \kappa$。}（《代数》定理
\ref{algebra-theorem-cohen-structure-theorem}）。

\medskip\noindent
设 $p$ 为奇素数，设 $k = \mathbf{F}_p(t)$，并令
$A = k[x, y]/(y^2 + x^p - t)$。则完备化 $A^\wedge$（即 $A$
在极大理想 $\mathfrak m = (y)$ 处的完备化）作为环同构于
$k(t^{1/p})[[z]]$，但不作为 $k$-代数同构。原因是 $A^\wedge$
不含有 $p$ 次幂为 $t$ 的元素（读者可以通过模 $y^2$ 计算看出这一点）。
这也说明 (\ref{varieties-equation-complete-local-ring}) 的任何同构都不可能是
$k$-代数同构。

\medskip\noindent
若 $\kappa/k$ 可分，则由《代数进阶》引理
\ref{more-algebra-lemma-lift-residue-field}，存在 $k$-代数同态
$\kappa \to \mathcal{O}_{X, x}^\wedge$，它在剩余域上诱导恒等映射。
设 $f_1, \ldots, f_n \in \mathfrak m_x$ 为一组生成元。考虑映射
$$
\kappa[[x_1, \ldots, x_n]] \longrightarrow \mathcal{O}_{X, x}^\wedge,\quad
x_i \longmapsto f_i
$$
由于两边关于 $(x_1, \ldots, x_n)$ 都完备（右边的完备性见《代数》引理
\ref{algebra-lemma-hathat-finitely-generated}），而且按构造，该映射模
$(x_1, \ldots, x_n)$ 是满射，所以由《代数》引理
\ref{algebra-lemma-completion-generalities}，该映射是满射。
\end{proof}

\begin{lemma}
\label{varieties-lemma-base-change-complete-local-ring}
设 $K/k$ 是域扩张。设 $X$ 是局部代数 $k$-概形。令 $Y = X_K$。
设 $y \in Y$ 是一点，其像为 $x \in X$。假设
$\dim(\mathcal{O}_{X, x}) = \dim(\mathcal{O}_{Y, y})$，并且
$\kappa(x)/k$ 可分。选取如 (\ref{varieties-equation-complete-local-ring}) 中的
$k$-代数同构
$$
\kappa(x)[[x_1, \ldots, x_n]]/(g_1, \ldots, g_m) \longrightarrow
\mathcal{O}_{X, x}^\wedge
$$
则存在如 (\ref{varieties-equation-complete-local-ring}) 中的 $K$-代数同构
$$
\kappa(y)[[x_1, \ldots, x_n]]/(g_1, \ldots, g_m) \longrightarrow
\mathcal{O}_{Y, y}^\wedge
$$
这里使用 $\kappa(x) \to \kappa(y)$，把 $g_j$ 看作 $\kappa(y)$
上的幂级数。
\end{lemma}

\begin{proof}
局部环映射 $\mathcal{O}_{X, x} \to \mathcal{O}_{Y, y}$ 诱导局部环映射
$\mathcal{O}_{X, x}^\wedge \to \mathcal{O}_{Y, y}^\wedge$。诱导映射
$$
\kappa(x) \to \kappa(x)[[x_1, \ldots, x_n]]/(g_1, \ldots, g_m)
\to \mathcal{O}_{X, x}^\wedge \to \mathcal{O}_{Y, y}^\wedge
$$
与到 $\kappa(y)$ 的投影复合后，就是典范同态
$\kappa(x) \to \kappa(y)$。由引理 \ref{varieties-lemma-change-fields-flat}，
剩余域 $\kappa(y)$ 是 $\kappa(x) \otimes_k K$ 在核 $\mathfrak p_0$ 处的局部化；
这里所指的同态为 $\kappa(x) \otimes_k K \to \kappa(y)$。
另一方面，由引理 \ref{varieties-lemma-change-fields-algebraic-unramified}，局部环
$(\kappa(x) \otimes_k K)_{\mathfrak p_0}$ 等于 $\kappa(y)$。因此映射
$$
\kappa(x) \otimes_k K \to \mathcal{O}_{Y, y}^\wedge
$$
典范地经过 $\kappa(y)$ 分解。我们得到交换图
$$
\xymatrix{
\kappa(y) \ar[rr] & & \mathcal{O}_{Y, y}^\wedge \\
\kappa(x) \ar[r] \ar[u] &
\kappa(x)[[x_1, \ldots, x_n]]/(g_1, \ldots, g_m) \ar[r] &
\mathcal{O}_{X, x}^\wedge \ar[u]
}
$$
设 $f_i \in \mathfrak m_x^\wedge \subset \mathcal{O}_{X, x}^\wedge$
为 $x_i$ 的像。注意，由于该映射是满射，有
$\mathfrak  m_x^\wedge = (f_1, \ldots, f_n)$。考虑映射
$$
\kappa(y)[[x_1, \ldots, x_n]] \longrightarrow \mathcal{O}_{Y, y}^\wedge,\quad
x_i \longmapsto f_i
$$
这里的 $f_i$ 实际上是指 $f_i$ 在 $\mathfrak m_y^\wedge$ 中的像。
由引理 \ref{varieties-lemma-change-fields-algebraic-unramified}，有
$\mathfrak m_x \mathcal{O}_{Y, y} = \mathfrak m_y$，所以右边关于
$(x_1, \ldots, x_n)$ 完备（使用《代数》引理
\ref{algebra-lemma-hathat-finitely-generated} 可知它是完备局部环）。
由于两边关于 $(x_1, \ldots, x_n)$ 都完备，而该映射模
$(x_1, \ldots, x_n)$ 是满射，故由《代数》引理
\ref{algebra-lemma-completion-generalities}，我们的映射是满射。
当然，幂级数 $g_1, \ldots, g_m$ 在此映射下映为零，因为它们在
$\mathcal{O}_{X, x}^\wedge$ 中已经映为零。因此有交换图
$$
\xymatrix{
\kappa(y)[[x_1, \ldots, x_n]]/(g_1, \ldots, g_m) \ar[r] &
\mathcal{O}_{Y, y}^\wedge \\
\kappa(x)[[x_1, \ldots, x_n]]/(g_1, \ldots, g_m) \ar[r] \ar[u] &
\mathcal{O}_{X, x}^\wedge \ar[u]
}
$$
还需要证明上方的水平箭头是同构。我们已经知道它是满射。由引理
\ref{varieties-lemma-change-fields-flat}，$\mathcal{O}_{X, x} \to \mathcal{O}_{Y, y}$
是平坦的，这蕴涵 $\mathcal{O}_{X, x}^\wedge \to \mathcal{O}_{Y, y}^\wedge$
是平坦的（《代数进阶》引理
\ref{more-algebra-lemma-flat-completion}）。因此可以应用《代数》引理
\ref{algebra-lemma-mod-injective}，其中
$R = \kappa(x)[[x_1, \ldots, x_n]]/(g_1, \ldots, g_m)$，
$S = \kappa(y)[[x_1, \ldots, x_n]]/(g_1, \ldots, g_m)$，
$M = \mathcal{O}_{Y, y}^\wedge$，并且 $N = S$，从而得出该映射是单射。
\end{proof}


\section{全局生成}
\label{varieties-section-global-generation}

\noindent
这里给出与准凝聚模全局生成有关的一些引理。

\begin{lemma}
\label{varieties-lemma-globally-generated-base-change}
设 $X \to \Spec(A)$ 是概形的态射。设 $A \subset A'$ 是忠实平坦环映射。
设 $\mathcal{F}$ 是准凝聚 $\mathcal{O}_X$-模。那么 $\mathcal{F}$ 全局生成，
当且仅当基变换 $\mathcal{F}_{A'}$ 全局生成。
\end{lemma}

\begin{proof}
更准确地说，令 $X_{A'} = X \times_{\Spec(A)} \Spec(A')$。令
$\mathcal{F}_{A'} = p^*\mathcal{F}$，其中 $p : X_{A'} \to X$ 是投影。
由《概形上同调》引理 \ref{coherent-lemma-flat-base-change-cohomology}，有
$H^0(X_{k'}, \mathcal{F}_{A'}) = H^0(X, \mathcal{F}) \otimes_A A'$。
因此，若 $s_i$（$i \in I$）是 $H^0(X, \mathcal{F})$ 作为 $A$-模的一组生成元，
则它们在 $H^0(X_{A'}, \mathcal{F}_{A'})$ 中的像是
$H^0(X_{A'}, \mathcal{F}_{A'})$ 作为 $A'$-模的一组生成元。
所以只需证明映射
$\alpha : \bigoplus_{i \in I} \mathcal{O}_X \to \mathcal{F}$，
$(f_i) \mapsto \sum f_i s_i$ 是满射，当且仅当 $p^*\alpha$ 是满射。
这可以在仿射开集 $U = \Spec(B)$ 上检验；它是 $X$ 的开集。此时
$\mathcal{F}|_U$ 对应于 $B$-模 $M$，而 $s_i|_U$ 对应于元素
$x_i \in M$。因此只需证明 $\bigoplus_{i \in I} B \to M$ 是满射，
当且仅当基变换
$\bigoplus_{i \in I} B \otimes_A A' \to M \otimes_A A'$ 是满射。
这成立是因为 $A \to A'$ 忠实平坦。
\end{proof}

\begin{lemma}
\label{varieties-lemma-very-ample-vanish-at-point}
设 $k$ 是无限域。设 $X$ 是 $k$ 上的有限型概形。设 $\mathcal{L}$
是 $X$ 上的甚丰沛可逆层。设 $n \geq 0$，并设
$x, x_1, \ldots, x_n \in X$ 是若干点，其中 $x$ 是 $k$-有理点，即
$\kappa(x) = k$，并且有 $x \not = x_i$，其中 $i = 1, \ldots, n$。
那么存在 $s \in H^0(X, \mathcal{L})$，它在 $x$ 处为零，但在 $x_i$
处不为零。
\end{lemma}

\begin{proof}
若 $n = 0$，结论显然，故假设 $n > 0$。由甚丰沛可逆层的定义，
引理立即归约到 $X = \mathbf{P}^r_k$ 的情形，其中某个 $r > 0$，并且
$\mathcal{L} = \mathcal{O}_X(1)$。写成
$\mathbf{P}^r_k = \text{Proj}(k[T_0, \ldots, T_r])$。令
$V = H^0(X, \mathcal{L}) = kT_0 \oplus \ldots \oplus kT_r$。
由于 $x$ 是 $k$-有理点，元素 $s \in V$ 中在 $x$ 处为零的那些元素组成余维为 $1$
的子空间 $W \subset V$，并且 $W$ 生成与 $x$ 对应的齐次素理想。
由于 $x_i \not = x$，相应的齐次素理想
$\mathfrak p_i \subset k[T_0, \ldots, T_r]$ 不包含 $W$。由于 $k$ 无限，
可知 $W \not = \bigcup W \cap \mathfrak q_i$，证明完成。
\end{proof}

\begin{lemma}
\label{varieties-lemma-generated-by-dim-plus-1-sections}
设 $k$ 是无限域。设 $X$ 是代数 $k$-概形。设 $\mathcal{L}$ 是可逆
$\mathcal{O}_X$-模。设 $V \to \Gamma(X, \mathcal{L})$ 是 $k$-向量空间的
线性映射，其像生成 $\mathcal{L}$。那么存在子空间 $W \subset V$，满足
$\dim_k(W) \leq \dim(X) + 1$，并且它生成 $\mathcal{L}$。
\end{lemma}

\begin{proof}
在整个证明中，我们将使用如下事实：对每个 $x \in X$，线性映射
$$
\psi_x : V \to \Gamma(X, \mathcal{L}) \to \mathcal{L}_x \to
\mathcal{L}_x \otimes_{\mathcal{O}_{X, x}} \kappa(x)
$$
非零。证明对 $\dim(X)$ 使用归纳法。

\medskip\noindent
基始情形是 $\dim(X) = 0$。此时 $X$ 只有有限多个点
$X = \{x_1, \ldots, x_n\}$（例如见引理
\ref{varieties-lemma-algebraic-scheme-dim-0}）。由于 $k$ 无限，存在向量
$v \in V$，使得 $\psi_{x_i}(v) \not = 0$ 对所有 $i$ 都成立。
于是 $W = k\cdot v$ 满足要求。

\medskip\noindent
假设 $\dim(X) > 0$。设 $X_i \subset X$ 是维数等于 $\dim(X)$ 的不可约分支。
由于 $X$ 是 Noether 概形，这样的分支只有有限多个。对每个 $i$，选取一点
$x_i \in X_i$。如上选取 $v \in V$，使得 $\psi_{x_i}(v) \not = 0$
对所有 $i$ 都成立。设 $Z \subset X$ 是 $v$ 在
$\Gamma(X, \mathcal{L})$ 中之像的零概形，见《除子》定义
\ref{divisors-definition-zero-scheme-s}。按构造，$\dim(Z) < \dim(X)$。
由归纳假设，可以找到 $W \subset V$，满足 $\dim(W) \leq \dim(X)$，
使得 $W$ 生成 $\mathcal{L}|_Z$。于是 $W + k\cdot v$ 生成 $\mathcal{L}$。
\end{proof}


\section{分离点与切向量}
\label{varieties-section-separating-points-tangent-vectors}

\noindent
本节就是下列结果。

\begin{lemma}
\label{varieties-lemma-separate-points-tangent-vectors}
设 $k$ 是代数闭域。设 $X$ 是固有 $k$-概形。设 $\mathcal{L}$ 是可逆
$\mathcal{O}_X$-模。设 $V \subset H^0(X, \mathcal{L})$ 是 $k$-向量子空间。
如果
\begin{enumerate}
\item 对每一对互异闭点 $x, y \in X$，存在截面 $s \in V$，它在 $x$
处为零而在 $y$ 处不为零；并且
\item 对每个闭点 $x \in X$ 和非零切向量 $\theta \in T_{X/k, x}$，
存在截面 $s \in V$，它在 $x$ 处为零，但经 $\theta$ 拉回后非零，
\end{enumerate}
那么 $\mathcal{L}$ 甚丰沛，并且典范态射
$\varphi_{\mathcal{L}, V} : X \to \mathbf{P}(V)$ 是闭浸入。
\end{lemma}

\begin{proof}
条件 (1) 尤其蕴涵 $V$ 的元素生成 $\mathcal{L}$（作为 $X$ 上的层）。因此由《构造》例
\ref{constructions-example-projective-space}，得到典范态射
$$
\varphi = \varphi_{\mathcal{L}, V} :
X
\longrightarrow
\mathbf{P}(V)
$$
由《态射》引理 \ref{morphisms-lemma-image-proper-scheme-closed}，态射
$\varphi$ 是固有的。由 (1)，映射 $\varphi$ 在闭点上是单射（略去计算）。
特别地，$\mathbf{P}(V)$ 的任意闭点上方的纤维都是单点（略去一个小细节）。
因此 $\varphi$ 是有限态射；例如可使用《概形上同调》引理
\ref{coherent-lemma-proper-finite-fibre-finite-in-neighbourhood}。
为完成证明，只需说明映射
$$
\varphi^\sharp :
\mathcal{O}_{\mathbf{P}(V)}
\longrightarrow
\varphi_*\mathcal{O}_X
$$
是满射。这可以在闭点的茎上检验。设 $x \in X$ 是闭点，其像为闭点
$p = \varphi(x) \in \mathbf{P}(V)$。由 (1)，有
$\varphi^{-1}(\{p\}) = \{x\}$；又因为 $\varphi$ 固有（因而是闭映射），
$\varphi^{-1}(U)$ 取遍 $x$ 的开邻域的一个基本系，其中 $U$ 取遍
$p$ 的开邻域的一个基本系。因此在 $p$ 处取茎，得到映射
$$
\varphi^\sharp_x :
\mathcal{O}_{\mathbf{P}(V), p}
\longrightarrow
\mathcal{O}_{X, x}
$$
特别地，$\mathcal{O}_{X, x}$ 是有限 $\mathcal{O}_{\mathbf{P}(V), p}$-模。
此外，$x$ 和 $p$ 的剩余域都等于 $k$（因为 $k$ 代数闭；使用 Hilbert
零点定理）。最后，条件 (2) 蕴涵映射
$$
T_{X/k, x} \longrightarrow T_{\mathbf{P}(V)/k, p}
$$
是单射，因为该映射的核中的任何非零 $\theta$ 都不可能满足 (2) 的结论。
就上述局部环映射而言，这意味着
$$
\mathfrak m_p/\mathfrak m_p^2 \longrightarrow
\mathfrak m_x/\mathfrak m_x^2
$$
是满射，见引理 \ref{varieties-lemma-tangent-space-rational-point}。现在应用《代数》引理
\ref{algebra-lemma-when-surjective-local} 即完成证明。
\end{proof}

\begin{lemma}
\label{varieties-lemma-variant-separate-points-tangent-vectors}
设 $k$ 是代数闭域。设 $X$ 是固有 $k$-概形。设 $\mathcal{L}$ 是可逆
$\mathcal{O}_X$-模。假设对每个闭子概形 $Z \subset X$，只要它的维数为 $0$、
次数为 $2$（相对于 $k$），映射
$$
H^0(X, \mathcal{L}) \longrightarrow H^0(Z, \mathcal{L}|_Z)
$$
就是满射。那么 $\mathcal{L}$ 在 $X$ 上相对于 $k$ 甚丰沛。
\end{lemma}

\begin{proof}
这是引理 \ref{varieties-lemma-separate-points-tangent-vectors} 的改述。事实上，给定互异闭点
$x, y \in X$，取 $Z = x \cup y$（视为闭子概形），便得到该引理的条件 (1)。
给定非零切向量 $\theta \in T_{X/k, x}$，态射
$\theta : \Spec(k[\epsilon]) \to X$ 是闭浸入。令 $Z = \Im(\theta)$，
便得到该引理的条件 (2)。
\end{proof}


\section{积的闭包}
\label{varieties-section-closure-of-products}

\noindent
以下给出若干关于闭包与积之间关系的结果。

\begin{lemma}
\label{varieties-lemma-closure-of-product}
设 $k$ 是域。设 $X$、$Y$ 是 $k$-概形，并设
$A \subset X$、$B \subset Y$ 是子集。令
$$
AB =
\{z \in X \times_k Y \mid \text{pr}_X(z) \in A, \ \text{pr}_Y(z) \in B\}
\subset X \times_k Y
$$
则在集合论意义下有
$$
\overline{A} \times_k \overline{B} = \overline{AB}
$$
\end{lemma}

\begin{proof}
包含关系 $\overline{AB} \subset \overline{A} \times_k \overline{B}$ 是显然的。
我们可将 $X$ 和 $Y$ 分别替换为既约闭子概形 $\overline{A}$ 和
$\overline{B}$。设 $W \subset X \times_k Y$ 是非空开子集。由《态射》引理
\ref{morphisms-lemma-scheme-over-field-universally-open}，子集
$U = \text{pr}_X(W)$ 是 $X$ 中的非空开集。因此 $A \cap U$ 非空；取
$a \in A \cap U$。记
$Y_a = \{a\} \times_k Y \cong Y_{\kappa(a)}$ 为投影
$\text{pr}_X : X \times_k Y \to X$ 在 $a$ 上方的纤维。再次由《态射》引理
\ref{morphisms-lemma-scheme-over-field-universally-open}，态射 $Y_a \to Y$
是开映射，因为 $\Spec(\kappa(a)) \to \Spec(k)$ 是万有开态射。因此非空开子集
$W_a = W \times_{X \times_k Y} Y_a$ 映到 $Y$ 的非空开子集。于是该像中存在
$b \in B$。故 $AB \cap W \not = \emptyset$，正是所需结论。
\end{proof}

\begin{lemma}
\label{varieties-lemma-closure-image-product-map}
设 $k$ 是域。设 $f : A \to X$、$g : B \to Y$ 是 $k$-概形之间的态射。
则在集合论意义下有
$$
\overline{f(A)} \times_k \overline{g(B)} =
\overline{(f \times g)(A \times_k B)}
$$
\end{lemma}

\begin{proof}
这由引理 \ref{varieties-lemma-closure-of-product} 得出，因为在该引理的记号下，
$f \times g$ 的像为 $f(A)g(B)$。
\end{proof}

\begin{lemma}
\label{varieties-lemma-scheme-theoretic-image-product-map}
设 $k$ 是域。设 $f : A \to X$、$g : B \to Y$ 是 $k$-概形之间的拟紧态射。
设 $Z \subset X$ 是 $f$ 的概形论像，见《态射》定义
\ref{morphisms-definition-scheme-theoretic-image}。类似地，设 $Z' \subset Y$
是 $g$ 的概形论像。则 $Z \times_k Z'$ 是 $f \times g$ 的概形论像。
\end{lemma}

\begin{proof}
回忆 $Z$ 是 $X$ 的最小闭子概形，并且 $f$ 分解通过它；$Z'$ 也类似。
设 $W \subset X \times_k Y$ 是 $f \times g$ 的概形论像。由于 $f \times g$
分解通过 $Z \times_k Z'$，故 $W \subset Z \times_k Z'$。

\medskip\noindent
为证明反向包含，设 $U \subset X$ 和 $V \subset Y$ 是仿射开集。由《态射》引理
\ref{morphisms-lemma-quasi-compact-scheme-theoretic-image}，概形 $Z \cap U$
是 $f|_{f^{-1}(U)} : f^{-1}(U) \to U$ 的概形论像；$Z' \cap V$ 和
$W \cap U \times_k V$ 也类似。因此可假设 $X$ 和 $Y$ 是仿射的。由于 $f$
和 $g$ 拟紧，这蕴涵 $A = \bigcup U_i$ 是仿射开集的有限并，而
$B = \bigcup V_j$ 也是仿射开集的有限并。于是可将 $A$ 替换为
$\coprod U_i$，将 $B$ 替换为 $\coprod V_j$；也就是说，还可假设 $A$ 和 $B$
也是仿射的。在此情形，$Z$ 由
$\Ker(\Gamma(X, \mathcal{O}_X) \to \Gamma(A, \mathcal{O}_A))$ 截出，
$Z'$ 和 $W$ 也类似。因此结论由等式
$$
\Gamma(A \times_k B, \mathcal{O}_{A \times_k B})
=
\Gamma(A, \mathcal{O}_A) \otimes_k \Gamma(B, \mathcal{O}_B)
$$
得出；该等式成立是因为 $A$ 和 $B$ 是仿射的。略去细节。
\end{proof}


\section{域上的光滑概形}
\label{varieties-section-smooth}

\noindent
下面给出两个刻画域上光滑概形的引理。

\begin{lemma}
\label{varieties-lemma-char-zero-differentials-free-smooth}
设 $k$ 是域。设 $X$ 是 $k$-概形。假设
\begin{enumerate}
\item $X$ 在 $k$ 上局部有限型，
\item $\Omega_{X/k}$ 局部自由，并且
\item $k$ 的特征为零。
\end{enumerate}
那么结构态射 $X \to \Spec(k)$ 是光滑的。
\end{lemma}

\begin{proof}
这由《代数》引理 \ref{algebra-lemma-characteristic-zero-local-smooth} 得出。
\end{proof}

\noindent
在正特征中，存在微分层自由的有限型非既约概形，例如
$\Spec(\mathbf{F}_p[t]/(t^p))$ 在 $\Spec(\mathbf{F}_p)$ 上的情形。若基域 $k$
的特征为 $p$ 且非完全，则存在既约概形 $X/k$，其微分层 $\Omega_{X/k}$
自由却不光滑，例如 $\Spec(k[t]/(t^p-a)$，其中 $a \in k$ 不是 $p$ 次幂。

\begin{lemma}
\label{varieties-lemma-char-p-differentials-free-smooth}
设 $k$ 是域。设 $X$ 是 $k$-概形。假设
\begin{enumerate}
\item $X$ 在 $k$ 上局部有限型，
\item $\Omega_{X/k}$ 局部自由，
\item $X$ 既约，并且
\item $k$ 完全。
\end{enumerate}
那么结构态射 $X \to \Spec(k)$ 是光滑的。
\end{lemma}

\begin{proof}
设 $x \in X$ 是一点。因为 $X$ 局部 Noether（见《态射》引理
\ref{morphisms-lemma-finite-type-noetherian}），只有有限多个不可约分支
$X_1, \ldots, X_n$ 通过 $x$（见《性质》引理
\ref{properties-lemma-Noetherian-topology} 和《拓扑》引理
\ref{topology-lemma-Noetherian}）。设 $\eta_i \in X_i$ 是泛点。由于 $X$ 既约，
有 $\mathcal{O}_{X, \eta_i} = \kappa(\eta_i)$，见《代数》引理
\ref{algebra-lemma-minimal-prime-reduced-ring}。此外，$\kappa(\eta_i)$ 是完全域
$k$ 的有限生成域扩张，因而在 $k$ 上可分生成（见《代数》第
\ref{algebra-section-separability} 节）。所以
$\Omega_{X/k, \eta_i} = \Omega_{\kappa(\eta_i)/k}$ 是自由的，其秩等于
$\kappa(\eta_i)$ 在 $k$ 上的超越次数。由《态射》引理
\ref{morphisms-lemma-dimension-fibre-at-a-point}，得到
$\dim_{\eta_i}(X_i) = \text{rank}_{\eta_i}(\Omega_{X/k})$。由于
$x \in X_1 \cap \ldots \cap X_n$，可见
$$
\text{rank}_x(\Omega_{X/k}) = \text{rank}_{\eta_i}(\Omega_{X/k}) = \dim(X_i).
$$
因此 $\dim_x(X) = \text{rank}_x(\Omega_{X/k})$，见《代数》引理
\ref{algebra-lemma-dimension-at-a-point-finite-type-over-field}。于是例如由《代数》引理
\ref{algebra-lemma-characterize-smooth-over-field}，$X \to \Spec(k)$ 在 $x$
处光滑。
\end{proof}

\begin{lemma}
\label{varieties-lemma-smooth-regular}
\begin{slogan}
域上光滑蕴涵正则
\end{slogan}
设 $X \to \Spec(k)$ 是光滑态射，其中 $k$ 是域。则 $X$ 是正则概形。
\end{lemma}

\begin{proof}
（另见引理 \ref{varieties-lemma-geometrically-regular-smooth}。）由《代数》引理
\ref{algebra-lemma-characterize-smooth-over-field}，每个局部环
$\mathcal{O}_{X, x}$ 都正则。又因为 $X$ 在 $k$ 上局部有限型，所以它局部
Noether。因此由《性质》引理 \ref{properties-lemma-characterize-regular}，
$X$ 正则。
\end{proof}

\begin{lemma}
\label{varieties-lemma-smooth-geometrically-normal}
设 $X \to \Spec(k)$ 是光滑态射，其中 $k$ 是域。则 $X$ 在 $k$ 上几何正则、
几何正规且几何既约。
\end{lemma}

\begin{proof}
（另见引理 \ref{varieties-lemma-geometrically-regular-smooth}。）设 $k'$ 是 $k$ 的有限纯不可分
扩张。只需证明 $X_{k'}$ 正则、正规且既约，见引理
\ref{varieties-lemma-geometrically-regular}、\ref{varieties-lemma-geometrically-normal} 和
\ref{varieties-lemma-check-only-finite-inseparable-extensions}。由《态射》引理
\ref{morphisms-lemma-base-change-smooth}，态射 $X_{k'} \to \Spec(k')$ 也是光滑的。
因此只需说明域上的光滑概形 $X$ 正则、正规且既约。由引理
\ref{varieties-lemma-smooth-regular} 可知 $X$ 正则。于是《性质》引理
\ref{properties-lemma-regular-normal} 保证 $X$ 正规。
\end{proof}

\begin{lemma}
\label{varieties-lemma-affine-space-over-field}
设 $k$ 是域。设 $d \geq 0$。设 $W \subset \mathbf{A}^d_k$ 是非空开集。
则存在闭点 $w \in W$，使得 $k \subset \kappa(w)$ 是有限可分扩张。
\end{lemma}

\begin{proof}
必要时缩小 $W$ 后，可假设 $W = \mathbf{A}^d_k \setminus V(f)$，其中
$f \in k[x_1, \ldots, x_d]$。若引理不成立，则
$f(a_1, \ldots, a_d) = 0$ 对所有
$(a_1, \ldots, a_d) \in (k^{sep})^d$ 都成立。
这是荒谬的，因为 $k^{sep}$ 是无限域。
\end{proof}

\begin{lemma}
\label{varieties-lemma-smooth-separable-closed-points-dense}
设 $k$ 是域。若 $X$ 在 $\Spec(k)$ 上光滑，则集合
$$
\{x \in X\text{ closed such that }k \subset \kappa(x)
\text{ is finite separable}\}
$$
在 $X$ 中稠密。
\end{lemma}

\begin{proof}
只需说明：给定非空光滑概形 $X$（定义在 $k$ 上），至少存在一个剩余域在 $k$ 上有限可分
的闭点。为此，选取图表
$$
\xymatrix{
X & U \ar[l] \ar[r]^-\pi & \mathbf{A}^d_k
}
$$
其中 $\pi$ 为平展态射，见《态射》引理
\ref{morphisms-lemma-smooth-etale-over-affine-space}。态射
$\pi : U \to \mathbf{A}^d_k$ 是开映射，见《态射》引理
\ref{morphisms-lemma-etale-open}。由引理 \ref{varieties-lemma-affine-space-over-field}，
可选取闭点 $w \in \pi(U)$，其剩余域在 $k$ 上有限可分。任选
$x \in U$ 使得 $\pi(x) = w$。由《态射》引理
\ref{morphisms-lemma-etale-over-field}，域扩张 $\kappa(x)/\kappa(w)$ 有限可分。
故 $\kappa(x)/k$ 有限可分。由《态射》引理
\ref{morphisms-lemma-algebraic-residue-field-extension-closed-point-fibre}，
点 $x$ 是 $X$ 的闭点。
\end{proof}

\begin{lemma}
\label{varieties-lemma-geometrically-reduced-dense-smooth-open}
设 $X$ 是域 $k$ 上局部有限型的既约概形。则 $X$ 在 $k$ 上几何既约，当且仅当
$X$ 含有一个在 $k$ 上光滑的稠密开集。
\end{lemma}

\begin{proof}
问题关于 $X$ 是局部的，故可假设 $X$ 拟紧。设
$X = X_1 \cup \ldots \cup X_n$ 是 $X$ 的不可约分支。先假设 $X$ 含有一个
在 $k$ 上光滑的稠密开集。那么 $X$ 在每个 $X_i$ 的泛点处几何既约（例如由引理
\ref{varieties-lemma-smooth-geometrically-normal}）。由引理
\ref{varieties-lemma-generic-points-geometrically-reduced}，可知 $X$ 几何既约。

\medskip\noindent
反之，假设 $X$ 几何既约。由于
$Z = \bigcup_{i \not = j} X_i \cap X_j$ 在 $X$ 中无处稠密，可将
$X$ 替换为 $X \setminus Z$。因为 $X \setminus Z$ 是不可约概形的不交并，问题归结为
$X$ 不可约的情形。由于 $X$ 不可约且既约，它是整概形，见《性质》引理
\ref{properties-lemma-characterize-integral}。设 $\eta \in X$ 是其泛点。那么函数域
$K = k(X) = \kappa(\eta) = \mathcal{O}_{X, \eta}$ 在 $k$ 上几何既约，因而在
$k$ 上可分，见《代数》引理
\ref{algebra-lemma-characterize-separable-field-extensions}。任取非空仿射开集
$U = \Spec(A) \subset X$，使得 $K = A_{(0)}$ 是 $A$ 的分式域。应用《代数》引理
\ref{algebra-lemma-separable-smooth}，可知 $A$ 在 $(0)$ 处在 $k$ 上光滑。
按定义，这意味着 $A$ 的某个主局部化在 $k$ 上光滑，结论得证。
\end{proof}

\begin{lemma}
\label{varieties-lemma-dense-smooth-open-variety-over-perfect-field}
设 $k$ 是完全域。设 $X$ 是局部代数的既约 $k$-概形，例如 $k$ 上的簇。则
$$
\{x \in X \mid X \to \Spec(k)\text{ is smooth at }x\} =
\{x \in X \mid \mathcal{O}_{X, x}\text{ 正则}\}
$$
且这是 $X$ 的稠密开子概形。
\end{lemma}

\begin{proof}
两个集合相等立即由《代数》引理 \ref{algebra-lemma-separable-smooth} 和各定义得出
（完全域的定义见《代数》定义 \ref{algebra-definition-perfect}）。该集合是开的，因为
概形态射的光滑点集合是开的，见《态射》定义
\ref{morphisms-definition-smooth}。最后给出两种论证说明它稠密：
(1) $X$ 的泛点属于该集合，因为泛点处的局部环是域（《代数》引理
\ref{algebra-lemma-minimal-prime-reduced-ring}），因而正则。
(2) 由引理 \ref{varieties-lemma-perfect-reduced}，$X$ 几何既约，因此可应用引理
\ref{varieties-lemma-geometrically-reduced-dense-smooth-open}。
\end{proof}

\begin{lemma}
\label{varieties-lemma-flat-under-smooth}
设 $k$ 是域。设 $f : X \to Y$ 是 $k$ 上局部有限型概形之间的态射。设
$x \in X$ 是一点，并令 $y = f(x)$。若 $X \to \Spec(k)$ 在 $x$ 处光滑，
且 $f$ 在 $x$ 处平坦，则 $Y \to \Spec(k)$ 在 $y$ 处光滑。特别地，若 $X$
在 $k$ 上光滑，并且 $f$ 平坦且满射，则 $Y$ 在 $k$ 上光滑。
\end{lemma}

\begin{proof}
只需说明 $Y$ 在 $y$ 处几何正则，见引理
\ref{varieties-lemma-geometrically-regular-smooth}。这由引理
\ref{varieties-lemma-flat-under-geometrically-regular} 得出（并将引理
\ref{varieties-lemma-geometrically-regular-smooth} 应用于 $(X, x)$）。
\end{proof}

\begin{lemma}
\label{varieties-lemma-variety-with-smooth-rational-point}
设 $k$ 是域。设 $X$ 是 $k$ 上的簇，并有一个 $k$-有理点 $x$，使得 $X$
在 $x$ 处光滑。则 $X$ 在 $k$ 上几何整。
\end{lemma}

\begin{proof}
设 $U \subset X$ 是 $X$ 的光滑点集。由假设，$U$ 非空，因而稠密且概形论稠密。
于是 $U_{\overline{k}} \subset X_{\overline{k}}$ 也稠密且概形论稠密（略去若干细节）。
因此只需说明 $U$ 几何整。由于 $U$ 有一个 $k$-有理点，由引理
\ref{varieties-lemma-geometrically-connected-if-connected-and-point}，它几何连通。另一方面，
$U_{\overline{k}}$ 既约且正规（引理
\ref{varieties-lemma-smooth-geometrically-normal}）。由于连通正规 Noether 概形是整概形
（《性质》引理 \ref{properties-lemma-normal-Noetherian}），证明完成。
\end{proof}

\begin{lemma}
\label{varieties-lemma-smooth-stratification}
设 $X$ 是域 $k$ 上的有限型概形。存在有限纯不可分扩张 $k'/k$、整数
$t \geq 0$ 和闭子概形
$$
X_{k'} \supset Z_0 \supset Z_1 \supset \ldots \supset Z_t = \emptyset
$$
使得 $Z_0 = (X_{k'})_{red}$，并且 $Z_i \setminus Z_{i + 1}$ 在 $k'$ 上
对所有 $i$ 都光滑。
\end{lemma}

\begin{proof}
可对 $\dim(X)$ 使用归纳法。由引理
\ref{varieties-lemma-finite-extension-geometrically-reduced}，可找到有限纯不可分扩张
$k'/k$，使得 $(X_{k'})_{red}$ 在 $k'$ 上几何既约。由引理
\ref{varieties-lemma-geometrically-reduced-dense-smooth-open}，存在无处稠密闭子概形
$X' \subset (X_{k'})_{red}$，使得 $(X_{k'})_{red} \setminus X'$ 在 $k'$ 上光滑。
于是 $\dim(X') < \dim(X)$。由归纳假设，存在有限纯不可分扩张 $k''/k'$、整数
$t' \geq 0$ 和闭子概形
$$
X'_{k''} \supset Y_0 \supset Y_1 \supset \ldots \supset Y_{t'} = \emptyset
$$
使得 $Y_0 = (X'_{k''})_{red}$，并且 $Y_i \setminus Y_{i + 1}$ 在 $k''$ 上
对所有 $i$ 都光滑。令 $t = t' + 1$，并考虑
$$
X_{k''} \supset Z_0 \supset Z_1 \supset \ldots \supset Z_t = \emptyset
$$
其中 $Z_0 = (X_{k''})_{red}$，且 $Z_i = Y_{i - 1}$（当 $i > 0$）；这是有意义的，
因为 $X'_{k''}$ 是 $X_{k''}$ 的闭子概形。略去所有所述性质均成立的验证。
\end{proof}


\section{簇的类型}
\label{varieties-section-types}

\noindent
本短节讨论簇的一些初等整体性质。

\begin{definition}
\label{varieties-definition-variety-type}
设 $k$ 是域。设 $X$ 是 $k$ 上的簇。
\begin{enumerate}
\item 若 $X$ 是仿射概形，则称 $X$ 为{\it 仿射簇}。这等价于要求 $X$
同构于 $\mathbf{A}^n_k$ 的闭子概形，其中 $n$ 适当选取。
\item 称 $X$ 为{\it 射影簇}，如果结构态射 $X \to \Spec(k)$ 是射影态射。
由《态射》引理 \ref{morphisms-lemma-characterize-locally-projective}，这成立当且仅当
$X$ 同构于 $\mathbf{P}^n_k$ 的闭子概形，其中 $n$ 适当选取。
\item 称 $X$ 为{\it 拟射影簇}，如果结构态射 $X \to \Spec(k)$ 是拟射影态射。
由《态射》引理 \ref{morphisms-lemma-characterize-locally-quasi-projective}，这成立当且仅当
$X$ 同构于 $\mathbf{P}^n_k$ 的局部闭子概形，其中 $n$ 适当选取。
\item {\it 固有簇}是态射 $X \to \Spec(k)$ 固有的簇。
\item {\it 光滑簇}是态射 $X \to \Spec(k)$ 光滑的簇。
\end{enumerate}
\end{definition}

\noindent
注意，射影簇是固有簇，见《态射》引理
\ref{morphisms-lemma-locally-projective-proper}。此外，仿射簇是拟射影的，因为
$\mathbf{A}^n_k$ 同构于 $\mathbf{P}^n_k$ 的开子概形，见《构造》引理
\ref{constructions-lemma-standard-covering-projective-space}。

\begin{lemma}
\label{varieties-lemma-regular-functions-proper-variety}
设 $X$ 是 $k$ 上的固有簇。则
\begin{enumerate}
\item $K = H^0(X, \mathcal{O}_X)$ 是域，并且是域 $k$ 的有限扩张；
\item 若 $X$ 几何既约，则 $K/k$ 可分；
\item 若 $X$ 几何不可约，则 $K/k$ 纯不可分；
\item 若 $X$ 几何整，则 $K = k$。
\end{enumerate}
\end{lemma}

\begin{proof}
这是引理 \ref{varieties-lemma-proper-geometrically-reduced-global-sections} 的特殊情形。
\end{proof}


\section{正规化}
\label{varieties-section-normalization}

\noindent
下面讨论与正规化有关的一些问题。

\begin{lemma}
\label{varieties-lemma-normalization-locally-algebraic}
设 $k$ 是域。设 $X$ 是 $k$ 上的局部代数概形。设
$\nu : X^\nu \to X$ 是正规化态射，见《态射》定义
\ref{morphisms-definition-normalization}。则
\begin{enumerate}
\item $\nu$ 有限且支配，而 $X^\nu$ 是 $k$ 上正规不可约局部代数概形的不交并；
\item $\nu$ 分解为 $X^\nu \to X_{red} \to X$，且第一个态射是
$X_{red}$ 的正规化态射；
\item 若 $X$ 是既约代数概形，则 $\nu$ 是双有理的；
\item 若 $X$ 是簇，则 $X^\nu$ 是簇，且 $\nu$ 是簇的有限双有理态射。
\end{enumerate}
\end{lemma}

\begin{proof}
由于 $X$ 是域上的局部有限型概形，可知 $X$ 局部 Noether（《态射》引理
\ref{morphisms-lemma-finite-type-noetherian}），因此每个拟紧开集只有有限多个不可约分支
（《性质》引理 \ref{properties-lemma-Noetherian-irreducible-components}）。所以可应用
《态射》定义 \ref{morphisms-definition-normalization}。正规化 $X^\nu$ 总是正规整概形
的不交并，且正规化态射 $\nu$ 总是支配的，见《态射》引理
\ref{morphisms-lemma-normalization-normal}。由于 $X$ 万有 Nagata（《态射》引理
\ref{morphisms-lemma-ubiquity-nagata}），可知 $\nu$ 有限（《态射》引理
\ref{morphisms-lemma-nagata-normalization}）。故 $X^\nu$ 也是局部代数的。至此已证明 (1)。

\medskip\noindent
(2) 是《态射》引理 \ref{morphisms-lemma-normalization-reduced}。

\medskip\noindent
(3) 是《态射》引理 \ref{morphisms-lemma-normalization-birational}。

\medskip\noindent
(4) 由 (1)、(2)、(3) 以及如下事实得出：$X^\nu$ 作为分离概形上的有限概形是分离的。
\end{proof}

\begin{lemma}
\label{varieties-lemma-normalize-projective}
设 $k$ 是域。设 $X$ 是 $k$ 上的固有概形。设
$\nu : X^\nu \to X$ 是正规化态射，见《态射》定义
\ref{morphisms-definition-normalization}。则 $X^\nu$ 在 $k$ 上固有。若 $X$
在 $k$ 上射影，则 $X^\nu$ 在 $k$ 上射影。
\end{lemma}

\begin{proof}
由引理 \ref{varieties-lemma-normalization-locally-algebraic}，态射 $\nu$ 有限。故由《态射》引理
\ref{morphisms-lemma-finite-proper} 和
\ref{morphisms-lemma-composition-proper}，$X^\nu$ 在 $k$ 上固有。
由《态射》引理 \ref{morphisms-lemma-finite-projective}，态射 $\nu$ 是射影的；
所以若 $X$ 在 $k$ 上射影，则由《态射》引理
\ref{morphisms-lemma-composition-projective}，$X^\nu$ 在 $k$ 上射影。
\end{proof}

\begin{lemma}
\label{varieties-lemma-relative-normalization-finite}
设 $k$ 是域。设 $f : Y \to X$ 是 $k$ 上局部代数概形之间的拟紧态射。
设 $X'$ 是 $X$ 在 $Y$ 中的正规化。若 $Y$ 既约，则 $X' \to X$ 有限。
\end{lemma}

\begin{proof}
由于 $Y$ 拟分离（由《性质》引理
\ref{properties-lemma-locally-Noetherian-quasi-separated} 和《态射》引理
\ref{morphisms-lemma-finite-type-noetherian}），态射 $f$ 拟分离（《概形》引理
\ref{schemes-lemma-compose-after-separated}）。因此可应用《态射》定义
\ref{morphisms-definition-normalization-X-in-Y}。结论由《态射》引理
\ref{morphisms-lemma-nagata-normalization-finite-general} 得出。这里用到：局部代数概形
局部 Noether（因而局部只有有限多个不可约分支），并且局部代数概形是 Nagata 的
（《态射》引理 \ref{morphisms-lemma-ubiquity-nagata}）。略去一些小细节。
\end{proof}

\begin{lemma}
\label{varieties-lemma-finite-extension-geometrically-normal}
设 $k$ 是域。设 $X$ 是代数 $k$-概形。则存在有限纯不可分扩张 $k'/k$，
使得正规化 $Y$（取自 $X_{k'}$）在 $k'$ 上几何正规。
\end{lemma}

\begin{proof}
令 $K = k^{perf}$ 为完全闭包。设 $Y_K$ 是 $X_K$ 的正规化，见引理
\ref{varieties-lemma-normalization-locally-algebraic}。由《极限》引理
\ref{limits-lemma-descend-finite-presentation}，存在有限子扩张 $K/k'/k$ 和有限表示态射
$\nu : Y \to X_{k'}$，其到 $K$ 的基变换是正规化态射
$\nu_K : Y_K \to X_K$。注意 $Y$ 在 $k'$ 上几何正规（引理
\ref{varieties-lemma-geometrically-normal}）。增大 $k'$ 后，可假设 $Y \to X_{k'}$ 有限
（《极限》引理 \ref{limits-lemma-descend-finite-finite-presentation}）。由于
$\nu_K : Y_K \to X_K$ 是正规化态射，它诱导双有理态射
$Y_K \to (X_K)_{red}$。故存在稠密开集 $V_K \subset X_K$，使得
$\nu_K^{-1}(V_K) \to V_K$ 是闭浸入（并诱导 $\nu_K^{-1}(V_K)$ 与
$V_{K, red}$ 的同构；例如见《态射》引理
\ref{morphisms-lemma-birational-isomorphism-over-dense-open}）。增大 $k'$ 后，可找到
$V_K$ 是某个稠密开集 $V \subset Y$ 的基变换，并且态射
$\nu^{-1}(V) \to V$ 是闭浸入（《极限》引理
\ref{limits-lemma-descend-opens} 和
\ref{limits-lemma-descend-closed-immersion-finite-presentation}）。由此容易看出
$\nu$ 是正规化态射，证明完成。
\end{proof}

\begin{lemma}
\label{varieties-lemma-normalization-and-change-of-fields}
设 $k$ 是域。设 $X$ 是局部代数 $k$-概形。设 $K/k$ 是域扩张。设
$\nu : X^\nu \to X$ 是 $X$ 的正规化，并设 $Y^\nu \to X_K$ 是基变换的
正规化。那么典范态射
$$
Y^\nu \longrightarrow X^\nu \times_{\Spec(k)} \Spec(K)
$$
在 $K/k$ 可分时是同构，一般情形下是万有同胚。
\end{lemma}

\begin{proof}
令 $Y = X_K$。设 $X^{(0)}$ 和 $Y^{(0)}$ 分别是 $X$ 和 $Y$ 的不可约分支泛点集。
那么投影态射 $\pi : Y \to X$ 满足 $\pi(Y^{(0)}) = X^{(0)}$。这是因为
$\pi$ 满射、开且保泛化，见《态射》引理
\ref{morphisms-lemma-scheme-over-field-universally-open} 和
\ref{morphisms-lemma-open-generizing}。若将 $X^{(0)}$ 和 $Y^{(0)}$ 分别视为（既约）
概形，则 $X^\nu$ 和 $Y^\nu$ 分别是 $X$ 和 $Y$ 在 $X^{(0)}$ 和
$Y^{(0})$ 中的正规化。因此《态射》引理
\ref{morphisms-lemma-functoriality-normalization} 给出典范态射
$Y^\nu \to X^\nu$，它位于 $Y \to X$ 上方；再由纤维积的万有性质给出
本引理中的典范态射。

\medskip\noindent
为证明该态射具有引理所述性质，可假设 $X = \Spec(A)$ 是仿射的。设
$Q(A_{red})$ 为 $A_{red}$ 的总分式环。那么 $X^\nu$ 是整闭包 $A'$ 的谱，
该整闭包取自 $A$ 在 $Q(A_{red})$ 中的整闭包，见《态射》引理
\ref{morphisms-lemma-normalization-reduced} 和
\ref{morphisms-lemma-description-normalization}。类似地，$Y^\nu$ 是整闭包 $B'$ 的谱，
该整闭包取自 $A \otimes_k K$ 在 $Q((A \otimes_k K)_{red})$ 中的整闭包。
有典范映射 $Q(A_{red}) \to Q((A \otimes_k K)_{red})$、典范映射
$A' \to B'$，而引理中的态射对应于诱导映射
$$
A' \otimes_k K \longrightarrow B'
$$
这个 $K$-代数映射的核由幂零元组成，因为
$Q(A_{red}) \otimes_k K \to Q((A \otimes_k K)_{red})$ 的核就是幂零元集合。

\medskip\noindent
若 $K/k$ 可分，则由引理 \ref{varieties-lemma-base-change-normal-by-separable}，
$A' \otimes_k K$ 正规。特别地，它既约，因而
$Q((A \otimes_k K)_{red}) = Q(A' \otimes_k K)$，并且由《代数》引理
\ref{algebra-lemma-characterize-reduced-ring-normal}，
$B' = A' \otimes_k K$。

\medskip\noindent
假设 $K/k$ 不可分。那么 $k$ 的特征为 $p > 0$。我们将证明：对每个
$b \in B'$，都存在幂 $q$，它是 $p$ 的幂，使得 $b^q$ 属于
$A' \otimes_k K$ 的像。
由《代数》引理 \ref{algebra-lemma-p-ring-map}，这将证明上述映射是万有同胚。
给定 $b$，存在子域 $F \subset K$，其中 $F/k$ 有限生成，使得 $b$ 属于
$Q((A \otimes_k F)_{red})$，且整于 $A \otimes_k F$。选取首一多项式
$P = T^d + \alpha_1 T^{d - 1} + \ldots + \alpha _d$，使得
$P(b) = 0$ 且 $\alpha_i \in A \otimes_k F$。选取超越基
$t_1, \ldots, t_r$，它是 $F$ 在 $k$ 上的超越基。设
$F/F'/k(t_1, \ldots, t_r)$ 是最大可分子扩张
（《域》引理 \ref{fields-lemma-separable-first}）。由于 $F/F'$ 有限纯不可分，
存在 $q$，使得 $\lambda^q \in F'$ 对所有 $\lambda \in F$ 都成立。于是 $b^q$
属于 $Q((A \otimes_k F')_{red})$，并满足多项式
$T^d + \alpha_1^q T^{d - 1} + \ldots + \alpha _d^q$，其中
$\alpha_i^q \in A \otimes_k F'$。由可分情形可知
$b^q \in A' \otimes_k F'$，证明完成。
\end{proof}

\begin{lemma}
\label{varieties-lemma-geometrically-normal-in-codim-1}
设 $k$ 是域。设 $X$ 是局部代数 $k$-概形。设 $\nu : X^\nu \to X$ 是
$X$ 的正规化。设 $x \in X$ 是一点，满足：(a) $\mathcal{O}_{X, x}$ 既约；
(b) $\dim(\mathcal{O}_{X, x}) = 1$；并且 (c) 对每个 $x' \in X^\nu$，
若 $\nu(x') = x$，则扩张 $\kappa(x')/k$ 可分。那么 $X$ 在 $x$ 处几何既约，
并且 $X^\nu$ 在 $x'$ 处几何正则，只要 $\nu(x') = x$。
\end{lemma}

\begin{proof}
以下不再特别说明地使用引理 \ref{varieties-lemma-normalization-locally-algebraic} 的结果。
设 $x' \in X^\nu$ 是 $x$ 上方的一点。由维数理论（第
\ref{varieties-section-algebraic-schemes} 节），有
$\dim(\mathcal{O}_{X^\nu, x'}) = 1$。由于 $X^\nu$ 正规，可知
$\mathcal{O}_{X^\nu, x'}$ 是离散赋值环（《性质》引理
\ref{properties-lemma-criterion-normal}）。因此 $\mathcal{O}_{X^\nu, x'}$
是正则局部 $k$-代数，其剩余域在 $k$ 上可分。故
$k \to \mathcal{O}_{X^\nu, x'}$ 在 $\mathfrak m_{x'}$-进拓扑下形式光滑，见
《更多代数》引理 \ref{more-algebra-lemma-regular-implies-fs}。于是由《更多代数》定理
\ref{more-algebra-theorem-regular-fs}，$\mathcal{O}_{X^\nu, x'}$ 在 $k$ 上几何正则。
再由引理 \ref{varieties-lemma-geometrically-regular-at-point}，$X^\nu$ 在 $x'$ 处几何正则。

\medskip\noindent
由于 $\mathcal{O}_{X, x}$ 既约，映射族
$\mathcal{O}_{X, x} \to \mathcal{O}_{X^\nu, x'}$ 是单射。由于
$\mathcal{O}_{X^\nu, x'}$ 是几何既约 $k$-代数，立即可知
$\mathcal{O}_{X, x}$ 是几何既约 $k$-代数。故由引理
\ref{varieties-lemma-geometrically-reduced-at-point}，$X$ 在 $x$ 处几何既约。
\end{proof}


\section{可逆函数群}
\label{varieties-section-units}

\noindent
群 $\mathcal{O}^*(X)/k^*$ 往往（但不总是）是有限生成 Abel 群，只要 $X$ 是 $k$
上的簇。我们用一系列引理证明这一点。一切都建立在下述特殊情形上。

\begin{lemma}
\label{varieties-lemma-open-in-normal-proper}
设 $k$ 是代数闭域。设 $\overline{X}$ 是 $k$ 上的固有簇。设
$X \subset \overline{X}$ 是开子概形。假设 $X$ 正规。则
$\mathcal{O}^*(X)/k^*$ 是有限生成 Abel 群。
\end{lemma}

\begin{proof}
由于命题只涉及 $X$，可以将 $\overline{X}$ 替换为 $k$ 上另一个固有簇。设
$\nu : \overline{X}^\nu \to
\overline{X}$ 是正规化态射。由引理
\ref{varieties-lemma-normalization-locally-algebraic}，$\nu$ 有限，且
$\overline{X}^\nu$ 是簇。由于 $X$ 正规，可知
$\nu^{-1}(X) \to X$ 是同构（略去一个小细节）。最后，因为有限态射是固有的
（《态射》引理 \ref{morphisms-lemma-finite-proper}），且固有态射的复合是固有的
（《态射》引理 \ref{morphisms-lemma-composition-proper}），可知
$\overline{X}^\nu$ 在 $k$ 上固有。因此可以且将假设 $\overline{X}$ 正规。

\medskip\noindent
以下不再特别说明地使用：对任意仿射开集 $U$（它是 $\overline{X}$ 的开集），环
$\mathcal{O}(U)$ 是有限生成 $k$-代数，它是 Noether 环、整环且正规；见《代数》引理
\ref{algebra-lemma-Noetherian-permanence}、《性质》定义
\ref{properties-definition-integral}、《性质》引理
\ref{properties-lemma-locally-Noetherian} 和
\ref{properties-lemma-locally-normal}，以及《态射》引理
\ref{morphisms-lemma-locally-finite-type-characterize}。

\medskip\noindent
设 $\xi_1, \ldots, \xi_r$ 是 $X$ 在 $\overline{X}$ 中之补集的泛点。它们只有有限多个，
因为 $\overline{X}$ 的底拓扑空间是 Noether 的（见《态射》引理
\ref{morphisms-lemma-finite-type-noetherian}、《性质》引理
\ref{properties-lemma-Noetherian-topology} 和《拓扑》引理
\ref{topology-lemma-Noetherian}）。对每个 $i$，局部环
$\mathcal{O}_i = \mathcal{O}_{X, \xi_i}$ 是正规 Noether 局部整环（因为它是 Noether 正规整环的
局部化）。设 $J \subset \{1, \ldots, r\}$ 为所有指标 $i$ 所成的集合，这些指标满足
$\dim(\mathcal{O}_i) = 1$。对 $j \in J$，局部环
$\mathcal{O}_j$ 是离散赋值环，见《代数》引理
\ref{algebra-lemma-characterize-dvr}。因此得到赋值
$$
v_j : k(\overline{X})^* \longrightarrow \mathbf{Z}
$$
并且 $v_j(f) \geq 0 \Leftrightarrow f \in \mathcal{O}_j$。

\medskip\noindent
将 $\mathcal{O}(X)$ 视为一个 $k$-子代数，包含在 $k(X) = k(\overline{X})$ 中。我们断言映射
$$
\mathcal{O}(X)^* \longrightarrow
\prod\nolimits_{j \in J} \mathbf{Z},
\quad
f \longmapsto \prod v_j(f)
$$
的核是 $k^*$。显然，该断言证明引理。事实上，设
$f \in \mathcal{O}(X)$ 是核中的一个元素。设
$U = \Spec(B) \subset \overline{X}$ 是任意仿射开集。则 $B$ 是 Noether 正规整环。
对每个高为一的素理想 $\mathfrak q \subset B$ 及其对应点 $\xi \in X$，可知要么
$\xi = \xi_j$ 对某个 $j \in J$ 成立，要么 $\xi \in X$。理由是，由《性质》引理
\ref{properties-lemma-codimension-local-ring}，
$\text{codim}(\overline{\{\xi\}}, \overline{X}) = 1$；所以若
$\xi \in \overline{X} \setminus X$，则它必是 $\overline{X} \setminus X$ 的泛点，从而等于某个
$\xi_j$，$j \in J$。两种情形下都有
$f \in \mathcal{O}_{X, \xi} = B_{\mathfrak q}$，因为 $f$ 在该映射的核中。因此
$f \in \bigcap_{\text{ht}(\mathfrak q) = 1} B_{\mathfrak q} = B$，见《代数》引理
\ref{algebra-lemma-normal-domain-intersection-localizations-height-1}。换言之，可知
$f \in \Gamma(\overline{X}, \mathcal{O}_{\overline{X}})$。但由于 $k$ 代数闭，由引理
\ref{varieties-lemma-regular-functions-proper-variety} 可得 $f \in k$。
\end{proof}

\noindent
接下来，我们用一些初等论证推广上述情形，仍然保持域代数闭。

\begin{lemma}
\label{varieties-lemma-units-integral-finite-type-algebraically-closed}
设 $k$ 是代数闭域。设 $X$ 是 $k$ 上的局部有限型整概形。则
$\mathcal{O}^*(X)/k^*$ 是有限生成 Abel 群。
\end{lemma}

\begin{proof}
由于 $X$ 是整的，限制映射 $\mathcal{O}(X) \to \mathcal{O}(U)$ 对任意非空开子概形
$U \subset X$ 都是单射。因此可假设 $X$ 仿射。选取闭浸入
$X \to \mathbf{A}^n_k$，并用 $\overline{X}$ 表示 $X$ 在 $\mathbf{P}^n_k$ 中经由通常浸入
$\mathbf{A}^n_k \to \mathbf{P}^n_k$ 所得的闭包。从而可假设 $X$ 是某个射影簇
$\overline{X}$ 的仿射开集。

\medskip\noindent
设 $\nu : \overline{X}^\nu \to \overline{X}$ 是正规化态射，见《态射》定义
\ref{morphisms-definition-normalization}。我们知道 $\nu$ 有限且支配，并且
$\overline{X}^\nu$ 是正规不可约概形，见《态射》引理
\ref{morphisms-lemma-normalization-normal}、
\ref{morphisms-lemma-Japanese-normalization} 和
\ref{morphisms-lemma-ubiquity-nagata}。因此 $\overline{X}^\nu$ 是固有簇，因为
$\overline{X} \to \Spec(k)$ 作为有限态射与固有态射的复合是固有的（见《态射》第
\ref{morphisms-section-proper} 节和第
\ref{morphisms-section-integral} 节中的结果）。还可知 $\nu$ 是满射，因为
$\nu$ 的像是闭的，并包含 $\overline{X}$ 的泛点。因此令
$X^\nu = \nu^{-1}(X)$，可知只需对 $X^\nu$ 证明结论。换言之，可假设
$X$ 是正规固有簇 $\overline{X}$ 的非空开集。该情形由引理
\ref{varieties-lemma-open-in-normal-proper} 处理。
\end{proof}

\noindent
前一引理立即给出下述稍微的推广。

\begin{lemma}
\label{varieties-lemma-units-general-algebraically-closed}
设 $k$ 是代数闭域。设 $X$ 是 $k$ 上局部有限型的连通既约概形，并且只有
有限多个不可约分支。则 $\mathcal{O}^*(X)/k^*$ 是有限生成 Abel 群。
\end{lemma}

\begin{proof}
设 $X = \bigcup X_i$ 是不可约分支的并。由引理
\ref{varieties-lemma-units-integral-finite-type-algebraically-closed}，可知
$\mathcal{O}(X_i)^*/k^*$ 是有限生成 Abel 群。设
$f \in \mathcal{O}(X)^*$ 属于映射
$$
\mathcal{O}(X)^* \longrightarrow \prod \mathcal{O}(X_i)^*/k^*.
$$
的核。则对每个 $i$，存在元素 $\lambda_i \in k$，使得
$f|_{X_i} = \lambda_i$。限制到 $X_i \cap X_j$ 可得 $\lambda_i = \lambda_j$，只要
$X_i \cap X_j \not = \emptyset$。由于 $X$ 连通，可知所有
$\lambda_i$ 相等，因而 $f \in k^*$。这证明了
$$
\mathcal{O}(X)^*/k^* \subset \prod \mathcal{O}(X_i)^*/k^*
$$
而引理得证，因为右端是有限多个有限生成 Abel 群的积。
\end{proof}

\begin{lemma}
\label{varieties-lemma-integral-closure-ground-field}
设 $k$ 是域。设 $X$ 是 $k$ 上连通且既约的概形。则 $k$ 在
$\Gamma(X, \mathcal{O}_X)$ 中的整闭包是域。
\end{lemma}

\begin{proof}
设 $k' \subset \Gamma(X, \mathcal{O}_X)$ 是 $k$ 的整闭包。则
$X \to \Spec(k)$ 经由 $\Spec(k')$ 分解，见《概形》引理
\ref{schemes-lemma-morphism-into-affine}。由于 $X$ 既约，可知 $k'$ 没有非零幂零元。
由于 $k \to k'$ 整，可知 $k'$ 的每个素理想既是极大理想，又是极小素理想，
而 $\Spec(k')$ 全不连通；见《代数》引理
\ref{algebra-lemma-integral-no-inclusion} 和
\ref{algebra-lemma-ring-with-only-minimal-primes}。由于 $X$ 连通，态射
$X \to \Spec(k')$ 为常值态射，设其像为对应于
$\mathfrak p \subset k'$ 的点。那么任意 $f \in k'$ 若满足
$f \not \in \mathfrak p$，其像都是 $\mathcal{O}_X$ 的可逆元。由 $k'$ 的定义，这进而迫使
$f$ 是 $k'$ 的单位。所以 $k'$ 是以 $\mathfrak p$ 为极大理想的局部环，见《代数》引理
\ref{algebra-lemma-characterize-local-ring}。我们已经看到 $k'$ 既约，因此这意味着 $k'$ 是域，
见《代数》引理 \ref{algebra-lemma-minimal-prime-reduced-ring}。
\end{proof}

\begin{proposition}
\label{varieties-proposition-units-general}
设 $k$ 是域。设 $X$ 是 $k$ 上的概形。假设 $X$ 在 $k$ 上局部有限型、连通、既约，
并且只有有限多个不可约分支。若除上述条件外，下列条件之一成立，则
$\mathcal{O}(X)^*/k^*$ 是有限生成 Abel 群：
\begin{enumerate}
\item $k$ 在 $\Gamma(X, \mathcal{O}_X)$ 中的整闭包是 $k$；
\item $X$ 有一个 $k$-有理点；或
\item $X$ 几何整。
\end{enumerate}
\end{proposition}

\begin{proof}
设 $\overline{k}$ 是 $k$ 的代数闭包。设 $Y$ 是
$(X_{\overline{k}})_{red}$ 的一个连通分支。注意，典范态射
$p : Y \to X$ 是开的（由《态射》引理
\ref{morphisms-lemma-scheme-over-field-universally-open}），也是闭的（由《态射》引理
\ref{morphisms-lemma-integral-universally-closed}）。因此 $p(Y) = X$，因为假设 $X$ 连通。
特别地，由于 $X$ 既约，这意味着
$\mathcal{O}(X) \subset \mathcal{O}(Y)$。由引理
\ref{varieties-lemma-galois-action-irreducible-components-locally-finite-type}，可知 $Y$ 只有有限多个不可约分支
（该引理中的域 $\overline{k}$ 是可分代数闭包，但这无关紧要，因为由引理
\ref{varieties-lemma-separably-closed-field-irreducible-components}，两个基变换
$X_{\overline{k}}$ 有相同的不可约分支）。因此引理
\ref{varieties-lemma-units-general-algebraically-closed} 可应用于 $Y$。这意味着：若
$\mathcal{O}(X)^*/k^*$ 不是有限生成 Abel 群，则存在元素
$f \in \mathcal{O}(X)$、$f \not \in k$，其像属于 $\overline{k}$，这里使用映射
$\mathcal{O}(X) \to \mathcal{O}(Y)$。在这种情形下，
$f$ 在 $k$ 上代数，从而整于 $k$。因此若条件 (1) 成立，则这不可能发生。
为完成证明，我们证明条件 (2) 和 (3) 都推出 (1)。

\medskip\noindent
设 $k \subset k' \subset \Gamma(X, \mathcal{O}_X)$ 为 $k$ 在
$\Gamma(X, \mathcal{O}_X)$ 中的整闭包。由引理
\ref{varieties-lemma-integral-closure-ground-field}，可知 $k'$ 是域。若
$e : \Spec(k) \to X$ 是 $k$-有理点，则
$e^\sharp : \Gamma(X, \mathcal{O}_X) \to k$ 是包含映射
$k \to \Gamma(X, \mathcal{O}_X)$ 的一个截面。特别地，$e^\sharp$ 限制到 $k'$ 后是域映射
$k' \to k$，且位于 $k$ 上；这显然表明 (2) 推出 (1)。

\medskip\noindent
若整闭包 $k'$，即 $k$ 在 $\Gamma(X, \mathcal{O}_X)$ 中的整闭包，非平凡，则可知：
要么 $X$ 不几何连通（当 $k'/k$ 不是纯不可分的时候），要么 $X$ 不几何既约
（当 $k'/k$ 是非平凡纯不可分的时候）。略去细节。因此 (3) 推出 (1)。
\end{proof}

\begin{lemma}
\label{varieties-lemma-units-variety}
设 $k$ 是域。设 $X$ 是 $k$ 上的簇。若下列条件中至少一个成立，则群
$\mathcal{O}(X)^*/k^*$ 是有限生成 Abel 群：
\begin{enumerate}
\item $k$ 在 $\Gamma(X, \mathcal{O}_X)$ 中整闭；
\item $k$ 在 $k(X)$ 中代数闭；
\item $X$ 在 $k$ 上几何整；或
\item $k$ 是域扩张 $\kappa(x)/k$ 的“交”，其中 $x$ 遍历 $x$ 的闭点。
\end{enumerate}
\end{lemma}

\begin{proof}
由命题 \ref{varieties-proposition-units-general} 可知 (1) 足够。我们略去 (2)、(3)、(4) 各自推出 (1) 的验证。
\end{proof}


\section{Künneth 公式，I}
\label{varieties-section-kunneth}

\noindent
本节证明底是域且我们考虑拟凝聚模上同调时的 Künneth 公式。
更一般的版本见《概形的导出范畴》第
\ref{perfect-section-kunneth} 节。

\begin{lemma}
\label{varieties-lemma-kunneth}
设 $k$ 是域。设 $X$ 和 $Y$ 是 $k$ 上的概形，并设 $\mathcal{F}$ 和
$\mathcal{G}$ 分别是拟凝聚 $\mathcal{O}_X$-模和 $\mathcal{O}_Y$-模。则有典范同构
$$
H^n(X \times_{\Spec(k)} Y, \text{pr}_1^*\mathcal{F}
\otimes_{\mathcal{O}_{X \times_{\Spec(k)} Y}} \text{pr}_2^*\mathcal{G}) =
\bigoplus\nolimits_{p + q = n}
H^p(X, \mathcal{F}) \otimes_k H^q(Y, \mathcal{G})
$$
只要 $X$ 和 $Y$ 拟紧且具有仿射对角态射\footnote{
$X$ 和 $Y$ 拟分离的情形将在下文引理
\ref{varieties-lemma-kunneth-general} 中讨论。}（例如 $X$ 和 $Y$ 都分离）。
\end{lemma}

\begin{proof}
在本证明中，没有额外标注的积与张量积都取在 $k$ 上。作为映射
$$
H^p(X, \mathcal{F}) \otimes H^q(Y, \mathcal{G})
\longrightarrow
H^n(X \times Y, \text{pr}_1^*\mathcal{F}
\otimes_{\mathcal{O}_{X \times Y}} \text{pr}_2^*\mathcal{G})
$$
我们先利用上同调的函子性得到映射
$H^p(X, \mathcal{F}) \to H^p(X \times Y, \text{pr}_1^*\mathcal{F})$ 和
$H^p(Y, \mathcal{G}) \to H^p(X \times Y, \text{pr}_2^*\mathcal{G})$，
然后使用杯积
$$
\cup :
H^p(X \times Y, \text{pr}_1^*\mathcal{F})
\otimes H^q(X \times Y, \text{pr}_2^*\mathcal{G})
\longrightarrow
H^n(X \times Y, \text{pr}_1^*\mathcal{F}
\otimes_{\mathcal{O}_{X \times Y}} \text{pr}_2^*\mathcal{G})
$$
当 $X$ 和 $Y$ 仿射时，结论由仿射概形上高次上同调群的消失
（《概形的上同调》引理
\ref{coherent-lemma-quasi-coherent-affine-cohomology-zero}）以及定义
（拟凝聚模的拉回和拟凝聚模的张量积的定义）得出。

\medskip\noindent
选取有限仿射开覆盖
$\mathcal{U} : X = \bigcup_{i \in I} U_i$ 和
$\mathcal{V} : Y = \bigcup_{j \in J} V_j$。它们确定仿射开覆盖
$\mathcal{W} : X \times Y = \bigcup_{(i, j) \in I \times J} U_i \times V_j$。
注意，$\mathcal{W}$ 同时是 $\text{pr}_1^{-1}\mathcal{U}$ 和
$\text{pr}_2^{-1}\mathcal{V}$ 的加细。所以由《上同调》引理
\ref{cohomology-lemma-functoriality-cech}，得到映射
$$
\check{\mathcal{C}}^\bullet(\mathcal{U}, \mathcal{F}) \to
\check{\mathcal{C}}^\bullet(\mathcal{W}, \text{pr}_1^*\mathcal{F})
\quad\text{且}\quad
\check{\mathcal{C}}^\bullet(\mathcal{V}, \mathcal{G}) \to
\check{\mathcal{C}}^\bullet(\mathcal{W}, \text{pr}_2^*\mathcal{G})
$$
它们与上同调上的拉回映射相容。在《上同调》等式
(\ref{cohomology-equation-needs-signs}) 中，我们已构造复形映射
$$
\text{Tot}(
\check{\mathcal{C}}^\bullet(\mathcal{W}, \text{pr}_1^*\mathcal{F})
\otimes
\check{\mathcal{C}}^\bullet(\mathcal{W}, \text{pr}_2^*\mathcal{G}))
\longrightarrow
\check{\mathcal{C}}^\bullet(\mathcal{W},
\text{pr}_1^*\mathcal{F} \otimes_{\mathcal{O}_{X \times Y}}
\text{pr}_2^*\mathcal{G})
$$
它定义上同调上的杯积。结合上述映射，得到复形映射
\begin{equation}
\label{varieties-equation-kunneth-on-cech}
\text{Tot}(
\check{\mathcal{C}}^\bullet(\mathcal{U}, \mathcal{F})
\otimes
\check{\mathcal{C}}^\bullet(\mathcal{V}, \mathcal{G}))
\longrightarrow
\check{\mathcal{C}}^\bullet(\mathcal{W},
\text{pr}_1^*\mathcal{F} \otimes_{\mathcal{O}_{X \times Y}}
\text{pr}_2^*\mathcal{G})
\end{equation}
我们提醒读者，该映射不是复形的同构。回忆，可用这些覆盖计算我们的拟凝聚层的上同调
（《概形的上同调》引理
\ref{coherent-lemma-affine-diagonal} 和
\ref{coherent-lemma-cech-cohomology-quasi-coherent}）。因此，
(\ref{varieties-equation-kunneth-on-cech}) 在上同调上重现本引理的映射。

\medskip\noindent
考虑拟凝聚模的短正合列
$0 \to \mathcal{F} \to \mathcal{F}' \to \mathcal{F}'' \to 0$。由于
(\ref{varieties-equation-kunneth-on-cech}) 的构造对 $\mathcal{F}$ 具有函子性，并且相关
{\v C}ech 复形的构造对变量 $\mathcal{F}$ 正合（因为我们在仿射开集上取截面），
得到短正合复形列之间的映射
\begingroup\scriptsize
$$
\xymatrix{
\text{Tot}(
\check{\mathcal{C}}^\bullet(\mathcal{U}, \mathcal{F})
\otimes
\check{\mathcal{C}}^\bullet(\mathcal{V}, \mathcal{G})) \ar[r] \ar[d] &
\text{Tot}(
\check{\mathcal{C}}^\bullet(\mathcal{U}, \mathcal{F}')
\otimes
\check{\mathcal{C}}^\bullet(\mathcal{V}, \mathcal{G})) \ar[r] \ar[d] &
\text{Tot}(
\check{\mathcal{C}}^\bullet(\mathcal{U}, \mathcal{F}'')
\otimes
\check{\mathcal{C}}^\bullet(\mathcal{V}, \mathcal{G})) \ar[d] \\
\check{\mathcal{C}}^\bullet(\mathcal{W},
\text{pr}_1^*\mathcal{F} \otimes_{\mathcal{O}_{X \times Y}}
\text{pr}_2^*\mathcal{G}) \ar[r] &
\check{\mathcal{C}}^\bullet(\mathcal{W},
\text{pr}_1^*\mathcal{F}' \otimes_{\mathcal{O}_{X \times Y}}
\text{pr}_2^*\mathcal{G}) \ar[r] &
\check{\mathcal{C}}^\bullet(\mathcal{W},
\text{pr}_1^*\mathcal{F}'' \otimes_{\mathcal{O}_{X \times Y}}
\text{pr}_2^*\mathcal{G})
}
$$
\endgroup
（我们略去了外侧的零项）。考察上同调长正合列，可知：若引理的结论对 $2$-out-of-$3$ of
$\mathcal{F}, \mathcal{F}', \mathcal{F}''$ 成立，则它对第三个也成立。

\medskip\noindent
注意，$X$ 对拟凝聚模具有有限上同调维数，见《概形的上同调》引理
\ref{coherent-lemma-vanishing-nr-affines}。对
$d(\mathcal{F}) = \max \{d \mid H^d(X, \mathcal{F}) \not = 0\}$
作归纳，我们将化归到 $d(\mathcal{F}) = 0$ 的情形。假设
$d(\mathcal{F}) > 0$。由《概形的上同调》引理
\ref{coherent-lemma-affine-diagonal-universal-delta-functor}，我们已经看到，存在嵌入
$\mathcal{F} \to \mathcal{F}'$，使得 $H^p(X, \mathcal{F}') = 0$ 对所有 $p \geq 1$ 成立。令
$\mathcal{F}'' = \Coker(\mathcal{F} \to \mathcal{F}')$，则
$d(\mathcal{F}'') < d(\mathcal{F})$。因此可用前段的结果看出，只需对
$\mathcal{F}'$ 和 $\mathcal{F}''$ 证明引理，由此完成归纳步骤。

\medskip\noindent
对 $\mathcal{G}$ 作同样的论证，可假设 $\mathcal{F}$ 和 $\mathcal{G}$ 都只在次数
$0$ 具有非零上同调。设 $V \subset Y$ 是仿射开集。考虑仿射开覆盖
$\mathcal{U}_V : X \times V = \bigcup_{i \in I} U_i \times V$。立即有
$$
\check{\mathcal{C}}^\bullet(\mathcal{U}, \mathcal{F}) \otimes \mathcal{G}(V) =
\check{\mathcal{C}}^\bullet(\mathcal{U}_V,
\text{pr}_1^*\mathcal{F} \otimes_{\mathcal{O}_{X \times Y}}
\text{pr}_2^*\mathcal{G})
$$
（复形相等）。我们得出
$$
R\text{pr}_{2, *}(\text{pr}_1^*\mathcal{F} \otimes_{\mathcal{O}_{X \times Y}}
\text{pr}_2^*\mathcal{G})
\cong
\Gamma(X, \mathcal{F}) \otimes_k \mathcal{G} \cong
\bigoplus\nolimits_{\alpha \in A} \mathcal{G}
$$
在 $Y$ 上成立。这里 $A$ 是 $k$-向量空间 $\Gamma(X, \mathcal{F})$ 的一组基。
$Y$ 上的上同调与直和可交换（《上同调》引理
\ref{cohomology-lemma-quasi-separated-cohomology-colimit}）。对 $\text{pr}_2$ 使用 Leray 谱序列
（经由《上同调》引理 \ref{cohomology-lemma-apply-Leray}），我们得出，
$H^n(X \times Y, \text{pr}_1^*\mathcal{F} \otimes_{\mathcal{O}_{X \times Y}}
\text{pr}_2^*\mathcal{G})$
在 $n > 0$ 时为零，且同构于
$H^0(X, \mathcal{F}) \otimes H^0(Y, \mathcal{G})$，当 $n = 0$ 时。这完成了证明
（除了还应检查该同构确实由次数 $0$ 上的杯积给出；我们略去验证）。
\end{proof}

\begin{lemma}
\label{varieties-lemma-kunneth-general}
设 $k$ 是域。设 $X$ 和 $Y$ 是 $k$ 上的概形，并设 $\mathcal{F}$ 和
$\mathcal{G}$ 分别是拟凝聚 $\mathcal{O}_X$-模和 $\mathcal{O}_Y$-模。则有典范同构
$$
H^n(X \times_{\Spec(k)} Y, \text{pr}_1^*\mathcal{F}
\otimes_{\mathcal{O}_{X \times_{\Spec(k)} Y}} \text{pr}_2^*\mathcal{G}) =
\bigoplus\nolimits_{p + q = n}
H^p(X, \mathcal{F}) \otimes_k H^q(Y, \mathcal{G})
$$
只要 $X$ 和 $Y$ 拟紧且拟分离。
\end{lemma}

\begin{proof}
若 $X$ 和 $Y$ 分离，或更一般地具有仿射对角态射，请参见引理
\ref{varieties-lemma-kunneth} 中的“更好”证明（它胜过本证明之处在于：它把这些映射识别为先拉回再作杯积）。
设 $X'$ 和 $Y'$ 分别是 $X$ 和 $Y$ 的无穷小加厚，其结构层分别为
$\mathcal{O}_{X'} = \mathcal{O}_X \oplus \mathcal{F}$
和 $\mathcal{O}_{Y'} = \mathcal{O}_Y \oplus \mathcal{G}$，其中 $\mathcal{F}$ 和
$\mathcal{G}$ 分别是平方为零的理想。那么
$$
\mathcal{O}_{X' \times Y'} =
\mathcal{O}_{X \times Y}
\oplus \text{pr}_1^*\mathcal{F}
\oplus \text{pr}_2^*\mathcal{G}
\oplus \text{pr}_1^*\mathcal{F}
\otimes_{\mathcal{O}_{X \times Y}} \text{pr}_2^*\mathcal{G}
$$
作为 $X \times Y$ 上的层成立。由此可知，只需证明
$$
H^n(X \times Y, \mathcal{O}_{X \times Y}) =
\bigoplus\nolimits_{p + q = n}
H^p(X, \mathcal{O}_X) \otimes_k H^q(Y, \mathcal{O}_Y)
$$
对 $k$ 上任意一对拟紧拟分离概形成立。略去一些细节。

\medskip\noindent
为证明此断言，我们使用《概形的上同调》引理
\ref{coherent-lemma-hypercoverings} 所给出形式的上同调与基变换。该引理告诉我们，存在一个下有界的
$k$-向量空间复形，也就是复形 $\mathcal{K}^\bullet$，它由 $\Spec(k)$ 上的拟凝聚模组成，
它万有地计算 $Y$ 在 $\Spec(k)$ 上的上同调。特别地，可知
$$
R\text{pr}_{1, *}(\mathcal{O}_{X \times Y}) \cong
(X \to \Spec(k))^*\mathcal{K}^\bullet
$$
在 $D(\mathcal{O}_X)$ 中成立。在同伦意义下，复形 $\mathcal{K}^\bullet$ 同构于
$\bigoplus_{q \geq 0} H^q(Y, \mathcal{O}_Y)[-q]$，因为域上的每个向量空间复形都如此。
我们得出
$$
R\text{pr}_{1, *}(\mathcal{O}_{X \times Y}) \cong
\bigoplus\nolimits_{q \geq 0}
H^q(Y, \mathcal{O}_Y)[-q] \otimes_k \mathcal{O}_X
$$
在 $D(\mathcal{O}_X)$ 中成立。于是
\begin{align*}
R\Gamma(X \times Y, \mathcal{O}_{X \times Y})
& =
R\Gamma(X, R\text{pr}_{1, *}(\mathcal{O}_{X \times Y})) \\
& =
R\Gamma(X, \bigoplus\nolimits_{q \geq 0}
H^q(Y, \mathcal{O}_Y)[-q] \otimes_k \mathcal{O}_X) \\
& =
\bigoplus\nolimits_{q \geq 0}
R\Gamma(X,  H^q(Y, \mathcal{O}_Y) \otimes \mathcal{O}_X)[-q] \\
& =
\bigoplus\nolimits_{q \geq 0}
R\Gamma(X, \mathcal{O}_X) \otimes_k H^q(Y, \mathcal{O}_Y)[-q] \\
& =
\bigoplus\nolimits_{p, q \geq 0}
H^p(X, \mathcal{O}_X)[-p] \otimes_k H^q(Y, \mathcal{O}_Y)[-q]
\end{align*}
如所求。第一个等式由 $\text{pr}_1$ 的 Leray 谱序列得出
（《上同调》引理 \ref{cohomology-lemma-before-Leray}）。
第二个等式由上面给出的总直像分解得出。
第三个等式成立，因为上同调总与有限直和可交换（并且由《概形的上同调》引理
\ref{coherent-lemma-vanishing-nr-affines-quasi-separated}，$Y$ 的上同调在充分高次数上消失）。
第四个等式成立，因为由《上同调》引理
\ref{cohomology-lemma-quasi-separated-cohomology-colimit}，$X$ 上的上同调与无限直和可交换。
最后一个等式由上面对域的导出范畴的评论得出。
\end{proof}

\section{簇的 Picard 群}
\label{varieties-section-picard-groups}

\noindent
本节汇集代数簇的 Picard 群的一些初等结果。

\begin{lemma}
\label{varieties-lemma-change-rings-pic-pre}
设 $A \to B$ 是忠实平坦环同态。设 $X$ 是 $A$ 上的拟紧拟分离概形。
设 $\mathcal{L}$ 是可逆 $\mathcal{O}_X$-模，其到 $X_B$ 的拉回平凡。则
$H^0(X, \mathcal{L})$ 和 $H^0(X, \mathcal{L}^{\otimes -1})$ 是可逆
$H^0(X, \mathcal{O}_X)$-模，并且乘法映射诱导同构
$$
H^0(X, \mathcal{L}) \otimes_{H^0(X, \mathcal{O}_X)}
H^0(X, \mathcal{L}^{\otimes -1}) \longrightarrow
H^0(X, \mathcal{O}_X)
$$
\end{lemma}

\begin{proof}
以 $\mathcal{L}_B$ 表示 $\mathcal{L}$ 到 $X_B$ 的拉回。
选取同构 $\mathcal{L}_B \to \mathcal{O}_{X_B}$。
令 $R = H^0(X, \mathcal{O}_X)$、$M = H^0(X, \mathcal{L})$，并将
$M$ 视为 $R$-模。对每个拟凝聚 $\mathcal{O}_X$-模 $\mathcal{F}$，记
$\mathcal{F}_B$ 为它在 $X_B$ 上的拉回，则有典范同构
$H^0(X_B, \mathcal{F}_B) = H^0(X, \mathcal{F}) \otimes_A B$；见
《概形的上同调》引理 \ref{coherent-lemma-flat-base-change-cohomology}。
因此有
$$
M \otimes_R (R \otimes_A B) =
M \otimes_A B = H^0(X_B, \mathcal{L}_B) \cong
H^0(X_B, \mathcal{O}_{X_B}) = R \otimes_A B
$$
由于 $R \to R \otimes_A B$ 忠实平坦（它是忠实平坦同态
$A \to B$ 的基变换），由《代数》命题
\ref{algebra-proposition-ffdescent-finite-projectivity} 可知，$M$ 是可逆
$R$-模。类似地，$N = H^0(X, \mathcal{L}^{\otimes -1})$ 是可逆 $R$-模。
要看关于张量积的断言成立，只需利用它在拉回到 $X_B$ 后成立，以及
$R \to R \otimes_A B$ 的忠实平坦性。略去一些细节。
\end{proof}

\begin{lemma}
\label{varieties-lemma-change-rings-pic}
设 $A \to B$ 是忠实平坦环同态。设 $X$ 是 $A$ 上的概形，并且
\begin{enumerate}
\item $X$ 拟紧且拟分离；
\item $R = H^0(X, \mathcal{O}_X)$ 是半局部环。
\end{enumerate}
则拉回映射 $\Pic(X) \to \Pic(X_B)$ 是单射。
\end{lemma}

\begin{proof}
设 $\mathcal{L}$ 是可逆 $\mathcal{O}_X$-模，其拉回 $\mathcal{L}'$ 在
$X_B$ 上平凡。令 $M = H^0(X, \mathcal{L})$ 和
$N = H^0(X, \mathcal{L}^{\otimes - 1})$。由引理
\ref{varieties-lemma-change-rings-pic-pre}，$R$-模 $M$ 和 $N$ 可逆。由于 $R$ 半局部，
$M \cong R$ 且 $N \cong R$；见《代数》引理
\ref{algebra-lemma-locally-free-semi-local-free}。选取生成元 $s \in M$ 和
$t \in N$。由引理 \ref{varieties-lemma-change-rings-pic-pre} 的最后一部分，
$st \in R = H^0(X, \mathcal{O}_X)$ 是单位。我们得出，$s$ 和 $t$ 分别给出
$\mathcal{L}$ 和 $\mathcal{L}^{\otimes -1}$ 在 $X$ 上的平凡化。
\end{proof}

\begin{lemma}
\label{varieties-lemma-change-fields-pic}
设 $k'/k$ 是域扩张。设 $X$ 是 $k$ 上的概形，并且
\begin{enumerate}
\item $X$ 拟紧且拟分离；
\item $R = H^0(X, \mathcal{O}_X)$ 是半局部环，例如当 $\dim_k R < \infty$ 时如此。
\end{enumerate}
则拉回映射 $\Pic(X) \to \Pic(X_{k'})$ 是单射。
\end{lemma}

\begin{proof}
这是引理 \ref{varieties-lemma-change-rings-pic} 的特殊情形。若 $\dim_k R < \infty$，则
$R$ 是 Artin 环，因而是半局部环（《代数》引理
\ref{algebra-lemma-finite-dimensional-algebra} 和
\ref{algebra-lemma-artinian-finite-nr-max}）。
\end{proof}

\begin{example}
\label{varieties-example-need-condition}
若不对概形 $X$（视为域 $k$ 上的概形）施加某种条件，引理
\ref{varieties-lemma-change-fields-pic} 并不成立。下面给出一个例子。
设 $k$ 是域。设 $t \in \mathbf{P}^1_k$ 是闭点。令
$X = \mathbf{P}^1 \setminus \{t\}$。则有满射
$$
\mathbf{Z} = \Pic(\mathbf{P}^1_k) \longrightarrow \Pic(X)
$$
第一个等式由《除子》引理 \ref{divisors-lemma-Pic-projective-space-UFD} 得出，
满射性由《除子》引理 \ref{divisors-lemma-extend-invertible-module} 得出
（因为 $\mathbf{P}^1_k$ 的维数为 $1$，且它在 $k$ 上光滑，故其所有局部环都是离散赋值环）。
若 $\mathcal{L}$ 位于所示映射的核中，则有
$\mathcal{L} \cong \mathcal{O}_{\mathbf{P}^1_k}(nt)$，其中 $n \in \mathbf{Z}$。我们把下述验证留给读者：
$\mathcal{O}_{\mathbf{P}^1_k}(t) \cong \mathcal{O}_{\mathbf{P}^1_k}(d)$，
其中 $d = [\kappa(t) : k]$。因此
$$
\Pic(X) = \mathbf{Z}/d\mathbf{Z}
$$
所以，若 $t$ 不是 $k$-有理点，则 $d > 1$，这个 Picard 群非零。
另一方面，若把底域 $k$ 扩张为任意域扩张 $k'$，且存在 $k$-嵌入
$\kappa(t) \to k'$，则 $\mathbf{P}^1_{k'} \setminus X_{k'}$ 有一个
$k'$-有理点 $t'$。因此
$\mathcal{O}_{\mathbf{P}^1_{k'}}(1) = \mathcal{O}_{\mathbf{P}^1_{k'}}(t')$
属于映射 $\mathbf{Z} \to \Pic(X_{k'})$ 的核；按与上面相同的方法可得
$\Pic(X_{k'}) = 0$。
\end{example}

\noindent
下列引理告诉我们，在正规簇上，“有理等价的可逆模”彼此同构。

\begin{lemma}
\label{varieties-lemma-rational-equivalence-for-Pic}
设 $k$ 是域。设 $X$ 是 $k$ 上的正规簇。设
$U \subset \mathbf{A}^n_k$ 是开子概形，并具有 $k$-有理点
$p, q \in U(k)$。对每个可逆模 $\mathcal{L}$（定义在
$X \times_{\Spec(k)} U$ 上），
其限制 $\mathcal{L}|_{X \times p}$ 和 $\mathcal{L}|_{X \times q}$ 同构。
\end{lemma}

\begin{proof}
$X \times_{\Spec(k)} U \to X$ 的纤维是域上的仿射 $n$-空间的开子概形。
因此，由《除子》引理 \ref{divisors-lemma-open-subscheme-UFD}，这些纤维的
Picard 群平凡。应用《除子》引理 \ref{divisors-lemma-in-image-pullback}，
可见 $\mathcal{L}$ 是某个可逆模 $\mathcal{N}$ 的拉回，而该模定义在 $X$ 上。
\end{proof}

\section{底域的唯一性}
\label{varieties-section-base-field}

\noindent
“设 $X$ 是 $k$ 上的概形”这一说法意指 $X$ 是一个概形，并配备态射
$X \to \Spec(k)$。现在可以问，域 $k$ 是否由概形 $X$ 唯一确定。
当然并非如此。例如，我们通常把 $\mathbf{A}^1_{\mathbf{C}}$ 看作复数域
$\mathbf{C}$ 上的概形，但也可以把它看作 $\mathbf{Q}$ 上的概形。然而，若要求态射
$X \to \Spec(k)$ 不通过
$X \to \Spec(k') \to \Spec(k)$ 分解，其中 $k'/k$ 是任意非平凡域扩张，
情形又如何？换言之，我们要求 $k$ 在某种意义下是使 $X$ 定义在 $k$ 上的极大域。

\medskip\noindent
下列例子表明，这仍不能保证 $k$ 的唯一性。考虑概形
$$
X =
\Spec\left(
\mathbf{Q}(x)[y]\left[\frac{1}{P(y)}, P \in \mathbf{Q}[y], P \not = 0\right]
\right)
$$
乍看之下，它似乎是 $\mathbf{Q}(x)$ 上的概形；再看便会发现，它也是
$\mathbf{Q}(y)$ 上的概形。此外，域 $\mathbf{Q}(x)$ 和 $\mathbf{Q}(y)$ 都是
$R = \Gamma(X, \mathcal{O}_X)$ 的子域，并且在 $R$ 的所有子域中极大
（略去细节）。特别地，$\mathbf{Q}(x)$ 和 $\mathbf{Q}(y)$ 在上述意义下都是极大的。
注意，态射 $X \to \Spec(\mathbf{Q}(x))$ 和
$X \to \Spec(\mathbf{Q}(y))$ 都“本质有限型”
（即对应的环同态本质有限型）。因此 $X$ 是有限维 Noether 概形，也就是说，
它并非完全病态。

\medskip\noindent
另一个可能妨碍唯一性的问题是，概形 $X$ 可能非既约。在这种情形，从 $X$ 到
给定域的谱可以有许多不同的态射。作为一个显式例子，考虑对偶数
$D = \mathbf{C}[y]/(y^2) = \mathbf{C} \oplus \epsilon \mathbf{C}$。
给定任意导子 $\theta : \mathbf{C} \to \mathbf{C}$（定义在 $\mathbf{Q}$ 上），便得到环同态
$$
\mathbf{C} \longrightarrow D, \quad
c \longmapsto c + \epsilon \theta(c).
$$
$\mathbf{C}$ 中使所有这些映射都相同的子域是 $\mathbf{Q}$ 的代数闭包。
这意味着，若 $X$ 非既约，把 $X$ 能定义于其上的所有域取交，所得可能是很小的域。

\medskip\noindent
与此有关的一个观察如下：给定域 $k$ 以及两个子域 $k_1, k_2$（它们都是 $k$ 的子域），
若 $k$ 在 $k_1$ 和 $k_2$ 上都有限，一般而言 $k$ 在 $k_1 \cap k_2$ 上
{\it 并不}有限。例如，取域 $k = \mathbf{Q}(t)$ 及其子域
$k_1 = \mathbf{Q}(t^2)$ 和 $\mathbf{Q}((t + 1)^2)$。在此情形，
$k_1 \cap k_2 = \mathbf{Q}$。因此在下文取域的交时必须谨慎。

\medskip\noindent
说完这些，我们现在证明：若 $X$ 在某个域上局部有限型，则仍有某种唯一性。
精确结果如下。

\begin{proposition}
\label{varieties-proposition-unique-base-field}
设 $X$ 是概形。设 $a : X \to \Spec(k_1)$ 和
$b : X \to \Spec(k_2)$ 是从 $X$ 到域的谱的态射。
假设 $a, b$ 局部有限型，并且 $X$ 既约且连通。则
$k_1' = k_2'$，其中 $k_i' \subset \Gamma(X, \mathcal{O}_X)$ 是
$k_i$ 在 $\Gamma(X, \mathcal{O}_X)$ 中的整闭包。
\end{proposition}

\begin{proof}
首先假设该引理在 $X$ 拟紧的情形成立（我们将在下面处理拟紧情形）。
由于 $X$ 在域上局部有限型，它局部 Noether；见《态射》引理
\ref{morphisms-lemma-finite-type-noetherian}。特别地，这意味着它局部连通，
开子集的连通分支是开集，并且拟紧开集的交拟紧；见《性质》引理
\ref{properties-lemma-Noetherian-topology}、《拓扑》引理
\ref{topology-lemma-locally-connected}、《拓扑》第
\ref{topology-section-noetherian} 节以及《拓扑》引理
\ref{topology-lemma-constructible-Noetherian-space}。
选取开覆盖 $X = \bigcup_{i \in I} U_i$，使每个 $U_i$ 都拟紧且连通。
对每个 $i$，令 $K_i \subset \mathcal{O}_X(U_i)$ 为 $k_1$ 和 $k_2$ 的整闭包。
对每对 $i, j \in I$，把
$$
U_i \cap U_j = \coprod U_{i, j, l}
$$
分解为其有限多个连通分支。记
$K_{i, j, l} \subset \mathcal{O}(U_{i, j, l})$ 为 $k_1$ 和 $k_2$ 的整闭包。
由引理 \ref{varieties-lemma-integral-closure-ground-field}，环 $K_i$ 和
$K_{i, j, l}$ 都是域。现在我们断言，$k_1'$ 和 $k_2'$ 都等于下列映射的核：
$$
\prod K_i \longrightarrow \prod K_{i, j, l}, \quad
(x_i)_i \longmapsto x_i|_{U_{i, j, l}} - x_j|_{U_{i, j, l}}
$$
这便证明了所求。的确，显然 $k_1'$ 包含于这个核。反过来，设
$(x_i)_i$ 位于该核中。由层条件，$(x_i)_i$ 对应于
$f \in \mathcal{O}(X)$。选取某个 $i_0 \in I$，并令 $P(T) \in k_1[T]$
为满足 $P(x_{i_0}) = 0$ 的首一多项式。于是我们断言 $P(f) = 0$；这证明了
$f \in k_1$。为证明这一点，须证明 $P(x_i) = 0$ 对所有 $i$ 成立。
选取 $i \in I$。由于 $X$ 连通，存在序列
$i_0, i_1, \ldots, i_n = i \in I$，使得
$U_{i_t} \cap U_{i_{t + 1}} \not = \emptyset$。这意味着，对每个 $t$，存在
$l_t$，使 $x_{i_t}$ 和 $x_{i_{t + 1}}$ 映到域
$K_{i_t, i_{t + 1}, l_t}$ 的同一个元素。因此若 $P(x_{i_t}) = 0$，则
$P(x_{i_{t + 1}}) = 0$。从 $P(x_{i_0}) = 0$ 出发归纳，便得到
$P(x_i) = 0$，如所求。

\medskip\noindent
为完成该引理的证明，我们在附加假设 $X$ 拟紧的条件下证明该引理。
由引理 \ref{varieties-lemma-integral-closure-ground-field}，把 $k_i$ 替换为 $k_i'$ 后，
可以假设 $k_i$ 在 $\Gamma(X, \mathcal{O}_X)$ 中整闭。这蕴涵
$\mathcal{O}(X)^*/k_i^*$ 是有限生成 Abel 群；见命题
\ref{varieties-proposition-units-general}。令 $k_{12} = k_1 \cap k_2$，视其为
$\mathcal{O}(X)$ 的子环。注意，$k_{12}$ 是域。由于有映射
$$
k_1^*/k_{12}^* \longrightarrow \mathcal{O}(X)^*/k_2^*
$$
可见 $k_1^*/k_{12}^*$ 也是有限生成 Abel 群。因此存在
$\alpha_1, \ldots, \alpha_n \in k_1^*$，使每个元素 $\lambda \in k_1$ 都具有形式
$$
\lambda = c \alpha_1^{e_1} \ldots \alpha_n^{e_n}
$$
其中 $e_i \in \mathbf{Z}$ 且 $c \in k_{12}$。特别地，环同态
$$
k_{12}[x_1, \ldots, x_n, \frac{1}{x_1 \ldots x_n}] \longrightarrow k_1, \quad
x_i \longmapsto \alpha_i
$$
是满射。由 Hilbert 零点定理，即《代数》定理
\ref{algebra-theorem-nullstellensatz}，可知 $k_1$ 是 $k_{12}$ 的有限扩张。
同理可知 $k_2$ 是 $k_{12}$ 的有限扩张。特别地，$k_1$ 和 $k_2$ 都包含于
整闭包 $k_{12}'$；这是 $k_{12}$ 在 $\Gamma(X, \mathcal{O}_X)$ 中的整闭包。但由引理
\ref{varieties-lemma-integral-closure-ground-field}，$k_{12}'$ 是域；并且由于我们选取
$k_i$ 在 $\Gamma(X, \mathcal{O}_X)$ 中整闭，故得到
$k_1 = k_{12} = k_2$，如所求。
\end{proof}

\section{自同构}
\label{varieties-section-automorphisms}

\noindent
本节讨论域上概形的自同构。关于曲线的（无穷小）自同构的一些信息，见
《代数曲线》第 \ref{curves-section-vector-fields} 节和
《曲线的模》第 \ref{moduli-curves-section-finite-aut} 节。

\begin{lemma}
\label{varieties-lemma-automorphism}
设 $X$ 是域 $k$ 上有限型的既约概形。设 $f : X \to X$ 是 $k$ 上的自同构，
并且它在 $X$ 的底拓扑空间上诱导恒等映射。则
\begin{enumerate}
\item 有 $f^*\mathcal{F} \cong \mathcal{F}$，对任意凝聚
$\mathcal{O}_X$-模皆然；
\item 若每个不可约分支的维数都满足 $\dim(Z) > 0$（其中 $Z \subset X$），
则 $f$ 是恒等映射。
\end{enumerate}
\end{lemma}

\begin{proof}
(1) 由 (2) 以及 $X$ 的维数为 $0$ 的连通分支都是域的谱这一事实得出。

\medskip\noindent
设 $Z \subset X$ 是不可约分支，并把它看作整的闭子概形。显然
$f(Z) \subset Z$，且 $f|_Z : Z \to Z$ 是 $k$ 上的自同构，并在 $Z$ 的
底拓扑空间上诱导恒等映射。由于 $X$ 既约，只须证明箭头
$f|_Z : Z \to Z$ 是恒等映射。这把问题归约到下一段讨论的情形。

\medskip\noindent
假设 $X$ 不可约且维数 $> 0$。选取非空仿射开集 $U \subset X$。
由于 $f(U) \subset U$，并且 $U \subset X$ 概形论稠密，只须证明
$f|_U : U \to U$ 是恒等映射。

\medskip\noindent
假设 $X = \Spec(A)$ 仿射、不可约、维数 $> 0$，并且 $k$ 是无限域。
设 $g \in A$ 非常数。集合
$$
S = \bigcup\nolimits_{\lambda \in k} V(g - \lambda)
$$
在 $X$ 中稠密，因为它是稠密子集 $\mathbf{A}^1_k(k)$ 在非常数态射
$g : X \to \mathbf{A}^1_k$ 下的逆像。若 $x \in S$，则
$g(x)$（即 $g$ 在 $\kappa(x)$ 中的像）属于 $k \to \kappa(x)$ 的像。因此
$f^\sharp : \kappa(x) \to \kappa(x)$ 固定 $g(x)$。所以
$f^\sharp(g)$ 在 $\kappa(x)$ 中的像等于 $g(x)$。我们得出
$$
S \subset V(g - f^\sharp(g))
$$
而由于 $X$ 既约且 $S$ 稠密，便得到 $g=f^\sharp(g)$。这证明
$f^\sharp = \text{id}_A$，因为 $A$ 作为 $k$-代数由上述元素 $g$ 生成
（略去细节；提示：常数函数之集是有限维 $k$-向量子空间，位于 $A$ 中）。
我们得出 $f = \text{id}_X$。

\medskip\noindent
假设 $X = \Spec(A)$ 仿射、不可约、维数 $> 0$，并且 $k$ 是有限域。
若对每个 $1$ 维整的闭子概形 $C \subset X$，限制 $f|_C : C \to C$ 都是
恒等映射，则 $f$ 是恒等映射。这把问题归约到 $X$ 为曲线的情形。
有限域上的曲线具有有限自同构群（略去细节）。因此 $f$ 具有有限阶，记为
$n$。然后如上选取非常数态射 $g : X \to \mathbf{A}^1_k$，并考虑
$$
S = \{x \in X\text{ closed such that }[\kappa(g(x)) : k]
\text{ is prime to }n\}
$$
如前论证，可知 $S$ 在 $X$ 中稠密。对闭点 $x \in X$，映射
$f^\sharp : \kappa(x) \to \kappa(x)$ 是阶整除 $n$ 的自同构。因此对
$x \in S$，该自同构在由 $g$ 在 $\kappa(x)$ 中的像生成的子域上作用平凡。
于是 $S \subset V(g - f^\sharp(g))$，按前述论证即得结论。
\end{proof}


\section{欧拉示性数}
\label{varieties-section-euler}

\noindent
本节证明域上固有概形上的凝聚层的欧拉示性数的一些初等性质。

\begin{definition}
\label{varieties-definition-euler-characteristic}
设 $k$ 是域。设 $X$ 是 $k$ 上的固有概形。设 $\mathcal{F}$ 是凝聚
$\mathcal{O}_X$-模。在此情形下，{\it $\mathcal{F}$ 的欧拉示性数}是整数
$$
\chi(X, \mathcal{F}) = \sum\nolimits_i (-1)^i \dim_k H^i(X, \mathcal{F}).
$$
关于这个公式的合理性，见下文。
\end{definition}

\noindent
在该定义的情形下，向量空间 $H^i(X, \mathcal{F})$ 中只有有限多个非零
（《概形的上同调》引理
\ref{coherent-lemma-quasi-coherence-higher-direct-images}），
并且这些空间都是有限维的（《概形的上同调》引理
\ref{coherent-lemma-proper-over-affine-cohomology-finite}）。因此
$\chi(X, \mathcal{F}) \in \mathbf{Z}$ 有良好定义。注意，这个定义依赖于域
$k$，而不只依赖于二元组 $(X, \mathcal{F})$。

\begin{lemma}
\label{varieties-lemma-euler-characteristic-additive}
设 $k$ 是域。设 $X$ 是 $k$ 上的固有概形。设
$0 \to \mathcal{F}_1 \to \mathcal{F}_2 \to \mathcal{F}_3 \to 0$
是 $X$ 上凝聚模的短正合列。则
$$
\chi(X, \mathcal{F}_2) = \chi(X, \mathcal{F}_1) + \chi(X, \mathcal{F}_3)
$$
\end{lemma}

\begin{proof}
考虑与该引理中的短正合列相伴的上同调长正合列
$$
0 \to H^0(X, \mathcal{F}_1) \to H^0(X, \mathcal{F}_2) \to
H^0(X, \mathcal{F}_3) \to H^1(X, \mathcal{F}_1) \to \ldots
$$
线性代数中的秩-零化度定理表明
$$
0 = \dim H^0(X, \mathcal{F}_1) - \dim H^0(X, \mathcal{F}_2)
+ \dim H^0(X, \mathcal{F}_3) - \dim H^1(X, \mathcal{F}_1) + \ldots
$$
这立即蕴涵该引理。
\end{proof}

\begin{lemma}
\label{varieties-lemma-chi-tensor-finite}
设 $k$ 是域。设 $X$ 是 $k$ 上的固有概形。设 $\mathcal{F}$ 是满足
$\dim(\text{Supp}(\mathcal{F})) \leq 0$ 的凝聚层。则
\begin{enumerate}
\item $\mathcal{F}$ 由全局截面生成，
\item $H^0(X, \mathcal{F}) =
\bigoplus_{x \in \text{Supp}(\mathcal{F})} \mathcal{F}_x$，
\item $H^i(X, \mathcal{F}) = 0$ 对每个 $i > 0$ 成立，
\item $\chi(X, \mathcal{F}) = \dim_k H^0(X, \mathcal{F})$，并且
\item 等式
$\chi(X, \mathcal{F} \otimes \mathcal{E}) = n\chi(X, \mathcal{F})$
对每个局部自由模 $\mathcal{E}$（其秩为 $n$）成立。
\end{enumerate}
\end{lemma}

\begin{proof}
由《概形的上同调》引理 \ref{coherent-lemma-coherent-support-closed}，
可知 $\mathcal{F} = i_*\mathcal{G}$，其中 $i : Z \to X$ 是
$\mathcal{F}$ 的概形论支集的包含映射，而 $\mathcal{G}$ 是凝聚
$\mathcal{O}_Z$-模。由概形论支集的定义，$Z$ 的底拓扑空间是
$\text{Supp}(\mathcal{F})$。由于 $Z$ 的维数为 $0$，由《性质》引理
\ref{properties-lemma-locally-Noetherian-dimension-0} 可知 $Z$ 是仿射概形。
所以 $\mathcal{G}$ 是全局生成的，并且 $\mathcal{G}$ 的高阶上同调群为零
（《概形的上同调》引理
\ref{coherent-lemma-quasi-coherent-affine-cohomology-zero}）。
事实上，由引理 \ref{varieties-lemma-algebraic-scheme-dim-0}，概形 $Z$ 是局部
Artin 环的谱的有限不交并。因此相应地有
$H^0(Z, \mathcal{G}) = \bigoplus_{z \in Z} \mathcal{G}_z$。
由《概形的上同调》引理 \ref{coherent-lemma-relative-affine-cohomology}，
$\mathcal{F}$ 与 $\mathcal{G}$ 的上同调相同。因此
$H^i(X, \mathcal{F}) = 0$ 对每个 $i > 0$ 成立，并且
$H^0(X, \mathcal{F}) = H^0(Z, \mathcal{G})$。
特别地，(3) 成立。对 $z \in Z$，若它对应于
$x \in \text{Supp}(\mathcal{F})$，则有
$\mathcal{G}_z = (i_*\mathcal{G})_x = \mathcal{F}_x$。
我们得出 (2) 成立。当然，(2) 蕴涵 (1)。由欧拉示性数
$\chi(X, \mathcal{F})$ 的定义和 (3)，得到 (4)。由投影公式
（《上同调》引理 \ref{cohomology-lemma-projection-formula}），有
$$
i_*(\mathcal{G} \otimes i^*\mathcal{E}) = \mathcal{F} \otimes \mathcal{E}.
$$
由于 $Z$ 的维数为 $0$，局部自由层 $i^*\mathcal{E}$ 同构于
$\mathcal{O}_Z^{\oplus n}$；按上述方式论证即知 (5) 成立。
\end{proof}

\begin{lemma}
\label{varieties-lemma-euler-characteristic-extend-base-field}
设 $k'/k$ 是域扩张。设 $X$ 是 $k$ 上的固有概形。设 $\mathcal{F}$ 是
$X$ 上的凝聚层。设 $\mathcal{F}'$ 是 $\mathcal{F}$ 在 $X_{k'}$ 上的拉回。
则 $\chi(X, \mathcal{F}) = \chi(X', \mathcal{F}')$。
\end{lemma}

\begin{proof}
这是因为由平坦基变换有
$$
H^i(X_{k'}, \mathcal{F}') = H^i(X, \mathcal{F}) \otimes_k k'
$$
见《概形的上同调》引理
\ref{coherent-lemma-flat-base-change-cohomology}。
\end{proof}

\begin{lemma}
\label{varieties-lemma-euler-characteristic-morphism}
设 $k$ 是域。设 $f : Y \to X$ 是 $k$ 上的固有概形之间的态射。设
$\mathcal{G}$ 是凝聚 $\mathcal{O}_Y$-模。则
$$
\chi(Y, \mathcal{G}) = \sum (-1)^i \chi(X, R^if_*\mathcal{G})
$$
\end{lemma}

\begin{proof}
该公式是有意义的：层 $R^if_*\mathcal{G}$ 是凝聚的，并且其中只有有限多个非零；
见《概形的上同调》命题
\ref{coherent-proposition-proper-pushforward-coherent} 和引理
\ref{coherent-lemma-quasi-coherence-higher-direct-images}。
由《上同调》引理 \ref{cohomology-lemma-Leray}，存在以
$$
E_2^{p, q} = H^p(X, R^qf_*\mathcal{G})
$$
为第二页并收敛到 $H^{p + q}(Y, \mathcal{G})$ 的谱序列。由 $X$ 上上同调的
有限性可知，只有有限多个 $E_2^{p, q}$ 非零，并且每个 $E_2^{p, q}$ 都是
有限维向量空间。由此，$E_r^{p, q}$ 对每个 $r \geq 2$ 亦如此，并且
$$
\sum (-1)^{p + q} \dim_k E_r^{p, q}
$$
与 $r$ 无关。对充分大的 $r$，有 $E_r^{p, q} = E_\infty^{p, q}$；而收敛是指
$H^n(Y, \mathcal{G})$ 上存在一个滤过，其分次片为
$E_\infty^{p, q}$，其中 $p + q = n$（这就是谱序列收敛的含义），故得出结论。
另见更一般的《同调》引理
\ref{homology-lemma-biregular-ss-relation-in-K0}。
\end{proof}

\section{射影空间}
\label{varieties-section-projective-space}

\noindent
域上的射影空间的一些结果。

\begin{lemma}
\label{varieties-lemma-projective-space-smooth}
\begin{slogan}
射影空间是光滑的。
\end{slogan}
设 $k$ 是域且 $n \geq 0$。则 $\mathbf{P}^n_k$ 是维数为 $n$ 的
$k$ 上光滑射影簇。
\end{lemma}

\begin{proof}
从略。
\end{proof}

\begin{lemma}
\label{varieties-lemma-intersection-in-affine-space}
设 $k$ 是域且 $n \geq 0$。设 $X, Y \subset \mathbf{A}^n_k$
为闭子集。假设 $X$ 和 $Y$ 都是等维的，
$\dim(X) = r$ 且 $\dim(Y) = s$。
则 $X \cap Y$ 的每个不可约分支的维数都 $\geq r + s - n$。
\end{lemma}

\begin{proof}
考虑闭子概形 $X \times Y \subset \mathbf{A}^{2n}_k$，
这里使用坐标 $x_1, \ldots, x_n, y_1, \ldots, y_n$。于是
$X \cap Y = X \times Y \cap V(x_1 - y_1, \ldots, x_n - y_n)$。
设 $t \in X \cap Y \subset X \times Y$ 为闭点。
由引理 \ref{varieties-lemma-dimension-product-locally-algebraic}，有
$\dim_t(X \times Y) = \dim(X) + \dim(Y)$。
因而由引理 \ref{varieties-lemma-dimension-locally-algebraic}，有
$\dim(\mathcal{O}_{X \times Y, t}) = r + s$。
由《代数》引理 \ref{algebra-lemma-one-equation}，可得
$$
\dim(\mathcal{O}_{X \cap Y, t}) =
\dim(\mathcal{O}_{X \times Y, t}/(x_1 - y_1, \ldots, x_n - y_n)) \geq
r + s - n
$$
再由引理 \ref{varieties-lemma-dimension-locally-algebraic} 即得结论。
\end{proof}

\begin{lemma}
\label{varieties-lemma-intersection-in-projective-space}
设 $k$ 是域且 $n \geq 0$。设 $X, Y \subset \mathbf{P}^n_k$
为非空闭子集。若 $\dim(X) = r$ 且 $\dim(Y) = s$，并且
$r + s \geq n$，则 $X \cap Y$ 非空，且
$\dim(X \cap Y) \geq r + s - n$。
\end{lemma}

\begin{proof}
记 $\mathbf{A}^n = \Spec(k[x_0, \ldots, x_n])$ 且
$\mathbf{P}^n = \text{Proj}(k[T_0, \ldots, T_n])$。
考虑态射
$\pi : \mathbf{A}^{n + 1} \setminus \{0\} \to \mathbf{P}^n$，
它把 $(x_0, \ldots, x_n)$ 送到点 $[x_0 : \ldots : x_n]$。
更准确地说，它是由二元组
$(\mathcal{O}_{\mathbf{A}^{n + 1} \setminus \{0\}}, (x_0, \ldots, x_n))$
相伴的态射；见《构造》引理 \ref{constructions-lemma-projective-space}。
在标准仿射开集 $D_+(T_i)$ 上，得到由环同态
$$
k\left[\frac{T_0}{T_i}, \ldots, \frac{T_n}{T_i}\right]
\longrightarrow
k\left[T_0, \ldots, T_n, \frac{1}{T_i}\right] \cong
k\left[\frac{T_0}{T_i}, \ldots, \frac{T_n}{T_i}\right]
\left[T_i, \frac{1}{T_i}\right]
$$
相伴的态射；该环同态是满的，并且是相对维数为 $1$、纤维不可约的光滑同态
（略去细节）。
因此 $\pi^{-1}(X)$ 和 $\pi^{-1}(Y)$ 是维数分别为
$r + 1$ 和 $s + 1$ 的非空闭子集。选取不可约分支
$V \subset \pi^{-1}(X)$，其维数为 $r + 1$；再选取不可约分支
$W \subset \pi^{-1}(Y)$，其维数为 $s + 1$。注意，这蕴涵 $V$ 和 $W$ 包含它们所遇到的
$\pi$ 的每条纤维（因为 $\pi$ 的纤维是维数为 $1$ 的不可约集，
且引理 \ref{varieties-lemma-dimension-fibres-locally-algebraic} 表明
$V \to \pi(V)$ 和 $W \to \pi(W)$ 的纤维维数都 $\geq 1$）。
令 $\overline{V}$ 和 $\overline{W}$ 分别为
$V$ 和 $W$ 在 $\mathbf{A}^{n + 1}$ 中的闭包。因为
$0 \in \mathbf{A}^{n + 1}$ 属于 $\pi$ 的每条纤维的闭包，所以
$0 \in \overline{V} \cap \overline{W}$。由引理
\ref{varieties-lemma-intersection-in-affine-space}，有
$\dim(\overline{V} \cap \overline{W}) \geq r + s - n + 1$。
再次使用引理 \ref{varieties-lemma-dimension-fibres-locally-algebraic} 如上论证，
便得出 $\pi(V \cap W) \subset X \cap Y$ 的维数至少为
$r + s - n$，如所需。
\end{proof}

\begin{lemma}
\label{varieties-lemma-equation-codim-1-in-projective-space}
设 $k$ 是域。设 $Z \subset \mathbf{P}^n_k$ 是没有嵌入点的闭子概形，
且 $Z$ 的每个不可约分支的维数都是 $n - 1$。则理想
$I(Z) \subset k[T_0, \ldots, T_n]$（与 $Z$ 对应）是主理想。
\end{lemma}

\begin{proof}
这是《除子》引理
\ref{divisors-lemma-equation-codim-1-in-projective-space} 的特殊情形。
\end{proof}

\section{射影空间上的凝聚层}
\label{varieties-section-coherent-Pn}

\noindent
本节证明域上的 $\mathbf{P}^n$ 上凝聚层之上同调的一些结果；这些结果可见
\cite{Mum}。以后讨论 Quot 概形和 Hilbert 概形时会用到它们。

\subsection{预备知识}
\label{varieties-subsection-preliminaries}

\noindent
设 $k$ 是域，$n \geq 1$，$d \geq 1$，并设
$s \in \Gamma(\mathbf{P}_k^n, \mathcal{O}(d))$
为非零截面。本节用 $\mathcal{O}(d)$ 表示射影空间上结构层的第
$d$ 次扭转（《构造》定义 \ref{constructions-definition-twist} 和
\ref{constructions-definition-projective-space}）。
由于 $\mathbf{P}^n_k$ 是簇，该截面是正则的；因而
$s$ 是 $\mathcal{O}(d)$ 的正则截面，并定义有效 Cartier 除子
$H = Z(s) \subset \mathbf{P}^n_k$；见《除子》第
\ref{divisors-section-effective-Cartier-divisors} 节。
这样的除子 $H$ 称为{it 超曲面}；若 $d = 1$，则称为{it 超平面}。

\begin{lemma}
\label{varieties-lemma-hyperplane}
设 $k$ 是域。设 $n \geq 1$。
设 $i : H \to \mathbf{P}^n_k$ 是超平面。
则存在同构
$$
\varphi : \mathbf{P}^{n - 1}_k \longrightarrow H
$$
使得 $i^*\mathcal{O}(1)$ 拉回为 $\mathcal{O}(1)$。
\end{lemma}

\begin{proof}
有 $\mathbf{P}^n_k = \text{Proj}(k[T_0, \ldots, T_n])$。
截面 $s$ 对应于 $T_0, \ldots, T_n$ 中次数为 $1$ 的齐次形式；见
《概形的上同调》第
\ref{coherent-section-cohomology-projective-space} 节。
设 $s = \sum a_i T_i$。
由《构造》引理 \ref{constructions-lemma-closed-in-projective-space}，有
$H = \text{Proj}(k[T_0, \ldots, T_n]/I)$，其中分次理想 $I$ 的定义是：
$I_d$ 等于映射
$\Gamma(\mathbf{P}^n_k, \mathcal{O}(d)) \to \Gamma(H, i^*\mathcal{O}(d))$
的核。由《除子》定义 \ref{divisors-definition-zero-scheme-s} 中
$Z(s)$ 的构造可知，在 $D_{+}(T_j)$ 上，$H$ 的理想由
$\sum a_i T_i/T_j$ 生成；这是多项式环
$k[T_0/T_j, \ldots, T_n/T_j]$ 中的元素。因此很清楚，$I$ 是由
$\sum a_i T_i$ 生成的理想。注意
$$
k[T_0, \ldots, T_n]/I = k[T_0, \ldots, T_n]/(\sum a_i T_i) \cong
k[S_0, \ldots, S_{n - 1}]
$$
为分次环同构。例如，若 $a_n \not = 0$，则令
$S_i$ 映到 $T_i$ 的类即可。由 $\text{Proj}$ 的函子性得到所需同构。
结构层扭转的相等性例如由《构造》引理
\ref{constructions-lemma-surjective-graded-rings-generated-degree-1-map-proj}
推出。
\end{proof}

\begin{lemma}
\label{varieties-lemma-exact-sequence-induction}
设 $k$ 是无限域。设 $n \geq 1$。
设 $\mathcal{F}$ 是 $\mathbf{P}^n_k$ 上的凝聚模。
则存在非零截面
$s \in \Gamma(\mathbf{P}^n_k, \mathcal{O}(1))$
和短正合列
$$
0 \to \mathcal{F}(-1) \to \mathcal{F} \to i_*\mathcal{G} \to 0
$$
其中 $i : H \to \mathbf{P}^n_k$ 是超平面 $H$，它与 $s$ 相伴，且
$\mathcal{G} = i^*\mathcal{F}$。
\end{lemma}

\begin{proof}
映射 $\mathcal{F}(-1) \to \mathcal{F}$ 来自《构造》等式
(\ref{constructions-equation-multiply-on-sheaf})，其中 $n = 1$、
$m = -1$，并取截面 $s$（属于 $\mathcal{O}(1)$）。
具体考察该映射限制到 $D_{+}(T_0)$ 后的形式。记
$D_{+}(T_0) = \Spec(k[x_1, \ldots, x_n])$，其中
$x_i = T_i/T_0$。将 $\mathcal{O}(1)|_{D_{+}(T_0)}$ 与
$\mathcal{O}$ 等同，所用截面为 $T_0$。因此，若
$s = \sum a_iT_i$，则在上述等同下
$s|_{D_{+}(T_0)} = a_0 + \sum a_ix_i$。再假设
$\mathcal{F}|_{D_{+}(T_0)}$ 对应于有限
$k[x_1, \ldots, x_n]$-模 $M$。通过等同
$\mathcal{F}(-1) = \mathcal{F} \otimes \mathcal{O}(-1)$
及所选的 $\mathcal{O}(1)$ 平凡化，可知
$\mathcal{F}(-1)$ 也对应于 $M$。所以，把
$\mathcal{F}(-1) \to \mathcal{F}$ 限制到 $D_{+}(T_0)$，得到映射
$$
M \xrightarrow{a_0 + \sum a_ix_i} M
$$
为说明该箭头是单射，只须选取 $a_0 + \sum a_ix_i$，使它不属于
$M$ 的任何相伴素理想；见《代数》引理
\ref{algebra-lemma-ass-zero-divisors}。由《代数》引理
\ref{algebra-lemma-finite-ass}，集合 $\text{Ass}(M)$（即
$M$ 的相伴素理想集）是有限的。注意，对
$\mathfrak p \in \text{Ass}(M)$，交集
$\mathfrak p \cap \{a_0 + \sum a_i x_i\}$ 是真 $k$-向量子空间。
我们得出，存在真子向量空间的有限族
$V_1, \ldots, V_m \subset \Gamma(\mathbf{P}^n_k, \mathcal{O}(1))$，
使得若取 $s$ 不属于 $\bigcup V_i$，则乘以 $s$ 在
$D_{+}(T_0)$ 上是单射。对限制到 $D_{+}(T_j)$ 的情形也同样，
其中 $j = 1, \ldots, n$。由于 $k$ 是无限域，有限个真子向量空间的并
绝不等于整个空间，故可选取 $s$ 使该映射为单射。
$\mathcal{F}(-1) \to \mathcal{F}$ 的余核被
$\Im(s : \mathcal{O}(-1) \to \mathcal{O})$ 零化；由
《除子》定义 \ref{divisors-definition-zero-scheme-s}，后者是 $H$ 的理想层。
因此，由《概形的上同调》引理
\ref{coherent-lemma-i-star-equivalence}，得到 $\mathcal{G}$，它在 $H$ 上。
\end{proof}

\begin{remark}
\label{varieties-remark-exact-sequence-induction}
设 $k$ 是无限域。设 $n \geq 1$。给定有限多个凝聚模
$\mathcal{F}_i$，它们在 $\mathbf{P}^n_k$ 上，可选取同一个
$s \in \Gamma(\mathbf{P}^n_k, \mathcal{O}(1))$，
使引理 \ref{varieties-lemma-exact-sequence-induction} 的陈述对每个模都成立。
为证明这一点，只需把该引理应用于 $\bigoplus \mathcal{F}_i$。
\end{remark}

\begin{remark}
\label{varieties-remark-exact-sequence-induction-cohomology}
在引理 \ref{varieties-lemma-hyperplane} 和
\ref{varieties-lemma-exact-sequence-induction} 的情形下，有
$H \cong \mathbf{P}^{n - 1}_k$，且 Serre 扭转满足
$\mathcal{O}_H(d) = i^*\mathcal{O}_{\mathbf{P}^n_k}(d)$。
对每个 $d \in \mathbf{Z}$，有短正合列
$$
0 \to \mathcal{F}(d - 1) \to \mathcal{F}(d) \to i_*(\mathcal{G}(d)) \to 0
$$
事实上，与 $\mathcal{O}_{\mathbf{P}^n_k}(d)$ 作张量积是正合函子，
并且由投影公式（《上同调》引理
\ref{cohomology-lemma-projection-formula}），有
$i_*(\mathcal{G}(d)) = i_*\mathcal{G} \otimes \mathcal{O}_{\mathbf{P}^n_k}(d)$。
得到相应的长正合列
$$
H^i(\mathbf{P}^n_k, \mathcal{F}(d - 1)) \to
H^i(\mathbf{P}^n_k, \mathcal{F}(d)) \to
H^i(H, \mathcal{G}(d)) \to
H^{i + 1}(\mathbf{P}^n_k, \mathcal{F}(d - 1))
$$
这是由上述事实以及下列等式推出的：由《概形的上同调》引理
\ref{coherent-lemma-relative-affine-cohomology}（闭浸入是仿射的），有
$H^i(\mathbf{P}^n_k, i_*\mathcal{G}(d)) = H^i(H, \mathcal{G}(d))$。
\end{remark}





\subsection{正则性}
\label{varieties-subsection-regularity}

\noindent
定义如下。

\begin{definition}
\label{varieties-definition-regularity}
设 $k$ 是域。设 $n \geq 0$。设 $\mathcal{F}$ 是
$\mathbf{P}^n_k$ 上的凝聚层。称 $\mathcal{F}$ 是{it $m$-正则的}，若
$$
H^i(\mathbf{P}^n_k, \mathcal{F}(m - i)) = 0
$$
对 $i = 1, \ldots, n$ 成立。
\end{definition}

\noindent
注意，$\mathcal{F} = \mathcal{O}(d)$ 是 $m$-正则的，当且仅当
$d \geq -m$。这由《概形的上同调》等式
(\ref{coherent-equation-identify}) 中的上同调群计算推出。
也就是说，$H^n(\mathbf{P}^n_k, \mathcal{O}(d)) = 0$，当且仅当
$d \geq -n$。

\begin{lemma}
\label{varieties-lemma-m-regular-extend-base-field}
设 $k'/k$ 是域扩张。设 $n \geq 0$。
设 $\mathcal{F}$ 是 $\mathbf{P}^n_k$ 上的凝聚层。
设 $\mathcal{F}'$ 是 $\mathcal{F}$ 在 $\mathbf{P}^n_{k'}$ 上的拉回。
则 $\mathcal{F}$ 是 $m$-正则的，当且仅当 $\mathcal{F}'$ 是
$m$-正则的。
\end{lemma}

\begin{proof}
这是因为由平坦基变换有
$$
H^i(\mathbf{P}^n_{k'}, \mathcal{F}') =
H^i(\mathbf{P}^n_k, \mathcal{F}) \otimes_k k'
$$
见《概形的上同调》引理
\ref{coherent-lemma-flat-base-change-cohomology}。
\end{proof}

\begin{lemma}
\label{varieties-lemma-m-regular}
在引理 \ref{varieties-lemma-exact-sequence-induction} 的情形下，若
$\mathcal{F}$ 是 $m$-正则的，则 $\mathcal{G}$ 是 $m$-正则的；
这里它在 $H \cong \mathbf{P}^{n - 1}_k$ 上。
\end{lemma}

\begin{proof}
回忆，由《概形的上同调》引理
\ref{coherent-lemma-relative-affine-cohomology}，有
$H^i(\mathbf{P}^n_k, i_*\mathcal{G}) = H^i(H, \mathcal{G})$。
因此由评注 \ref{varieties-remark-exact-sequence-induction-cohomology} 可知，
对 $i \geq 1$ 得到
$$
H^i(\mathbf{P}^n_k, \mathcal{F}(m - i)) \to
H^i(H, \mathcal{G}(m - i)) \to
H^{i + 1}(\mathbf{P}^n_k, \mathcal{F}(m - 1 - i))
$$
该引理随即成立。
\end{proof}

\begin{lemma}
\label{varieties-lemma-m-regular-up}
设 $k$ 是域。设 $n \geq 0$。
设 $\mathcal{F}$ 是 $\mathbf{P}^n_k$ 上的凝聚层。
若 $\mathcal{F}$ 是 $m$-正则的，则 $\mathcal{F}$ 是
$(m + 1)$-正则的。
\end{lemma}

\begin{proof}
对 $n$ 作归纳证明。若 $n = 0$，则每个层对所有 $m$ 都是 $m$-正则的，
无需证明。由引理 \ref{varieties-lemma-m-regular-extend-base-field}，可把 $k$
替换为无限扩域，并假设 $k$ 是无限域。
于是可应用引理 \ref{varieties-lemma-exact-sequence-induction}。
由引理 \ref{varieties-lemma-m-regular}，知道 $\mathcal{G}$ 是
$m$-正则的。对 $n$ 使用归纳假设，便知 $\mathcal{G}$ 是
$(m + 1)$-正则的。考察与下列正合列相伴的上同调长正合列
$$
0 \to \mathcal{F}(m - i) \to \mathcal{F}(m + 1 - i)
\to i_*\mathcal{G}(m + 1 - i) \to 0
$$
如评注 \ref{varieties-remark-exact-sequence-induction-cohomology} 所述，读者容易推出：
对 $i \geq 1$，$H^i(\mathbf{P}^n_k, \mathcal{F}(m + 1 - i))$ 消失；
这是由 $H^i(\mathbf{P}^n_k, \mathcal{F}(m - i))$ 和
$H^i(\mathbf{P}^n_k, \mathcal{G}(m + 1 - i))$ 的（已知）消失得出的。
\end{proof}

\begin{lemma}
\label{varieties-lemma-m-regular-multiply}
设 $k$ 是域。设 $n \geq 0$。
设 $\mathcal{F}$ 是 $\mathbf{P}^n_k$ 上的凝聚层。
若 $\mathcal{F}$ 是 $m$-正则的，则乘法映射
$$
H^0(\mathbf{P}^n_k, \mathcal{F}(m)) \otimes_k
H^0(\mathbf{P}^n_k, \mathcal{O}(1))
\longrightarrow
H^0(\mathbf{P}^n_k, \mathcal{F}(m + 1))
$$
是满射。
\end{lemma}

\begin{proof}
设 $k'/k$ 是域扩张。设 $\mathcal{F}'$ 如引理
\ref{varieties-lemma-m-regular-extend-base-field} 中所定义。由《概形的上同调》引理
\ref{coherent-lemma-flat-base-change-cohomology}，该引理的线性映射基变换到
$k'$ 后，就是层 $\mathcal{F}'$ 的同一线性映射。由于
$k \to k'$ 忠实平坦，只须在 $k'$ 上证明该引理；也就是说，可以假设
$k$ 是无限域。

\medskip\noindent
假设 $k$ 是无限域。对 $n$ 作归纳证明。情形 $n = 0$ 是平凡的，因为
$\mathcal{O}(1) \cong \mathcal{O}$ 由 $T_0$ 生成。对 $n > 0$，应用
引理 \ref{varieties-lemma-exact-sequence-induction}，并将该正合列与
$\mathcal{O}(m + 1)$ 作张量积，得到
$$
0 \to \mathcal{F}(m) \xrightarrow{s} \mathcal{F}(m + 1) \to
i_*\mathcal{G}(m + 1) \to 0
$$
见评注 \ref{varieties-remark-exact-sequence-induction-cohomology}。
设 $t \in H^0(\mathbf{P}^n_k, \mathcal{F}(m + 1))$。
由归纳假设，像 $\overline{t} \in H^0(H, \mathcal{G}(m + 1))$
是 $\sum g_i \otimes \overline{s}_i$ 的像，其中
$\overline{s}_i \in \Gamma(H, \mathcal{O}(1))$ 且
$g_i \in H^0(H, \mathcal{G}(m))$。由于 $\mathcal{F}$ 是 $m$-正则的，
有 $H^1(\mathbf{P}^n_k, \mathcal{F}(m - 1)) = 0$；因此，与短正合列
$$
0 \to \mathcal{F}(m - 1) \xrightarrow{s} \mathcal{F}(m) \to
i_*\mathcal{G}(m) \to 0
$$
相伴的上同调长正合列表明，可把 $g_i$ 提升为
$f_i \in H^0(\mathbf{P}^n_k, \mathcal{F}(m))$。
也可以把 $\overline{s}_i$ 提升为
$s_i \in H^0(\mathbf{P}^n_k, \mathcal{O}(1))$（例如见引理
\ref{varieties-lemma-hyperplane} 的证明）。减去
$\sum f_i \otimes s_i$ 的像后，可把 $t$ 取为满足
$\overline{t} = 0$。但这恰好意味着 $t$ 是某个
$f \otimes s$ 的像，其中
$f \in H^0(\mathbf{P}^n_k, \mathcal{F}(m))$，如所需。
\end{proof}

\begin{lemma}
\label{varieties-lemma-m-regular-globally-generated}
设 $k$ 是域。设 $n \geq 0$。
设 $\mathcal{F}$ 是 $\mathbf{P}^n_k$ 上的凝聚层。
若 $\mathcal{F}$ 是 $m$-正则的，则 $\mathcal{F}(m)$ 是全局生成的。
\end{lemma}

\begin{proof}
对所有 $d \gg 0$，层 $\mathcal{F}(d)$ 都是全局生成的。
这例如由《概形的上同调》引理
\ref{coherent-lemma-coherent-projective} 的第一部分推出。
选取 $d \geq m$，使 $\mathcal{F}(d)$ 全局生成。
选取一组基 $f_1, \ldots, f_r \in H^0(\mathbf{P}^n_k, \mathcal{F})$。
由引理 \ref{varieties-lemma-m-regular-multiply}，每个元素
$f \in H^0(\mathbf{P}^n_k, \mathcal{F}(d))$ 都可写成
$f = \sum P_if_i$，其中某些
$P_i \in k[T_0, \ldots, T_n]$ 是次数为 $d - m$ 的齐次多项式。
由于截面 $f$ 生成 $\mathcal{F}(d)$，可知截面 $f_i$ 生成
$\mathcal{F}(m)$。
\end{proof}

\subsection{Hilbert 多项式}
\label{varieties-subsection-hilbert}

\noindent
下面的引理将被更一般的引理
\ref{varieties-lemma-numerical-polynomial-from-euler} 所取代。

\begin{lemma}
\label{varieties-lemma-hilbert-polynomial}
设 $k$ 是域。设 $n \geq 0$。设 $\mathcal{F}$ 是
$\mathbf{P}^n_k$ 上的凝聚层。函数
$$
d \longmapsto \chi(\mathbf{P}^n_k, \mathcal{F}(d))
$$
是多项式。
\end{lemma}

\begin{proof}
对 $n$ 作归纳证明。若 $n = 0$，则
$\mathbf{P}^n_k = \Spec(k)$ 且 $\mathcal{F}(d) = \mathcal{F}$。
因此在此情形下，该函数是常数，即次数为 $0$ 的多项式。
假设 $n > 0$。由引理
\ref{varieties-lemma-euler-characteristic-extend-base-field}，可假设 $k$ 是无限域。
应用引理 \ref{varieties-lemma-exact-sequence-induction}。
把引理 \ref{varieties-lemma-euler-characteristic-additive} 应用于扭转后的正合列
$0 \to \mathcal{F}(d - 1) \to \mathcal{F}(d) \to i_*\mathcal{G}(d) \to 0$，
得到
$$
\chi(\mathbf{P}^n_k, \mathcal{F}(d)) -
\chi(\mathbf{P}^n_k, \mathcal{F}(d - 1)) =
\chi(H, \mathcal{G}(d))
$$
见评注 \ref{varieties-remark-exact-sequence-induction-cohomology}。
由于 $H \cong \mathbf{P}^{n - 1}_k$，由归纳假设，右端是多项式。
最后注意：任何函数 $f : \mathbf{Z} \to \mathbf{Z}$，若映射
$d \mapsto f(d) - f(d - 1)$ 是多项式，则其自身也是多项式；这便完成证明。
略去该事实的证明（提示：与《代数》引理
\ref{algebra-lemma-numerical-polynomial} 比较）。
\end{proof}

\begin{definition}
\label{varieties-definition-hilbert-polynomial}
设 $k$ 是域。设 $n \geq 0$。设 $\mathcal{F}$ 是
$\mathbf{P}^n_k$ 上的凝聚层。函数
$d \mapsto \chi(\mathbf{P}^n_k, \mathcal{F}(d))$ 称为
$\mathcal{F}$ 的{it Hilbert 多项式}。
\end{definition}

\noindent
Hilbert 多项式的系数属于 $\mathbf{Q}$，一般并不属于 $\mathbf{Z}$。
例如，$\mathcal{O}_{\mathbf{P}^n_k}$ 的 Hilbert 多项式为
$$
d \longmapsto {d + n \choose n} = \frac{d^n}{n!} + \ldots
$$
这由下面的引理和下列事实推出：
$$
H^0(\mathbf{P}^n_k, \mathcal{O}_{\mathbf{P}^n_k}(d)) = k[T_0, \ldots, T_n]_d
$$
是次数为 $d$ 的部分，其在 $k$ 上的维数是 ${d + n \choose n}$。

\begin{lemma}
\label{varieties-lemma-hilbert-polynomial-H0}
设 $k$ 是域。设 $n \geq 0$。设 $\mathcal{F}$ 是
$\mathbf{P}^n_k$ 上的凝聚层，其 Hilbert 多项式为
$P \in \mathbf{Q}[t]$。则
$$
P(d) = \dim_k H^0(\mathbf{P}^n_k, \mathcal{F}(d))
$$
对所有 $d \gg 0$ 成立。
\end{lemma}

\begin{proof}
这由 $\mathcal{F}$ 的充分高次扭转之上同调消失推出。见
《概形的上同调》引理
\ref{coherent-lemma-coherent-projective}。
\end{proof}




\subsection{商的有界性}
\label{varieties-subsection-boundedness}

\noindent
本小节用 Hilbert 多项式来界定 $\mathbf{P}^n$ 上给定凝聚层之商的正则性。

\begin{lemma}
\label{varieties-lemma-bound-quotients-free}
设 $k$ 是域。设 $n \geq 0$。设 $r \geq 1$。设 $P \in \mathbf{Q}[t]$。
存在整数 $m$，它依赖于 $n$、$r$ 和 $P$，并具有如下性质：若
$$
0 \to \mathcal{K} \to \mathcal{O}^{\oplus r} \to \mathcal{F} \to 0
$$
是 $\mathbf{P}^n_k$ 上凝聚层的短正合列，并且 $\mathcal{F}$ 的
Hilbert 多项式为 $P$，则 $\mathcal{K}$ 是 $m$-正则的。
\end{lemma}

\begin{proof}
对 $n$ 作归纳证明。若 $n = 0$，则
$\mathbf{P}^n_k = \Spec(k)$，任何凝聚模都是 $0$-正则的，
而任何满射在全局截面上也是满射。
假设 $n > 0$。考虑引理中的正合列。
设 $P' \in \mathbf{Q}[t]$ 为多项式
$P'(t) = P(t) - P(t - 1)$。设 $m'$ 为对
$n - 1$、$r$ 和 $P'$ 有效的整数。
由引理 \ref{varieties-lemma-m-regular-extend-base-field} 和
\ref{varieties-lemma-euler-characteristic-extend-base-field}，可把 $k$ 替换为域扩张，
故可假设 $k$ 是无限域。把引理
\ref{varieties-lemma-exact-sequence-induction} 应用于凝聚层 $\mathcal{F}$。
$\mathcal{F}' = i^*\mathcal{F}$ 的 Hilbert 多项式为
$P'$（见引理 \ref{varieties-lemma-hilbert-polynomial} 的证明）。
由于 $i^*$ 右正合，可知 $\mathcal{F}'$ 是
$\mathcal{O}_H^{\oplus r} = i^*\mathcal{O}^{\oplus r}$ 的商。
因而归纳假设适用于 $\mathcal{F}'$；它在
$H \cong \mathbf{P}^{n - 1}_k$ 上（引理 \ref{varieties-lemma-hyperplane}）。
注意，映射 $\mathcal{K}(-1) \to \mathcal{K}$ 是单射，因为
$\mathcal{K} \subset \mathcal{O}^{\oplus r}$；其余核为
$i_*\mathcal{H}$，其中 $\mathcal{H} = i^*\mathcal{K}$。
由蛇形引理（《同调》引理 \ref{homology-lemma-snake}），得到行、列都正合的交换图
$$
\xymatrix{
& 0 \ar[d] & 0 \ar[d] & 0 \ar[d] \\
0 \ar[r] &
\mathcal{K}(-1) \ar[r] \ar[d] &
\mathcal{O}^{\oplus r}(-1) \ar[r] \ar[d] &
\mathcal{F}(-1) \ar[d] \ar[r] & 0 \\
0 \ar[r] &
\mathcal{K} \ar[r] \ar[d] &
\mathcal{O}^{\oplus r} \ar[r] \ar[d] &
\mathcal{F} \ar[d] \ar[r] & 0\\
0 \ar[r] &
i_*\mathcal{H} \ar[r] \ar[d] &
i_*\mathcal{O}_H^{\oplus r} \ar[r] \ar[d] &
i_*\mathcal{F}' \ar[r] \ar[d] & 0 \\
& 0 & 0 & 0
}
$$
因此，归纳假设适用于正合列
$0 \to \mathcal{H} \to \mathcal{O}_H^{\oplus r} \to \mathcal{F}' \to 0$；
该列在 $H \cong \mathbf{P}^{n - 1}_k$ 上（引理
\ref{varieties-lemma-hyperplane}），且 $\mathcal{H}$ 是 $m'$-正则的。
回忆，这蕴涵 $\mathcal{H}$ 是 $d$-正则的，对所有 $d \geq m'$ 都如此
（引理 \ref{varieties-lemma-m-regular-up}）。

\medskip\noindent
设 $i \geq 2$ 且 $d \geq m'$。由与上图左列相伴的上同调长正合列以及
$H^{i - 1}(H, \mathcal{H}(d))$ 的消失可知，映射
$$
H^i(\mathbf{P}^n_k, \mathcal{K}(d - 1))
\longrightarrow
H^i(\mathbf{P}^n_k, \mathcal{K}(d))
$$
是单射。由于这些群对 $d \gg 0$ 为零（《概形的上同调》引理
\ref{coherent-lemma-coherent-projective}），可知
$H^i(\mathbf{P}^n_k, \mathcal{K}(d))$ 对所有 $d \geq m'$ 和
$i \geq 2$ 都为零。

\medskip\noindent
还须控制 $H^1$。首先注意，所有映射
$$
H^1(\mathbf{P}^n_k, \mathcal{K}(m' - 1)) \to
H^1(\mathbf{P}^n_k, \mathcal{K}(m')) \to
H^1(\mathbf{P}^n_k, \mathcal{K}(m' + 1)) \to \ldots
$$
都是满射，因为 $H^1(H, \mathcal{H}(d))$ 对
$d \geq m'$ 消失。假设 $d > m'$，且
$$
H^1(\mathbf{P}^n_k, \mathcal{K}(d - 1))
\longrightarrow
H^1(\mathbf{P}^n_k, \mathcal{K}(d))
$$
是单射。则
$H^0(\mathbf{P}^n_k, \mathcal{K}(d)) \to H^0(H, \mathcal{H}(d))$
是满射。考虑交换图
$$
\xymatrix{
H^0(\mathbf{P}^n_k, \mathcal{K}(d)) \otimes_k
H^0(\mathbf{P}^n_k, \mathcal{O}(1))
\ar[r] \ar[d] &
H^0(\mathbf{P}^n_k, \mathcal{K}(d + 1)) \ar[d] \\
H^0(H, \mathcal{H}(d)) \otimes_k
H^0(H, \mathcal{O}_H(1))
\ar[r] &
H^0(H, \mathcal{H}(d + 1))
}
$$
由引理 \ref{varieties-lemma-m-regular-multiply}，下方水平箭头是满射。
因此右侧竖直箭头是满射。我们得出
$$
H^1(\mathbf{P}^n_k, \mathcal{K}(d))
\longrightarrow
H^1(\mathbf{P}^n_k, \mathcal{K}(d + 1))
$$
是单射。归纳可知
$$
H^1(\mathbf{P}^n_k, \mathcal{K}(d - 1)) \to
H^1(\mathbf{P}^n_k, \mathcal{K}(d)) \to
H^1(\mathbf{P}^n_k, \mathcal{K}(d + 1)) \to \ldots
$$
全都是单射；由这些群最终消失，得出
$H^1(\mathbf{P}^n_k, \mathcal{K}(d - 1)) = 0$。
所以，群 $H^1(\mathbf{P}^n_k, \mathcal{K}(d))$ 的维数在
$d \geq m'$ 时严格递减，直至变为零。
由此，$\mathcal{K}$ 的正则性被下式所界：
$m' + \dim_k H^1(\mathbf{P}^n_k, \mathcal{K}(m'))$。
另一方面，由高阶上同调群的消失，有
$$
\dim_k H^1(\mathbf{P}^n_k, \mathcal{K}(m')) = 
- \chi(\mathbf{P}^n_k, \mathcal{K}(m')) +
\dim_k H^0(\mathbf{P}^n_k, \mathcal{K}(m'))
$$
注意，$H^0$ 的维数被
$H^0(\mathbf{P}^n_k, \mathcal{O}^{\oplus r}(m'))$ 的维数所界；
后者至多为 $r{n + m' \choose n}$（若 $m' > 0$），否则为零。
最后，项 $\chi(\mathbf{P}^n_k, \mathcal{K}(m'))$ 等于
$r{n + m' \choose n} - P(m')$。这给出了所需类型的界，从而完成引理的证明。
\end{proof}







\section{Frobenius 态射}
\label{varieties-section-frobenius}

\noindent
设 $p$ 为素数。若 $X$ 是概形，则称
``$X$ 具有特征 $p$''，或称 ``$X$ 是特征 $p$ 的''，
或称 ``$X$ 处于特征 $p$''，若 $p$ 在 $\mathcal{O}_X$ 中为零。

\begin{definition}
\label{varieties-definition-absolute-frobenius}
设 $p$ 为素数。设 $X$ 是特征 $p$ 的概形。
$X$ 的\emph{绝对 Frobenius} 是态射 $F_X : X \to X$，
它在底层拓扑空间上为恒等映射，并且其
$F_X^\sharp : \mathcal{O}_X \to \mathcal{O}_X$ 由 $g \mapsto g^p$ 给出。
\end{definition}

\noindent
这是有意义的，因为对于任意环 $A$（其特征为 $p$），映射
$F_A : A \to A$，$a \mapsto a^p$ 是一个环自同态，并且在
$\Spec(A)$ 上诱导恒等映射。此外，若 $A$ 是局部环，则 $F_A$ 是局部同态。
因此，我们看到 $X$ 的绝对 Frobenius 是 $X$ 的一个自同态，
且属于概形范畴。事实证明，绝对 Frobenius 在特征 $p$
的概形范畴上定义了恒等函子的一个自映射。

\begin{lemma}
\label{varieties-lemma-frobenius-endomorphism-identity}
设 $p > 0$ 为素数。设 $f : X \to Y$ 是特征 $p$ 的概形之间的态射。
则下图交换：
$$
\xymatrix{
X \ar[d]_f \ar[r]_{F_X} & X \ar[d]^f \\
Y \ar[r]^{F_Y} & Y
}
$$
\end{lemma}

\begin{proof}
这直接由如下平凡的代数事实得到：若
$\varphi : A \to B$ 是特征 $p$ 的环同态，则
$\varphi(a^p) = \varphi(a)^p$。
\end{proof}

\begin{lemma}
\label{varieties-lemma-frobenius}
设 $p > 0$ 为素数。设 $X$ 是特征 $p$ 的概形。
则绝对 Frobenius $F_X : X \to X$ 是泛同胚，是整的，且
诱导纯不可分的剩余域扩张。
\end{lemma}

\begin{proof}
这由 Frobenius 自同态 $F_A : A \to A$ 的相应结果得到，其中 $A$
是特征 $p > 0$ 的环。例如参见《代数》一节
\ref{algebra-section-universal-homeomorphism} 的讨论，特别是引理
\ref{algebra-lemma-p-ring-map}。
\end{proof}

\noindent
若我们在固定的底概形上工作，则有 Frobenius 态射的相对版本。

\begin{definition}
\label{varieties-definition-relative-frobenius}
设 $p > 0$ 为素数。设 $S$ 是特征 $p$ 的概形。
设 $X$ 是 $S$ 上的概形。定义
$$
X^{(p)} = X^{(p/S)} = X \times_{S, F_S} S
$$
并将其视为 $S$ 上的概形。应用引理
\ref{varieties-lemma-frobenius-endomorphism-identity}，可知存在唯一的
态射 $F_{X/S} : X \longrightarrow X^{(p)}$，它在 $S$ 上且使下图交换：
$$
\xymatrix{
X \ar[rr]_{F_{X/S}} \ar[rrd] \ar@/^1em/[rrrr]^{F_X}
& & X^{(p)} \ar[rr] \ar[d] & & X \ar[d] \\
& & S \ar[rr]^{F_S} & & S
}
$$
其中右方的方形是笛卡尔的。态射 $F_{X/S}$ 称为
$X/S$ 的\emph{相对 Frobenius 态射}。
\end{definition}

\noindent
注意，$X \mapsto X^{(p)}$ 是一个函子；它是绝对 Frobenius 态射
$F_S : S \to S$ 的基变换函子。关于相对 Frobenius 态射，
我们有与此前相同的引理。

\begin{lemma}
\label{varieties-lemma-relative-frobenius-endomorphism-identity}
设 $p > 0$ 为素数。设 $S$ 是特征 $p$ 的概形。
设 $f : X \to Y$ 是 $S$ 上概形之间的态射。
则下图交换：
$$
\xymatrix{
X \ar[d]_f \ar[r]_{F_{X/S}} & X^{(p)} \ar[d]^{f^{(p)}} \\
Y \ar[r]^{F_{Y/S}} & Y^{(p)}
}
$$
\end{lemma}

\begin{proof}
这由引理 \ref{varieties-lemma-frobenius-endomorphism-identity} 及定义得到。
\end{proof}

\begin{lemma}
\label{varieties-lemma-relative-frobenius}
设 $p > 0$ 为素数。设 $S$ 是特征 $p$ 的概形。
设 $X$ 是 $S$ 上的概形。则相对 Frobenius
$F_{X/S} : X \to X^{(p)}$ 是泛同胚，是整的，且
诱导纯不可分的剩余域扩张。
\end{lemma}

\begin{proof}
由引理 \ref{varieties-lemma-frobenius}，态射 $F_X : X \to X$ 以及
基变换 $h : X^{(p)} \to X$（它是 $F_S$ 的基变换）都是泛同胚。
由于 $h \circ F_{X/S} = F_X$，我们得到 $F_{X/S}$ 是泛同胚
（《态射》引理 \ref{morphisms-lemma-permanence-universal-homeomorphism}）。
由《态射》引理
\ref{morphisms-lemma-universal-homeomorphism} 和
\ref{morphisms-lemma-universally-injective}，我们也得到
$F_{X/S}$ 具有其余性质。
\end{proof}

\begin{lemma}
\label{varieties-lemma-relative-frobenius-omega}
设 $p > 0$ 为素数。设 $S$ 是特征 $p$ 的概形。
设 $X$ 是 $S$ 上的概形。则 $\Omega_{X/S} = \Omega_{X/X^{(p)}}$。
\end{lemma}

\begin{proof}
这转化为如下代数事实。设 $A \to B$ 是特征 $p$ 的环同态。
置 $B' = B \otimes_{A, F_A} A$，并考虑环映射
$F_{B/A} : B' \to B$，$b \otimes a \mapsto b^pa$。
我们要断言的是 $\Omega_{B/A} = \Omega_{B/B'}$。
这是正确的，因为有 $\text{d}(b^pa) = 0$；其中
$\text{d} : B \to \Omega_{B/A}$ 是泛导子，因而 $\text{d}$ 是一个 $B'$-导子。
\end{proof}

\begin{lemma}
\label{varieties-lemma-relative-frobenius-finite}
设 $p > 0$ 为素数。设 $S$ 是特征 $p$ 的概形。
设 $X$ 是 $S$ 上的概形。若 $X \to S$ 局部有限型，则 $F_{X/S}$ 是有限的。
\end{lemma}

\begin{proof}
这转化为如下代数事实。设 $A \to B$ 是特征 $p$ 的有限型环同态。
置 $B' = B \otimes_{A, F_A} A$，并考虑环映射
$F_{B/A} : B' \to B$，$b \otimes a \mapsto b^pa$。
我们要断言的是 $F_{B/A}$ 是有限的。
事实上，若 $x_1, \ldots, x_n \in B$ 是在 $A$ 上的生成元，
则 $x_i$ 在 $B'$ 上整，因为 $x_i^p = F_{B/A}(x_i \otimes 1)$。
因此由《代数》引理
\ref{algebra-lemma-characterize-finite-in-terms-of-integral}，
$F_{B/A} : B' \to B$ 是有限的。
\end{proof}

\begin{lemma}
\label{varieties-lemma-geometrically-reduced-p}
设 $k$ 是特征 $p > 0$ 的域。设 $X$ 是 $k$ 上的概形。
则 $X$ 几何约化，当且仅当 $X^{(p)}$ 约化。
\end{lemma}

\begin{proof}
考虑绝对 Frobenius $F_k : k \to k$。则 $F_k(k) = k^p$；
换言之，$F_k : k \to k$ 同构于 $k$ 到 $k^{1/p}$ 的嵌入。
因此引理由引理 \ref{varieties-lemma-geometrically-reduced} 得到。
\end{proof}

\begin{lemma}
\label{varieties-lemma-inseparable-deg-p-smooth}
设 $k$ 是特征 $p > 0$ 的域。设 $X$ 是 $k$ 上的簇。
以下条件等价：
\begin{enumerate}
\item $X^{(p)}$ 约化；
\item $X$ 几何约化；
\item 存在非空开子集 $U \subset X$，其到 $k$ 光滑。
\end{enumerate}
在此情形，$X^{(p)}$ 是 $k$ 上的簇，且 $F_{X/k} : X \to X^{(p)}$
是次数为 $p^{\dim(X)}$ 的有限主导态射。
\end{lemma}

\begin{proof}
我们已经在引理 \ref{varieties-lemma-geometrically-reduced-p} 中证明了
(1) 与 (2) 等价。我们已经在引理
\ref{varieties-lemma-geometrically-reduced-dense-smooth-open} 中证明了
(2) 蕴含 (3)。若 (3) 成立，则 $U$ 几何约化
（例如见引理 \ref{varieties-lemma-geometrically-regular-smooth}），
从而由引理 \ref{varieties-lemma-generic-points-geometrically-reduced}，
$X$ 几何约化。由此可见 (1)、(2) 与 (3) 等价。

\medskip\noindent
假设 (1)、(2) 与 (3) 成立。由于 $F_{X/k}$ 是同胚
（引理 \ref{varieties-lemma-relative-frobenius}），我们看到 $X^{(p)}$ 是一个簇。
随后由引理 \ref{varieties-lemma-relative-frobenius-finite}，$F_{X/k}$ 是有限的。
它是主导的，因为它是满射。为计算次数
（《态射》定义 \ref{morphisms-definition-degree}），
只需计算 $F_{U/k} : U \to U^{(p)}$ 的次数
（因为由引理 \ref{varieties-lemma-relative-frobenius-endomorphism-identity}，
$F_{U/k} = F_{X/k}|_U$）。
缩小一点 $U$ 后，不妨设存在 'etale 态射
$h : U \to \mathbf{A}^n_k$，见《态射》引理
\ref{morphisms-lemma-smooth-etale-over-affine-space}。
当然 $n = \dim(U)$，因为 $\mathbf{A}^n_k \to \Spec(k)$
是相对维数为 $n$ 的光滑态射，'etale 态射 $h$ 是相对维数为 $0$ 的光滑态射，
而 $U \to \Spec(k)$ 是相对维数为 $\dim(U)$ 的光滑态射，且相对维数正确相加
（《态射》引理 \ref{morphisms-lemma-composition-relative-dimension-d}）。
注意，$h$ 是簇之间的一般有限主导态射，因而 $\deg(h)$ 有定义。
由引理 \ref{varieties-lemma-relative-frobenius-endomorphism-identity}，有交换图
$$
\xymatrix{
U \ar[rr]_{F_{X/k}} \ar[d]_h & & U^{(p)} \ar[d]^{h^{(p)}} \\
\mathbf{A}^n_k \ar[rr]^{F_{\mathbf{A}^n_k/k}} & &
(\mathbf{A}^n_k)^{(p)}
}
$$
由于 $h^{(p)}$ 是 $h$ 的基变换，它也是 'etale 的，
因此 $h^{(p)}$ 也是簇之间的一般有限主导态射。$h^{(p)}$ 的次数是扩张
$k(X^{(p)})/k((\mathbf{A}^n_k)^{(p)})$ 的次数；
该次数与扩张 $k(X)/k(\mathbf{A}^n_k)$ 的次数相同，因为 $h^{(p)}$ 是 $h$ 的基变换
（省略细节）。
由次数的乘法性（《态射》引理
\ref{morphisms-lemma-degree-composition}），只需证明
$F_{\mathbf{A}^n_k/k}$ 的次数为 $p^n$。为此注意
$(\mathbf{A}^n_k)^{(p)} = \mathbf{A}^n_k$，且
$F_{\mathbf{A}^n_k/k}$ 由将坐标映射为其 $p$ 次幂的映射给出。
\end{proof}

\begin{remark}
\label{varieties-remark-n-fold-relative-frobenius}
设 $p > 0$ 为素数。设 $S$ 是特征 $p$ 的概形。
设 $X$ 是 $S$ 上的概形。对 $n \geq 1$，
$$
X^{(p^n)} = X^{(p^n/S)} = X \times_{S, F_S^n} S
$$
并将其视为 $S$ 上的概形。注意，$X \mapsto X^{(p^n)}$
是一个函子。应用引理 \ref{varieties-lemma-frobenius-endomorphism-identity}，
可知 $F_{X/S, n} = (F_X^n, \text{id}_S) : X \longrightarrow X^{(p^n)}$
是一个在 $S$ 上的态射，并且适合下列交换图：
$$
\xymatrix{
X \ar[rr]_{F_{X/S, n}} \ar[rrd] \ar@/^1em/[rrrr]^{F_X^n}
& & X^{(p^n)} \ar[rr] \ar[d] & & X \ar[d] \\
& & S \ar[rr]^{F_S^n} & & S
}
$$
其中右方的方形是笛卡尔的。态射 $F_{X/S, n}$ 有时称为
\emph{$n$ 重相对 Frobenius 态射}（$X/S$）。
这是有意义的，因为我们有公式
$$
F_{X/S, n} =
F_{X^{(p^{n - 1})}/S} \circ \ldots \circ F_{X^{(p)}/S} \circ F_{X/S}
$$
这表明 $F_{X/S, n}$ 是 $n$ 个相对 Frobenius 的复合。由于
$$
F_{X^{(p^m)}/S} = F_{X^{(p^{m - 1})}/S}^{(p)} = \ldots = F_{X/S}^{(p^m)}
$$
（细节略去），我们还得到
$$
F_{X/S, n} =
F_{X/S}^{(p^{n - 1})} \circ \ldots \circ F_{X/S}^{(p)} \circ F_{X/S}
$$
\end{remark}




\section{粘合一维环}
\label{varieties-section-glueing-dim-1}

\noindent
本节给出一些代数预备，用于证明：分离概形的有限个余维 $1$ 点包含于某个仿射开集。

\begin{situation}
\label{varieties-situation-glue}
给定环的交换图
$$
\xymatrix{
A \ar[r] & K \\
R \ar[u] \ar[r] & B \ar[u]
}
$$
其中 $K$ 是域，$A$、$B$ 是 $K$ 的子环，且分式域都是 $K$。最后，
$R = A \times_K B = A \cap B$。
\end{situation}

\begin{lemma}
\label{varieties-lemma-glue-valuation-ring}
在情形 \ref{varieties-situation-glue} 中，假设 $B$ 是赋值环。
那么对每个单元 $u$（来自 $A$），或者 $u \in R$，或者 $u^{-1} \in R$。
\end{lemma}

\begin{proof}
事实上，若 $c$（$u$ 在 $K$ 中的像）属于 $B$，则 $u \in R$。
否则 $c^{-1} \in B$
（《代数》引理 \ref{algebra-lemma-valuation-ring-x-or-x-inverse}），
从而 $u^{-1} \in R$。
\end{proof}

\noindent
下面的引理解释稍后反复出现的条件
“$A \otimes B \to K$ 是满射”的含义。

\begin{lemma}
\label{varieties-lemma-glue-separated}
在情形 \ref{varieties-situation-glue} 中，假设 $A$ 是维数为 $1$ 的 Noether 环。
以下条件等价：
\begin{enumerate}
\item $A \otimes B \to K$ 不是满射；
\item 存在离散赋值环 $\mathcal{O} \subset K$，它包含 $A$ 与 $B$。
\end{enumerate}
\end{lemma}

\begin{proof}
(2) 蕴含 (1) 是显然的。另一方面，若 $A \otimes B \to K$
不是满射，则其像 $C \subset K$ 不是域，因而 $C$ 有非零极大理想
$\mathfrak m$。取 $\mathcal{O} \subset K$ 为支配 $C_\mathfrak m$ 的赋值环。
由《代数》引理
\ref{algebra-lemma-krull-akizuki} 应用于 $A \subset \mathcal{O}$，
环 $\mathcal{O}$ 是 Noether 环。因此，由《代数》引理
\ref{algebra-lemma-valuation-ring-Noetherian-discrete}，
$\mathcal{O}$ 是离散赋值环。
\end{proof}

\begin{lemma}
\label{varieties-lemma-semi-local}
在情形 \ref{varieties-situation-glue} 中，假设
\begin{enumerate}
\item $A$ 是维数为 $1$ 的 Noether 半局部整环；
\item $B$ 是离散赋值环；
\end{enumerate}
则有以下两种可能：
\begin{enumerate}
\item[(a)] 若 $A^*$ 不包含于 $R$，则
$\Spec(A) \to \Spec(R)$ 与 $\Spec(B) \to \Spec(R)$
都是覆盖 $\Spec(R)$ 的开浸入，并且 $K = A \otimes_R B$；
\item[(b)] 若 $A^*$ 包含于 $R$，则 $B$ 支配 $A$ 在某个极大理想处的
局部环，并且 $A \otimes B \to K$ 不是满射。
\end{enumerate}
\end{lemma}

\begin{proof}
假设 (a) 成立，则存在一个单元 $u$，它属于 $A$，其在 $K$ 中的像属于 $B$ 的极大理想。
于是 $u$ 是 $R$ 中的非零因子，并且对每个 $a \in A$，存在某个 $n$ 使
$u^n a \in R$。因此 $A = R_u$。

\medskip\noindent
令 $\mathfrak m_A$ 为 $A$ 的 Jacobson 根。取非零元
$x \in \mathfrak m_A$。由于 $\dim(A) = 1$，可知 $K = A_x$。
将 $x$ 替换为某个 $x^n u^m$（其中 $n \geq 1$ 且 $m \in \mathbf{Z}$）后，
可假设 $x$ 映到 $B$ 的单元。于是对每个 $b \in B$，都有 $x^nb$
属于 $R$ 的像，对某个 $n$ 成立。因此 $B = R_x$。

\medskip\noindent
设 $z \in R$。若 $z \not \in \mathfrak m_A$ 且 $z$ 不映到
$\mathfrak m_B$ 中的元素，则 $z$ 可逆。因此 $x + u$ 在 $R$ 中可逆。
于是 $\Spec(R) = D(x) \cup D(u)$。上面已经看到
$D(u) = \Spec(A)$ 且 $D(x) = \Spec(B)$。

\medskip\noindent
情形 (b)。若 $x \in \mathfrak m_A$，则 $1 + x$ 是单元，因而
$1 + x \in R$，即 $x \in R$。所以
$\mathfrak m_A \subset R \subset A$。事实上，在此情形 $A$ 在 $R$ 上整。
将
$A/\mathfrak m_A = \kappa_1 \times \ldots \times \kappa_n$
写成域的乘积。设
$x = (c_1, \ldots, c_r, 0, \ldots, 0)$
是一个其中 $c_i \not = 0$ 的元素。则
$$
x^2 - x(c_1, \ldots, c_r, 1, \ldots, 1)  = 0
$$
由于 $R$ 包含所有单元，可知 $A/\mathfrak m_A$ 在 $R$ 在其中的像上整，
从而 $A$ 在 $R$ 上整。由此得到 $R \subset A \subset B$，因为 $B$ 整闭。
此外，若非零元 $x \in \mathfrak m_A$，则
$K = A_x = \bigcup x^{-n}A = \bigcup x^{-n}R$。
因此 $x^{-1} \not \in B$，即 $x \in \mathfrak m_B$。
我们得到 $\mathfrak m_A \subset \mathfrak m_B$。
于是 $A \cap \mathfrak m_B$ 是 $A$ 的极大理想，证明完成。
\end{proof}

\begin{lemma}
\label{varieties-lemma-semi-local-dimension-one-conductor}
设 $B$ 是维数为 $1$ 的 Noether 半局部整环。
设 $B'$ 是 $B$ 在其分式域中的整闭包。
则 $B'$ 是半局部 Dedekind 整环。
设 $x$ 是 $B'$ 的 Jacobson 根中的非零元。
则对每个 $y \in B'$，存在某个 $n$ 使得 $x^n y \in B$。
\end{lemma}

\begin{proof}
令 $\mathfrak m_B$ 为 $B$ 的 Jacobson 根。
$B'$ 的结构由《代数》引理
\ref{algebra-lemma-integral-closure-Dedekind} 给出。
给定引理陈述中的 $x, y \in B'$，考虑子环
$B \subset A \subset B'$，它由 $x$ 与 $y$ 生成。则 $A$ 在 $B$ 上有限（《代数》引理
\ref{algebra-lemma-characterize-finite-in-terms-of-integral}）。
由于 $B$ 与 $A$ 的分式域相同，可知有限模 $A/B$ 支集于 $B$ 的闭点集。
因此有 $\mathfrak m_B^n A \subset B$，其中 $n$ 适当。
此外，$\Spec(B') \to \Spec(A)$ 是满射（《代数》引理
\ref{algebra-lemma-integral-overring-surjective}），所以 $A$ 也是半局部的。
同样可知 $x$ 属于 Jacobson 根 $\mathfrak m_A$，这是 $A$ 的 Jacobson 根。
注意 $\mathfrak m_A = \sqrt{\mathfrak m_B A}$。
因此有 $x^m y \in \mathfrak m_B A$，其中 $m$ 适当。
于是 $x^{nm} y \in B$。
\end{proof}

\begin{lemma}
\label{varieties-lemma-semi-local-both-side}
在情形 \ref{varieties-situation-glue} 中，假设
\begin{enumerate}
\item $A$ 是维数为 $1$ 的 Noether 半局部整环；
\item $B$ 是维数为 $1$ 的 Noether 半局部整环；
\item $A \otimes B \to K$ 是满射。
\end{enumerate}
则 $\Spec(A) \to \Spec(R)$ 与 $\Spec(B) \to \Spec(R)$
是覆盖 $\Spec(R)$ 的开浸入，并且 $K = A \otimes_R B$。
\end{lemma}

\begin{proof}
特殊情形：$B$ 在 $K$ 中整闭。这意味着 $B$ 是 Dedekind 整环
（《代数》引理 \ref{algebra-lemma-characterize-Dedekind}），
所以它在极大理想处的所有局部化都是离散赋值环。
令 $\mathfrak m_1, \ldots, \mathfrak m_r$ 为 $B$ 的极大理想。置
$$
R_1 = A \times_K B_{\mathfrak m_1}
$$
注意到 $A \otimes_{R_1} B_{\mathfrak m_1} \to K$ 是满射，
由引理 \ref{varieties-lemma-semi-local} 可知，$A$ 与
$B_{\mathfrak m_1}$ 定义覆盖 $\Spec(R_1)$ 的开子概形，并且
$K = A \otimes_{R_1} B_{\mathfrak m_1}$。特别地，$R_1$ 是维数为 $1$ 的
Noether 半局部环。归纳定义
$$
R_{i + 1} = R_i \times_K B_{\mathfrak m_{i + 1}}
$$
其中 $i = 1, \ldots, r - 1$。注意，$R = R_r$，因为
$B = B_{\mathfrak m_1} \cap \ldots \cap B_{\mathfrak m_r}$
（见《代数》引理
\ref{algebra-lemma-normal-domain-intersection-localizations-height-1}）。
由这一归纳过程，$R \to A$ 定义开浸入
$\Spec(A) \to \Spec(R)$。另一方面，极大理想
$\mathfrak n_i$（属于 $R$ 且不在这个开集内）对应于极大理想
$\mathfrak m_i$（属于 $B$）；事实上环同态
$R \to B$ 定义同构
$R_{\mathfrak n_i} \to B_{\mathfrak m_i}$（细节略去；提示：
在每一步中我们恰好向 $\Spec(R_i)$ 添加了一个极大理想）。
所以 $\Spec(B) \to \Spec(R)$ 是开浸入，如所需。

\medskip\noindent
一般情形。令 $B' \subset K$ 为 $B$ 的整闭包。
见引理 \ref{varieties-lemma-semi-local-dimension-one-conductor}。
于是特殊情形适用于 $R' = A \times_K B'$。
取 $x \in R'$，使它不属于 $A$ 的极大理想，而属于 $B'$ 的极大理想
（见《代数》引理 \ref{algebra-lemma-chinese-remainder}）。
由引理 \ref{varieties-lemma-semi-local-dimension-one-conductor}，
存在整数 $n$ 使得 $x^n \in R = A \times_K B$。
将 $x$ 替换为 $x^n$，于是 $x \in R$。对每个 $y \in R'$，
存在整数 $n$ 使得 $x^n y \in R$。另一方面，显然 $R'_x = A$。
因此 $R_x = A$。
交换 $A$ 与 $B$ 的角色，还可找到 $y \in R$ 使得 $B = R_y$。
注意，同时对 $x$ 与 $y$ 取逆后只剩下 $(0)$ 这一素理想。
因此 $K = R_{xy} = R_x \otimes_R R_y$。
证明完成。
\end{proof}

\begin{lemma}
\label{varieties-lemma-glue-a-bunch-of-local-rings}
设 $K$ 为域。设 $A_1, \ldots, A_r \subset K$ 是维数为 $1$ 的
Noether 半局部环，其分式域为 $K$。若
$A_i \otimes A_j \to K$ 对所有 $i \not = j$ 都是满射，
则存在 Noether 半局部整环 $A \subset K$，其维数为 $1$，并且它包含于
$A_1, \ldots, A_r$，并且满足：
\begin{enumerate}
\item $A \to A_i$ 诱导开浸入 $j_i : \Spec(A_i) \to \Spec(A)$；
\item $\Spec(A)$ 是各开集 $j_i(\Spec(A_i))$ 的并；
\item $\Spec(A)$ 的每个闭点恰好属于这些开集中的一个。
\end{enumerate}
\end{lemma}

\begin{proof}
事实上，可以取 $A = A_1 \cap \ldots \cap A_r$。首先注意，一旦证明了
(1) 和 (2)，则 (3) 由假设 $A_i \otimes A_j \to K$ 是满射推出，
因为若
$\mathfrak m \in j_i(\Spec(A_i)) \cap j_j(\Spec(A_j))$，则
$A_i \otimes A_j \to K$ 的像落在 $A_\mathfrak m$ 中。
为证明 (1) 和 (2)，对 $r$ 归纳。
若 $r > 1$，由归纳假设，(1) 和 (2) 对
$B = A_2 \cap \ldots \cap A_r$ 成立。再应用引理
\ref{varieties-lemma-semi-local-both-side}，可知它们对
$A = A_1 \cap B$ 也成立。
\end{proof}

\begin{lemma}
\label{varieties-lemma-create-globally-generated}
设 $A$ 为分式域是 $K$ 的整环。设 $B_1, \ldots, B_r \subset K$
为 Noether $1$-维半局部整环，其分式域均是 $K$。若
$A \otimes B_i \to K$ 对 $i = 1, \ldots, r$ 都是满射，
则存在 $x \in A$，使得 $x^{-1}$ 属于 $B_i$ 的 Jacobson 根，
对 $i = 1, \ldots, r$ 都成立。
\end{lemma}

\begin{proof}
设 $B_i'$ 为 $B_i$ 在 $K$ 中的整闭包。假设能找到非零元
$x \in A$，使得 $x^{-1}$ 属于 $B'_i$ 的 Jacobson 根，
对 $i = 1, \ldots, r$ 都成立。则由引理
\ref{varieties-lemma-semi-local-dimension-one-conductor}，
将 $x$ 替换为某个幂后，得到 $x^{-1} \in B_i$。
由于 $\Spec(B'_i) \to \Spec(B_i)$ 是满射，可知
$x^{-1}$ 也属于 $B_i$ 的 Jacobson 根。因此可以假设每个
$B_i$ 都是半局部 Dedekind 整环。

\medskip\noindent
若 $B_i$ 不是局部环，则从列表中删去 $B_i$，并加入有限个局部环
$(B_i)_\mathfrak m$。因此可以假设 $B_i$ 是离散赋值环，
对 $i = 1, \ldots, r$ 都成立。

\medskip\noindent
设 $v_i : K \to \mathbf{Z}$（$i = 1, \ldots, r$）为相应的离散赋值
（见《代数》引理 \ref{algebra-lemma-characterize-Dedekind}）。
我们要找非零元 $x \in A$，使得 $v_i(x) < 0$，
对 $i = 1, \ldots, r$ 都成立。对此对 $r$ 归纳。

\medskip\noindent
若 $r = 1$ 且结论不成立，则 $A \subset B$，并且映射
$A \otimes B \to K$ 不是满射，矛盾。

\medskip\noindent
若 $r > 1$，则由归纳假设，可以找到非零元 $x \in A$，
使得 $v_i(x) < 0$ 对 $i = 1, \ldots, r - 1$ 都成立。
若 $v_r(x) < 0$，则已经完成；所以可假设 $v_r(x) \geq 0$。
由基本情形，可找到非零元 $y \in A$，使得 $v_r(y) < 0$。
将 $x$ 替换为某个幂后，可假设 $v_i(x) < v_i(y)$，
对 $i = 1, \ldots, r - 1$ 都成立。于是 $x + y$ 就是所需元素。
\end{proof}

\begin{lemma}
\label{varieties-lemma-power-equal}
设 $A$ 是维数为 $1$ 的 Noether 局部环。
令 $L = \prod A_\mathfrak p$，其中乘积取遍 $A$ 的极小素理想。
设 $a_1, a_2 \in \mathfrak m_A$ 映到 $L$ 中的同一元素。
则有 $a_1^n = a_2^n$，对某个 $n > 0$ 成立。
\end{lemma}

\begin{proof}
写成 $a_1 = a_2 + x$。则 $x$ 在 $L$ 中的像为零。因此 $x$ 是
$A$ 的幂零元，因为 $\bigcap \mathfrak p$ 是 $(0)$ 的根基，而
零化理想 $I$（即 $x$ 的零化理想）包含极大理想的某个幂，这是由于对所有极小素理想都有
$\mathfrak p \not \in V(I)$。设 $x^k = 0$ 且
$\mathfrak m^n \subset I$。则
\begingroup
\small
$$
a_1^{k + n} = a_2^{k + n} + {n + k \choose 1} a_2^{n + k - 1} x +
{n + k \choose 2} a_2^{n + k - 2} x^2 + \ldots +
{n + k \choose k - 1} a_2^{n + 1} x^{k - 1} = a_2^{n + k}
$$
\endgroup
因为 $a_2 \in \mathfrak m_A$。
\end{proof}

\begin{lemma}
\label{varieties-lemma-power-works}
设 $A$ 是维数为 $1$ 的 Noether 局部环。
令 $L = \prod A_\mathfrak p$ 且 $I = \bigcap \mathfrak p$，
其中乘积与交均取遍 $A$ 的极小素理想。设 $f \in L$ 是形如
$f = i + a$ 的元素，其中 $a \in \mathfrak m_A$ 且
$i \in IL$。则 $f$ 的某个幂属于 $A \to L$ 的像。
\end{lemma}

\begin{proof}
由于 $A$ 是 Noether 环，有 $I^t = 0$，对某个 $t > 0$ 成立。
假设已知 $f = a + i$，其中 $i \in I^kL$。
则 $f^n = a^n + na^{n - 1}i \bmod I^{k + 1}L$。
所以只需证明：$na^{n - 1}i$ 属于 $I^k \to I^kL$ 的像，
对某个 $n \gg 0$ 成立。
为此，取 $g \in A$，使得 $\mathfrak m_A = \sqrt{(g)}$
（《代数》引理 \ref{algebra-lemma-height-1}）。
于是，例如由《代数》命题
\ref{algebra-proposition-dimension-zero-ring} 可知 $L = A_g$。
另一方面，存在某个 $n$ 使得 $a^n \in (g)$。
因此，可以清除 $L$ 中元素的分母，方法是乘以 $a$ 的高次幂。
\end{proof}

\begin{lemma}
\label{varieties-lemma-good-intersection}
设 $A$ 是维数为 $1$ 的 Noether 局部环。
令 $L = \prod A_\mathfrak p$，其中乘积取遍 $A$ 的极小素理想。
设 $K \to L$ 为整的环同态。
则存在 $a \in \mathfrak m_A$ 与 $x \in K$，它们映到
$L$ 中的同一元素，并且 $\mathfrak m_A = \sqrt{(a)}$。
\end{lemma}

\begin{proof}
由引理 \ref{varieties-lemma-power-works}，可以把 $A$ 替换为
$A/(\bigcap \mathfrak p)$，并假设 $A$ 是约化环
（略去一些细节）。
还可以把 $K$ 替换为 $K \to L$ 的像。
于是 $K$ 是约化环。映射 $\Spec(L) \to \Spec(K)$ 是满射且是闭映射
（细节略去）。因此 $\Spec(K)$ 是有限离散空间。
所以 $K$ 是有限多个域的积。

\medskip\noindent
令 $\mathfrak p_j$（$j = 1, \ldots, m$）为 $A$ 的极小素理想。
令 $L_j$ 为 $A_j$ 的分式域，从而
$L = \prod_{j = 1, \ldots, m} L_j$。
令 $A_j$ 为 $A/\mathfrak p_j$ 的正规化。则
$A_j$ 是至少有一个极大理想的半局部 Dedekind 整环，见《代数》引理
\ref{algebra-lemma-integral-closure-Dedekind}。
令 $n$ 为 $A_1, \ldots, A_m$ 的极大理想个数之和。
对这样的极大理想 $\mathfrak m \subset A_j$，考虑函数
$$
v_{\mathfrak m} : L \to L_j \to \mathbf{Z} \cup \{\infty\}
$$
其中第二个箭头是与离散赋值环 $(A_j)_{\mathfrak m}$ 相应的离散赋值，
并通过把 $0$ 映到 $\infty$ 来扩张。这样得到 $n$ 个函数
$v_1, \ldots, v_n : L \to \mathbf{Z} \cup \{\infty\}$。
我们将找出 $x \in K$，使得 $v_i(x) < 0$，
对所有 $i = 1, \ldots, n$ 成立。

\medskip\noindent
首先断言：对每个 $i$，存在 $x \in K$ 使得 $v_i(x) < 0$。
事实上，假设 $v_i$ 对应于 $\mathfrak m \subset A_j$。
若 $v_i(x) \geq 0$ 对所有 $x \in K$ 都成立，则 $K$ 映入
$(A_j)_{\mathfrak m}$，这是分式域 $L_j$（$A_j$ 的分式域）中的子环。
$K$ 在 $L_j$ 中的像是域；$L_j$ 在该域上是代数的，这是由《代数》引理
\ref{algebra-lemma-integral-under-field} 得到的。
结合起来，这与《代数》引理
\ref{algebra-lemma-valuation-ring-cap-field-finite} 矛盾。

\medskip\noindent
假设已经找到 $x \in K$，使得
$v_1(x) < 0, \ldots, v_r(x) < 0$，其中某个 $r < n$。
若 $v_{r + 1}(x) < 0$，则 $x$ 对 $r + 1$ 也适用。
否则，由上一段的结果，可以选取某个 $y \in K$，使
$v_{r + 1}(y) < 0$。将 $x$ 替换为某个 $x^n$（其中 $n > 0$）后，
可以假设 $v_i(x) < v_i(y)$ 对 $i = 1, \ldots, r$ 成立。
于是由赋值的性质，
$v_j(x + y) = v_j(x) < 0$ 对 $j = 1, \ldots, r$ 成立；类似地，
$v_{r + 1}(x + y) = v_{r + 1}(y) < 0$。
归纳可得 $x \in K$，使得 $v_i(x) < 0$
对 $i = 1, \ldots, n$ 成立。

\medskip\noindent
特别地，元素 $x \in K$ 在 $K$ 的每个因子中的投影均非零
（回忆 $K$ 是有限多个域的积；如果 $x$ 的某个分量为零，
则某个值 $v_i(x)$ 会是 $\infty$）。因此 $x$ 可逆，而
$x^{-1} \in K$ 满足
$\infty > v_i(x^{-1}) > 0$，对所有 $i$ 都成立。由引理
\ref{varieties-lemma-semi-local-dimension-one-conductor}，可知对某个
$e < 0$，元素 $x^e \in K$ 映到
$\mathfrak m_A/\mathfrak p_j \subset A/\mathfrak p_j$ 中，
对所有 $j = 1, \ldots, m$ 都成立。
注意，映射
$\mathfrak m_A \to \prod \mathfrak m_A/\mathfrak p_j$
的余核被 $\mathfrak m_A$ 的某个幂零化。因此，把 $e$ 替换为一个
更小的 $e$ 后，可找到 $a \in \mathfrak m_A$，使它在
$\mathfrak m_A/\mathfrak p_j$ 中的像等于 $x^e$ 的像。
于是二元组 $(a, x^e)$ 满足引理的结论。
\end{proof}

\begin{lemma}
\label{varieties-lemma-localization-semi-local}
设 $A$ 是环。设 $\mathfrak p_1, \ldots, \mathfrak p_r$
是 $A$ 的有限个素理想。令
$S = A \setminus \bigcup \mathfrak p_i$。则 $S$ 是乘法系，
而 $S^{-1}A$ 是半局部环，其极大理想对应于集合
$\{\mathfrak p_i\}$ 的极大元。
\end{lemma}

\begin{proof}
若 $a, b \in A$ 且 $a, b \in S$，则
$a, b \not \in \mathfrak p_i$，从而
$ab \not \in \mathfrak p_i$，所以 $ab \in S$。又有 $1 \in S$。
因此 $S$ 是 $A$ 的乘法子集。由《代数》引理
\ref{algebra-lemma-spec-localization} 中对 $\Spec(S^{-1}A)$ 的描述，
以及《代数》引理 \ref{algebra-lemma-silly}，可知
$S^{-1}A$ 的素理想对应于 $A$ 中包含于某个
$\mathfrak p_i$ 的素理想。
所以 $S^{-1}A$ 的极大理想一一对应于集合
$\{\mathfrak p_1, \ldots, \mathfrak p_r\}$ 中关于包含关系的极大元。
\end{proof}

\section{一维 Noether 概形}
\label{varieties-section-dimension-one}

\noindent
本节的主要结果是：维数为 $1$ 的 Noether 分离概形具有
丰沛可逆层。见命题
\ref{varieties-proposition-dim-1-noetherian-separated-has-ample}。

\begin{lemma}
\label{varieties-lemma-affine}
设 $X$ 是概形，其所有局部环都是维数
$\leq 1$ 的 Noether 环。设 $U \subset X$ 是逆紧开集，并以
$j : U \to X$ 表示包含态射。则
$R^pj_*\mathcal{F} = 0$，其中 $p > 0$；
这对每个拟凝聚 $\mathcal{O}_U$-模 $\mathcal{F}$ 都成立。
\end{lemma}

\begin{proof}
可以在茎上检验 $R^pj_*\mathcal{F}$ 的消失性。
$R^qj_*$ 的形成与平坦基变换可交换，见《概形上同调》引理
\ref{coherent-lemma-flat-base-change-cohomology}。
因此可假设 $X$ 是维数 $\leq 1$ 的 Noether 局部环的谱。
在此情形，$X$ 有一个闭点 $x$，以及有限个其他点
$x_1, \ldots, x_n$；它们特化到 $x$，但彼此之间不发生特化
（见《代数》引理
\ref{algebra-lemma-Noetherian-irreducible-components}）。
若 $x \in U$，则 $U = X$，结论显然。否则，在必要时重编号各点后，
有 $U = \{x_1, \ldots, x_r\}$，对某个 $r$ 成立。
于是 $U$ 是仿射的
（《概形》引理 \ref{schemes-lemma-scheme-finite-discrete-affine}）。
因此结论由《概形上同调》引理
\ref{coherent-lemma-relative-affine-vanishing} 得到。
\end{proof}

\begin{lemma}
\label{varieties-lemma-open-in-affine-curve-affine}
设 $X$ 是仿射概形，其所有局部环都是维数 $\leq 1$ 的 Noether 环。
则任意拟紧开集 $U \subset X$ 都是仿射的。
\end{lemma}

\begin{proof}
以 $j : U \to X$ 表示包含态射。设 $\mathcal{F}$
是拟凝聚 $\mathcal{O}_U$-模。
由引理 \ref{varieties-lemma-affine}，高阶正像
$R^pj_*\mathcal{F}$ 均为零。$\mathcal{O}_X$-模 $j_*\mathcal{F}$
是拟凝聚的
（《概形》引理 \ref{schemes-lemma-push-forward-quasi-coherent}）。
因此，由《概形上同调》引理
\ref{coherent-lemma-quasi-coherent-affine-cohomology-zero}，
它的高阶上同调群均为零。
由 Leray 谱序列（《上同调》引理
\ref{cohomology-lemma-apply-Leray}），有
$H^p(U, \mathcal{F}) = 0$，对所有 $p > 0$ 都成立。
所以 $U$ 是仿射的，例如可由《概形上同调》引理
\ref{coherent-lemma-quasi-compact-h1-zero-covering} 得到。
\end{proof}

\begin{lemma}
\label{varieties-lemma-complement-codim-1-closed-points}
设 $X$ 是概形，$U \subset X$ 是开集。假设
\begin{enumerate}
\item $U$ 是 $X$ 的逆紧开集；
\item $X \setminus U$ 是离散的；
\item 对 $x \in X \setminus U$，局部环
$\mathcal{O}_{X, x}$ 是维数 $\leq 1$ 的 Noether 环。
\end{enumerate}
则 (1) 存在可逆 $\mathcal{O}_X$-模 $\mathcal{L}$
以及一个截面 $s$，使得 $U = X_s$；并且 (2) 映射
$\Pic(X) \to \Pic(U)$ 是满射。
\end{lemma}

\begin{proof}
令 $X \setminus U = \{x_i; i \in I\}$。
选择仿射开集 $U_i \subset X$，使得 $x_i \in U_i$ 且
$x_j \not \in U_i$，其中 $j \not = i$。这由条件 (2) 可以做到。
写 $U_i = \Spec(A_i)$。令 $\mathfrak m_i \subset A_i$
为与 $x_i$ 对应的极大理想。由对局部环的假设，只有有限多个素理想
$\mathfrak q \subset \mathfrak m_i$ 满足
$\mathfrak q \not = \mathfrak m_i$（见《代数》引理
\ref{algebra-lemma-Noetherian-irreducible-components}）。
所以，由素理想回避（《代数》引理
\ref{algebra-lemma-silly}），可以找到 $f_i \in \mathfrak m_i$，
使它不属于上述任何素理想。于是
$V(f_i) = \{\mathfrak m_i\} \amalg Z_i$，其中某个闭子集
$Z_i \subset U_i$；这是因为 $Z_i$ 是
$V(f_i)$ 中逆紧且对特化封闭的开子集，见《代数》引理
\ref{algebra-lemma-constructible-stable-specialization-closed}。
缩小 $U_i$ 后，可假设 $V(f_i) = \{x_i\}$。于是
$$
\mathcal{U} : X = U \cup \bigcup U_i
$$
是 $X$ 的开覆盖。考虑 $2$-余循环，它取值于
$\mathcal{O}_X^*$：由 $f_i$ 在 $U \cap U_i$ 上给出，并由
$f_i/f_j$ 在 $U_i \cap U_j$ 上给出。它定义一个线丛
$\mathcal{L}$；而截面 $s$ 由 $1$ 在 $U$ 上定义、由 $f_i$
在 $U_i$ 上定义，它满足引理的陈述。

\medskip\noindent
设 $\mathcal{N}$ 是可逆 $\mathcal{O}_U$-模。
设 $N_i$ 为可逆 $(A_i)_{f_i}$-模，使得
$\mathcal{N}|_{U \cap U_i}$ 等于 $\tilde N_i$。
注意 $(A_{\mathfrak m_i})_{f_i}$ 是 Artin 环
（因为它是零维 Noether 环，见《代数》引理
\ref{algebra-lemma-Noetherian-dimension-0}）。
所以它是局部环的乘积
（《代数》引理 \ref{algebra-lemma-artinian-finite-length}），
因而其 Picard 群平凡。因此，缩小 $U_i$
（即把 $A_i$ 替换为某个 $(A_i)_g$，其中 $g \in A_i$ 且
$g \not \in \mathfrak m_i$）后，
可假设 $N_i = (A_i)_{f_i}$，即
$\mathcal{N}|_{U \cap U_i}$ 平凡。在这种情形，如何将
$\mathcal{N}$ 扩张为 $X$ 上的可逆层是显然的
（在每个 $U_i$ 上用平凡可逆模扩张）。
\end{proof}

\begin{lemma}
\label{varieties-lemma-find-globally-generated}
设 $X$ 是整分离概形。设 $U \subset X$ 是非空仿射开集，使得
$X \setminus U$ 是有限点集
$x_1, \ldots, x_r$，其中 $\mathcal{O}_{X, x_i}$
是维数为 $1$ 的 Noether 环。
则存在全局生成的可逆 $\mathcal{O}_X$-模
$\mathcal{L}$ 以及一个截面 $s$，使得 $U = X_s$。
\end{lemma}

\begin{proof}
写 $U = \Spec(A)$，并令 $K$ 为 $X$ 的函数域。
写 $B_i = \mathcal{O}_{X, x_i}$ 且
$\mathfrak m_i = \mathfrak m_{x_i}$。
由于 $x_i \not \in U$，可知开集
$U \times_X \Spec(B_i)$（属于 $\Spec(B_i)$）只有一个点，即
$U \times_X \Spec(B_i) = \Spec(K)$。
由于 $X$ 是分离的，可知 $\Spec(K)$ 是
$U \times \Spec(B_i)$ 的闭子概形，即映射
$A \otimes B_i \to K$ 是满射。
由引理 \ref{varieties-lemma-create-globally-generated}，可以找到非零元
$f \in A$，使得 $f^{-1} \in \mathfrak m_i$ 对
$i = 1, \ldots, r$ 都成立。
选择开集 $x_i \in U_i \subset X$，使得
$f^{-1} \in \mathcal{O}(U_i)$。于是
$$
\mathcal{U} : X = U \cup \bigcup U_i
$$
是 $X$ 的开覆盖。考虑 $2$-余循环，它取值于
$\mathcal{O}_X^*$：由 $f$ 在 $U \cap U_i$ 上给出，并由
$1$ 在 $U_i \cap U_j$ 上给出。它定义一个线丛
$\mathcal{L}$，并且该线丛有两个截面：
\begin{enumerate}
\item 截面 $s$ 由 $1$ 在 $U$ 上定义，由
$f^{-1}$ 在 $U_i$ 上定义；它满足引理的陈述；
\item 截面 $t$ 由 $f$ 在 $U$ 上定义，并由
$1$ 在 $U_i$ 上定义。
\end{enumerate}
注意 $X_t \supset U_1 \cup \ldots \cup U_r$。
因此 $s, t$ 生成 $\mathcal{L}$，引理得证。
\end{proof}

\begin{lemma}
\label{varieties-lemma-enough-globally-generated-ample}
设 $X$ 是拟紧概形。若对每个 $x \in X$ 都存在二元组
$(\mathcal{L}, s)$，其中 $\mathcal{L}$ 是全局生成的可逆层，
$s$ 是全局截面，且 $x \in X_s$、$X_s$ 是仿射的，
则 $X$ 具有丰沛可逆层。
\end{lemma}

\begin{proof}
由于 $X$ 拟紧，可以找到有限多个二元组
$(\mathcal{L}_i, s_i)$（$i = 1, \ldots, n$），使得
$\mathcal{L}_i$ 全局生成，$X_{s_i}$ 是仿射的，并且
$X = \bigcup X_{s_i}$。
再由 $X$ 拟紧，对每个 $i$ 可以找到有限多个截面
$t_{i, j}$，它们属于 $\mathcal{L}_i$（$j = 1, \ldots, m_i$），使得
$X = \bigcup X_{t_{i, j}}$。置 $t_{i, 0} = s_i$。
考虑可逆层
$$
\mathcal{L} = \mathcal{L}_1 
\otimes_{\mathcal{O}_X} \ldots
\otimes_{\mathcal{O}_X} \mathcal{L}_n
$$
及其全局截面
$$
\tau_J = t_{1, j_1} \otimes \ldots \otimes t_{n, j_n}
$$
由《性质》引理 \ref{properties-lemma-affine-cap-s-open}，
开集 $X_{\tau_J}$ 在只要有 $j_i = 0$（对某个 $i$）时就是仿射的。
不难看出这些开集覆盖 $X$。所以按定义，$\mathcal{L}$ 是丰沛的。
\end{proof}

\begin{lemma}
\label{varieties-lemma-dim-1-noetherian-integral-separated-has-ample}
设 $X$ 是维数为 $1$ 的 Noether 整分离概形。
则 $X$ 具有丰沛可逆层。
\end{lemma}

\begin{proof}
选择仿射开覆盖 $X = U_1 \cup \ldots \cup U_n$。
由于 $X$ 是 Noether 的，每个集合 $X \setminus U_i$ 都是有限的。
因此，由引理 \ref{varieties-lemma-find-globally-generated}，
可以找到二元组 $(\mathcal{L}_i, s_i)$，
其中 $\mathcal{L}_i$ 是全局生成的可逆层，
$s_i$ 是全局截面，且 $U_i = X_{s_i}$。
由引理 \ref{varieties-lemma-enough-globally-generated-ample}，
得到 $X$ 具有丰沛可逆层。
\end{proof}

\begin{lemma}
\label{varieties-lemma-surjective-pic-birational-finite}
设 $f : X \to Y$ 是概形的有限态射。假设存在开集
$V \subset Y$，使得 $f^{-1}(V) \to V$ 是同构，并且
$Y \setminus V$ 是离散空间。则每个可逆
$\mathcal{O}_X$-模都是某个可逆 $\mathcal{O}_Y$-模的拉回。
\end{lemma}

\begin{proof}
我们将使用 $\Pic(X) = H^1(X, \mathcal{O}_X^*)$，见《上同调》引理
\ref{cohomology-lemma-h1-invertible}。
考虑阿贝尔层 $\mathcal{O}_X^*$ 与 $f$ 的 Leray 谱序列，
见《上同调》引理 \ref{cohomology-lemma-Leray}。
考虑诱导映射
$$
H^1(X, \mathcal{O}_X^*) \longrightarrow H^0(Y, R^1f_*\mathcal{O}_X^*)
$$
《除子》引理
\ref{divisors-lemma-finite-trivialize-invertible-upstairs}
恰好说明此映射为零。因此 Leray 谱序列给出
$H^1(X, \mathcal{O}_X^*) = H^1(Y, f_*\mathcal{O}_X^*)$。
接着考虑映射
$$
f^\sharp : \mathcal{O}_Y^* \longrightarrow f_*\mathcal{O}_X^*
$$
由假设，此映射的核与余核都支集于闭子集
$T = Y \setminus V$（属于 $Y$）上。
因为按假设 $T$ 是离散拓扑空间，所以 $Y$ 上支集于
$T$ 的任意阿贝尔层的高阶上同调群都为零
（这由《上同调》引理
\ref{cohomology-lemma-cohomology-and-closed-immersions}、
《模》引理 \ref{modules-lemma-i-star-exact}，以及如下事实推出：
$H^i(T, \mathcal{F}) = 0$，对任意 $i > 0$ 与任意阿贝尔层
$\mathcal{F}$（在 $T$ 上）都成立）。
把上述映射分解为短正合列
$$
0 \to \Ker(f^\sharp) \to \mathcal{O}_Y^* \to \Im(f^\sharp) \to 0,\quad
0 \to \Im(f^\sharp) \to f_*\mathcal{O}_X^* \to \Coker(f^\sharp) \to 0
$$
先得到 $H^1(Y, \mathcal{O}_Y^*) \to H^1(Y, \Im(f^\sharp))$
是满射，再得到
$H^1(Y, \Im(f^\sharp)) \to H^1(Y, f_*\mathcal{O}_X^*)$ 是满射。
结合以上各点，得到
$H^1(Y, \mathcal{O}_Y^*) \to
H^1(X, \mathcal{O}_X^*)$ 是满射，如所需。
\end{proof}

\begin{lemma}
\label{varieties-lemma-glue-invertible-sheaves}
设 $X$ 是概形。设 $Z_1, \ldots, Z_n \subset X$ 为闭子概形。
设 $\mathcal{L}_i$ 为 $Z_i$ 上的可逆层。假设
\begin{enumerate}
\item $X$ 是约化的；
\item 在集合论意义下 $X = \bigcup Z_i$；
\item $Z_i \cap Z_j$ 是离散拓扑空间，其中 $i \not = j$。
\end{enumerate}
则存在可逆层 $\mathcal{L}$（在 $X$ 上），其在 $Z_i$ 上的限制是
$\mathcal{L}_i$。此外，若给定截面
$s_i \in \Gamma(Z_i, \mathcal{L}_i)$，它们在
$Z_i \cap Z_j$ 的各点处不消失，则可选择 $\mathcal{L}$，使得存在
$s \in \Gamma(X, \mathcal{L})$，并且有
$s|_{Z_i} = s_i$，对所有 $i$ 都成立。
\end{lemma}

\begin{proof}
$\mathcal{L}$ 的存在性可以由引理
\ref{varieties-lemma-surjective-pic-birational-finite} 推出；
但我们也将给出直接证明，并用该证明说明关于截面的陈述成立。
置 $T = \bigcup_{i \not = j} Z_i \cap Z_j$。由于 $X$ 约化，作为概形有
$$
X \setminus T = \bigcup (Z_i \setminus T)
$$
。假设 (3) 蕴含 $T$ 是 $X$ 的离散子集。
所以对每个 $t \in T$，可以找到开集 $U_t \subset X$，
使得 $t \in U_t$，但 $t' \not \in U_t$；这里
$t' \in T$ 且 $t' \not = t$。
必要时缩小 $U_t$，可假设存在同构
$\varphi_{t, i} : \mathcal{L}_i|_{U_t \cap Z_i} \to
\mathcal{O}_{U_t \cap Z_i}$。此外，对每个 $i$ 选择开覆盖
$$
Z_i \setminus T = \bigcup\nolimits_j U_{ij}
$$
使得存在同构
$\varphi_{i, j} : \mathcal{L}_i|_{U_{ij}} \cong \mathcal{O}_{U_{ij}}$。
注意
$$
\mathcal{U} : X = \bigcup U_t \cup \bigcup U_{ij}
$$
是 $X$ 的开覆盖。我们断言，可以用这些同构
$\varphi_{t, i}$ 和 $\varphi_{i, j}$ 为此覆盖定义一个
$2$-余循环，它取值于 $\mathcal{O}_X^*$，并定义引理陈述中的
$\mathcal{L}$。

\medskip\noindent
事实上，若 $i \not = i'$，则
$U_{i, j} \cap U_{i', j'} = \emptyset$，无事可做。
对 $U_{i, j} \cap U_{i, j'}$，由证明开头的说明，有
$\mathcal{O}_X(U_{i, j} \cap U_{i, j'}) =
\mathcal{O}_{Z_i}(U_{i, j} \cap U_{i, j'})$。
所以 $\mathcal{L}_i$ 的转移函数（更准确地说，
$\varphi_{i, j} \circ \varphi_{i, j'}^{-1}$）
定义了余循环在这一交集上的值。
对 $U_t \cap U_{i, j}$ 可以作同样处理。
最后，当 $t \not = t'$ 时，有
$$
U_t \cap U_{t'} = \coprod (U_t \cap U_{t'}) \cap Z_i
$$
，而且交集 $U_t \cap U_{t'} \cap Z_i$ 包含于
$Z_i \setminus T$。所以，由与前面相同的理由，有
$$
\mathcal{O}_X(U_t \cap U_{t'}) =
\prod \mathcal{O}_{Z_i}(U_t \cap U_{t'} \cap Z_i)
$$
，并且可以用 $\mathcal{L}_i$ 的转移函数（更准确地说，
$\varphi_{t, i} \circ \varphi_{t', i}^{-1}$）定义余循环在
$U_t \cap U_{t'}$ 上的值。这就完成了 $\mathcal{L}$ 的存在性证明。

\medskip\noindent
给定引理最后一个断言中的截面 $s_i$。在上述论证中，选择
$U_t$，使得 $s_i|_{U_t \cap Z_i}$ 不消失，并选择
$\varphi_{t, i}$，使得
$\varphi_{t, i}(s_i|_{U_t \cap Z_i}) = 1$。
于是，使用 $1$（在 $U_t$ 上）以及
$\varphi_{i, j}(s_i|_{U_{i, j}})$（在 $U_{i, j}$ 上），
将定义 $\mathcal{L}$ 的一个截面，
它限制为 $s_i$（在 $Z_i$ 上）。
\end{proof}

\begin{remark}
\label{varieties-remark-conductor}
设 $A$ 是约化环。设 $I, J$ 是 $A$ 的理想，并且
$V(I) \cup V(J) = \Spec(A)$。置 $B = A/J$。
则 $I \to IB$ 是 $A$-模的同构。事实上，有
$IB = I + J/J = I/(I \cap J)$；而 $I \cap J$ 为零，因为
$A$ 约化，并且
$\Spec(A) = V(I) \cup V(J) = V(I \cap J)$。
所以对任意投射 $A$-模 $P$，也有 $IP = I(P/JP)$。
\end{remark}

\begin{lemma}
\label{varieties-lemma-dim-1-noetherian-reduced-separated-has-ample}
设 $X$ 是维数为 $1$ 的 Noether 约化分离概形。
则 $X$ 具有丰沛可逆层。
\end{lemma}

\begin{proof}
令 $Z_i$（$i = 1, \ldots, n$）为 $X$ 的不可约分支。
把它们视为 $X$ 的约化闭子概形。
由引理 \ref{varieties-lemma-dim-1-noetherian-integral-separated-has-ample}，
存在丰沛可逆层 $\mathcal{L}_i$（在 $Z_i$ 上）。
置 $T = \bigcup_{i \not = j} Z_i \cap Z_j$。由于 $X$ 是维数
$1$ 的 Noether 概形，集合 $T$ 是有限的，并由 $X$ 的闭点组成。
对每个 $i$，在必要时把 $\mathcal{L}_i$ 替换为某个幂后，
可以选择 $s_i \in \Gamma(Z_i, \mathcal{L}_i)$，使得
$(Z_i)_{s_i}$ 是仿射的并包含 $T \cap Z_i$，见《性质》引理
\ref{properties-lemma-ample-finite-set-in-principal-affine}。

\medskip\noindent
由引理 \ref{varieties-lemma-glue-invertible-sheaves}，可以找到可逆层
$\mathcal{L}$（在 $X$ 上）以及 $s \in \Gamma(X, \mathcal{L})$，使得
$(\mathcal{L}, s)|_{Z_i} = (\mathcal{L}_i, s_i)$。
注意 $X_s$ 包含 $T$，而且在集合论意义下等于这些仿射闭子概形
$(Z_i)_{s_i}$ 的并。因此，由《极限》引理
\ref{limits-lemma-affines-glued-in-closed-affine}，它是仿射的。
为完成证明，只需对每个 $x \in X$（其中 $x \not \in T$）
找到整数 $m > 0$ 和截面
$t \in \Gamma(X, \mathcal{L}^{\otimes m})$，使得 $X_t$ 是仿射的且
$x \in X_t$。由于 $x \not \in T$，可知 $x \in Z_i$，其中
$i$ 唯一；设 $i = 1$。
令 $Z \subset X$ 为约化闭子概形，其底拓扑空间是
$Z_2 \cup \ldots \cup Z_n$。
令 $\mathcal{I} \subset \mathcal{O}_X$ 为 $Z$ 的理想层。
以 $\mathcal{I}_1 \subset \mathcal{O}_{Z_1}$ 表示该理想层在包含态射
$Z_1 \to X$ 下的逆像。注意
$$
\Gamma(X, \mathcal{I}\mathcal{L}^{\otimes m}) =
\Gamma(Z_1, \mathcal{I}_1 \mathcal{L}_1^{\otimes m})
$$
，见注 \ref{varieties-remark-conductor}。所以只需找到 $m > 0$ 以及
$t \in \Gamma(Z_1, \mathcal{I}_1 \mathcal{L}_1^{\otimes m})$，
使得 $x \in (Z_1)_t$ 且它是仿射的。
由于 $\mathcal{L}_1$ 丰沛，并且 $x$ 不在
$Z_1 \cap T = V(\mathcal{I}_1)$ 中，可以找到截面
$t_1 \in \Gamma(Z_1, \mathcal{I}_1 \mathcal{L}_1^{\otimes m_1})$，
使得 $x \in (Z_1)_{t_1}$，见《性质》命题
\ref{properties-proposition-characterize-ample}。
由于 $\mathcal{L}_1$ 丰沛，可以找到截面
$t_2 \in \Gamma(Z_1, \mathcal{L}_1^{\otimes m_2})$，
使得 $x \in (Z_1)_{t_2}$ 且 $(Z_1)_{t_2}$ 是仿射的，见《性质》定义
\ref{properties-definition-ample}。
置 $m = m_1 + m_2$ 且 $t = t_1 t_2$。则由构造，
$t \in \Gamma(Z_1, \mathcal{I}_1 \mathcal{L}_1^{\otimes m})$，
且 $x \in (Z_1)_t$；并且由《性质》引理
\ref{properties-lemma-affine-cap-s-open}，$(Z_1)_t$ 是仿射的。
\end{proof}

\begin{lemma}
\label{varieties-lemma-lift-line-bundle-from-reduction-dimension-1}
设 $i : Z \to X$ 是概形的闭浸入。
若 $X$ 的底拓扑空间是 Noether 的，并且
$\dim(X) \leq 1$，则 $\Pic(X) \to \Pic(Z)$ 是满射。
\end{lemma}

\begin{proof}
考虑 $X$ 上阿贝尔层的短正合列
$$
0 \to (1 + \mathcal{I}) \cap \mathcal{O}_X^* \to
\mathcal{O}^*_X \to i_*\mathcal{O}^*_Z \to 0
$$
，其中 $\mathcal{I}$ 是与 $Z$ 对应的拟凝聚理想层。
由于 $\dim(X) \leq 1$，有
$H^2(X, \mathcal{F}) = 0$；这对任意阿贝尔层 $\mathcal{F}$ 成立，见《上同调》命题
\ref{cohomology-proposition-vanishing-Noetherian}。
因此映射
$H^1(X, \mathcal{O}^*_X) \to H^1(X, i_*\mathcal{O}_Z^*)$
是满射。由《上同调》引理
\ref{cohomology-lemma-cohomology-and-closed-immersions}，有
$H^1(X, i_*\mathcal{O}_Z^*) = H^1(Z, \mathcal{O}_Z^*)$。
再由《上同调》引理
\ref{cohomology-lemma-h1-invertible}，这就证明了引理。
\end{proof}

\begin{proposition}
\label{varieties-proposition-dim-1-noetherian-separated-has-ample}
设 $X$ 是维数为 $1$ 的 Noether 分离概形。
则 $X$ 具有丰沛可逆层。
\end{proposition}

\begin{proof}
令 $Z \subset X$ 为 $X$ 的约化。由引理
\ref{varieties-lemma-dim-1-noetherian-reduced-separated-has-ample}，
概形 $Z$ 具有丰沛可逆层。
所以，由引理 \ref{varieties-lemma-lift-line-bundle-from-reduction-dimension-1}，
存在可逆 $\mathcal{O}_X$-模 $\mathcal{L}$（在 $X$ 上），
它在 $Z$ 上的限制是丰沛的。
于是，应用《概形上同调》引理
\ref{coherent-lemma-ample-on-reduction}，可知 $\mathcal{L}$ 是丰沛的。
\end{proof}

\begin{remark}
\label{varieties-remark-useless-generalization}
事实上，若 $X$ 是概形，其约化是维数为 $1$ 的 Noether 分离概形，
则 $X$ 具有丰沛可逆层。证明这一点的论证与命题
\ref{varieties-proposition-dim-1-noetherian-separated-has-ample} 的证明相同，
只是使用《极限》引理
\ref{limits-lemma-ample-on-reduction}
来代替《概形上同调》引理
\ref{coherent-lemma-ample-on-reduction}。
\end{remark}

\noindent
下面的引理实际上对拟有限分离态射也成立；读者可以使用 Zariski 主定理
（《关于态射的更多内容》引理
\ref{more-morphisms-lemma-quasi-finite-separated-pass-through-finite}）
以及引理 \ref{varieties-lemma-complement-codim-1-closed-points} 看出这一点。

\begin{lemma}
\label{varieties-lemma-surjection-on-pic-quasi-finite}
设 $f : X \to Y$ 是概形的态射。假设
$Y$ 是维数 $\leq 1$ 的 Noether 概形，$f$ 是有限的，并且
存在稠密开集 $V \subset Y$，使得
$f^{-1}(V) \to V$ 是闭浸入。则每个可逆
$\mathcal{O}_X$-模都是某个可逆 $\mathcal{O}_Y$-模的拉回。
\end{lemma}

\begin{proof}
把 $f$ 分解为 $X \to Z \to Y$，其中 $Z$ 是 $f$ 的概形论像。
则 $X \to Z$ 在 $V \cap Z$ 上是同构，并可应用引理
\ref{varieties-lemma-surjective-pic-birational-finite}。另一方面，
引理 \ref{varieties-lemma-lift-line-bundle-from-reduction-dimension-1}
适用于 $Z \to Y$。略去一些细节。
\end{proof}

\section{\texorpdfstring{\ensuremath{\delta}}{delta} 不变量}
\label{varieties-section-delta-invariant}

\noindent
本节定义奇异点的 $\delta$-不变量，其中该点位于约化
$1$ 维 Nagata 概形上。

\begin{lemma}
\label{varieties-lemma-pre-pre-delta-invariant}
设 $(A, \mathfrak m)$ 是 Noether $1$ 维局部环。
设 $f \in \mathfrak m$。下列条件等价：
\begin{enumerate}
\item $\mathfrak m = \sqrt{(f)}$；
\item $f$ 不包含在 $A$ 的任何极小素理想中；
\item 有 $A_f = \prod_{\mathfrak p\text{ minimal}} A_\mathfrak p$（作为 $A$-代数）。
\end{enumerate}
这样的 $f \in \mathfrak m$ 总是存在。若 $\text{depth}(A) = 1$（例如
$A$ 约化），则 (1)--(3) 还等价于：
\begin{enumerate}
\item[(4)] $f$ 是非零因子；
\item[(5)] $A_f$ 是 $A$ 的总分式环。
\end{enumerate}
若 $A$ 约化，则 (1)--(5) 还等价于：
\begin{enumerate}
\item[(6)] $A_f$ 是 $A$ 在各极小素理想处的剩余域之积。
\end{enumerate}
\end{lemma}

\begin{proof}
$A$ 的谱有有限多个素理想
$\mathfrak p_1, \ldots, \mathfrak p_n$ 位于 $\mathfrak m$ 之外，
并且它们全都极小；见
《代数》引理 \ref{algebra-lemma-Noetherian-irreducible-components}。
于是 (1) 与 (2) 的等价性由《代数》引理
\ref{algebra-lemma-Zariski-topology} 得出。
显然，(3) 蕴含 (2)。反之，若 (2) 成立，则
$A_f$ 的谱是子集 $\{\mathfrak p_1, \ldots, \mathfrak p_n\}$，
它位于 $\Spec(A)$ 中并赋予诱导拓扑；见
《代数》引理 \ref{algebra-lemma-spec-localization}。
这是有限离散拓扑空间。
故由《代数》命题 \ref{algebra-proposition-dimension-zero-ring}，
$A_f = \prod_{\mathfrak p\text{ minimal}} A_\mathfrak p$。
《代数》引理 \ref{algebra-lemma-height-1} 断言了这样的 $f$ 的存在性。

\medskip\noindent
假设 $A$ 的深度为 $1$。（由《代数》引理
\ref{algebra-lemma-bound-depth}，这是最大可能值；若 $A$ 约化，
则由《代数》引理 \ref{algebra-lemma-criterion-reduced} 可知它成立。）
于是 $\mathfrak m$ 不是 $A$ 的相伴素理想。
$A$ 的每个极小素理想都是相伴素理想
（《代数》命题
\ref{algebra-proposition-minimal-primes-associated-primes}）。
因此，由《代数》引理 \ref{algebra-lemma-ass-zero-divisors}，
$A$ 的非零因子集合恰为不包含在任何极小素理想中的元素集合。
所以 (4) 等价于 (2)。
由《代数》引理
\ref{algebra-lemma-total-ring-fractions-no-embedded-points}，
(5) 等价于 (3)。

\medskip\noindent
$A_\mathfrak p$ 是域，其中 $\mathfrak p \subset A$ 极小；见
《代数》引理 \ref{algebra-lemma-minimal-prime-reduced-ring}。
故 (3) 等价于 (6)。
\end{proof}

\begin{lemma}
\label{varieties-lemma-pre-delta-invariant}
设 $(A, \mathfrak m)$ 是约化 Nagata $1$ 维局部环。
设 $A'$ 是 $A$ 在 $A$ 的总分式环中的整闭包。则
$A'$ 是正规 Nagata 环，$A \to A'$ 有限，并且
$A'/A$ 作为 $A$-模具有有限长度。
\end{lemma}

\begin{proof}
总分式环在 $A$ 上本质有限型；因此 $A \to A'$ 有限，因为
$A$ 是 Nagata 环；见《代数》引理
\ref{algebra-lemma-nagata-in-reduced-finite-type-finite-integral-closure}。
例如，由《代数》引理
\ref{algebra-lemma-characterize-reduced-ring-normal} 和
\ref{algebra-lemma-Noetherian-irreducible-components} 可知环 $A'$ 正规。
例如，由《代数》引理 \ref{algebra-lemma-quasi-finite-over-nagata}
可知环 $A'$ 是 Nagata 环。
按引理 \ref{varieties-lemma-pre-pre-delta-invariant} 选取 $f \in \mathfrak m$。
由于 $A' \subset A_f$，显然 $A_f = A'_f$。因此有限
$A$-模 $A'/A$ 的支撑包含在 $\{\mathfrak m\}$ 中。
由《代数》引理 \ref{algebra-lemma-support-point}，它具有有限长度。
\end{proof}

\begin{definition}
\label{varieties-definition-delta-invariant-algebra}
设 $A$ 是维数为 $1$ 的约化 Nagata 局部环。
{\it $\delta$-不变量（关于 $A$）}定义为 $\text{length}_A(A'/A)$，
其中 $A'$ 如引理 \ref{varieties-lemma-pre-delta-invariant} 所述。
\end{definition}

\noindent
下面证明关于这个不变量之行为的一些引理。

\begin{lemma}
\label{varieties-lemma-delta-invariant-is-zero}
设 $A$ 是维数为 $1$ 的约化 Nagata 局部环。
$\delta$-不变量（关于 $A$）为 $0$，当且仅当
$A$ 是离散赋值环。
\end{lemma}

\begin{proof}
若 $A$ 是离散赋值环，则 $A$ 正规，并且环
$A'$ 等于 $A$。反之，若
$\delta$-不变量（关于 $A$）为 $0$，则 $A$ 在其总分式环中整闭，从而
$A$ 正规
（《代数》引理 \ref{algebra-lemma-characterize-reduced-ring-normal}）；
再由《代数》引理 \ref{algebra-lemma-characterize-dvr}，
这迫使 $A$ 成为离散赋值环。
\end{proof}

\begin{lemma}
\label{varieties-lemma-normalization-same-after-completion}
设 $A$ 是维数为 $1$ 的约化 Nagata 局部环。
令 $A \to A'$ 如引理 \ref{varieties-lemma-pre-delta-invariant} 所述。
令 $A^h$、$A^{sh}$ 以及相应的 $A^\wedge$
分别为 $A$ 的 Hensel 化、严格 Hensel 化以及完备化。
则 $A^h$、$A^{sh}$ 以及相应的 $A^\wedge$ 都是约化 Nagata
局部环，维数为 $1$，并且
$A' \otimes_A A^h$、$A' \otimes_A A^{sh}$ 以及相应的 $A' \otimes_A A^\wedge$
分别是 $A^h$、$A^{sh}$ 以及相应的 $A^\wedge$
在其总分式环中的整闭包。
\end{lemma}

\begin{proof}
注意 $A^\wedge$ 约化；见《代数续篇》引理
\ref{more-algebra-lemma-completion-reduced}。
由《代数续篇》引理 \ref{more-algebra-lemma-henselization-reduced}，
环 $A^h$ 与 $A^{sh}$ 约化。
由《代数续篇》引理
\ref{more-algebra-lemma-completion-dimension} 和
\ref{more-algebra-lemma-henselization-dimension}，
$A$、$A^h$、$A^{sh}$ 与 $A^\wedge$ 的维数相同。

\medskip\noindent
回顾：Noether 局部环是 Nagata 环，当且仅当
$A$ 的形式纤维几何约化；见《代数续篇》引理
\ref{more-algebra-lemma-Nagata-local-ring}。
这个性质由 $A^h$ 与 $A^{sh}$ 继承；见《代数续篇》第
\ref{more-algebra-section-properties-formal-fibres} 节中的材料，
尤其是引理 \ref{more-algebra-lemma-henselization-P-ring}。
由《代数》引理 \ref{algebra-lemma-Noetherian-complete-local-Nagata}，
完备化是 Nagata 环。

\medskip\noindent
现在证明关于整闭包的陈述。继续之前，按引理
\ref{varieties-lemma-pre-pre-delta-invariant} 选取 $f \in \mathfrak m$。
$f$ 在 $A^h$、$A^{sh}$ 与 $A^\wedge$ 中的像
显然都在相应环中满足性质 (1)--(6)。

\medskip\noindent
由于 $A \to A'$ 有限，可知
$A' \otimes_A A^h$ 与 $A' \otimes_A A^{sh}$
分别是若干 Hensel 局部环之积，而这些局部环分别在 $A^h$ 与
$A^{sh}$ 上有限；见《代数》引理
\ref{algebra-lemma-finite-over-henselian}。
其中每个局部环都是 $A'$ 在某个极大理想
$\mathfrak m' \subset A'$ 处的 Hensel 化，该理想位于 $\mathfrak m$
上方；见《代数》引理
\ref{algebra-lemma-quasi-finite-henselization} 或
\ref{algebra-lemma-quasi-finite-strict-henselization}。
故由《代数续篇》引理
\ref{more-algebra-lemma-henselization-normal}，这些局部环都是正规整环。
于是 $A' \otimes_A A^h$ 与 $A' \otimes_A A^{sh}$
是正规环。由于 $A^h \to A' \otimes_A A^h$ 与
$A^{sh} \to A' \otimes_A A^{sh}$ 有限（因而整），并且
$A' \otimes_A A^h \subset (A^h)_f = Q(A^h)$ 与
$A' \otimes_A A^{sh} \subset (A^{sh})_f = Q(A^{sh})$，
我们断定 $A' \otimes_A A^h$ 与 $A' \otimes_A A^{sh}$
就是所需的整闭包。

\medskip\noindent
对于完备化，论证完全相同。首先，由《代数》引理
\ref{algebra-lemma-completion-finite-extension} 有
$$
A' \otimes_A A^\wedge = (A')^\wedge =
\prod\nolimits (A'_{\mathfrak m'})^\wedge
$$
局部环 $A'_{\mathfrak m'}$ 正规且维数为 $1$
（例如由《代数》引理 \ref{algebra-lemma-finite-in-codim-1}，或由
《代数》第 \ref{algebra-section-homomorphism-dimension} 节中的讨论）。
故由《代数》引理 \ref{algebra-lemma-characterize-dvr}，
$A'_{\mathfrak m'}$ 是离散赋值环。
再由《代数续篇》引理 \ref{more-algebra-lemma-completion-dvr}，
$(A'_{\mathfrak m'})^\wedge$ 是离散赋值环。
于是 $A' \otimes_A A^\wedge$ 是正规环，并且
可以完全按前面的方式得出结论。
\end{proof}

\begin{lemma}
\label{varieties-lemma-delta-same-after-completion}
设 $A$ 是维数为 $1$ 的约化 Nagata 局部环。
$\delta$-不变量（关于 $A$）与 $\delta$-不变量（关于
$A$ 的 Hensel 化、严格 Hensel 化或完备化）相同。
\end{lemma}

\begin{proof}
我们证明完备化 $B = A^\wedge$ 的情形；其他情形的证明完全相同。
令 $A'$ 以及相应的 $B'$ 分别为 $A$ 以及相应的 $B$
在其总分式环中的整闭包。
由引理 \ref{varieties-lemma-normalization-same-after-completion}，
$B' = A' \otimes_A B$。
故 $B'/B = (A'/A) \otimes_A B$。
现在，等式由《代数》引理 \ref{algebra-lemma-pullback-module}
以及 $B \otimes_A \kappa_A = \kappa_B$ 这一事实得出。
\end{proof}

\begin{definition}
\label{varieties-definition-delta-invariant}
设 $k$ 是域。设 $X$ 是局部代数 $k$-概形。
设 $x \in X$ 是一点，使得 $\mathcal{O}_{X, x}$
约化且 $\dim(\mathcal{O}_{X, x}) = 1$。
{\it $\delta$-不变量（$X$ 在 $x$ 处）}定义为
$\delta$-不变量（关于 $\mathcal{O}_{X, x}$），如定义
\ref{varieties-definition-delta-invariant-algebra} 所述。
\end{definition}

\noindent
这个定义有意义，因为局部代数概形的局部环由
《代数》命题 \ref{algebra-proposition-ubiquity-nagata} 是 Nagata 环。
当然，更一般地，只要概形的一点 $x \in X$ 的局部环
$\mathcal{O}_{X, x}$ 是约化 Nagata 环且维数为 $1$，我们都可以作此定义。
由引理 \ref{varieties-lemma-delta-same-after-completion}，
$\delta$-不变量（$X$ 在 $x$ 处）为
$$
\delta\text{-invariant of }X\text{ 在 }x =
\delta\text{-invariant of }\mathcal{O}_{X, x}^h =
\delta\text{-invariant of }\mathcal{O}_{X, x}^\wedge
$$
由此可见，$\delta$-不变量是该点之完备局部环的不变量。

\begin{lemma}
\label{varieties-lemma-delta-invariant-and-change-of-fields}
设 $k$ 是域。设 $X$ 是局部代数 $k$-概形。
设 $K/k$ 是域扩张，并令 $Y = X_K$。
设 $y \in Y$ 的像为 $x \in X$。
假设 $X$ 在 $x$ 处几何约化，并且
$\dim(\mathcal{O}_{X, x}) = \dim(\mathcal{O}_{Y, y}) = 1$。
则
$$
\delta\text{-invariant of }X\text{ 在 }x \leq
\delta\text{-invariant of }Y\text{ 在 }y
$$
\end{lemma}

\begin{proof}
置 $A = \mathcal{O}_{X, x}$ 与 $B = \mathcal{O}_{Y, y}$。
由引理 \ref{varieties-lemma-geometrically-reduced-at-point} 可知
$A$ 几何约化。
因此，$B$ 是 $A \otimes_k K$ 的一个局部化。
令 $A \to A'$ 如引理 \ref{varieties-lemma-pre-delta-invariant} 所述。则
$$
B' = B \otimes_{(A \otimes_k K)} (A' \otimes_k K)
$$
在 $B$ 上有限，并且 $B \to B'$ 在总分式环上诱导同构。
事实上，选取满足引理 \ref{varieties-lemma-pre-pre-delta-invariant}
中 (1)--(6) 的 $f \in \mathfrak m_A$；由于
$\dim(B) = 1$，可知 $f \in \mathfrak m_B$
对 $B$ 起同样作用；又有 $B_f = B'_f$，因为 $A_f = A'_f$。
令 $B''$ 是 $B$ 在其总分式环中的整闭包，如引理
\ref{varieties-lemma-pre-delta-invariant} 所述。
则 $B' \subset B''$。因此，$\delta$-不变量（$Y$ 在 $y$ 处）为
$\text{length}_B(B''/B)$，并且
\begin{align*}
\text{length}_B(B''/B)
& \geq
\text{length}_B(B'/B) \\
& =
\text{length}_B((A'/A) \otimes_A B) \\
& =
\text{length}_B(B/\mathfrak m_A B) \text{length}_A(A'/A)
\end{align*}
这是因为 $A \to B$ 平坦（它是 $A \to A \otimes_k K$ 的一个局部化），
并可应用《代数》引理 \ref{algebra-lemma-pullback-module}。
由于 $\text{length}_A(A'/A)$ 是 $\delta$-不变量（$X$ 在 $x$ 处），
又由于 $\text{length}_B(B/\mathfrak m_A B) \geq 1$，引理得证。
\end{proof}

\begin{lemma}
\label{varieties-lemma-delta-invariant-and-change-of-fields-better}
设 $k$ 是域。设 $X$ 是局部代数 $k$-概形。
设 $K/k$ 是域扩张，并令 $Y = X_K$。
设 $y \in Y$ 的像为 $x \in X$。
假设引理 \ref{varieties-lemma-geometrically-normal-in-codim-1}
中的假设 (a)、(b)、(c) 对 $x \in X$ 成立，并且
$\dim(\mathcal{O}_{Y, y}) = 1$。
则 $\delta$-不变量（$X$ 在 $x$ 处）
等于 $\delta$-不变量（$Y$ 在 $y$ 处）。
\end{lemma}

\begin{proof}
置 $A = \mathcal{O}_{X, x}$ 与 $B = \mathcal{O}_{Y, y}$。
由引理 \ref{varieties-lemma-geometrically-normal-in-codim-1}
可知 $A$ 几何约化。
因此，$B$ 是 $A \otimes_k K$ 的一个局部化。
令 $A \to A'$ 如引理 \ref{varieties-lemma-pre-delta-invariant} 所述。
由引理 \ref{varieties-lemma-geometrically-normal-in-codim-1}
可知 $A' \otimes_k K$ 正规。
因此
$$
B' = B \otimes_{(A \otimes_k K)} (A' \otimes_k K)
$$
正规、在 $B$ 上有限，并且 $B \to B'$ 在总分式环上诱导同构。
事实上，选取满足引理 \ref{varieties-lemma-pre-pre-delta-invariant}
中 (1)--(6) 的 $f \in \mathfrak m_A$；由于
$\dim(B) = 1$，可知 $f \in \mathfrak m_B$
对 $B$ 起同样作用；又有 $B_f = B'_f$，因为 $A_f = A'_f$。
由此，$B \to B'$ 对 $B$ 而言如引理
\ref{varieties-lemma-pre-delta-invariant} 所述。因此只需证明
$\text{length}_A(A'/A) =
\text{length}_B(B'/B) = \text{length}_B((A'/A) \otimes_A B)$。
由于 $A \to B$ 平坦（它是 $A \to A \otimes_k K$ 的一个局部化），
并且 $\mathfrak m_B = \mathfrak m_A B$（因为
$B/\mathfrak m_A B$ 由上述说明是零维的，又是
$K \otimes_k \kappa(x)$ 的一个局部化；后者约化，因为
$\kappa(x)$ 在 $k$ 上可分），结论由《代数》引理
\ref{algebra-lemma-pullback-module} 得出。
\end{proof}





\section{分枝数}
\label{varieties-section-number-of-branches}

\noindent
我们已在《性质》第 \ref{properties-section-unibranch} 节中
定义了概形在一点处的分枝数。

\begin{lemma}
\label{varieties-lemma-number-of-branches}
设 $X$ 是概形。假设 $X$ 的每个拟紧开集都只有有限多个不可约分支。
设 $\nu : X^\nu \to X$ 是 $X$ 的正规化。设 $x \in X$。
\begin{enumerate}
\item $X$ 在 $x$ 处的分枝数等于 $x$ 在 $X^\nu$ 中的逆像数。
\item $X$ 在 $x$ 处的几何分枝数为
$\sum_{\nu(x^\nu) = x} [\kappa(x^\nu) : \kappa(x)]_s$。
\end{enumerate}
\end{lemma}

\begin{proof}
首先注意，对 $X$ 的假设恰好意味着正规化有定义；见《态射》定义
\ref{morphisms-definition-normalization}。
于是茎 $A' = (\nu_*\mathcal{O}_{X^\nu})_x$ 是
$A = \mathcal{O}_{X, x}$ 在 $A_{red}$ 的总分式环中的整闭包；见
《态射》引理 \ref{morphisms-lemma-stalk-normalization}。
由于 $\nu$ 是整态射，可知 $X^\nu$ 中位于 $x$ 上方的点
对应于 $A'$ 中位于极大理想 $\mathfrak m$ 上方的素理想，
而后者是 $A$ 的极大理想。
由于 $A \to A'$ 是整的，这也就是 $A'$ 的极大理想
（《代数》引理 \ref{algebra-lemma-integral-no-inclusion} 与
\ref{algebra-lemma-integral-going-up}）。
因此，本引理由其代数版本得出：
《代数续篇》引理 \ref{more-algebra-lemma-number-of-branches-1}。
\end{proof}

\begin{lemma}
\label{varieties-lemma-geometric-branches-and-change-of-fields}
设 $k$ 是域。设 $X$ 是局部代数 $k$-概形。
设 $K/k$ 是域扩张。设 $y \in X_K$ 是一点，其像 $x$ 位于
$X$ 中。则 $X$ 在 $x$ 处的几何分枝数等于
$X_K$ 在 $y$ 处的几何分枝数。
\end{lemma}

\begin{proof}
记 $Y = X_K$，并令 $X^\nu$ 以及相应的 $Y^\nu$ 为
$X$ 以及相应的 $Y$ 的正规化。考虑交换图
$$
\xymatrix{
Y^\nu \ar[r] \ar[d] & X^\nu_K \ar[r] \ar[d]_{\nu_K} & X^\nu \ar[d]_\nu \\
Y \ar@{=}[r] & Y \ar[r] & X
}
$$
由引理 \ref{varieties-lemma-normalization-and-change-of-fields}，左上方的水平箭头
是泛同胚。因此它诱导纯不可分剩余域扩张；见《态射》引理
\ref{morphisms-lemma-universal-homeomorphism} 与
\ref{morphisms-lemma-universally-injective}。
故由引理 \ref{varieties-lemma-number-of-branches}，$Y$ 在 $y$ 处的几何分枝数为
$\sum_{\nu_K(y') = y} [\kappa(y') : \kappa(y)]_s$。
类似地，$\sum_{\nu(x') = x} [\kappa(x') : \kappa(x)]_s$
是 $X$ 在 $x$ 处的几何分枝数。
利用《概形》引理 \ref{schemes-lemma-points-fibre-product}，
我们的陈述由下列代数事实得出：给定域扩张 $l/\kappa$
与代数域扩张 $m/\kappa$，则
$$
\sum\nolimits_{m \otimes_\kappa l \to m'} [m' : l']_s = [m : \kappa]_s
$$
其中求和遍历 $m \otimes_\kappa l$ 的商域。
可以用初等方法证明这一点，也可以把引理
\ref{varieties-lemma-separably-closed-field-connected-components} 应用于
\begingroup
\small
$$
\Spec(m \otimes_\kappa l) \times_{\Spec(l)} \Spec(\overline{l}) =
\Spec(m) \otimes_{\Spec(\kappa)} \Spec(\overline{l})
\longrightarrow
\Spec(m) \times_{\Spec(\kappa)} \Spec(\overline{\kappa})
$$
\endgroup
因为可以把 $[m : \kappa]_s$ 解释为右端的连通分支数，而把和
$\sum_{m \otimes_\kappa l \to m'} [m' : l']_s$
解释为左端的连通分支数。
\end{proof}

\begin{lemma}
\label{varieties-lemma-geometrically-unibranch-and-change-of-fields}
设 $k$ 是域。设 $X$ 是局部代数 $k$-概形。
设 $K/k$ 是域扩张。设 $y \in X_K$ 是一点，其像 $x$ 位于
$X$ 中。则 $X$ 在 $x$ 处几何单枝，当且仅当
$X_K$ 在 $y$ 处几何单枝。
\end{lemma}

\begin{proof}
这立即由引理 \ref{varieties-lemma-geometric-branches-and-change-of-fields}
和《代数续篇》引理
\ref{more-algebra-lemma-number-of-branches-1} 得出。
\end{proof}

\begin{definition}
\label{varieties-definition-wedge}
设 $A$ 以及 $A_i$（$1 \leq i \leq n$）都是局部环。称
{\it $A$ 是 $A_1, \ldots, A_n$ 的楔合}，若存在同构
$$
\kappa_{A_1} \to \kappa_{A_2} \to \ldots \to \kappa_{A_n}
$$
且 $A$ 同构于由那些 $n$-元组
$(a_1, \ldots, a_n) \in A_1 \times \ldots \times A_n$ 构成的环，
其中这些元组映到 $\kappa_{A_n}$ 的同一个元素。
\end{definition}

\noindent
若给定基环 $\Lambda$，并且 $A$ 与 $A_i$ 都是 $\Lambda$-代数，
则还要求 $\kappa_{A_i} \to \kappa_{A_{i + 1}}$
是 $\Lambda$-代数同构，并要求 $A$ 作为 $\Lambda$-代数同构于如下
$\Lambda$-代数：它由那些 $n$-元组
$(a_1, \ldots, a_n) \in A_1 \times \ldots \times A_n$ 构成，
且这些元组映到 $\kappa_{A_n}$ 的同一个元素。
特别地，若 $\Lambda = k$ 是域，且映射
$k \to \kappa_{A_i}$ 都是同构，则诸同构
$\kappa_{A_i} \to \kappa_{A_{i + 1}}$ 有唯一选择，
并且我们常说{\it $A_1, \ldots, A_n$ 的楔合}。

\begin{lemma}
\label{varieties-lemma-delta-number-branches-inequality-sh}
设 $(A, \mathfrak m)$ 是严格 Hensel 的
$1$ 维约化 Nagata 局部环。则
$$
\delta\text{-invariant of }A \geq \text{number of geometric branches of }A - 1
$$
若等号成立，则 $A$ 是 $n \geq 1$ 个严格 Hensel 离散赋值环的楔合。
\end{lemma}

\begin{proof}
几何分枝数等于 $A$ 的分枝数（这由《代数续篇》定义
\ref{more-algebra-definition-number-of-branches} 立即得出）。
令 $A \to A'$ 如引理 \ref{varieties-lemma-pre-delta-invariant} 所述。
注意，$A$ 的分枝数等于 $A'$ 的极大理想数；见
《代数续篇》引理 \ref{more-algebra-lemma-number-of-branches-1}。
有满射
$$
A'/A \longrightarrow
\left(\prod\nolimits_{\mathfrak m'} \kappa(\mathfrak m')\right)/
\kappa(\mathfrak m)
$$
由于 $\dim_{\kappa(\mathfrak m)} \prod \kappa(\mathfrak m')$
是 $\geq$ 分枝数，不等式显然成立。

\medskip\noindent
若等号成立，则 $\kappa(\mathfrak m') = \kappa(\mathfrak m)$
对所有 $\mathfrak m' \subset A'$ 都成立，并且上面的显示箭头是同构。
由于 $A$ 是 Hensel 环且 $A \to A'$ 有限，可知 $A'$ 是若干
Hensel 局部环之积；见《代数》引理
\ref{algebra-lemma-finite-over-henselian}。
这些因子就是局部环 $A'_{\mathfrak m'}$；由于
$A'$ 正规，这些因子是离散赋值环
（《代数》引理 \ref{algebra-lemma-characterize-dvr}）。
由于显示箭头是同构，可知 $A$ 确实是这些局部环的楔合。
\end{proof}

\begin{lemma}
\label{varieties-lemma-delta-number-branches-inequality}
设 $(A, \mathfrak m)$ 是 $1$ 维约化 Nagata 局部环。则
$$
\delta\text{-invariant of }A \geq \text{number of geometric branches of }A - 1
$$
\end{lemma}

\begin{proof}
可以把 $A$ 替换为 $A$ 的严格 Hensel 化，而不改变
$\delta$-不变量（引理 \ref{varieties-lemma-delta-same-after-completion}），
也不改变 $A$ 的几何分枝数（这由定义立即得出；见
《代数续篇》定义 \ref{more-algebra-definition-number-of-branches}）。
因此可以假设 $A$ 严格 Hensel，并应用引理
\ref{varieties-lemma-delta-number-branches-inequality-sh}。
\end{proof}






\section{一维概形的正规化}
\label{varieties-section-normalization-one-dimensional}

\noindent
由 Krull--Akizuki 结果，维数为 $1$ 的 Noether 概形的正规化态射
具有意想不到的良好性质。

\begin{lemma}
\label{varieties-lemma-normalize-noetherian-dim-1}
设 $X$ 是维数为 $1$ 的局部 Noether 概形。
设 $\nu : X^\nu \to X$ 是正规化。则
\begin{enumerate}
\item $\nu$ 是整的、满的，并在不可约分支上诱导双射；
\item 存在分解 $X^\nu \to X_{red} \to X$，且态射
$X^\nu \to X_{red}$ 是 $X_{red}$ 的正规化；
\item $X^\nu \to X_{red}$ 是双有理的；
\item 对每个闭点 $x \in X$，茎
$(\nu_*\mathcal{O}_{X^\nu})_x$ 是 $\mathcal{O}_{X, x}$
在 $(\mathcal{O}_{X, x})_{red} = \mathcal{O}_{X_{red}, x}$
的总分式环中的整闭包；
\item $\nu$ 的纤维有限，且剩余域扩张有限；
\item $X^\nu$ 是若干整正规局部 Noether 概形的不交并，
且每个仿射开集都是有限个 Dedekind 整环之积的谱。
\end{enumerate}
\end{lemma}

\begin{proof}
许多结论实际上都是正规化态射的一般性质；见《态射》引理
\ref{morphisms-lemma-normalization-reduced}、
\ref{morphisms-lemma-stalk-normalization}、
\ref{morphisms-lemma-normalization-normal} 和
\ref{morphisms-lemma-normalization-birational}。
并不明显的是：纤维有限，所诱导的剩余域扩张有限，以及
$X^\nu$ 局部看起来是 Dedekind 整环的谱（因而是 Noether 的）。
为证明这一点，可以假设 $X = \Spec(A)$ 仿射、Noether、维数为
$1$，并且 $A$ 约化。
然后可以使用《态射》引理
\ref{morphisms-lemma-description-normalization} 中的描述，
把问题约化到 $A$ 是维数为 $1$ 的 Noether 整环的情形。
此时所需性质由 Krull--Akizuki 定理得出，其形式见《代数》引理
\ref{algebra-lemma-integral-closure-Dedekind}。
\end{proof}

\noindent
当然，当 $X$ 不约化时，下列引理也有一个变体。

\begin{lemma}
\label{varieties-lemma-prepare-delta-invariant}
设 $X$ 是维数为 $1$ 的约化 Nagata 概形。设 $\nu : X^\nu \to X$
是正规化。设 $x \in X$ 表示一个闭点。则
\begin{enumerate}
\item $\nu : X^\nu \to X$ 有限、满且双有理；
\item $\mathcal{O}_X \subset \nu_*\mathcal{O}_{X^\nu}$，并且
$\nu_*\mathcal{O}_{X^\nu}/\mathcal{O}_X$ 是摩天层的直和，
其值为 $\mathcal{Q}_x$，位于 $x$ 处；后者非零当且仅当
$x$ 是 $X$ 的奇异点；
\item $A' = (\nu_*\mathcal{O}_{X^\nu})_x$ 是
$A = \mathcal{O}_{X, x}$ 在其总分式环中的整闭包；
\item $\mathcal{Q}_x = A'/A$ 具有有限长度，且该长度等于
$\delta$-不变量（$X$ 在 $x$ 处）；
\item $A'$ 是半局部环，它是有限个 Dedekind 整环之积；
\item $A^\wedge$ 是维数为 $1$ 的约化 Noether 完备局部环；
\item $(A')^\wedge$ 是 $A^\wedge$ 在其总分式环中的整闭包；
\item $(A')^\wedge$ 是有限个完备离散赋值环之积；
\item $A'/A \cong (A')^\wedge/A^\wedge$。
\end{enumerate}
\end{lemma}

\begin{proof}
我们可以并且将要使用引理 \ref{varieties-lemma-normalize-noetherian-dim-1}
的全部结论。
$\nu$ 的有限性由《态射》引理
\ref{morphisms-lemma-nagata-normalization} 得出。
由于 $X$ 约化、Nagata 且维数为 $1$，由《代数续篇》命题
\ref{more-algebra-proposition-ubiquity-J-2} 可知正则轨迹是
稠密开集 $U \subset X$。
由于正则概形正规，这说明 $\nu$ 在 $U$ 上是同构。
由于 $\dim(X) \leq 1$，这意味着 $\nu$ 在一离散闭点集
$x \in X$ 上不是同构。
特别地，有短正合列
$$
0 \to \mathcal{O}_X \to \nu_*\mathcal{O}_{X^\nu} \to
\bigoplus\nolimits_{x \in X \setminus U} \mathcal{Q}_x \to 0
$$
由于引理 \ref{varieties-lemma-normalize-noetherian-dim-1} 已描述
$\nu_*\mathcal{O}_{X^\nu}$ 的茎，我们断定 $Q_x = A'/A$
的长度确实等于 $\delta$-不变量（$X$ 在 $x$ 处）。
例如由引理 \ref{varieties-lemma-delta-invariant-is-zero} 可知，
$Q_x \not = 0$ 恰好当 $x$ 是奇异点。
$A'$ 是若干半局部 Dedekind 整环之积这一描述也由引理
\ref{varieties-lemma-normalize-noetherian-dim-1} 得出。
$A$、$A'$ 与 $(A')^\wedge$ 之间的关系，我们已在引理
\ref{varieties-lemma-normalization-same-after-completion}（及其证明）中见过。
\end{proof}




\section{寻找仿射开集}
\label{varieties-section-finding-affine-opens}

\noindent
我们继续《性质》第 \ref{properties-section-finding-affine-opens} 节中
开始的讨论。事实证明，在分离概形上，可以找到包含给定有限个余维
$1$ 的点的仿射开集。见命题
\ref{varieties-proposition-finite-set-of-points-of-codim-1-in-affine}。

\medskip\noindent
我们将在《下降》引理
\ref{descent-lemma-characterize-open-immersion} 中改进下列引理。

\begin{lemma}
\label{varieties-lemma-characterize-open-immersion}
设 $f : X \to Y$ 是概形态射。以 $X^0$ 表示
$X$ 的不可约分支的泛点集合。若
\begin{enumerate}
\item $f$ 分离；
\item 存在开覆盖 $X = \bigcup U_i$，使得
$f|_{U_i} : U_i \to Y$ 是开浸入；
\item 若 $\xi, \xi' \in X^0$ 且 $\xi \not = \xi'$，则
$f(\xi) \not = f(\xi')$，
\end{enumerate}
则 $f$ 是开浸入。
\end{lemma}

\begin{proof}
假设 $y = f(x) = f(x')$。选取特化 $y_0 \leadsto y$，
其中 $y_0$ 是 $Y$ 的某个不可约分支的泛点。
由于 $f$ 在源上局部为同构，可以选取特化
$x_0 \leadsto x$ 与 $x'_0 \leadsto x'$，它们映到
$y_0 \leadsto y$。
注意 $x_0, x'_0 \in X^0$。故由假设 (3)，$x_0 = x'_0$。
由于 $f$ 分离，我们断定 $x = x'$。因此 $f$ 是开浸入。
\end{proof}

\begin{lemma}
\label{varieties-lemma-local-isomorphism}
设 $X \to S$ 是概形态射。设 $x \in X$ 是一点，其像为
$s \in S$。若
\begin{enumerate}
\item $\mathcal{O}_{X, x} = \mathcal{O}_{S, s}$；
\item $S$ 约化；
\item $X \to S$ 是有限型的；
\item $S$ 有有限多个不可约分支，
\end{enumerate}
则存在开邻域 $U$（关于 $x$），使得 $f|_U$ 是开浸入。
\end{lemma}

\begin{proof}
可以移除 $S$ 中不包含 $s$ 的（有限多个）不可约分支。
可以用仿射开邻域替换 $S$，该邻域包含 $s$。
可以用仿射开邻域替换 $X$，该邻域包含 $x$。
记 $S = \Spec(A)$、$X = \Spec(B)$。令
$\mathfrak q \subset B$ 以及相应的 $\mathfrak p \subset A$
分别为对应于 $x$ 以及相应的 $s$ 的素理想。
由于 $A$ 约化，且 $A$ 的所有极小素理想都包含在
$\mathfrak p$ 中，可知 $A \subset A_\mathfrak p$。
由于 $X \to S$ 是有限型的，$B$ 在 $A$ 上有限型。
设 $b_1, \ldots, b_n \in B$ 是生成 $B$（在 $A$ 上）的元素。
由于 $A_\mathfrak p = B_\mathfrak q$，可以找到
$f \in A$、$f \not \in \mathfrak p$ 以及 $a_i \in A$，
使得 $b_i$ 与 $a_i/f$ 在 $B_\mathfrak q$ 中有相同的像。
因此可以找到 $g \in B$、$g \not \in \mathfrak q$，使得
$g(fb_i - a_i) = 0$ 于 $B$ 中。由此，
$A_f \to B_{fg}$ 的像包含 $b_1, \ldots, b_n$ 的像，
特别也包含 $g$ 的像。
选取 $n \geq 0$ 与 $f' \in A$，使得 $f'/f^n$
映到 $g$ 在 $B_{fg}$ 中的像。由于
$A_\mathfrak p = B_\mathfrak q$，可知 $f' \not \in \mathfrak p$。
我们断定 $A_{ff'} \to B_{fg}$ 满。
最后，由于 $A_{ff'} \subset A_\mathfrak p = B_\mathfrak q$
（见上文），映射 $A_{ff'} \to B_{fg}$ 单，因而是同构。
\end{proof}

\begin{lemma}
\label{varieties-lemma-points-in-affine}
设 $f : T \to X$ 是概形态射。以 $X^0$ 以及相应的 $T^0$
表示不可约分支的泛点集合。
设 $t_1, \ldots, t_m \in T$ 是有限点集，其像为
$x_j = f(t_j)$。若
\begin{enumerate}
\item $T$ 仿射；
\item $X$ 拟分离；
\item $X^0$ 有限；
\item $f(T^0) \subset X^0$，且 $f : T^0 \to X^0$ 单；
\item $\mathcal{O}_{X, x_j} = \mathcal{O}_{T, t_j}$，
\end{enumerate}
则存在 $X$ 的仿射开集，包含 $x_1, \ldots, x_r$。
\end{lemma}

\begin{proof}
使用《极限》命题 \ref{limits-proposition-affine}，
可立即约化到 $X$ 与 $T$ 都约化的情形。略去细节。

\medskip\noindent
假设 $X$ 与 $T$ 约化。可以把 $T = \lim_{i \in I} T_i$
写成 $X$ 上有限表示概形的有向极限，转移态射为仿射态射；见
《极限》引理 \ref{limits-lemma-relative-approximation}。
选取 $i \in I$ 使得 $T_i$ 仿射；见《极限》引理
\ref{limits-lemma-limit-affine}。
记 $T_i = \Spec(R_i)$、$T = \Spec(R)$。
令 $R' \subset R$ 是 $R_i \to R$ 的像。
则 $T' = \Spec(R')$ 仿射、约化、在 $X$ 上有限型，且
$T \to T'$ 支配。对 $j = 1, \ldots, r$，令
$t'_j \in T'$ 为 $t_j$ 的像。考虑局部环映射
$$
\mathcal{O}_{X, x_j} \to
\mathcal{O}_{T', t'_j} \to
\mathcal{O}_{T, t_j}
$$
以 $(T')^0$ 表示 $T'$ 的不可约分支的泛点集合。
设 $\xi \leadsto t'_j$ 是一个特化，其中 $\xi \in (T')^0$。
由于 $T \to T'$ 支配，可以选取 $\eta \in T^0$ 映到
$\xi$（警告：事先并不知道 $\eta$ 特化到 $t_j$）。
把假设 (3) 应用于 $\eta$，可知像 $\theta$（即
$\xi$ 在 $X$ 中的像）对应于 $\mathcal{O}_{X, x_j}$
的一个极小素理想。
通过 (5) 中的同构提升 $\xi$，得到特化
$\eta' \leadsto t_j$，其中 $\eta' \in T^0$ 映到
$\theta \leadsto x_j$。
(4) 的单射性表明 $\eta = \eta'$。因此
$\mathcal{O}_{T', t'_j}$ 的每个极小素理想都位于
$\mathcal{O}_{T, t_j}$ 的一个极小素理想之下。我们断定
$\mathcal{O}_{T', t'_j} \to \mathcal{O}_{T, t_j}$ 单，
因而上面两个映射都是同构。

\medskip\noindent
由引理 \ref{varieties-lemma-local-isomorphism}，存在开集
$U \subset T'$，它包含所有点 $t'_j$，使得
$U \to X$ 是引理 \ref{varieties-lemma-characterize-open-immersion}
意义下的局部同构。
由该引理可知 $U \to X$ 是开浸入。
最后，由《性质》引理
\ref{properties-lemma-ample-finite-set-in-affine}，
可以找到开集 $W \subset U \subset T'$，它包含所有 $t'_j$。
$W$ 在 $X$ 中的像就是所需的仿射开集。
\end{proof}

\begin{lemma}
\label{varieties-lemma-finite-set-codim-1-points-in-affine}
设 $X$ 是整分离概形。设 $x_1, \ldots, x_r \in X$
是有限点集，使得 $\mathcal{O}_{X, x_i}$ 是维数
$\leq 1$ 的 Noether 环。则存在 $X$ 的仿射开子概形，
它包含所有 $x_1, \ldots, x_r$。
\end{lemma}

\begin{proof}
设 $K$ 是 $X$ 的有理函数域。
置 $A_i = \mathcal{O}_{X, x_i}$。则 $A_i \subset K$，且 $K$
是 $A_i$ 的分式域。由于 $X$ 分离，并且
$x_i \not = x_j$，不可能存在赋值环 $\mathcal{O} \subset K$
同时支配 $A_i$ 与 $A_j$。事实上，考虑图
$$
\xymatrix{
\Spec(\mathcal{O}) \ar[r] \ar[d] & \Spec(A_1) \ar[d] \\
\Spec(A_2) \ar[r] & X
}
$$
并应用分离性的赋值判据
（《概形》引理 \ref{schemes-lemma-separated-implies-valuative}），
就会得到 $x_i = x_j$。因此，由引理
\ref{varieties-lemma-glue-separated} 可知，$A_i \otimes A_j \to K$
对所有 $i \not = j$ 都满。
由引理 \ref{varieties-lemma-glue-a-bunch-of-local-rings} 可知
$A = A_1 \cap \ldots \cap A_r$ 是 Noether 半局部环，
恰有 $r$ 个极大理想
$\mathfrak m_1, \ldots, \mathfrak m_r$，使得
$A_i = A_{\mathfrak m_i}$。此外，
$$
\Spec(A) = \Spec(A_1) \cup \ldots \cup \Spec(A_r)
$$
是开覆盖，并且该覆盖任意两块的交都是 $\Spec(K)$。
因此，给定态射 $\Spec(A_i) \to X$ 可粘合为概形态射
$$
\Spec(A) \longrightarrow X
$$
它把 $\mathfrak m_i$ 映到 $x_i$，并诱导局部环同构。
于是结论由引理 \ref{varieties-lemma-points-in-affine} 得出。
\end{proof}

\begin{lemma}
\label{varieties-lemma-extra-silly}
设 $A$ 是环，$I \subset A$ 是理想，
$\mathfrak p_1, \ldots, \mathfrak p_r$ 是 $A$ 的素理想，并且
$\overline{f} \in A/I$ 是元素。若
$I \not \subset \mathfrak p_i$ 对所有 $i$ 都成立，则存在
$f \in A$、$f \not \in \mathfrak p_i$，它映到
$\overline{f}$（在 $A/I$ 中）。
\end{lemma}

\begin{proof}
可以假设诸 $\mathfrak p_i$ 之间没有包含关系
（移除较小的素理想即可）。首先选取任意 $f \in A$ 提升
$\overline{f}$。令 $S$ 为那些 $s \in \{1, \ldots, r\}$ 的集合，
使得 $f \in \mathfrak p_s$。若 $S$ 为空，则已经完成。否则，
考虑理想 $J = I \prod_{i \not \in S} \mathfrak p_i$。
注意，$J$ 不包含在 $\mathfrak p_s$ 中，其中 $s \in S$，
因为诸 $\mathfrak p_i$ 之间没有包含关系，并且
$I$ 不包含在任何 $\mathfrak p_i$ 中。
故由《代数》引理 \ref{algebra-lemma-silly}，可以选取
$g \in J$、$g \not \in \mathfrak p_s$，其中 $s \in S$。
于是 $f + g$ 是引理所述问题的一个解。
\end{proof}

\begin{lemma}
\label{varieties-lemma-finite-set-codim-1-points-in-affine-per-component}
设 $X$ 是概形。设 $T \subset X$ 是有限点集。假设
\begin{enumerate}
\item $X$ 有有限多个不可约分支 $Z_1, \ldots, Z_t$；
\item $Z_i \cap T$ 包含在与 $Z_i$ 对应的赋予约化诱导结构的
子概形的一个仿射开集中。
\end{enumerate}
则存在 $X$ 的仿射开子概形，它包含 $T$。
\end{lemma}

\begin{proof}
使用《极限》命题 \ref{limits-proposition-affine}，
可立即约化到 $X$ 约化的情形。略去细节。在证明的其余部分，
对 $X$ 的每个闭子集都赋予诱导约化闭子概形结构。

\medskip\noindent
我们归纳证明，可以找到仿射开集
$U \subset Z_1 \cup \ldots \cup Z_r$，它包含
$T \cap (Z_1 \cup \ldots \cup Z_r)$。当 $r = 1$ 时，
这由假设成立。
设 $r > 1$，并令 $U \subset Z_1 \cup \ldots \cup Z_{r -  1}$
是包含 $T \cap (Z_1 \cup \ldots \cup Z_{r - 1})$ 的仿射开集。
设 $V \subset X_r$ 是包含 $T \cap Z_r$ 的仿射开集
（由假设存在）。则 $U \cap V$ 包含
$T \cap ( Z_1 \cup \ldots \cup Z_{r - 1} ) \cap Z_r$。
因此
$$
\Delta = (U \cap Z_r) \setminus (U \cap V)
$$
不包含 $T$ 的任何元素。注意 $\Delta$ 是 $U$ 的闭子集。
由素理想回避（《代数》引理 \ref{algebra-lemma-silly}），
可以找到标准开集 $U'$（属于 $U$），它包含 $T \cap U$ 并避开
$\Delta$，即 $U' \cap Z_r \subset U \cap V$。
把 $U$ 替换为 $U'$ 后，可以假设 $U \cap V$ 在 $U$ 中闭。

\medskip\noindent
与上面相同的论证还表明，集合
$\Delta' = (U \cap (Z_1 \cup \ldots \cup Z_{r - 1})) \setminus (U \cap V)$
不包含 $T$ 的任何元素。利用这一点，可以找到
$h \in \mathcal{O}(V)$，使得 $D(h) \subset V$ 包含
$T \cap V$，并且 $U \cap D(h) \subset U \cap V$。
利用 $U \cap V$ 在 $U$ 中闭这一事实，可以应用引理
\ref{varieties-lemma-extra-silly} 找到 $g \in \mathcal{O}(U)$，
它在 $U \cap V$ 上的限制等于 $h$ 在 $U \cap V$ 上的限制，
并且 $T \cap U \subset D(g)$。然后可以把 $U$ 替换为 $D(g)$，
把 $V$ 替换为 $D(h)$，从而达到 $U \cap V$ 在 $U$ 与 $V$
中都闭的情形。此时概形 $U \cup V$ 由《极限》引理
\ref{limits-lemma-affines-glued-in-closed-affine} 是仿射的。
这证明了归纳步骤，从而证明了引理。
\end{proof}

\noindent
下面是我们可以从上述材料中得出的结论。

\begin{proposition}
\label{varieties-proposition-finite-set-of-points-of-codim-1-in-affine}
设 $X$ 是分离概形，使得每个拟紧开集只有有限多个不可约分支。
设 $x_1, \ldots, x_r \in X$ 是点，使得
$\mathcal{O}_{X, x_i}$ 是维数 $\leq 1$ 的 Noether 环。
则存在 $X$ 的仿射开子概形，它包含所有 $x_1, \ldots, x_r$。
\end{proposition}

\begin{proof}
可以替换 $X$，所用拟紧开集包含 $x_1, \ldots, x_r$，
因而可以假设 $X$ 有有限多个不可约分支。
由引理 \ref{varieties-lemma-finite-set-codim-1-points-in-affine-per-component}，
可约化到 $X$ 为整概形的情形。这就是引理
\ref{varieties-lemma-finite-set-codim-1-points-in-affine}。
\end{proof}




\section{曲线}
\label{varieties-section-curves}

\noindent
在 Stacks 计划中，我们采用如下曲线定义。

\begin{definition}
\label{varieties-definition-curve}
设 $k$ 为域。所谓一条{\it 曲线}，是指维数为 $1$ 的、域 $k$ 上的簇。
\end{definition}

\noindent
域 $k$ 上曲线的两个标准例子是仿射直线 $\mathbf{A}^1_k$
和射影直线 $\mathbf{P}^1_k$。
\par\noindent
概形 $X = \Spec(k[x, y]/(f))$
是曲线，当且仅当 $f \in k[x, y]$ 不可约。

\medskip\noindent
曲线的定义具有与簇的定义相同的问题；见定义 \ref{varieties-definition-variety}
之后的讨论。此外，这意味着每条曲线都带有一个指定的定义域。
例如，$X = \Spec(\mathbf{C}[x])$ 是域 $\mathbf{C}$ 上的曲线，
但我们也可以把它看成域 $\mathbf{R}$ 上的曲线。概形
$\Spec(\mathbf{Z})$ 不是曲线，尽管概形 $\Spec(\mathbf{Z})$
与 $\mathbf{A}^1_{\mathbf{F}_p}$ 在许多方面表现相似。

\begin{lemma}
\label{varieties-lemma-proper-minus-point}
设 $X$ 是维数 $> 0$ 的、域 $k$ 上的分离不可约概形。
设 $x \in X$ 为闭点。开子概形 $X \setminus \{x\}$
在 $k$ 上不固有。
\end{lemma}

\begin{proof}
由于 $X$ 不可约，$U = X \setminus \{x\}$ 在 $X$ 中不是闭集。
特别地，浸入 $U \to X$ 不固有。
由《态射》引理 \ref{morphisms-lemma-image-proper-scheme-closed}
（这里使用了 $X$ 分离），$U \to \Spec(k)$ 也不固有。
\end{proof}

\begin{lemma}
\label{varieties-lemma-dim-1-quasi-projective}
设 $X$ 是域 $k$ 上分离的有限型概形。
若 $\dim(X) \leq 1$，则 $X$ 在 $k$ 上 H-拟射影。
\end{lemma}

\begin{proof}
由命题 \ref{varieties-proposition-dim-1-noetherian-separated-has-ample}，
概形 $X$ 具有丰沛可逆层 $\mathcal{L}$。
由《态射》引理 \ref{morphisms-lemma-quasi-projective-finite-type-over-S}，
可知 $X$ 同构于 $\mathbf{P}^n_k$ 在 $\Spec(k)$ 上的一个局部闭子概形。
这正是在 $k$ 上 H-拟射影的定义；见《态射》定义
\ref{morphisms-definition-quasi-projective}。
\end{proof}

\begin{lemma}
\label{varieties-lemma-dim-1-proper-projective}
设 $X$ 是域 $k$ 上的固有概形。
若 $\dim(X) \leq 1$，则 $X$ 在 $k$ 上 H-射影。
\end{lemma}

\begin{proof}
由引理 \ref{varieties-lemma-dim-1-quasi-projective}，可知
$X$ 是 $\mathbf{P}^n_k$ 的局部闭子概形，其中 $k$ 为某个域。
由于 $X$ 在 $k$ 上固有，可知 $X$ 是 $\mathbf{P}^n_k$
的闭子概形（《态射》引理
\ref{morphisms-lemma-image-proper-scheme-closed}）。
\end{proof}

\begin{lemma}
\label{varieties-lemma-dim-1-projective-completion}
设 $X$ 是 $k$ 上分离的有限型概形。
若 $\dim(X) \leq 1$，则存在开浸入
$j : X \to \overline{X}$，满足下列性质：
\begin{enumerate}
\item $\overline{X}$ 在 $k$ 上 H-射影，即
$\overline{X}$ 是 $\mathbf{P}^d_k$ 的闭子概形，对某个 $d$ 成立；
\item $j(X) \subset \overline{X}$ 稠密且概形论稠密；
\item $\overline{X} \setminus X = \{x_1, \ldots, x_n\}$，
其中 $x_i \in \overline{X}$ 是闭点。
\end{enumerate}
\end{lemma}

\begin{proof}
由引理 \ref{varieties-lemma-dim-1-quasi-projective}，可假设
$X$ 是 $\mathbf{P}^d_k$ 的局部闭子概形，其中 $d$ 为某个整数。令
$\overline{X} \subset \mathbf{P}^d_k$ 为 $X \to \mathbf{P}^d_k$
的概形论像；见《态射》定义
\ref{morphisms-definition-scheme-theoretic-image}。
《态射》引理 \ref{morphisms-lemma-quasi-compact-immersion}
中的描述给出性质 (1) 与 (2)。
例如考察泛点并应用引理 \ref{varieties-lemma-dimension-locally-algebraic}，可知
$\dim(X) = 1 \Rightarrow \dim(\overline{X}) = 1$。
由于 $\overline{X}$ 是 Noether 概形，因而
$\overline{X} \setminus X = \{x_1, \ldots, x_n\}$
是有限个闭点组成的集合。
\end{proof}

\begin{lemma}
\label{varieties-lemma-reduced-dim-1-projective-completion}
设 $X$ 是 $k$ 上分离的有限型概形。
若 $X$ 约化且 $\dim(X) \leq 1$，则存在开浸入
$j : X \to \overline{X}$，使得
\begin{enumerate}
\item $\overline{X}$ 在 $k$ 上 H-射影，即
$\overline{X}$ 是 $\mathbf{P}^d_k$ 的闭子概形，对某个 $d$ 成立；
\item $j(X) \subset \overline{X}$ 稠密且概形论稠密；
\item $\overline{X} \setminus X = \{x_1, \ldots, x_n\}$，
其中 $x_i \in \overline{X}$ 是闭点；
\item 局部环 $\mathcal{O}_{\overline{X}, x_i}$ 都是离散赋值环，
对 $i = 1, \ldots, n$ 成立。
\end{enumerate}
\end{lemma}

\begin{proof}
取引理 \ref{varieties-lemma-dim-1-projective-completion} 中的
$j : X \to \overline{X}$。考虑正规化 $X'$，即 $\overline{X}$ 在
$X$ 中的正规化。由引理
\ref{varieties-lemma-relative-normalization-finite}，态射
$X' \to \overline{X}$ 有限。
由《态射》引理 \ref{morphisms-lemma-finite-projective}，
$X' \to \overline{X}$ 射影。由《态射》引理
\ref{morphisms-lemma-projective-over-quasi-projective-is-H-projective}，
可知 $X' \to \overline{X}$ H-射影。
由《态射》引理 \ref{morphisms-lemma-H-projective-composition}，
可知 $X' \to \Spec(k)$ H-射影。
令 $\{x'_1, \ldots, x'_m\} \subset X'$ 为
$\{x_1, \ldots, x_n\} = \overline{X} \setminus X$ 的逆像。
则都有 $\dim(\mathcal{O}_{X', x'_i}) = 1$，
对所有 $1 \leq i \leq m$ 成立。
因此，由《态射》引理
\ref{morphisms-lemma-relative-normalization-normal-codim-1}，
局部环 $\mathcal{O}_{X', x'}$ 都是离散赋值环。
于是 $X \to X'$ 与 $\{x'_1, \ldots, x'_m\}$ 即为所求。
\end{proof}

\begin{lemma}
\label{varieties-lemma-dim-1-projective-completion-after-insep}
设 $X$ 是 $k$ 上分离的有限型概形，且 $\dim(X) \leq 1$。
则存在交换图
$$
\xymatrix{
\overline{Y}_1 \amalg \ldots \amalg \overline{Y}_n \ar[rd] &
Y_1 \amalg \ldots \amalg Y_n \ar[r]_-\nu \ar[d] \ar[l]^j &
X_{k'} \ar[r] \ar[d] &
X \ar[d]^f \\
& \Spec(k'_1) \amalg \ldots \amalg \Spec(k'_n) \ar[r] &
\Spec(k') \ar[r] &
\Spec(k)
}
$$
满足下列性质：
\begin{enumerate}
\item $k'/k$ 是域的有限纯不可分扩张；
\item $\nu$ 是 $X_{k'}$ 的正规化；
\item $j$ 是像稠密的开浸入；
\item $k'_i/k'$ 是有限可分扩张，对 $i = 1, \ldots, n$ 成立；
\item $\overline{Y}_i$ 光滑、射影、几何不可约，维数 $\leq 1$，
其基域为 $k'_i$。
\end{enumerate}
\end{lemma}

\begin{proof}
由于可以把 $X$ 替换为它的约化，我们可以并且确实假设 $X$ 约化。
选取引理 \ref{varieties-lemma-reduced-dim-1-projective-completion}
中的 $X \to \overline{X}$。
若能对 $\overline{X}$ 证明本引理，则对 $X$ 的结论也随之成立
（略去细节）。因此可以并且确实假设 $X$ 射影。

\medskip\noindent
选取有限纯不可分扩张 $k'/k$，使 $X_{k'}$ 的正规化
在 $k'$ 上几何正规；见引理
\ref{varieties-lemma-finite-extension-geometrically-normal}。
以 $Y = (X_{k'})^\nu$ 记该正规化；关于正规化的性质，
见第 \ref{varieties-section-normalization} 节。
$Y$ 几何正则，因为在维数 $\leq 1$ 时正规与正则等价；见《性质》引理
\ref{properties-lemma-normal-dimension-1-regular}。
因而由引理 \ref{varieties-lemma-geometrically-regular-smooth}，
$Y$ 在 $k'$ 上光滑。
设 $Y = Y_1 \amalg \ldots \amalg Y_n$ 为 $Y$ 的不可约分支分解。
置 $k'_i = \Gamma(Y_i, \mathcal{O}_{Y_i})$。
由引理 \ref{varieties-lemma-proper-geometrically-reduced-global-sections}，
这些都是 $k'$ 的有限可分扩张。
最后应用引理 \ref{varieties-lemma-baby-stein} 即完成证明。
\end{proof}

\begin{lemma}
\label{varieties-lemma-regular-point-on-curve}
设 $k$ 为域，$X$ 为域 $k$ 上的曲线，$x \in X$ 为闭点。
我们把 $x$ 看作 $X$ 的一个（约化）闭子概形，其理想层为
$\mathcal{I}$。下列条件等价：
\begin{enumerate}
\item $\mathcal{O}_{X, x}$ 正则；
\item $\mathcal{O}_{X, x}$ 正规；
\item $\mathcal{O}_{X, x}$ 是离散赋值环；
\item $\mathcal{I}$ 是可逆 $\mathcal{O}_X$-模；
\item $x$ 是 $X$ 上的有效 Cartier 除子。
\end{enumerate}
若 $k$ 完全，或者 $\kappa(x)$ 在 $k$ 上可分，
则这些条件还等价于
\begin{enumerate}
\item[(6)] $X \to \Spec(k)$ 在 $x$ 处光滑。
\end{enumerate}
\end{lemma}

\begin{proof}
由于 $X$ 是曲线，局部环 $\mathcal{O}_{X, x}$ 是维数为 $1$ 的
Noether 局部整环（引理 \ref{varieties-lemma-dimension-locally-algebraic}）。
条件 (4) 与 (5) 依定义等价，并且都等价于
$\mathcal{I}_x = \mathfrak m_x \subset \mathcal{O}_{X, x}$
可以由一个元素生成（《除子》引理
\ref{divisors-lemma-effective-Cartier-in-points}）。
因此，(1)、(2)、(3)、(4)、(5) 的等价性由《代数》引理
\ref{algebra-lemma-characterize-dvr} 得出。
当 $k$ 完全时，最后的陈述由引理
\ref{varieties-lemma-dense-smooth-open-variety-over-perfect-field} 得出。
若 $\kappa(x)/k$ 可分，则该等价性由《代数》引理
\ref{algebra-lemma-separable-smooth} 得出。
\end{proof}

\begin{remark}
\label{varieties-remark-smooth-projective-compactification}
设 $k$ 为域，$X$ 为域 $k$ 上的正则曲线。
由引理 \ref{varieties-lemma-regular-point-on-curve} 与
\ref{varieties-lemma-reduced-dim-1-projective-completion}，
存在非奇异射影曲线 $\overline{X}$，它是 $X$ 的紧化；换言之，存在开浸入
$j : X \to \overline{X}$，其补集由有限个闭点组成。
若 $k$ 完全，则 $X$ 与 $\overline{X}$ 都在 $k$ 上光滑，
而 $\overline{X}$ 是 $X$ 的光滑射影紧化。
\end{remark}

\noindent
注意，若仿射概形 $X$ 在 $k$ 上且在 $k$ 上固有，
则 $X$ 在 $k$ 上有限（《态射》引理
\ref{morphisms-lemma-finite-proper}），因而维数为 $0$
（《代数》引理 \ref{algebra-lemma-finite-dimensional-algebra} 与
命题 \ref{algebra-proposition-dimension-zero-ring}）。
因此，维数 $> 0$ 的、域 $k$ 上的概形不可能既仿射又在 $k$ 上固有。
所以，下面引理中的两种可能相互排斥。

\begin{lemma}
\label{varieties-lemma-curve-affine-projective}
设 $X$ 是域 $k$ 上的曲线。则 $X$ 或者是仿射概形，
或者 $X$ 在 $k$ 上 H-射影。
\end{lemma}

\begin{proof}
选取引理 \ref{varieties-lemma-reduced-dim-1-projective-completion} 中的
$X \to \overline{X}$，使得
$\overline{X} \setminus X = \{x_1, \ldots, x_r\}$。
于是 $\overline{X}$ 也是曲线。
若 $r = 0$，则 $X = \overline{X}$ 在 $k$ 上 H-射影。
因此可以假设 $r \geq 1$，我们的目标是证明 $X$ 仿射。
由引理 \ref{varieties-lemma-open-in-affine-curve-affine}，
只须证明 $\overline{X} \setminus \{x_1\}$ 仿射。
这把问题约化为下一段中的断言。

\medskip\noindent
设 $X$ 为域 $k$ 上的 H-射影曲线。设 $x \in X$ 为闭点，
使得 $\mathcal{O}_{X, x}$ 是离散赋值环。我们断言
$U = X \setminus \{x\}$ 仿射。
由引理 \ref{varieties-lemma-regular-point-on-curve}，
点 $x$ 定义了 $X$ 上的有效 Cartier 除子。
对 $n \geq 1$，以 $nx = x + \ldots + x$ 表示 $n$ 重和；见《除子》定义
\ref{divisors-definition-sum-effective-Cartier-divisors}。
以 $\mathcal{O}_{nx}$ 表示 $nx$ 的结构层，把它看作 $X$ 上的凝聚模。
由于局部概形 $nx$ 上的每个可逆模都是平凡的，
《除子》注 \ref{divisors-remark-ses-regular-section}
中的第一个短正合列在这里成为
$$
0 \to \mathcal{O}_X
\xrightarrow{1} \mathcal{O}_X(nx) \to \mathcal{O}_{nx} \to 0
$$
注意，$\dim_k H^0(X, \mathcal{O}_{nx}) \geq n$。
事实上，由引理 \ref{varieties-lemma-chi-tensor-finite}，有
$H^0(X, \mathcal{O}_{nx}) = \mathcal{O}_{X, x}/(\pi^n)$，其中
$\pi$ 是 $\mathcal{O}_{X, x}$ 中的素元，并且幂 $\pi^i$ 映到
关于 $k$ 线性无关的 $\mathcal{O}_{X, x}/(\pi^n)$ 中元素，
对 $i = 0, 1, \ldots, n - 1$ 成立。
由《概形上同调》引理
\ref{coherent-lemma-proper-over-affine-cohomology-finite}，
有 $\dim_k H^1(X, \mathcal{O}_X) < \infty$。
若 $n > \dim_k H^1(X, \mathcal{O}_X)$，
则由同调长正合列可知，存在
$s \in \Gamma(X, \mathcal{O}_X(nx))$，
它不是 $\mathcal{O}_X$ 的截面。
若取具有此性质的最小 $n$，则 $s$ 映到茎
$\left(\mathcal{O}_X(nx)\right)_x$ 的一个生成元；
否则它会定义 $\mathcal{O}_X((n - 1)x) \subset \mathcal{O}_X(nx)$
的一个截面。对这个 $n$，可知 $s_0 = 1$ 与 $s_1 = s$
生成可逆模 $\mathcal{L} = \mathcal{O}_X(nx)$。

\medskip\noindent
考虑《构造》第 \ref{constructions-section-projective-space} 节中的
相应态射
$f = \varphi_{\mathcal{L}, (s_0, s_1)} : X \to \mathbf{P}^1_k$。
注意，$D_{+}(T_0)$ 的逆像是 $U = X \setminus \{x\}$，
因为截面 $s_0$（作为 $\mathcal{L}$ 的截面）仅在 $x$ 处消失。
特别地，$f$ 不是常值态射，即 $\Im(f)$ 含有不止一个点。
因此，$f$ 必把泛点 $\eta$（$X$ 的泛点）映到
$\mathbf{P}^1_k$ 的泛点。所以，若
$y \in \mathbf{P}^1_k$ 是闭点，则 $f^{-1}(\{y\})$
是 $X$ 中不含 $\eta$ 的闭集，因而是有限集。
最后，$f$ 固有\footnote{事实上，由《态射》引理
\ref{morphisms-lemma-projective-is-quasi-projective-proper}，
H-射影簇是固有簇。由《态射》引理
\ref{morphisms-lemma-image-proper-scheme-closed}，
源为固有簇的簇之间的态射是固有态射。}。
由《概形上同调》引理
\ref{coherent-lemma-proper-finite-fibre-finite-in-neighbourhood}\footnote{可以避免使用这个依赖形式函数定理的引理。
事实上，$X$ 射影，因此只须证明：固有态射 $f : X \to Y$
若连接两个在 $k$ 上拟射影的概形且纤维有限，则是有限态射。为此，选取 $X$ 中包含
$f$ 在某点 $y$ 上之纤维的仿射开集；这里使用了拟射影概形的任意有限点集
在域 $k$ 上都包含于某个仿射开集这一事实。把 $Y$ 缩小到 $y$ 的一个小仿射
邻域后，问题约化为仿射概形之间的固有态射。由《态射》引理
\ref{morphisms-lemma-integral-universally-closed}，这种态射有限。}，
可知 $f$ 有限。因此 $U = f^{-1}(D_{+}(T_0))$ 仿射。
\end{proof}

\noindent
下面的引理与引理 \ref{varieties-lemma-proper-minus-point} 合用说明：
给定分离概形 $X$，其维数为 $1$ 且在 $k$ 上有限型，
那么 $X \setminus Z$ 是仿射概形，只要闭子集 $Z$
与 $X$ 的每个不可约分支相交。

\begin{lemma}
\label{varieties-lemma-dim-1-nonproper-affine}
设 $X$ 是 $k$ 上分离的有限型概形。
若 $\dim(X) \leq 1$，并且 $X$ 的不可约分支中没有维数为 $1$
的固有分支，则 $X$ 仿射。
\end{lemma}

\begin{proof}
令 $X = \bigcup X_i$ 为 $X$ 的不可约分支分解。
我们把 $X_i$ 看作整概形（赋予约化诱导概形结构；见《概形》定义
\ref{schemes-definition-reduced-induced-scheme}）。
特别地，$X_i$ 或者是单点集（因而仿射），或者是曲线，
从而由引理 \ref{varieties-lemma-curve-affine-projective} 仿射。
于是 $\coprod X_i \to X$ 是有限满射，而 $\coprod X_i$ 仿射。
因此，由《概形上同调》引理
\ref{coherent-lemma-image-affine-finite-morphism-affine-Noetherian}，
可知 $X$ 仿射。
\end{proof}




\section{曲线上的次数}
\label{varieties-section-divisors-curves}

\noindent
我们开始定义域上维数为 $1$ 的固有概形之可逆层，乃至更一般的
局部自由层的次数。在第 \ref{varieties-section-euler} 节中，
我们定义了凝聚层 $\mathcal{F}$（位于固有概形 $X$ 上，基域为
$k$）的欧拉示性数，公式为
$$
\chi(X, \mathcal{F}) = \sum (-1)^i \dim_k H^i(X, \mathcal{F}).
$$

\begin{definition}
\label{varieties-definition-degree-invertible-sheaf}
设 $k$ 为域，$X$ 为维数 $\leq 1$ 的固有概形，其基域为 $k$，
$\mathcal{L}$ 为可逆 $\mathcal{O}_X$-模。$\mathcal{L}$ 的{\it 次数}定义为
$$
\deg(\mathcal{L}) = \chi(X, \mathcal{L}) - \chi(X, \mathcal{O}_X)
$$
更一般地，若 $\mathcal{E}$ 是秩为 $n$ 的局部自由层，
则定义 $\mathcal{E}$ 的{\it 次数}为
$$
\deg(\mathcal{E}) = \chi(X, \mathcal{E}) - n\chi(X, \mathcal{O}_X)
$$
\end{definition}

\noindent
注意，这依赖于三元组 $\mathcal{E}/X/k$。
若 $X$ 不连通，而 $\mathcal{E}$ 有限局部自由（但秩不恒定），
则可以修改定义，对 $\mathcal{E}$ 在 $X$ 各连通分支上的限制之次数求和。
若 $\mathcal{E}$ 只是凝聚层，则有若干不同的方式扩展此定义\footnote{若
$X$ 是固有曲线，$\mathcal{F}$ 是 $X$ 上的凝聚层，
则次数常定义为 $\chi(X, \mathcal{F}) - r\chi(X, \mathcal{O}_X)$，其中
$r = \dim_{\kappa(\xi)} \mathcal{F}_\xi$ 是 $\mathcal{F}$
在泛点 $\xi$（$X$ 的泛点）处的秩。}。
下面用一系列引理证明，该定义具有次数应有的一切性质。

\begin{lemma}
\label{varieties-lemma-degree-base-change}
设 $k'/k$ 为域扩张。设 $X$ 是维数 $\leq 1$ 的固有概形，其基域为 $k$。
设 $\mathcal{E}$ 是局部自由 $\mathcal{O}_X$-模，其秩恒为 $n$。
则 $\mathcal{E}/X/k$ 的次数等于
$\mathcal{E}_{k'}/X_{k'}/k'$ 的次数。
\end{lemma}

\begin{proof}
更明确地，置 $X_{k'} = X \times_{\Spec(k)} \Spec(k')$。
令 $\mathcal{E}_{k'} = p^*\mathcal{E}$，其中
$p : X_{k'} \to X$ 是投影。由《概形上同调》引理
\ref{coherent-lemma-flat-base-change-cohomology}，有
$H^i(X_{k'}, \mathcal{E}_{k'}) = H^i(X, \mathcal{E}) \otimes_k k'$
以及
$H^i(X_{k'}, \mathcal{O}_{X_{k'}}) = H^i(X, \mathcal{O}_X) \otimes_k k'$。
因此欧拉示性数不变，次数也不变。
\end{proof}

\begin{lemma}
\label{varieties-lemma-degree-additive}
设 $k$ 为域。设 $X$ 是维数 $\leq 1$ 的固有概形，其基域为 $k$。
设 $0 \to \mathcal{E}_1 \to \mathcal{E}_2 \to \mathcal{E}_3 \to 0$
是有限恒秩局部自由 $\mathcal{O}_X$-模的短正合列。则
$$
\deg(\mathcal{E}_2) = \deg(\mathcal{E}_1) + \deg(\mathcal{E}_3)
$$
\end{lemma}

\begin{proof}
这立即由欧拉示性数的可加性（引理
\ref{varieties-lemma-euler-characteristic-additive}）与秩的可加性得出。
\end{proof}

\begin{lemma}
\label{varieties-lemma-degree-birational-pullback}
设 $k$ 为域。设 $f : X' \to X$ 是维数 $\leq 1$
的固有概形之间的双有理态射，这些概形定义在 $k$ 上。则
$$
\deg(f^*\mathcal{E}) = \deg(\mathcal{E})
$$
对每个恒秩有限局部自由层都成立。更一般地，只须 $f$
在维数为 $1$ 的不可约分支之间诱导双射，并在相应泛点处的局部环之间诱导同构。
\end{lemma}

\begin{proof}
态射 $f$ 固有（《态射》引理
\ref{morphisms-lemma-image-proper-scheme-closed}），
且其纤维维数 $\leq 0$。因此 $f$ 有限（《概形上同调》引理
\ref{coherent-lemma-proper-finite-fibre-finite-in-neighbourhood}）。
从而
$$
Rf_*f^*\mathcal{E} = f_*f^*\mathcal{E} =
\mathcal{E} \otimes_{\mathcal{O}_X} f_*\mathcal{O}_{X'}
$$
由于 $f$ 在所有维数为 $1$ 的不可约分支的泛点处的局部环上诱导同构，
可知核与余核
$$
0 \to \mathcal{K} \to \mathcal{O}_X \to f_*\mathcal{O}_{X'}
\to \mathcal{Q} \to 0
$$
的支撑维数 $\leq 0$。注意，由于 $\mathcal{E}$ 局部自由，
把该列与 $\mathcal{E}$ 张量后仍是正合列。于是得到
\begin{align*}
\chi(X, \mathcal{E}) - \chi(X', f^*\mathcal{E})
& =
\chi(X, \mathcal{E}) - \chi(X, f_*f^*\mathcal{E}) \\
& =
\chi(X, \mathcal{E}) - \chi(X, \mathcal{E} \otimes f_*\mathcal{O}_{X'}) \\
& =
\chi(X, \mathcal{K} \otimes \mathcal{E}) -
\chi(X, \mathcal{Q} \otimes \mathcal{E}) \\
& =
n\chi(X, \mathcal{K}) -
n\chi(X, \mathcal{Q}) \\
& =
n\chi(X, \mathcal{O}_X) - n\chi(X, f_*\mathcal{O}_{X'}) \\
& =
n\chi(X, \mathcal{O}_X) - n\chi(X', \mathcal{O}_{X'})
\end{align*}
这正是所需结论。第一个等号来自 $f$ 有限；见《概形上同调》引理
\ref{coherent-lemma-relative-affine-cohomology}。
第二个等号来自投影公式；见《同调》引理
\ref{cohomology-lemma-projection-formula}。
第三个等号来自欧拉示性数的可加性；见引理
\ref{varieties-lemma-euler-characteristic-additive}。
第四个等号来自引理 \ref{varieties-lemma-chi-tensor-finite}。
\end{proof}

\begin{lemma}
\label{varieties-lemma-degree-on-proper-curve}
设 $k$ 为域。设 $X$ 是 $k$ 上泛点为 $\xi$ 的固有曲线。
设 $\mathcal{E}$ 是局部自由 $\mathcal{O}_X$-模，其秩为 $n$，
$\mathcal{F}$ 是凝聚 $\mathcal{O}_X$-模。则
$$
\chi(X, \mathcal{E} \otimes \mathcal{F}) =
r \deg(\mathcal{E}) + n \chi(X, \mathcal{F})
$$
其中 $r = \dim_{\kappa(\xi)} \mathcal{F}_\xi$ 是 $\mathcal{F}$ 的秩。
\end{lemma}

\begin{proof}
令 $\mathcal{P}$ 为凝聚层 $\mathcal{F}$（位于 $X$ 上）的如下性质：
引理中的公式成立。我们断言，《概形上同调》引理
\ref{coherent-lemma-property} 的假设 (1) 与 (2) 对 $\mathcal{P}$ 成立。
事实上，(1) 成立，因为欧拉示性数与秩 $r$
在凝聚层的短正合列中都可加。(2) 也成立：若 $Z = X$，
则可取 $\mathcal{G} = \mathcal{O}_X$，而
$\mathcal{P}(\mathcal{O}_X)$ 依次数定义为真。
若 $i : Z \to X$ 是一个闭点的包含态射，则可取
$\mathcal{G} = i_*\mathcal{O}_Z$；而 $\mathcal{P}$ 成立，
这由引理 \ref{varieties-lemma-chi-tensor-finite} 以及此时 $r = 0$ 这一事实得出。
\end{proof}

\noindent
设 $k$ 为域。设 $X$ 是 $k$ 上维数 $\leq 1$ 的有限型概形。
设 $C_i \subset X$（$i = 1, \ldots, t$）为维数 $1$ 的不可约分支。
我们赋予 $C_i$ 约化诱导概形结构，从而把它看作概形。
令 $\xi_i \in C_i$ 为泛点。$C_i$ 在 $X$ 中的{\it 重数}定义为长度
$$
m_i = \text{length}_{\mathcal{O}_{X, \xi_i}} \mathcal{O}_{X, \xi_i}
$$
这是有意义的，因为 $\mathcal{O}_{X, \xi_i}$ 是零维 Noether 局部环，
因而作为自身上的模具有有限长度（《代数》命题
\ref{algebra-proposition-dimension-zero-ring}）。
更多信息见《Chow 同调》第
\ref{chow-section-cycle-of-closed-subscheme} 节。
事实上，局部自由层的次数只依赖于它在各不可约分支上的限制。

\begin{lemma}
\label{varieties-lemma-degree-in-terms-of-components}
设 $k$ 为域。设 $X$ 是维数 $\leq 1$ 的固有概形，其基域为 $k$。
设 $\mathcal{E}$ 是局部自由 $\mathcal{O}_X$-模，其秩为 $n$。则
$$
\deg(\mathcal{E}) = \sum m_i \deg(\mathcal{E}|_{C_i})
$$
其中 $C_i \subset X$（$i = 1, \ldots, t$）是赋予约化诱导概形结构的
维数 $1$ 的不可约分支，而 $m_i$ 是 $C_i$ 在 $X$ 中的重数。
\end{lemma}

\begin{proof}
先注意，该陈述有意义，因为 $C_i \to \Spec(k)$ 固有且维数为 $1$；
下面还将使用如下事实：由《态射》引理
\ref{morphisms-lemma-closed-immersion-proper} 与
\ref{morphisms-lemma-composition-proper}，$X$ 的任意闭子概形都在 $k$ 上固有。
考虑开子概形 $U_i = X \setminus (\bigcup_{j \not = i} C_j)$，
并令 $X_i \subset X$ 为 $U_i$ 的概形论闭包。
注意，$X_i \cap U_i = U_i$（概形论意义下），而
$X_i \cap U_j = \emptyset$（集合论意义下），对 $i \not = j$ 成立；这由《态射》引理
\ref{morphisms-lemma-quasi-compact-immersion} 中对概形论闭包的描述得出。
因此可把引理 \ref{varieties-lemma-degree-birational-pullback} 应用于态射
$X' = \coprod X_i \to X$。显然，$C_i \subset X_i$（概形论意义下），
且 $C_i$ 在 $X_i$ 中的重数等于 $C_i$ 在 $X$ 中的重数，
故问题约化为下一段所讨论的情形。

\medskip\noindent
假设 $X$ 不可约，泛点为 $\xi$。设
$C = X_{red}$ 的重数为 $m$。
我们必须证明 $\deg(\mathcal{E}) = m \deg(\mathcal{E}|_C)$。
令 $\mathcal{I} \subset \mathcal{O}_X$ 为定义闭子概形 $C$ 的理想。
取最小的 $e \geq 0$，使 $\mathcal{I}^{e + 1} = 0$
（《概形上同调》引理
\ref{coherent-lemma-power-ideal-kills-sheaf}）。
我们对 $e$ 作归纳。若 $e = 0$，则 $X = C$，结论显然。
否则，置 $\mathcal{F} = \mathcal{I}^e$，将其看作凝聚
$\mathcal{O}_C$-模（《概形上同调》引理
\ref{coherent-lemma-i-star-equivalence}）。
令 $X' \subset X$ 为凝聚理想 $\mathcal{I}^e$ 截出的闭子概形，
并令 $m'$ 为 $C$ 在 $X'$ 中的重数。取短正合列
$$
0 \to \mathcal{F} \to \mathcal{O}_X \to \mathcal{O}_{X'} \to 0
$$
在 $\xi$ 处的茎，可得（使用《代数》引理
\ref{algebra-lemma-length-additive}、
\ref{algebra-lemma-dimension-is-length} 与
\ref{algebra-lemma-length-independent}）
$$
m = \text{length}_{\mathcal{O}_{X, \xi}} \mathcal{O}_{X, \xi}
= \dim_{\kappa(\xi)} \mathcal{F}_\xi +
\text{length}_{\mathcal{O}_{X', \xi}} \mathcal{O}_{X', \xi}
= r + m'
$$
其中 $r$ 是 $\mathcal{F}$ 作为 $C$ 上凝聚层的秩。
与 $\mathcal{E}$ 张量后，得到短正合列
$$
0 \to \mathcal{E}|_C \otimes \mathcal{F} \to \mathcal{E} \to
\mathcal{E} \otimes \mathcal{O}_{X'} \to 0
$$
由归纳假设，有
$\deg(\mathcal{E}|_{X'}) = m' \deg(\mathcal{E}|_C)$。
由引理 \ref{varieties-lemma-degree-on-proper-curve}，有
$\chi(\mathcal{E}|_C \otimes \mathcal{F}) =
r \deg(\mathcal{E}|_C) + n \chi(\mathcal{F})$。
综合以上各式，即得结论。
\end{proof}




\begin{lemma}
\label{varieties-lemma-degree-tensor-product}
设 $k$ 为域，$X$ 为维数 $\leq 1$ 的固有概形，其基域为 $k$，
$\mathcal{E}$、$\mathcal{V}$ 为恒定有限秩的局部自由
$\mathcal{O}_X$-模。则
$$
\deg(\mathcal{E} \otimes \mathcal{V}) =
\text{rank}(\mathcal{E}) \deg(\mathcal{V}) +
\text{rank}(\mathcal{V}) \deg(\mathcal{E})
$$
\end{lemma}

\begin{proof}
由引理 \ref{varieties-lemma-degree-in-terms-of-components} 与初等算术，
问题约化为固有曲线的情形。此情形由引理
\ref{varieties-lemma-degree-on-proper-curve} 得出。
\end{proof}

\begin{lemma}
\label{varieties-lemma-degree-and-det}
设 $k$ 为域，$X$ 为维数 $\leq 1$ 的固有概形，其基域为 $k$，
$\mathcal{E}$ 为局部自由 $\mathcal{O}_X$-模，秩为 $n$。则
$$
\deg(\mathcal{E}) = \deg(\wedge^n(\mathcal{E})) = \deg(\det(\mathcal{E}))
$$
\end{lemma}

\begin{proof}
由引理 \ref{varieties-lemma-degree-in-terms-of-components} 与初等算术，
问题约化为固有曲线的情形。
此时存在修正 $f : X' \to X$，使得
$f^*\mathcal{E}$ 具有一个逐次商均为可逆模的滤过；见《除子》引理
\ref{divisors-lemma-filter-after-modification}。
由引理 \ref{varieties-lemma-degree-birational-pullback}，可在 $X'$ 上处理。
因此，可以假设存在由局部自由 $\mathcal{O}_X$-模组成的滤过
$$
0 = \mathcal{E}_0 \subset \mathcal{E}_1 \subset \mathcal{E}_2 \subset
\ldots \subset \mathcal{E}_n = \mathcal{E}
$$
使得 $\mathcal{L}_i = \mathcal{E}_i/\mathcal{E}_{i - 1}$ 可逆。
由《模》引理 \ref{modules-lemma-det-ses} 及归纳法，可得
$\det(\mathcal{E}) = \mathcal{L}_1 \otimes \ldots \otimes \mathcal{L}_n$。
于是等式由引理 \ref{varieties-lemma-degree-tensor-product}
与可加性（引理 \ref{varieties-lemma-degree-additive}）得出。
\end{proof}

\begin{lemma}
\label{varieties-lemma-degree-effective-Cartier-divisor}
设 $k$ 为域，$X$ 为维数 $\leq 1$ 的固有概形，其基域为 $k$。
设 $D$ 是 $X$ 上的有效 Cartier 除子。
则 $D$ 在 $\Spec(k)$ 上有限，次数为
$\deg(D) = \dim_k \Gamma(D, \mathcal{O}_D)$。对于局部自由层
$\mathcal{E}$，若其秩为 $n$，则有
$$
\deg(\mathcal{E}(D)) = n\deg(D) + \deg(\mathcal{E})
$$
其中 $\mathcal{E}(D) = \mathcal{E} \otimes_{\mathcal{O}_X} \mathcal{O}_X(D)$。
\end{lemma}

\begin{proof}
由于 $D$ 在 $X$ 中无处稠密（《除子》引理
\ref{divisors-lemma-complement-effective-Cartier-divisor}），
可知 $\dim(D) \leq 0$。因此由引理
\ref{varieties-lemma-algebraic-scheme-dim-0}，$D$ 在 $k$ 上有限。
由于 $k$ 为域，态射 $D \to \Spec(k)$ 有限局部自由，
因而具有次数（《态射》定义
\ref{morphisms-definition-finite-locally-free}）；该次数显然等于
$\dim_k \Gamma(D, \mathcal{O}_D)$，如引理所述。
由《除子》定义
\ref{divisors-definition-invertible-sheaf-effective-Cartier-divisor}，
存在短正合列
$$
0 \to \mathcal{O}_X \to \mathcal{O}_X(D) \to i_*i^*\mathcal{O}_X(D) \to 0
$$
其中 $i : D \to X$ 是闭浸入。与 $\mathcal{E}$ 张量，得到短正合列
$$
0 \to \mathcal{E} \to \mathcal{E}(D) \to i_*i^*\mathcal{E}(D) \to 0
$$
引理中的等式由欧拉示性数的可加性（引理
\ref{varieties-lemma-euler-characteristic-additive}）与引理
\ref{varieties-lemma-chi-tensor-finite} 得出。
\end{proof}

\begin{lemma}
\label{varieties-lemma-divisible}
设 $k$ 为域。设 $X$ 是 $k$ 上约化且连通的固有概形。
令 $\kappa = H^0(X, \mathcal{O}_X)$。则 $\kappa/k$
是有限域扩张，并且 $w = [\kappa : k]$ 整除下列各数：
\begin{enumerate}
\item $\deg(\mathcal{E})$，其中所取的是任意局部自由
$\mathcal{O}_X$-模 $\mathcal{E}$；
\item $[\kappa(x) : k]$，其中 $x \in X$ 为任意闭点；以及
\item $\deg(D)$，其中 $D \subset X$ 为任意零维闭子概形。
\end{enumerate}
\end{lemma}

\begin{proof}
关于 $\kappa$ 的断言见引理
\ref{varieties-lemma-proper-geometrically-reduced-global-sections}。
对每个拟凝聚 $\mathcal{O}_X$-模，$k$-向量空间
$H^i(X, \mathcal{F})$ 都是 $\kappa$-向量空间。
上述整除性很容易由这一陈述与各定义得出。
\end{proof}



\begin{lemma}
\label{varieties-lemma-degree-pullback-map-proper-curves}
设 $k$ 为域。设 $f : X \to Y$ 是 $k$ 上固有曲线之间的非常值态射。
设 $\mathcal{E}$ 是局部自由 $\mathcal{O}_Y$-模。则
$$
\deg(f^*\mathcal{E}) = \deg(X/Y) \deg(\mathcal{E})
$$
\end{lemma}

\begin{proof}
$X$ 在 $Y$ 上的次数定义见《态射》定义
\ref{morphisms-definition-degree}。
因此 $f_*\mathcal{O}_X$ 是凝聚 $\mathcal{O}_Y$-模，其秩为
$\deg(X/Y)$；换言之，
$\deg(X/Y) = \dim_{\kappa(\xi)} (f_*\mathcal{O}_X)_\xi$，其中
$\xi$ 是 $Y$ 的泛点。于是得到
\begin{align*}
\chi(X, f^*\mathcal{E})
& =
\chi(Y, f_*f^*\mathcal{E}) \\
& =
\chi(Y, \mathcal{E} \otimes f_*\mathcal{O}_X) \\
& =
\deg(X/Y) \deg(\mathcal{E}) + n \chi(Y, f_*\mathcal{O}_X) \\
& =
\deg(X/Y) \deg(\mathcal{E}) + n \chi(X, \mathcal{O}_X)
\end{align*}
这正是所需结论。第一个等号来自 $f$ 有限；见《概形上同调》引理
\ref{coherent-lemma-relative-affine-cohomology}。
第二个等号来自投影公式；见《同调》引理
\ref{cohomology-lemma-projection-formula}。
第三个等号来自引理 \ref{varieties-lemma-degree-on-proper-curve}。
\end{proof}

\noindent
下面是以上次数结果的一个平凡但重要的推论。

\begin{lemma}
\label{varieties-lemma-check-invertible-sheaf-trivial}
设 $k$ 为域。设 $X$ 是 $k$ 上的固有曲线。
设 $\mathcal{L}$ 是可逆 $\mathcal{O}_X$-模。
\begin{enumerate}
\item 若 $\mathcal{L}$ 有非零截面，则
$\deg(\mathcal{L}) \geq 0$。
\item 若 $\mathcal{L}$ 有一个在某点处消失的非零截面 $s$，则
$\deg(\mathcal{L}) > 0$。
\item 若 $\mathcal{L}$ 与 $\mathcal{L}^{-1}$ 都有非零截面，则
$\mathcal{L} \cong \mathcal{O}_X$。
\item 若 $\deg(\mathcal{L}) \leq 0$ 且 $\mathcal{L}$ 有非零截面，则
$\mathcal{L} \cong \mathcal{O}_X$。
\item 若 $\mathcal{N} \to \mathcal{L}$ 是可逆
$\mathcal{O}_X$-模之间的非零映射，则
$\deg(\mathcal{L}) \geq \deg(\mathcal{N})$；若等号成立，则该映射是同构。
\end{enumerate}
\end{lemma}

\begin{proof}
设 $s$ 是 $\mathcal{L}$ 的非零截面。由于 $X$ 是曲线，
可知 $s$ 是正则截面。因此存在有效 Cartier 除子 $D \subset X$
以及同构 $\mathcal{L} \to \mathcal{O}_X(D)$，它把 $s$
映到典范截面 $1$，即 $\mathcal{O}_X(D)$ 的典范截面；见《除子》引理
\ref{divisors-lemma-characterize-OD}。
于是由引理 \ref{varieties-lemma-degree-effective-Cartier-divisor}，
$\deg(\mathcal{L}) = \deg(D)$。
由于 $\deg(D) \geq 0$，且其 $= 0$ 当且仅当
$D = \emptyset$，这证明了 (1) 与 (2)。
在情形 (3) 中，可知 $\deg(\mathcal{L}) = 0$ 且
$D = \emptyset$。情形 (4) 同理。为证明 (5)，把 (1) 与 (4)
应用于可逆层
$$
\mathcal{L} \otimes_{\mathcal{O}_X} \mathcal{N}^{\otimes -1} =
\SheafHom_{\mathcal{O}_X}(\mathcal{N}, \mathcal{L})
$$
即可；由引理 \ref{varieties-lemma-degree-tensor-product}，该层的次数为
$\deg(\mathcal{L}) - \deg(\mathcal{N})$。
\end{proof}

\begin{lemma}
\label{varieties-lemma-no-sections-dual-nef}
设 $k$ 为域。设 $X$ 是 $k$ 上约化、连通且等维数的固有概形，
其维数为 $1$。设 $\mathcal{L}$ 是可逆 $\mathcal{O}_X$-模。
若 $\deg(\mathcal{L}|_C) \leq 0$ 对每个不可约分支 $C$
均成立（这些分支均取自 $X$），则或者
$H^0(X, \mathcal{L}) = 0$，或者
$\mathcal{L} \cong \mathcal{O}_X$。
\end{lemma}

\begin{proof}
设 $s \in H^0(X, \mathcal{L})$ 非零。由于 $X$ 约化，存在
不可约分支 $C$，它是 $X$ 的分支，并且 $s|_C \not = 0$。
然而，若 $s|_C$ 非零，则由引理
\ref{varieties-lemma-check-invertible-sheaf-trivial}，$s$
在 $C$ 上处处不消失。这又蕴含 $s$ 在 $X$ 的每个与 $C$
相交的不可约分支上都处处不消失。
由于 $X$ 连通，可知 $s$ 无处消失，引理得证。
\end{proof}

\begin{lemma}
\label{varieties-lemma-ample-curve}
设 $k$ 为域。设 $X$ 是 $k$ 上的固有曲线。
设 $\mathcal{L}$ 是可逆 $\mathcal{O}_X$-模。
则 $\mathcal{L}$ 丰沛，当且仅当 $\deg(\mathcal{L}) > 0$。
\end{lemma}

\begin{proof}
若 $\mathcal{L}$ 丰沛，则存在 $n > 0$ 以及截面
$s \in H^0(X, \mathcal{L}^{\otimes n})$，使 $X_s$ 仿射。
由于 $X$ 不仿射（否则由《态射》引理
\ref{morphisms-lemma-finite-proper}，$X$ 将是有限概形），
可知 $s$ 在某点处消失。因此由引理
\ref{varieties-lemma-check-invertible-sheaf-trivial}，有
$\deg(\mathcal{L}^{\otimes n}) > 0$。
由引理 \ref{varieties-lemma-degree-tensor-product}，可得
$\deg(\mathcal{L}) = 1/n\deg(\mathcal{L}^{\otimes n}) > 0$。

\medskip\noindent
假设 $\deg(\mathcal{L}) > 0$。则
$$
\dim_k H^0(X, \mathcal{L}^{\otimes n}) \geq \chi(X, \mathcal{L}^n)
= n\deg(\mathcal{L}) + \chi(X, \mathcal{O}_X)
$$
随 $n$ 线性增长。因此，对任意有限个闭点
$x_1, \ldots, x_t$（均在 $X$ 上），可以找到 $n$，使得
$\dim_k H^0(X, \mathcal{L}^{\otimes n}) > \sum \dim_k \kappa(x_i)$。
（回忆，由 Hilbert 零点定理，域扩张
$\kappa(x_i)/k$ 有限；例如见《态射》引理
\ref{morphisms-lemma-closed-point-fibre-locally-finite-type}。）
因此可以找到非零截面 $s \in H^0(X, \mathcal{L}^{\otimes n})$，
它在 $x_1, \ldots, x_t$ 处消失。特别地，若选取
$x_1, \ldots, x_t$，使 $X \setminus \{x_1, \ldots, x_t\}$ 仿射，
则 $X_s$ 也仿射（例如由《性质》引理
\ref{properties-lemma-affine-cap-s-open}；不过，若选取该有限集，使得
$\mathcal{L}|_{X \setminus \{x_1, \ldots, x_t\}}$ 平凡，
则结论是立即的）。所以可以找到 $n > 0$ 与非零截面
$s \in H^0(X, \mathcal{L}^{\otimes n})$，使 $X_s$ 仿射。

\medskip\noindent
我们将证明：对每个拟凝聚理想层 $\mathcal{I}$，
都存在 $m > 0$，使
$H^1(X, \mathcal{I} \otimes \mathcal{L}^{\otimes m})$ 为零。
由《概形上同调》引理
\ref{coherent-lemma-vanshing-gives-ample}，这将完成证明。
为此，考虑映射
$$
\mathcal{I} \xrightarrow{s}
\mathcal{I} \otimes \mathcal{L}^{\otimes n} \xrightarrow{s}
\mathcal{I} \otimes \mathcal{L}^{\otimes 2n} \xrightarrow{s} \ldots
$$
由于 $\mathcal{I}$ 无挠，这些映射是单射，并且在 $X_s$ 上是同构；
因而这些映射的余核的 $H^1$ 为零（例如由《概形上同调》引理
\ref{coherent-lemma-coherent-support-dimension-0}）。
可知向量空间的映射
$$
H^1(X, \mathcal{I}) \to
H^1(X, \mathcal{I} \otimes \mathcal{L}^{\otimes n}) \to
H^1(X, \mathcal{I} \otimes \mathcal{L}^{\otimes 2n}) \to \ldots
$$
均为满射。另一方面，$H^1(X, \mathcal{I})$ 的维数有限，
且由《概形上同调》引理
\ref{coherent-lemma-section-affine-open-kills-classes}，
每个元素最终都映到零。因此，对某个 $e > 0$，有
$H^1(X, \mathcal{I} \otimes \mathcal{L}^{\otimes en})$ 为零。
证明完成。
\end{proof}

\begin{lemma}
\label{varieties-lemma-ampleness-in-terms-of-degrees-components}
设 $k$ 为域。设 $X$ 是维数 $\leq 1$ 的固有概形，其基域为 $k$。
设 $\mathcal{L}$ 是可逆 $\mathcal{O}_X$-模。
设 $C_i \subset X$（$i = 1, \ldots, t$）为维数 $1$ 的不可约分支。
下列条件等价：
\begin{enumerate}
\item $\mathcal{L}$ 丰沛；以及
\item $\deg(\mathcal{L}|_{C_i}) > 0$ 对 $i = 1, \ldots, t$ 均成立。
\end{enumerate}
\end{lemma}

\begin{proof}
设 $x_1, \ldots, x_r \in X$ 为孤立闭点。
把 $x_i = \Spec(\kappa(x_i))$ 看作概形。考虑概形态射
$$
f : C_1 \amalg \ldots \amalg C_t \amalg x_1 \amalg \ldots \amalg x_r
\longrightarrow X
$$
这是 $k$ 上固有概形之间的有限满射（略去细节）。
因此，$\mathcal{L}$ 丰沛，当且仅当 $f^*\mathcal{L}$ 丰沛
（《概形上同调》引理
\ref{coherent-lemma-surjective-finite-morphism-ample}）。
于是结论由引理 \ref{varieties-lemma-ample-curve} 得出。
\end{proof}

\begin{lemma}
\label{varieties-lemma-general-degree-g-line-bundle}
设 $k$ 为代数闭域。设 $X$ 是 $k$ 上的固有曲线。则存在
\begin{enumerate}
\item 可逆 $\mathcal{O}_X$-模 $\mathcal{L}$，使
$\dim_k H^0(X, \mathcal{L}) = 1$ 且 $H^1(X, \mathcal{L}) = 0$；以及
\item 可逆 $\mathcal{O}_X$-模 $\mathcal{N}$，使
$\dim_k H^0(X, \mathcal{N}) = 0$ 且 $H^1(X, \mathcal{N}) = 0$。
\end{enumerate}
\end{lemma}

\begin{proof}
选取闭浸入 $i : X \to \mathbf{P}^n_k$
（引理 \ref{varieties-lemma-dim-1-proper-projective}）。
置 $\mathcal{L} = i^*\mathcal{O}_{\mathbf{P}^n}(d)$，其中
$d \gg 0$，
可知存在可逆层 $\mathcal{L}$，满足
$H^0(X, \mathcal{L}) \not = 0$ 且
$H^1(X, \mathcal{L}) = 0$（关于消失，见《概形上同调》引理
\ref{coherent-lemma-vanshing-gives-ample}；关于非消失，见该处所列参考）。
我们对 $t = \dim_k H^0(X, \mathcal{L})$ 作降归纳，以完成 (1) 的证明。
基础情形 $t = 1$ 显然。假设 $t > 1$。

\medskip\noindent
令 $U \subset X$ 为引理
\ref{varieties-lemma-dense-smooth-open-variety-over-perfect-field}
中所研究的非奇异点组成的非空开集。
设 $s \in H^0(X, \mathcal{L})$ 非零。存在闭点
$x \in U$，使得 $s$ 在 $x$ 处不消失。
设 $\mathcal{I}$ 为引理 \ref{varieties-lemma-regular-point-on-curve}
中 $i : x \to X$ 的理想层。考察短正合列
$$
0 \to \mathcal{I} \otimes_{\mathcal{O}_X} \mathcal{L} \to
\mathcal{L} \to i_*i^*\mathcal{L} \to 0
$$
注意，$H^0(X, i_*i^*\mathcal{L}) = H^0(x, i^*\mathcal{L})$
的维数为 $1$，因为 $x$ 是 $k$-有理点（$k$ 代数闭）。
由于 $s$ 在 $x$ 处不消失，可知
$$
H^0(X, \mathcal{L}) \longrightarrow H^0(X, i_*i^*\mathcal{L})
$$
满射。因此
$\dim_k H^0(X, \mathcal{I} \otimes_{\mathcal{O}_X} \mathcal{L}) = t - 1$。
最后，同调长正合列还表明
$H^1(X, \mathcal{I} \otimes_{\mathcal{O}_X} \mathcal{L}) = 0$，
从而完成归纳步骤的证明。

\medskip\noindent
为得到 (2) 中的可逆层，取 (1) 中的可逆层
$\mathcal{L}$，再作一次上一段的论证即可。
\end{proof}

\begin{lemma}
\label{varieties-lemma-vanishing-degree-2g-and-1-line-bundle}
设 $k$ 为代数闭域。设 $X$ 是 $k$ 上的固有曲线。
置 $g = \dim_k H^1(X, \mathcal{O}_X)$。
对每个可逆 $\mathcal{O}_X$-模 $\mathcal{L}$，若
$\deg(\mathcal{L}) \geq 2g - 1$，则
$H^1(X, \mathcal{L}) = 0$。
\end{lemma}

\begin{proof}
令 $\mathcal{N}$ 为引理
\ref{varieties-lemma-general-degree-g-line-bundle} 第 (2) 部分中找到的可逆模。
$\mathcal{N}$ 的次数为
$\chi(X, \mathcal{N}) - \chi(X, \mathcal{O}_X) = 0 - (1 - g) = g - 1$。
因此 $\mathcal{L} \otimes \mathcal{N}^{\otimes - 1}$ 的次数为
$\deg(\mathcal{L}) - (g - 1) \geq g$。从而
$\chi(X, \mathcal{L} \otimes \mathcal{N}^{\otimes -1}) \geq g + 1 - g = 1$。
所以存在非零整体截面 $s$，其零点概形是有效 Cartier 除子
$D$，且次数为 $\deg(\mathcal{L}) - (g - 1)$。
这给出短正合列
$$
0 \to \mathcal{N} \xrightarrow{s} \mathcal{L} \to i_*(\mathcal{L}|_D) \to 0
$$
其中 $i : D \to X$ 是包含态射。可知
$H^0(X, \mathcal{L})$ 同构地映到
$H^0(D, \mathcal{L}|_D)$，后者的维数为
$\deg(\mathcal{L}) - (g - 1)$。结论由次数定义得出。
\end{proof}



\section{数值相交}
\label{varieties-section-num}

\noindent
本节利用固有概形上凝聚层的欧拉示性数作一些运算，从而得到
可逆模的数值相交数。我们的主要工具是下述引理。

\begin{lemma}
\label{varieties-lemma-numerical-polynomial-from-euler}
设 $k$ 为域。设 $X$ 是 $k$ 上的固有概形。设 $\mathcal{F}$
是凝聚 $\mathcal{O}_X$-模。设
$\mathcal{L}_1, \ldots, \mathcal{L}_r$ 是可逆 $\mathcal{O}_X$-模。
映射
$$
(n_1, \ldots, n_r) \longmapsto
\chi(X, \mathcal{F} \otimes
\mathcal{L}_1^{\otimes n_1} \otimes \ldots \otimes
\mathcal{L}_r^{\otimes n_r})
$$
是关于 $n_1, \ldots, n_r$ 的数值多项式，其总次数不超过
$\mathcal{F}$ 的支撑的维数。
\end{lemma}

\begin{proof}
我们对 $\dim(\text{Supp}(\mathcal{F}))$ 作归纳。若该数为零，
则由引理 \ref{varieties-lemma-chi-tensor-finite}，此函数为常值函数，
其值为 $\dim_k \Gamma(X, \mathcal{F})$。假设
$\dim(\text{Supp}(\mathcal{F})) > 0$。

\medskip\noindent
若 $\mathcal{F}$ 有嵌入伴随点，则可考虑《除子》引理
\ref{divisors-lemma-remove-embedded-points} 中构造的短正合列
$0 \to \mathcal{K} \to \mathcal{F} \to \mathcal{F}' \to 0$。
由于 $\mathcal{K}$ 的支撑维数严格较小，由归纳假设，结论对
$\mathcal{K}$ 成立，且所得总次数严格较小。
由欧拉示性数的可加性（引理
\ref{varieties-lemma-euler-characteristic-additive}），只需对
$\mathcal{F}'$ 证明结论。因此可以假设 $\mathcal{F}$ 没有嵌入伴随点。

\medskip\noindent
若 $i : Z \to X$ 是闭浸入，且 $\mathcal{F} = i_*\mathcal{G}$，
则关于 $X$、$\mathcal{F}$、
$\mathcal{L}_1, \ldots, \mathcal{L}_r$ 的结论等价于关于
$Z$、$\mathcal{G}$、
$i^*\mathcal{L}_1, \ldots, i^*\mathcal{L}_r$ 的结论
（因为二者的上同调相同；见《概形上同调》引理
\ref{coherent-lemma-relative-affine-cohomology}）。
应用《除子》引理 \ref{divisors-lemma-no-embedded-points-endos}，
可以假设 $X$ 没有嵌入分支，并且
$X = \text{Supp}(\mathcal{F})$。

\medskip\noindent
选取正则亚纯截面 $s$，它是 $\mathcal{L}_1$ 的截面；见《除子》引理
\ref{divisors-lemma-regular-meromorphic-section-exists}。
令 $\mathcal{I} \subset \mathcal{O}_X$ 为 $s$ 的分母理想，
并考虑《除子》引理
\ref{divisors-lemma-make-maps-regular-section} 中的映射
$$
\mathcal{I}\mathcal{F} \to \mathcal{F},\quad
\mathcal{I}\mathcal{F} \to \mathcal{F} \otimes \mathcal{L}_1
$$
它们都是单射，余核分别为 $\mathcal{Q}$、$\mathcal{Q}'$，
这些余核支撑在
$X = \text{Supp}(\mathcal{F})$ 的无处稠密闭子概形上。
与可逆模
$\mathcal{L}_1^{\otimes n_1} \otimes \ldots \otimes \mathcal{L}_r^{\otimes n_r}$
张量是正合的；因此再次使用可加性可知
\begin{align*}
&\chi(X, \mathcal{F} \otimes \mathcal{L}_1^{\otimes n_1} \otimes \ldots \otimes
\mathcal{L}_r^{\otimes n_r}) -
\chi(X, \mathcal{F} \otimes \mathcal{L}_1^{\otimes n_1 + 1}
\otimes \ldots \otimes \mathcal{L}_r^{\otimes n_r}) \\
& =
\chi(\mathcal{Q} \otimes \mathcal{L}_1^{\otimes n_1} \otimes \ldots \otimes
\mathcal{L}_r^{\otimes n_r}) -
\chi(\mathcal{Q}' \otimes \mathcal{L}_1^{\otimes n_1} \otimes \ldots \otimes
\mathcal{L}_r^{\otimes n_r})
\end{align*}
因此，引理中的函数 $P(n_1, \ldots, n_r)$ 具有如下性质：
$$
P(n_1 + 1, n_2, \ldots, n_r) - P(n_1, \ldots, n_r)
$$
是总次数 $<$ $\mathcal{F}$ 支撑维数的数值多项式。
当然，由对称性，对任意 $i \in \{1, \ldots, r\}$，同一结论也适用于
$$
P(n_1, \ldots, n_{i - 1}, n_i + 1, n_{i + 1}, \ldots, n_r)
- P(n_1, \ldots, n_r)
$$
一个简单的算术论证表明，$P$ 是总次数不超过
$\dim(\text{Supp}(\mathcal{F}))$ 的数值多项式。
\end{proof}

\noindent
下述引理粗略地说明，首项系数只依赖于凝聚模在其支撑泛点处的长度。

\begin{lemma}
\label{varieties-lemma-numerical-polynomial-leading-term}
设 $k$ 为域。设 $X$ 是 $k$ 上的固有概形。设
$\mathcal{F}$ 是凝聚 $\mathcal{O}_X$-模。设
$\mathcal{L}_1, \ldots, \mathcal{L}_r$ 是可逆 $\mathcal{O}_X$-模。
令 $d = \dim(\text{Supp}(\mathcal{F}))$。
设 $Z_i \subset X$ 是 $\text{Supp}(\mathcal{F})$ 的维数为 $d$
的不可约分支。令 $\xi_i \in Z_i$ 为泛点，并置
$m_i = \text{length}_{\mathcal{O}_{X, \xi_i}}(\mathcal{F}_{\xi_i})$。
则
$$
\chi(X, \mathcal{F} \otimes \mathcal{L}_1^{\otimes n_1} \otimes \ldots \otimes
\mathcal{L}_r^{\otimes n_r}) -
\sum\nolimits_i
m_i\ \chi(Z_i, \mathcal{L}_1^{\otimes n_1} \otimes \ldots \otimes
\mathcal{L}_r^{\otimes n_r}|_{Z_i})
$$
是关于 $n_1, \ldots, n_r$ 的数值多项式，其总次数 $< d$。
\end{lemma}

\begin{proof}
考虑二元组 $(\xi , Z)$，其中 $Z \subset X$ 是维数为 $d$
的整闭子概形，而 $\xi$ 是其泛点。有限
$\mathcal{O}_{X, \xi}$-模 $\mathcal{F}_\xi$ 的支撑包含在
$\{\xi\}$ 中；因此长度
$m_Z = \text{length}_{\mathcal{O}_{X, \xi}}(\mathcal{F}_\xi)$
有限（《代数》引理 \ref{algebra-lemma-support-point}），
并且除非 $Z = Z_i$ 对某个 $i$ 成立，否则它为零。
所以引理中的表达式可写成
$$
E(\mathcal{F}) =
\chi(X, \mathcal{F} \otimes \mathcal{L}_1^{\otimes n_1} \otimes \ldots \otimes
\mathcal{L}_r^{\otimes n_r}) -
\sum\nolimits
m_Z\ \chi(Z, \mathcal{L}_1^{\otimes n_1} \otimes \ldots \otimes
\mathcal{L}_r^{\otimes n_r}|_Z)
$$
其中求和遍历所有整闭子概形 $Z \subset X$，其维数为 $d$。
赋值 $\mathcal{F} \mapsto E(\mathcal{F})$ 在短正合列
$0 \to \mathcal{F} \to \mathcal{F}' \to \mathcal{F}'' \to 0$
中可加；这些列由凝聚 $\mathcal{O}_X$-模组成，其支撑维数
$\leq d$。
这由欧拉示性数的可加性（引理
\ref{varieties-lemma-euler-characteristic-additive}）与长度的可加性
（《代数》引理 \ref{algebra-lemma-length-additive}）得出。
应用《概形上同调》引理
\ref{coherent-lemma-coherent-filter}，取滤过
$$
0 = \mathcal{F}_0 \subset \mathcal{F}_1 \subset
\ldots \subset \mathcal{F}_m = \mathcal{F}
$$
其中各项均为凝聚子层，并且对每个 $j = 1, \ldots, m$，
都存在整闭子概形 $V_j \subset X$ 与非零理想层
$\mathcal{I}_j \subset \mathcal{O}_{V_j}$，使得
$$
\mathcal{F}_j/\mathcal{F}_{j - 1}
\cong (V_j \to X)_* \mathcal{I}_j
$$
于是 $V_j \subset \text{Supp}(\mathcal{F})$，故
$\dim(V_j) \leq d$。
由上述可加性，只须对每个次商
$\mathcal{F}_j/\mathcal{F}_{j - 1}$ 证明结论。
所以只须考虑
$\mathcal{F} = (V \to X)_*\mathcal{I}$，其中
$V \subset X$ 是维数 $\leq d$ 的整闭子概形，而
$\mathcal{I} \subset \mathcal{O}_V$ 是非零凝聚理想层。
若 $\dim(V) < d$，更一般地，若 $\mathcal{F}$ 的支撑维数
$< d$，则由引理 \ref{varieties-lemma-numerical-polynomial-from-euler}，
$E(\mathcal{F})$ 中第一项的总次数 $< d$，第二项为零。
若 $\dim(V) = d$，则可使用短正合列
$$
0 \to (V \to X)_*\mathcal{I} \to (V \to X)_*\mathcal{O}_V
\to (V \to X)_*(\mathcal{O}_V/\mathcal{I}) \to 0
$$
结论对中间的层成立，因为求和中唯一出现的 $Z$ 是 $Z = V$，
且 $m_Z = 1$；此外，由投影公式
（《同调》第 \ref{cohomology-section-projection-formula} 节）
与《概形上同调》引理
\ref{coherent-lemma-relative-affine-cohomology}，有
$$
H^i(X, ((V \to X)_*\mathcal{O}_V) \otimes 
 \mathcal{L}_1^{\otimes n_1} \otimes \ldots \otimes
\mathcal{L}_r^{\otimes n_r}) =
H^i(V,  \mathcal{L}_1^{\otimes n_1} \otimes \ldots \otimes
\mathcal{L}_r^{\otimes n_r}|_V)
$$
因此在此情形中，实际上 $E(\mathcal{F}) = 0$。
结论对右端的层也成立，因为其支撑维数 $< d$。
所以结论对左端的层成立，引理得证。
\end{proof}

\begin{definition}
\label{varieties-definition-intersection-number}
设 $k$ 为域。设 $X$ 是 $k$ 上的固有概形。设
$i : Z \to X$ 是维数为 $d$ 的闭子概形。设
$\mathcal{L}_1, \ldots, \mathcal{L}_d$ 是可逆
$\mathcal{O}_X$-模。定义{\it 相交数}
$(\mathcal{L}_1 \cdots \mathcal{L}_d \cdot Z)$
为下述数值多项式中 $n_1 \ldots n_d$ 的系数：
$$
\chi(X, i_*\mathcal{O}_Z \otimes \mathcal{L}_1^{\otimes n_1} \otimes
\ldots \otimes \mathcal{L}_d^{\otimes n_d}) =
\chi(Z, \mathcal{L}_1^{\otimes n_1} \otimes
\ldots \otimes \mathcal{L}_d^{\otimes n_d}|_Z)
$$
在特殊情形
$\mathcal{L}_1 = \ldots = \mathcal{L}_d = \mathcal{L}$ 下，
记作 $(\mathcal{L}^d \cdot Z)$。
\end{definition}

\noindent
定义中的显示等式由投影公式
（《同调》第 \ref{cohomology-section-projection-formula} 节）
与《概形上同调》引理
\ref{coherent-lemma-relative-affine-cohomology} 得出。
下面证明关于这些相交数的几个引理。

\begin{lemma}
\label{varieties-lemma-intersection-number-integer}
在定义 \ref{varieties-definition-intersection-number} 的情形中，相交数
$(\mathcal{L}_1 \cdots \mathcal{L}_d \cdot Z)$ 是整数。
\end{lemma}

\begin{proof}
任意次数为 $e$、关于 $n_1, \ldots, n_d$ 的数值多项式，
都可以唯一地写成 $\mathbf{Z}$-线性组合，所用函数为
${n_1 \choose k_1}{n_2 \choose k_2} \ldots {n_d \choose k_d}$，其中
$k_1 + \ldots + k_d \leq e$。取 $e = d$ 即可。
细节留作练习。
\end{proof}

\begin{lemma}
\label{varieties-lemma-intersection-number-additive}
在定义 \ref{varieties-definition-intersection-number} 的情形中，相交数
$(\mathcal{L}_1 \cdots \mathcal{L}_d \cdot Z)$ 可加：
若 $\mathcal{L}_i = \mathcal{L}_i' \otimes \mathcal{L}_i''$，则
$$
(\mathcal{L}_1 \cdots \mathcal{L}_i \cdots \mathcal{L}_d \cdot Z) =
(\mathcal{L}_1 \cdots \mathcal{L}_i' \cdots \mathcal{L}_d \cdot Z) +
(\mathcal{L}_1 \cdots \mathcal{L}_i'' \cdots \mathcal{L}_d \cdot Z)
$$
\end{lemma}

\begin{proof}
这是因为，由引理 \ref{varieties-lemma-numerical-polynomial-from-euler}，函数
\begingroup\small
$$
(n_1, \ldots, n_{i - 1}, n_i', n_i'', n_{i + 1}, \ldots, n_d)
\mapsto
\chi(Z, \mathcal{L}_1^{\otimes n_1} \otimes
\ldots \otimes (\mathcal{L}_i')^{\otimes n_i'} \otimes
(\mathcal{L}_i'')^{\otimes n_i''} \otimes \ldots \otimes
\mathcal{L}_d^{\otimes n_d}|_Z)
$$
\endgroup
是总次数至多为 $d$ 的数值多项式，变量个数为 $d + 1$。
\end{proof}


\begin{lemma}
\label{varieties-lemma-intersection-number-in-terms-of-components}
在定义 \ref{varieties-definition-intersection-number} 的情形中，
设 $Z_i \subset Z$ 是维数为 $d$ 的不可约分支。令
$m_i = \text{length}_{\mathcal{O}_{X, \xi_i}}(\mathcal{O}_{Z, \xi_i})$，
其中 $\xi_i \in Z_i$ 为泛点。则
$$
(\mathcal{L}_1 \cdots \mathcal{L}_d \cdot Z) =
\sum m_i(\mathcal{L}_1 \cdots \mathcal{L}_d \cdot Z_i)
$$
\end{lemma}

\begin{proof}
这立即由引理 \ref{varieties-lemma-numerical-polynomial-leading-term} 与各定义得出。
\end{proof}

\begin{lemma}
\label{varieties-lemma-intersection-number-and-pullback}
设 $k$ 为域。设 $f : Y \to X$ 是 $k$ 上固有概形之间的态射。
设 $Z \subset Y$ 是维数为 $d$ 的整闭子概形，并设
$\mathcal{L}_1, \ldots, \mathcal{L}_d$ 是可逆 $\mathcal{O}_X$-模。
则
$$
(f^*\mathcal{L}_1 \cdots f^*\mathcal{L}_d \cdot Z) =
\deg(f|_Z : Z \to f(Z)) (\mathcal{L}_1 \cdots \mathcal{L}_d \cdot f(Z))
$$
其中 $\deg(Z \to f(Z))$ 的定义见《态射》定义
\ref{morphisms-definition-degree}；也可将其取为 $0$，条件是
$\dim(f(Z)) < d$。
\end{lemma}

\begin{proof}
左端由下述函数中 $n_1 \ldots n_d$ 的系数计算：
$$
\chi(Y, \mathcal{O}_Z \otimes f^*\mathcal{L}_1^{\otimes n_1} \otimes
\ldots \otimes f^*\mathcal{L}_d^{\otimes n_d}) =
\sum (-1)^i
\chi(X, R^if_*\mathcal{O}_Z \otimes
\mathcal{L}_1^{\otimes n_1} \otimes \ldots \otimes
\mathcal{L}_d^{\otimes n_d})
$$
该等式由引理 \ref{varieties-lemma-euler-characteristic-morphism} 与投影公式
（《同调》引理 \ref{cohomology-lemma-projection-formula}）得出。
若 $f(Z)$ 的维数 $< d$，则由引理
\ref{varieties-lemma-numerical-polynomial-from-euler}，右端是总次数
$< d$ 的多项式，故结论成立。假设 $\dim(f(Z)) = d$。
令 $\xi \in f(Z)$ 为泛点。由维数理论（见引理
\ref{varieties-lemma-dimension-locally-algebraic} 与
\ref{varieties-lemma-dimension-fibres-locally-algebraic}），
$Z$ 的泛点是 $Z$ 中唯一映到 $\xi$ 的点。
于是 $f : Z \to f(Z)$ 在 $f(Z)$ 的某个非空开集上有限；见《态射》引理
\ref{morphisms-lemma-generically-finite}。
因此 $\deg(f : Z \to f(Z))$ 有定义；事实上，它等于
$f_*\mathcal{O}_Z$ 在 $\xi$ 处的茎作为
$\mathcal{O}_{X, \xi}$-模的长度。此外，
$R^if_*\mathcal{O}_X$ 在 $\xi$ 处的茎对 $i > 0$ 为零，
因为刚刚已经看到
$f|_Z$ 在 $\xi$ 的一个邻域上有限
（因而《概形上同调》引理
\ref{coherent-lemma-finite-pushforward-coherent} 给出该消失）。
所以各项
$\chi(X, R^if_*\mathcal{O}_Z \otimes
\mathcal{L}_1^{\otimes n_1} \otimes \ldots \otimes
\mathcal{L}_d^{\otimes n_d})$（其中 $i > 0$）的总次数 $< d$，并且由引理
\ref{varieties-lemma-numerical-polynomial-leading-term}，模掉一个总次数
$< d$ 的多项式后，有
$$
\chi(X, f_*\mathcal{O}_Z \otimes
\mathcal{L}_1^{\otimes n_1} \otimes \ldots \otimes
\mathcal{L}_d^{\otimes n_d})
=
\deg(f : Z \to f(Z)) \chi(f(Z),
\mathcal{L}_1^{\otimes n_1} \otimes \ldots \otimes
\mathcal{L}_d^{\otimes n_d}|_{f(Z)})
$$
所需结论随即得出。
\end{proof}

\begin{lemma}
\label{varieties-lemma-numerical-intersection-effective-Cartier-divisor}
设 $k$ 为域。设 $X$ 在 $k$ 上固有。设 $Z \subset X$
是维数为 $d$ 的闭子概形。设
$\mathcal{L}_1, \ldots, \mathcal{L}_d$ 是可逆
$\mathcal{O}_X$-模。假设存在有效 Cartier 除子
$D \subset Z$，使得 $\mathcal{L}_1|_Z \cong \mathcal{O}_Z(D)$。
则
$$
(\mathcal{L}_1 \cdots \mathcal{L}_d \cdot Z) =
(\mathcal{L}_2 \cdots \mathcal{L}_d \cdot D)
$$
\end{lemma}

\begin{proof}
可以把 $X$ 替换为 $Z$，并把 $\mathcal{L}_i$ 替换为
$\mathcal{L}_i|_Z$。因此可以假设
$X = Z$ 且 $\mathcal{L}_1 = \mathcal{O}_X(D)$。
此时 $\mathcal{L}_1^{-1}$ 是 $D$ 的理想层，并可考虑短正合列
$$
0 \to \mathcal{L}_1^{\otimes -1} \to \mathcal{O}_X \to \mathcal{O}_D \to 0
$$
置
$P(n_1, \ldots, n_d) =
\chi(X, \mathcal{L}_1^{\otimes n_1} \otimes \ldots \otimes
\mathcal{L}_d^{\otimes n_d})$
以及
$Q(n_1, \ldots, n_d) =
\chi(D, \mathcal{L}_1^{\otimes n_1} \otimes \ldots \otimes
\mathcal{L}_d^{\otimes n_d}|_D)$。
由可加性可得
$$
P(n_1, \ldots, n_d) - P(n_1 - 1, n_2, \ldots, n_d) =
Q(n_1, \ldots, n_d)
$$
由于 $P$ 的总次数至多为 $d$，可知 $n_1 \ldots n_d$
在 $P$ 中的系数等于 $n_2 \ldots n_d$ 在 $Q$ 中的系数。
\end{proof}

\begin{lemma}
\label{varieties-lemma-ample-positive}
设 $k$ 为域。设 $X$ 在 $k$ 上固有。设 $Z \subset X$
是维数为 $d$ 的闭子概形。若
$\mathcal{L}_1, \ldots, \mathcal{L}_d$ 都丰沛，则
$(\mathcal{L}_1 \cdots \mathcal{L}_d \cdot Z)$ 为正。
\end{lemma}

\begin{proof}
我们对 $d$ 作归纳。$d = 0$ 的情形由引理
\ref{varieties-lemma-chi-tensor-finite} 得出。假设 $d > 0$。
由引理 \ref{varieties-lemma-intersection-number-in-terms-of-components}，
可以假设 $Z$ 是整闭子概形。事实上，可以把 $X$ 替换为 $Z$，
并把 $\mathcal{L}_i$ 替换为 $\mathcal{L}_i|_Z$，
从而约化到 $Z = X$ 是维数为 $d$ 的固有簇的情形。
由引理 \ref{varieties-lemma-intersection-number-additive}，可以把
$\mathcal{L}_1$ 替换为它的某个正张量幂。
所以可以假设存在非零截面
$s \in \Gamma(X, \mathcal{L}_1)$，使得
$X_s$ 仿射（这里使用丰沛可逆层的定义；见《性质》定义
\ref{properties-definition-ample}）。
注意，$X$ 不仿射，因为固有且仿射蕴含有限
（《态射》引理 \ref{morphisms-lemma-finite-proper}），
这与 $d > 0$ 矛盾。因此 $s$ 有非空消失概形
$Z(s) \subset X$。由于 $X$ 是簇，$s$ 是
$\mathcal{L}_1$ 的正则截面，所以 $Z(s)$ 是有效 Cartier 除子；
因而 $Z(s)$ 的余维为 $1$（这是在 $X$ 中而言），所以 $Z(s)$ 的维数为
$d - 1$（这里使用《除子》第
\ref{divisors-section-effective-Cartier-divisors}、
\ref{divisors-section-effective-Cartier-invertible}、
\ref{divisors-section-Noetherian-effective-Cartier} 节中的材料，
以及引理 \ref{varieties-lemma-dimension-locally-algebraic} 所用的维数理论）。
由引理 \ref{varieties-lemma-numerical-intersection-effective-Cartier-divisor}，有
$$
(\mathcal{L}_1 \cdots \mathcal{L}_d \cdot X) =
(\mathcal{L}_2 \cdots \mathcal{L}_d \cdot Z(s))
$$
由归纳假设，右端为正，证明完成。
\end{proof}

\begin{definition}
\label{varieties-definition-degree}
设 $k$ 为域。设 $X$ 是 $k$ 上的固有概形。设
$\mathcal{L}$ 是丰沛可逆 $\mathcal{O}_X$-模。
对任意闭子概形 $Z$，其关于 $\mathcal{L}$ 的{\it 次数}记作
$\deg_\mathcal{L}(Z)$，并定义为相交数
$(\mathcal{L}^d \cdot Z)$，其中 $d = \dim(Z)$。
\end{definition}

\noindent
由引理 \ref{varieties-lemma-ample-positive}，子概形的次数总是正整数。
注意，$\deg_\mathcal{L}(Z) = d$ 当且仅当
$$
\chi(Z, \mathcal{L}^{\otimes n}|_Z) = \frac{d}{\dim(Z)!} n^{\dim(Z)} + l.o.t
$$
这可由下式看出：
$$
(n_1 + \ldots + n_{\dim(Z)})^{\dim(Z)} =
\dim(Z)!\ n_1 \ldots n_{\dim(Z)} + \text{other terms}
$$

\begin{lemma}
\label{varieties-lemma-degree-finite-morphism-in-terms-degrees}
设 $k$ 为域。设 $f : Y \to X$ 是 $k$ 上固有簇之间的有限支配态射。
设 $\mathcal{L}$ 是丰沛可逆 $\mathcal{O}_X$-模。则
$$
\deg_{f^*\mathcal{L}}(Y) = \deg(f) \deg_\mathcal{L}(X)
$$
其中 $\deg(f)$ 的定义见《态射》定义
\ref{morphisms-definition-degree}。
\end{lemma}

\begin{proof}
该陈述有意义，因为由《态射》引理
\ref{morphisms-lemma-pullback-ample-tensor-relatively-ample}，
$f^*\mathcal{L}$ 丰沛。说明这一点后，结论就是引理
\ref{varieties-lemma-intersection-number-and-pullback} 的特殊情形。
\end{proof}

\noindent
最后，我们把与曲线的相交数同第
\ref{varieties-section-divisors-curves} 节中引入的曲线上可逆模次数联系起来。

\begin{lemma}
\label{varieties-lemma-intersection-numbers-and-degrees-on-curves}
设 $k$ 为域。设 $X$ 是 $k$ 上的固有概形。
设 $Z \subset X$ 是维数 $\leq 1$ 的闭子概形。
设 $\mathcal{L}$ 是可逆 $\mathcal{O}_X$-模。则
$$
(\mathcal{L} \cdot Z) = \deg(\mathcal{L}|_Z)
$$
其中 $\deg(\mathcal{L}|_Z)$ 的定义见定义
\ref{varieties-definition-degree-invertible-sheaf}。若 $\mathcal{L}$ 丰沛，则
$\deg_\mathcal{L}(Z) = \deg(\mathcal{L}|_Z)$。
\end{lemma}

\begin{proof}
这是因为函数
$n \mapsto \chi(Z, \mathcal{L}|_Z^{\otimes n})$ 的次数为 $1$，
故其首项系数等于相邻两值之差。
\end{proof}

\begin{proposition}[渐近 Riemann--Roch]
\label{varieties-proposition-asymptotic-riemann-roch}
设 $k$ 为域。设 $X$ 是 $k$ 上维数为 $d$ 的固有概形。
设 $\mathcal{L}$ 是丰沛可逆 $\mathcal{O}_X$-模。则
$$
\dim_k \Gamma(X, \mathcal{L}^{\otimes n}) \sim c n^d + l.o.t.
$$
其中 $c = \deg_\mathcal{L}(X)/d!$ 是正的常数。
\end{proposition}

\begin{proof}
这由各定义、引理 \ref{varieties-lemma-ample-positive}，以及《概形上同调》引理
\ref{coherent-lemma-vanshing-gives-ample} 中的高次上同调消失得出。
\end{proof}

\section{嵌入维数}
\label{varieties-section-embedding-dimension}

\noindent
嵌入维数有若干种定义方式，但对于代数闭域上的代数概形的闭点，
所有定义都与下述定义等价。

\begin{definition}
\label{varieties-definition-embed-dim}
设 $k$ 是代数闭域。设 $X$ 是局部代数 $k$-概形，并设 $x \in X$ 是闭点。
定义{\it $X$ 在 $x$ 处的嵌入维数}为
$\dim_k \mathfrak m_x/\mathfrak m_x^2$。
\end{definition}

\noindent
下面列出关于嵌入维数的一些事实。
设 $k, X, x$ 如定义 \ref{varieties-definition-embed-dim} 所述。
\begin{enumerate}
\item $X$ 在 $x$ 处的嵌入维数，等于切空间
$T_{X/k, x}$（定义 \ref{varieties-definition-tangent-space}）作为 $k$-向量空间的维数。
\item $X$ 在 $x$ 处的嵌入维数，是使得存在如下满射的最小整数
$d \geq 0$：
$$
k[[x_1, \ldots, x_d]] \longrightarrow \mathcal{O}_{X, x}^\wedge
$$
这里要求它是 $k$-代数满射。
\item $X$ 在 $x$ 处的嵌入维数，是最小整数 $d \geq 0$，使得存在开邻域
$U \subset X$（它是 $x$ 的邻域）以及闭浸入 $U \to Y$，其中 $Y$ 是维数为 $d$
的 $k$ 上光滑簇。
\item $X$ 在 $x$ 处的嵌入维数，是最小整数 $d \geq 0$，使得存在开邻域
$U \subset X$（它是 $x$ 的邻域）以及非分歧态射 $U \to \mathbf{A}^d_k$。
\item 若给定闭嵌入 $X \to Y$，且 $Y$ 在 $k$ 上光滑，则 $X$ 在 $x$ 处的嵌入维数
是满足下述条件的最小整数 $d \geq 0$：存在闭子概形 $Z \subset Y$，使得
$X \subset Z$，$Z \to \Spec(k)$ 在 $x$ 处光滑，并且 $\dim_x(Z) = d$。
\end{enumerate}
若以后需要这些事实，我们会陈述精确结果并给出证明。

\medskip\noindent
非代数闭的基域或非闭点。
设 $k$ 是域，且 $X$ 是局部代数 $k$-概形。
若 $x \in X$ 是一点，则定义 $X$ 在 $x$ 处的嵌入维数有若干种选择。具体说，可以采用
\begin{enumerate}
\item $\dim_{\kappa(x)}(\mathfrak m_x/\mathfrak m_x^2)$，
\item $\dim_{\kappa(x)}(T_{X/k, x}) =
\dim_{\kappa(x)}(\Omega_{X/k, x} \otimes_{\mathcal{O}_{X, x}} \kappa(x))$
（引理 \ref{varieties-lemma-tangent-space-cotangent-space}），
\item 最小整数 $d \geq 0$，使得存在开邻域 $U \subset X$（它是 $x$ 的邻域）
以及闭浸入 $U \to Y$，其中 $Y$ 是维数为 $d$ 的 $k$ 上光滑簇。
\end{enumerate}
在特征零的情形，(1) $=$ (2) 在 $x$ 是闭点时成立；
更一般地，若 $\kappa(x)$ 是 $k$ 上的可分代数扩张，这一等式也成立，
见引理 \ref{varieties-lemma-tangent-space-rational-point}。
几何定义 (3) 似乎最贴近“嵌入维数”这一措辞所唤起的几何直观。
由于可以证明 (3) 与 (2) 定义同一个数
（这由引理 \ref{varieties-lemma-immersion-into-affine} 得出），
我们将采用这个定义。在术语中，我们会明确说明所取的是相对于基域的嵌入维数。

\begin{definition}
\label{varieties-definition-embedding-dimension}
设 $k$ 是域。设 $X$ 是局部代数 $k$-概形。
设 $x \in X$ 是一点。定义{\it $X/k$ 在 $x$ 处的嵌入维数}为
$\dim_{\kappa(x)}(T_{X/k, x})$。
\end{definition}

\noindent
若 $(A, \mathfrak m, \kappa)$ 是 Noether 局部环，
有时将 $A$ 的{\it 嵌入维数}定义为 $\mathfrak m/\mathfrak m^2$ 在 $\kappa$ 上的维数。
上文已经看到，若 $A$ 具有域 $k$ 上的代数结构，采用
$\dim_\kappa(\Omega_{A/k} \otimes_A \kappa)$ 可能更合适。
我们将此量称为 $A/k$ 的{\it 嵌入维数}。采用这一术语，有
$$
\text{embed dim of }X/k\text{ 在 }x =
\text{embed dim of }\mathcal{O}_{X, x}/k =
\text{embed dim of }\mathcal{O}_{X, x}^\wedge/k
$$
其中 $k, X, x$ 如定义 \ref{varieties-definition-embedding-dimension} 所述。

\section{Bertini 定理}
\label{varieties-section-bertini}

\noindent
本节证明如下类型的结果：给定一个光滑射影簇 $X$，其基域为 $k$，
存在一个光滑的丰沛除子 $H \subset X$。

\begin{lemma}
\label{varieties-lemma-generate-over-complement}
设 $k$ 是域。设 $X$ 是 $k$ 上的本征概形。设 $\mathcal{L}$ 是丰沛可逆
$\mathcal{O}_X$-模。设 $Z \subset X$ 是闭子概形。则存在整数 $n_0$，使得对所有
$n \geq n_0$，令 $V_n$ 表示映射
$\Gamma(X, \mathcal{L}^{\otimes n}) \to \Gamma(Z, \mathcal{L}^{\otimes n}|_Z)$
的核；它生成 $\mathcal{L}^{\otimes n}|_{X \setminus Z}$，并且典范态射
$$
X \setminus Z \longrightarrow \mathbf{P}(V_n)
$$
是 $k$ 上概形的浸入。
\end{lemma}

\begin{proof}
设 $\mathcal{I} \subset \mathcal{O}_X$ 是 $Z$ 的拟凝聚理想层。注意，通过包含关系
$\mathcal{I} \otimes_{\mathcal{O}_X} \mathcal{L}^{\otimes n} \subset
\mathcal{L}^{\otimes n}$，有
$V_n = \Gamma(X, \mathcal{I} \otimes_{\mathcal{O}_X} \mathcal{L}^{\otimes n})$。
选取 $n_1$，使得对 $n \geq n_1$，层
$\mathcal{I} \otimes \mathcal{L}^{\otimes n}$ 是全局生成的，见
《性质》，命题 \ref{properties-proposition-characterize-ample}。
于是 $V_n$ 生成 $\mathcal{L}^{\otimes n}|_{X \setminus Z}$，只要 $n \geq n_1$。

\medskip\noindent
对 $n \geq n_1$，将限制映射记为
$\psi_n : V_n \to
\Gamma(X \setminus Z, \mathcal{L}^{\otimes n}|_{X \setminus Z})$。
由《构造》，例 \ref{constructions-example-projective-space}，得到典范态射
$$
\varphi = \varphi_{\mathcal{L}^{\otimes n}|_{X \setminus Z}, \psi_n} :
X \setminus Z
\longrightarrow
\mathbf{P}(V_n)
$$
。选取 $n_2$，使得对所有 $n \geq n_2$，可逆层
$\mathcal{L}^{\otimes n}$ 在 $X$ 上很丰沛。
我们断言 $n_0 = n_1 + n_2$ 可用。

\medskip\noindent
断言的证明。设 $n \geq n_0$，并写成 $n = n_1 + n'$。
对 $x \in X \setminus Z$，可选取 $s_1 \in V_1$，使它在 $x$ 处不消失。
置 $V' = \Gamma(X, \mathcal{L}^{\otimes n'})$。
由对 $n$ 与 $n'$ 的选取可知，相应态射
$\varphi' : X \to \mathbf{P}(V')$ 是闭浸入。因此，若选取
$s' \in \Gamma(X, \mathcal{L}^{\otimes n'})$ 在 $x$ 处不消失，则
$X_{s'} = (\varphi')^{-1}(D_+(s'))$（见《构造》，
引理 \ref{constructions-lemma-invertible-map-into-proj}）是仿射的，
且 $X_{s'} \to D_+(s')$ 是闭浸入。
于是 $s = s_1 \otimes s' \in V_n$ 在 $x$ 处不消失。
若 $D_+(s) \subset \mathbf{P}(V_n)$ 表示该射影空间中相应的开仿射空间，则
$\varphi^{-1}(D_+(s)) = X_s \subset X \setminus Z$（见上引文）。
开集 $X_s = X_{s'} \cap X_{s_1}$ 是仿射的，见
《性质》，引理 \ref{properties-lemma-affine-cap-s-open}。
考虑定义态射 $X_s \to D_+(s)$ 的环映射
$$
\text{Sym}(V)_{(s)} \longrightarrow \mathcal{O}_X(X_s)
$$
。因为 $X_{s'} \to D_+(s')$ 是闭浸入，以下元素的像
$$
\frac{s_1 \otimes t'}{s_1 \otimes s'}
$$
（其中 $t' \in V'$）生成映射
$\mathcal{O}_X(X_{s'}) \to \mathcal{O}_X(X_s)$ 的像。
由于 $X_s \to X_{s'}$ 是开浸入，这说明 $X_s \to D_+(s)$
是仿射概形的浸入（见下文）。从而由《态射》，
引理 \ref{morphisms-lemma-check-immersion}，$\varphi_n$ 是浸入。

\medskip\noindent
设 $a : A' \to A$ 与 $c : B \to A$ 是环映射，且
$\Spec(a)$ 是浸入并满足 $\Im(a) \subset \Im(c)$。
置 $B' = A' \times_A B$，其投影为 $b : B' \to B$ 与
$c' : B' \to A'$。由假设，$c'$ 是满射，因而
$\Spec(c')$ 是闭浸入。于是 $\Spec(c') \circ \Spec(a)$ 是浸入
（《概形》，引理 \ref{schemes-lemma-composition-immersion}）。
由于 $\Spec(c)$ 分解浸入
$\Spec(c') \circ \Spec(a) = \Spec(b) \circ \Spec(c)$，
故它必为浸入，见《态射》，引理 \ref{morphisms-lemma-immersion-permanence}。
\end{proof}

\begin{situation}
\label{varieties-situation-family-divisors}
设 $k$ 是域，
设 $X$ 是 $k$ 上的概形，
设 $\mathcal{L}$ 是可逆 $\mathcal{O}_X$-模，
设 $V$ 是有限维 $k$-向量空间，并且
设 $\psi : V \to \Gamma(X, \mathcal{L})$ 是 $k$-线性映射。
设 $\dim(V) = r$，且 $v_1, \ldots, v_r$ 是 $V$ 的一组基。
于是得到“泛除子”
$$
H_{univ} = Z(s_{univ}) \subset  \mathbf{A}^r \times_k X
$$
，它是截面
$$
s_{univ} = \sum\nolimits_{i = 1, \ldots, r} x_i \psi(v_i) \in
\Gamma(\mathbf{A}^r \times_k X, \text{pr}_2^*\mathcal{L})
$$
的零概形（《除子》，定义 \ref{divisors-definition-zero-scheme-s}）。
对域扩张 $k'/k$，$k'$-点 $v \in \mathbf{A}^r_k(k')$
对应于向量 $(a_1, \ldots, a_r)$，其分量均属于 $k'$。
因此，一方面可以把 $v$ 看作元素
$v = \sum_{i = 1, \ldots, r} a_i v_i \in V \otimes_k k'$，
另一方面可以向 $v$ 指派截面
$$
\psi(v) = \sum\nolimits_{i = 1, \ldots, r} a_i \psi(v_i) \in
\Gamma(X_{k'}, \mathcal{L}|_{X_{k'}})
$$
。按此记号，$H_{univ}$ 在 $v \in V \otimes k'$ 上的纤维显然是
$\psi(v)$ 的零概形。写成公式即
$$
H_v = H_{univ, v} = Z(\psi(v))
$$
。如式中所示，将这个公共值记作 $H_v$。最后，在此情形中，令 $P$ 是向量
$v \in V \otimes_k k'$ 的一个性质，其中 $k'/k$ 是任意域扩张\footnote{例如，
可以考虑 $H_v$ 在 $k'$ 上光滑，或在 $k'$ 上几何不可约这一条件。}。
称 $P$ 对{\it 一般的} $v \in V \otimes_k k'$ {\it 成立}，如果存在非空 Zariski 开集
$U \subset \mathbf{A}^r_k$，使得只要 $v$ 对应于一个 $k'$-点且该点属于 $U$，
则对任意 $k'/k$ 都有 $P(v)$ 成立。
\end{situation}

\begin{lemma}
\label{varieties-lemma-bertini}
在情形 \ref{varieties-situation-family-divisors} 中，假设
\begin{enumerate}
\item $X$ 在 $k$ 上光滑，
\item $\psi : V \to \Gamma(X, \mathcal{L})$ 的像生成 $\mathcal{L}$，
\item 相应态射
$\varphi_{\mathcal{L}, \psi} : X \to \mathbf{P}(V)$ 是浸入。
\end{enumerate}
则对一般的 $v \in V \otimes_k k'$，概形 $H_v$ 在 $k'$ 上光滑。
\end{lemma}

\begin{proof}
（注意，$X$ 作为射影空间的局部闭子概形，是分离且有限型的。）
采用引理陈述上方引入的记号。考虑投影
$$
\xymatrix{
\mathbf{A}^r_k \times_k X \ar[d] &
H_{univ} \ar[l] \ar[ld]^p \ar[r] \ar[rd]_q &
\mathbf{A}^r_k \times_k X \ar[d] \\
X & &
\mathbf{A}^r_k
}
$$
。设 $\Sigma \subset H_{univ}$ 是态射
$q : H_{univ} \to \mathbf{A}^r_k$ 的奇异轨迹，即 $q$ 不光滑的点集。
由于态射的光滑轨迹按定义是开的，$\Sigma$ 是闭的。
光滑态射的纤维是光滑的，故只需证明 $q(\Sigma)$ 包含在
$\mathbf{A}^r_k$ 的真闭子集中。
因为 $\Sigma$（赋予约化诱导概形结构）是 $k$ 上有限型概形，
只需证明 $\dim(\Sigma) < r$。这由引理
\ref{varieties-lemma-dimension-fibres-locally-algebraic} 得出。
以更大的域代替 $k$ 不改变维数
（《态射》，引理 \ref{morphisms-lemma-dimension-fibre-after-base-change}），
因此不妨假设 $k$ 是代数闭的。
由维数理论（引理 \ref{varieties-lemma-dimension-fibres-locally-algebraic}），
只需证明：对闭点 $x \in X \setminus Z$，有
$p^{-1}(\{x\}) \cap \Sigma$ 的维数 $< r - \dim(X')$，
其中 $X'$ 是 $X$ 中包含 $x$ 的唯一不可约分支。
由于 $X$ 在 $k$ 上光滑且 $x$ 是闭点，有
$\dim(X') = \dim \mathfrak m_x/\mathfrak m_x^2$
（《态射》，引理 \ref{morphisms-lemma-smooth-omega-finite-locally-free}
及《代数》，引理 \ref{algebra-lemma-rank-omega}）。
因此，只要对所有闭点 $x \in X$ 都有
$$
\dim p^{-1}(x) \cap \Sigma < r - \dim \mathfrak m_x/\mathfrak m_x^2
$$
，结论即成立。

\medskip\noindent
由于 $V$ 全局生成 $\mathcal{L}$，对每个不可约分支
$X'$（它属于 $X$），$\mathbf{A}^r$ 都有一个非空 Zariski 开集，使得 $q$ 在该开集上的纤维不包含
$X'$。（例如，若 $x' \in X'$ 是闭点，则可取对应于下述向量的开集：
$v \in V$，且 $\psi(v)$ 在 $x'$ 处不消失。这个开集是
$\mathbf{A}^r_k$ 中一个超平面的补集。）
令 $U \subset \mathbf{A}^r$ 为这些开集的（非空）交。
则 $q^{-1}(U) \to U$ 的纤维是 $U \times_k X \to U$ 的纤维上的有效 Cartier 除子
（因为整概形上可逆模的非零截面是正则截面）。因此由《除子》，
引理 \ref{divisors-lemma-fibre-Cartier}，态射 $q^{-1}(U) \to U$ 是平坦的。
从而，对闭点 $x \in X$ 与 $v \in V = \mathbf{A}^r_k(k)$，若
$(x, v) \in H_{univ}$，即 $x \in H_v$，则 $q$ 在 $(x, v)$ 处光滑当且仅当纤维
$H_v$ 在 $x$ 处光滑，见《态射》，引理 \ref{morphisms-lemma-smooth-at-point}。

\medskip\noindent
考虑像 $\psi(v)_x$；它位于茎 $\mathcal{L}_x$ 中，来自与 $v \in V$ 对应的截面。有
$$
x \in H_v \Leftrightarrow \psi(v)_x \in \mathfrak m_x\mathcal{L}_x
$$
。若此条件成立，则有
$$
H_v\text{ singular at }x \Leftrightarrow
\psi(v)_x \in \mathfrak m_x^2\mathcal{L}_x
$$
。事实上，$\psi(v)_x$ 不属于 $\mathfrak m_x^2\mathcal{L}_x$
$\Leftrightarrow$
$H_v \subset X$ 在 $x$ 处的局部方程不属于 $\mathfrak m_x^2$
$\Leftrightarrow$
$\mathcal{O}_{H_v, x}$ 是正则的
（《代数》，引理 \ref{algebra-lemma-regular-ring-CM}）
$\Leftrightarrow$
$H_v$ 在点 $x$ 处相对于 $k$ 光滑
（《代数》，引理 \ref{algebra-lemma-separable-smooth}）。
由此，$p^{-1}(x) \cap \Sigma$ 的闭点对应于这些 $v \in V$：它们满足
$\psi(v)_x \in \mathfrak m_x^2\mathcal{L}_x$。
然而，由于 $\varphi_{\mathcal{L}, \psi}$ 是浸入，映射
$$
V \longrightarrow \mathcal{L}_x/\mathfrak m_x^2\mathcal{L}_x
$$
是满射（略去一个小细节）。由上可知，将轨迹 $p^{-1}(x) \cap \Sigma$
视为 $V$ 的子空间时，其闭点构成此映射的核，因而其维数为
$r - \dim \mathfrak m_x/\mathfrak m_x^2 - 1$，如所欲证。
\end{proof}


\section{Enriques--Severi--Zariski 定理}
\label{varieties-section-vanishing-negative}

\noindent
本节证明若干如下形式的结果：以一个“非常负”的可逆模扭转会使低次上同调消失。
我们还推出，维数 $\geq 2$ 的正规本征概形的超曲面截面是连通的。

\begin{lemma}
\label{varieties-lemma-vanishin-h0-negative}
设 $k$ 是域。设 $X$ 是 $k$ 上的本征概形。设 $\mathcal{L}$ 是丰沛可逆
$\mathcal{O}_X$-模。设 $\mathcal{F}$ 是凝聚 $\mathcal{O}_X$-模。
若 $\text{Ass}(\mathcal{F})$ 不含任何闭点，则
$\Gamma(X, \mathcal{F} \otimes_{\mathcal{O}_X} \mathcal{L}^{\otimes n}) = 0$，
只要 $n \ll 0$。
\end{lemma}

\begin{proof}
对凝聚 $\mathcal{O}_X$-模 $\mathcal{F}$，令 $\mathcal{P}(\mathcal{F})$
表示如下性质：存在 $n_0 \in \mathbf{Z}$，使得对 $n \leq n_0$，
每个截面 $s$（它是
$\mathcal{F} \otimes_{\mathcal{O}_X} \mathcal{L}^{\otimes n}$ 的截面）
的支撑都只由闭点组成。
由于 $\text{Ass}(\mathcal{F}) =
\text{Ass}(\mathcal{F} \otimes_{\mathcal{O}_X} \mathcal{L}^{\otimes n})$，
只需证明 $\mathcal{P}$ 对 $X$ 上所有凝聚模都成立。
为此，证明《概形的上同调》，引理
\ref{coherent-lemma-property-higher-rank-cohomological}
中的条件 (1)、(2) 和 (3) 均成立。

\medskip\noindent
为验证条件 (1)，设
$$
0 \to \mathcal{F}_1 \to \mathcal{F} \to \mathcal{F}_2 \to 0
$$
是凝聚 $\mathcal{O}_X$-模的短正合列，并且已有 $\mathcal{P}$
对 $\mathcal{F}_i$ 成立，其中 $i = 1, 2$。
令 $n_1, n_2$ 为所得界限。
令 $\mathcal{F}'_2 \subset \mathcal{F}_2$ 为支撑是有限闭点集的最大凝聚子模。
令 $\mathcal{I} \subset \mathcal{O}_X$ 为 $\mathcal{F}'_2$ 的零化子。
由于 $\mathcal{L}$ 丰沛，可找到 $e > 0$，使
$\mathcal{I} \otimes_{\mathcal{O}_X} \mathcal{L}^{\otimes e}$
是全局生成的。置 $n_0 = \min(n_2, n_1 - e)$。
取 $n \leq n_0$，并令 $t$ 是
$\mathcal{F} \otimes \mathcal{L}^{\otimes n}$ 的一个全局截面。
截面 $t$ 在
$\mathcal{F}_2 \otimes \mathcal{L}^{\otimes n}$ 中的像落在
$\mathcal{F}'_2 \otimes \mathcal{L}^{\otimes n}$ 中，这是因为 $n \leq n_2$。
因此，对任意
$s \in \Gamma(X, \mathcal{I} \otimes_{\mathcal{O}_X} \mathcal{L}^{\otimes e})$，
乘积 $t \otimes s$ 属于
$\mathcal{F}_1 \otimes \mathcal{L}^{\otimes n + e}$。
$t \otimes s$ 的支撑包含在
$\text{Ass}(\mathcal{F}_1)$ 中的有限闭点集中，因为 $n + e \leq n_1$。
依照对 $e$ 的选取，可选取 $s$ 在任何不属于
$\mathcal{F}'_2$ 支撑的点处可逆；故 $t$ 的支撑包含在
$\text{Ass}(\mathcal{F}_1)$ 中的有限闭点集与
$\text{Ass}(\mathcal{F}_2)$ 中的有限闭点集之并中。
这就完成了条件 (1) 的证明。

\medskip\noindent
条件 (2) 是立即的。

\medskip\noindent
对条件 (3)，取 $\mathcal{G} = \mathcal{O}_Z$。
在这种情形下，若 $Z$ 是 $X$ 的闭点，则无须证明。
若 $\dim(Z) > 0$，则证明
$\Gamma(Z, \mathcal{L}^{\otimes n}|_Z) = 0$，只要 $n < 0$。
事实上，令 $s$ 是 $\mathcal{L}|_Z$ 的某个负次幂的非零截面。
选取非零截面 $t$，它属于 $\mathcal{L}|_Z$ 的某个正次幂
（这是可能的，因为 $\mathcal{L}$ 丰沛，见《性质》，命题
\ref{properties-proposition-characterize-ample}）。
则
$s^{\deg(t)} \otimes t^{\deg(s)}$
是 $\mathcal{O}_Z$ 的非零全局截面（因为 $Z$ 是整的），
从而是单位（引理 \ref{varieties-lemma-proper-geometrically-reduced-global-sections}）。
这说明 $t$ 是 $\mathcal{L}$ 的某个正次幂的平凡化截面。
于是函数 $n \mapsto \dim_k \Gamma(X, \mathcal{L}^{\otimes n})$
在一个无穷正整数集上有界；由于 $\dim(Z) > 0$，
这与渐近 Riemann--Roch（命题
\ref{varieties-proposition-asymptotic-riemann-roch}）矛盾。
\end{proof}

\begin{lemma}[Enriques--Severi--Zariski]
\label{varieties-lemma-vanishin-h1-negative}
设 $k$ 是域。设 $X$ 是 $k$ 上的本征概形。设 $\mathcal{L}$ 是丰沛可逆
$\mathcal{O}_X$-模。设 $\mathcal{F}$ 是凝聚 $\mathcal{O}_X$-模。
假设对闭点 $x \in X$ 有 $\text{depth}(\mathcal{F}_x) \geq 2$。
则
$H^1(X, \mathcal{F} \otimes_{\mathcal{O}_X} \mathcal{L}^{\otimes m}) = 0$，
只要 $m \ll 0$。
\end{lemma}

\begin{proof}
选取闭浸入 $i : X \to \mathbf{P}^n_k$，使得
$i^*\mathcal{O}(1) \cong \mathcal{L}^{\otimes e}$，其中 $e > 0$
（见《态射》，引理
\ref{morphisms-lemma-finite-type-over-affine-ample-very-ample}）。
只需在 $\mathbf{P}^n_k$ 上对
$$
\mathcal{G} =
i_*(\mathcal{F} \oplus \mathcal{F} \otimes \mathcal{L} \oplus \ldots
\oplus \mathcal{F} \otimes \mathcal{L}^{\otimes e - 1})
\quad\text{且}\quad \mathcal{O}(1)
$$
证明引理。事实上，由《概形的上同调》，引理
\ref{coherent-lemma-relative-affine-cohomology}，有
$$
H^1(\mathbf{P}^n_k, \mathcal{G}(m)) =
\bigoplus\nolimits_{j = 0, \ldots, e - 1}
H^1(X, \mathcal{F} \otimes \mathcal{L}^{\otimes j + me})
$$
。另外，若 $y \in \mathbf{P}^n_k$ 是闭点，则
$\text{depth}(\mathcal{G}_y) = \infty$，如果 $y \not \in i(X)$；
并且 $\text{depth}(\mathcal{G}_y) = \text{depth}(\mathcal{F}_x)$，
如果 $y = i(x)$。这是因为在后一种情形下，有
$\mathcal{G}_y \cong \mathcal{F}_x^{\oplus e}$，这是作为
$\mathcal{O}_{\mathbf{P}^n_k, x}$-模而言；并可用例如
《代数》，引理 \ref{algebra-lemma-depth-goes-down-finite}
得到该等式。

\medskip\noindent
假设 $X = \mathbf{P}^n_k$、$\mathcal{L} = \mathcal{O}(1)$，且
$k$ 是无限域。选取 $s \in H^0(\mathbf{P}^1_k, \mathcal{O}(1))$，
它确定正合列
$$
0 \to \mathcal{F}(-1) \xrightarrow{s} \mathcal{F} \to \mathcal{G} \to 0
$$
，如引理 \ref{varieties-lemma-exact-sequence-induction} 中所述。
由于映射 $\mathcal{F}(-1) \to \mathcal{F}$ 在仿射局部由
$\mathcal{F}$ 上的非零因子相乘给出，所以对闭点
$x \in \mathbf{P}^n_k$ 有 $\text{depth}(\mathcal{G}_x) \geq 1$，
见《代数》，引理 \ref{algebra-lemma-depth-drops-by-one}。
因此由引理 \ref{varieties-lemma-vanishin-h0-negative}，
$H^0(\mathcal{G}(m)) = 0$，只要 $m \ll 0$。
考察扭转后的上同调长正合列
（见注 \ref{varieties-remark-exact-sequence-induction-cohomology}），可知数列
$$
\dim H^1(\mathbf{P}^n_k, \mathcal{F}(m))
$$
在 $m \leq m_0$ 时稳定，其中 $m_0$ 是某个整数。
令 $N$ 为这些空间在 $m \leq m_0$ 时的公共维数。需要证明 $N = 0$。

\medskip\noindent
对 $d > 0$ 与 $m \leq m_0$，考虑双线性映射
$$
H^0(\mathbf{P}^n_k, \mathcal{O}(d)) \times
H^1(\mathbf{P}^n_k, \mathcal{F}(m - d))
\longrightarrow
H^1(\mathbf{P}^n_k, \mathcal{F}(m))
$$
。由线性代数，存在余维 $\leq N^2$ 的子空间
$V_m \subset H^0(\mathbf{P}^n_k, \mathcal{O}(d))$，使得由
$s' \in V_m$ 给出的乘法零化
$H^1(\mathbf{P}^n_k, \mathcal{F}(m - d))$。
注意，对 $m' < m \leq m_0$，图表
$$
\xymatrix{
H^0(\mathbf{P}^n_k, \mathcal{O}(d)) \times
H^1(\mathbf{P}^n_k, \mathcal{F}(m' - d)) \ar[r] \ar[d]^{1 \times s^{m' - m}} &
H^1(\mathbf{P}^n_k, \mathcal{F}(m')) \ar[d]^{s^{m' - m}}\\
H^0(\mathbf{P}^n_k, \mathcal{O}(d)) \times
H^1(\mathbf{P}^n_k, \mathcal{F}(m - d)) \ar[r] &
H^1(\mathbf{P}^n_k, \mathcal{F}(m))
}
$$
可交换，且竖直映射均为同构。因此 $V_m = V$ 与
$m \leq m_0$ 无关。对 $x \in \text{Ass}(\mathcal{F})$，
置 $Z = \overline{\{x\}}$。当 $d$ 足够大时，线性映射
$$
H^0(\mathbf{P}^n_k, \mathcal{O}(d)) \to H^0(Z, \mathcal{O}(d)|_Z)
$$
的秩 $> N^2$，因为 $\dim(Z) \geq 1$
（例如，这由渐近 Riemann--Roch 与 $\mathcal{O}(1)$ 的丰沛性得出；略去细节）。
因此可找到 $s' \in V$，使得 $s'$ 在
$\mathcal{F}$ 的任何伴随点都不消失（使用伴随点集有限这一事实）。
于是得到
$$
0 \to \mathcal{F}(-d) \xrightarrow{s'} \mathcal{F} \to \mathcal{G}' \to 0
$$
；与前面一样可知，乘以 $s'$ 在
$H^1(\mathbf{P}^n_k, \mathcal{F}(m - d))$ 上是单射，只要 $m \ll 0$。
这与对 $s'$ 的选取矛盾，除非 $N = 0$，故结论成立。

\medskip\noindent
还须处理 $k$ 是有限域的情形。
此时，令 $K/k$ 为任意无限代数域扩张。
记 $\mathcal{F}_K$ 与 $\mathcal{L}_K$ 为
$\mathcal{F}$ 与 $\mathcal{L}$ 到
$X_K = \Spec(K) \times_{\Spec(k)} X$ 上的拉回。
由《概形的上同调》，引理
\ref{coherent-lemma-flat-base-change-cohomology}，有
$$
H^1(X_K, \mathcal{F}_K \otimes \mathcal{L}_K^{\otimes m}) =
H^1(X, \mathcal{F} \otimes \mathcal{L}^{\otimes m}) \otimes_k K
$$
。另一方面，闭点 $x_K$（它属于 $X_K$）映到闭点
$x$（它属于 $X$），因为 $K/k$ 是代数扩张。
环映射 $\mathcal{O}_{X, x} \to \mathcal{O}_{X_K, x_K}$
是平坦的（引理 \ref{varieties-lemma-change-fields-flat}）。因此由《代数》，引理
\ref{algebra-lemma-apply-grothendieck-module}，有
$$
\text{depth}(\mathcal{F}_{x_K}) =
\text{depth}(\mathcal{F}_x \otimes_{\mathcal{O}_{X, x}}
\mathcal{O}_{X_K, x_K}) \geq
\text{depth}(\mathcal{F}_x)
$$
（事实上这里等号成立，但我们不需要它）。
所以 $k$ 上的结果由无限域 $K$ 上的结果得出，证明完成。
\end{proof}

\begin{lemma}
\label{varieties-lemma-connectedness-ample-divisor}
设 $k$ 是域。设 $X$ 是 $k$ 上的本征概形。设 $\mathcal{L}$ 是丰沛可逆
$\mathcal{O}_X$-模。设 $s \in \Gamma(X, \mathcal{L})$。假设
\begin{enumerate}
\item $s$ 是正则截面
（《除子》，定义 \ref{divisors-definition-regular-section}），
\item 对每个闭点 $x \in X$，有
$\text{depth}(\mathcal{O}_{X, x}) \geq 2$，并且
\item $X$ 是连通的。
\end{enumerate}
则 $Z(s)$（即 $s$ 的零概形）是连通的。
\end{lemma}

\begin{proof}
由于 $s$ 是正则截面，
$s^n \in \Gamma(X, \mathcal{L}^{\otimes n})$ 对所有 $n > 1$ 也是正则截面。
此外，包含态射 $Z(s) \to Z(s^n)$ 在底拓扑空间上是双射。
因此，若 $Z(s)$ 不连通，则 $Z(s^n)$ 也不连通。
现在考虑典范短正合列
$$
0 \to \mathcal{L}^{\otimes -n} \xrightarrow{s^n}
\mathcal{O}_X \to \mathcal{O}_{Z(s^n)} \to 0
$$
。考虑 $k$-代数 $R_n = \Gamma(X, \mathcal{O}_{Z(s^n)})$。
若 $Z(s)$ 不连通，即 $Z(s^n)$ 不连通，则或者
$R_n$ 为零（当 $Z(s^n) = \emptyset$ 时），或者
$R_n$ 含有非平凡幂等元（当
$Z(s^n) = U \amalg V$，其中 $U, V \subset Z(s^n)$ 开且非空时；
读者可查阅引理
\ref{varieties-lemma-proper-geometrically-reduced-global-sections}）。
所以映射 $\Gamma(X, \mathcal{O}_X) \to R_n$ 不可能是同构。
从而，要么 $H^0(X, \mathcal{L}^{\otimes -n})$ 非零，要么
$H^1(X, \mathcal{L}^{\otimes -n})$ 对无穷多个正整数 $n$ 非零。
这与引理 \ref{varieties-lemma-vanishin-h0-negative} 或
\ref{varieties-lemma-vanishin-h1-negative} 矛盾，证明完成。
\end{proof}
